% !TeX program = xelatex
% Standalone English volume. Use XeLaTeX until the references stabilize.
\documentclass[a4paper,11pt]{article}
\usepackage[T1]{fontenc}
\usepackage{lmodern}
\usepackage[english]{babel}
\usepackage[left=26mm,right=26mm,top=25mm,bottom=25mm,headheight=15pt,headsep=7mm,footskip=11mm]{geometry}
\usepackage{amsmath,amssymb,mathtools,amsthm}
\newtheorem{proposition}{Proposition}[section]
\newtheorem{lemma}[proposition]{Lemma}
\newtheorem{theorem}[proposition]{Theorem}
\newtheorem{corollary}[proposition]{Corollary}
\newtheorem{maintheorem}[proposition]{Main Theorem}
\theoremstyle{definition}
\newtheorem{definition}[proposition]{Definition}
\theoremstyle{remark}
\newtheorem{remark}[proposition]{Remark}
\makeatletter
\renewcommand*\l@section{\@dottedtocline{1}{0em}{2.4em}}
\renewcommand*\l@subsection{\@dottedtocline{2}{1.45em}{3.45em}}
\renewcommand*\l@subsubsection{\@dottedtocline{3}{3.8em}{4.55em}}
\makeatother
\usepackage{xcolor,fancyhdr,titlesec,needspace,xurl}
\usepackage{longtable,multicol}
\usepackage{hyperref}
\definecolor{ink}{HTML}{1A3542}
\definecolor{muted}{HTML}{4F606A}
\hypersetup{unicode=true,colorlinks=true,linkcolor=ink,filecolor=ink,urlcolor=ink,citecolor=ink,bookmarksnumbered=true,
 pdfdisplaydoctitle=true,pdflang={en},
 pdftitle={The generalized Riemann hypothesis for Dirichlet L-functions. Edition 930, Volume XIV},
 pdfsubject={A natural-language reconstruction of the proposed proof},
 pdfauthor={A mathematical reconstruction}}
\titleformat{\section}{\Large\bfseries\color{ink}}{\thesection}{.65em}{}
\titlespacing*{\section}{0pt}{.8em}{.65em}
\setlength{\parindent}{1.2em}
\linespread{1.045}
\setlength{\parskip}{.22em plus .05em minus .03em}
\setlength{\emergencystretch}{2em}
\setcounter{topnumber}{2}
\setcounter{tocdepth}{3}\setcounter{secnumdepth}{3}
\widowpenalty=4500
\clubpenalty=4500
\pagestyle{fancy}\fancyhf{}
\fancyhead[L]{\small\color{muted}Dirichlet GRH \textperiodcentered\ Volume I}
\fancyhead[R]{\small\color{muted}Dirichlet zeros and power intervals}
\fancyfoot[C]{\small\thepage}
\renewcommand{\headrulewidth}{.3pt}
\newcommand{\lead}[1]{\par\smallskip\noindent{\bfseries\color{ink}#1.}\enspace}
\newcommand{\sour}[2]{\unskip\penalty150\hspace{.35em}{\small\hyperlink{ref#1}{[#2]}}}
\titleformat{\subsection}{\large\bfseries}{\thesubsection}{.6em}{}
\titlespacing*{\subsection}{0pt}{1.15em plus .15em}{.4em}
\setlength{\abovedisplayskip}{9pt plus 2pt minus 3pt}
\setlength{\belowdisplayskip}{9pt plus 2pt minus 3pt}
\newcommand{\RR}{\mathbb R}\newcommand{\CC}{\mathbb C}\newcommand{\QQ}{\mathbb Q}
\newcommand{\ZZ}{\mathbb Z}\newcommand{\ord}{\operatorname{ord}}
\newcommand{\Frob}{\operatorname{Frob}}\newcommand{\Gal}{\operatorname{Gal}}
\newcommand{\len}{\operatorname{length}}\newcommand{\Res}{\operatorname*{Res}}
\newcommand{\wmean}{\operatorname{av}}\newcommand{\Ker}{\operatorname{ker}}
\newcommand{\Lh}{\mathcal L}\newcommand{\Gg}{\mathcal G}
\newcommand{\real}{\operatorname{Re}}\newcommand{\imag}{\operatorname{Im}}
\newcommand{\abs}[1]{\left|#1\right|}\newcommand{\eps}{\varepsilon}
\newcommand{\anchoritem}[2]{\par\Needspace{2.8\baselineskip}\hypertarget{ref#1}{}\noindent{\bfseries\color{ink}#2.}\enspace}
\newcommand{\srcloc}[2]{{\small\nolinkurl{#1.lean}}, line~#2}
\numberwithin{equation}{section}
\raggedbottom

\newcommand{\Sel}{\mathcal S}
\newcommand{\Cl}{\operatorname{Cl}}
\newcommand{\Aut}{\operatorname{Aut}}
\newcommand{\GL}{\operatorname{GL}}
\newcommand{\OO}{\mathcal O}
\hypersetup{bookmarksdepth=3}


\addto\captionsenglish{\renewcommand{\partname}{Volume}}


\makeatletter
\expandafter\def\csname r@cor:actualdenom803\endcsname{{52.21}{80}{}{proposition.52.21}{Dirichlet-GRH-v930-Volume-X.pdf}}
\expandafter\def\csname r@cor:actualmarks714\endcsname{{12.10}{38}{}{proposition.12.10}{Dirichlet-GRH-v930-Volume-II.pdf}}
\expandafter\def\csname r@cor:fieldrelations712\endcsname{{12.4}{32}{}{proposition.12.4}{Dirichlet-GRH-v930-Volume-II.pdf}}
\expandafter\def\csname r@cor:fouriermean354\endcsname{{120.9}{20}{}{proposition.120.9}{Dirichlet-GRH-v930-Volume-XXIII.pdf}}
\expandafter\def\csname r@cor:original200\endcsname{{83.7}{57}{}{proposition.83.7}{Dirichlet-GRH-v930-Volume-XVI.pdf}}
\expandafter\def\csname r@cor:quotientchainheight734\endcsname{{13.14}{73}{}{proposition.13.14}{Dirichlet-GRH-v930-Volume-II.pdf}}
\expandafter\def\csname r@eq:gap\endcsname{{6.3}{26}{}{equation.6.3}{Dirichlet-GRH-v930-Volume-I.pdf}}
\expandafter\def\csname r@eq:letters111\endcsname{{9.1}{117}{}{equation.9.1}{Dirichlet-GRH-v930-Volume-I.pdf}}
\expandafter\def\csname r@eq:region198\endcsname{{20.12}{33}{}{equation.20.12}{Dirichlet-GRH-v930-Volume-IV.pdf}}
\expandafter\def\csname r@eq:src109\endcsname{{87.4}{37}{}{equation.87.4}{Dirichlet-GRH-v930-Volume-XVII.pdf}}
\expandafter\def\csname r@eq:tail114\endcsname{{19.15}{17}{}{equation.19.15}{Dirichlet-GRH-v930-Volume-IV.pdf}}
\expandafter\def\csname r@eq:weil161\endcsname{{93.1}{21}{}{equation.93.1}{Dirichlet-GRH-v930-Volume-XVIII.pdf}}
\expandafter\def\csname r@eq:window\endcsname{{{W}}{2}{}{AMS.2}{Dirichlet-GRH-v930-Volume-I.pdf}}
\expandafter\def\csname r@lem:augmentedindex687\endcsname{{131.16}{27}{}{proposition.131.16}{Dirichlet-GRH-v930-Volume-XXV.pdf}}
\expandafter\def\csname r@lem:bilinearsource710\endcsname{{142.11}{44}{}{proposition.142.11}{Dirichlet-GRH-v930-Volume-XXVII.pdf}}
\expandafter\def\csname r@lem:dimension239\endcsname{{94.22}{55}{}{proposition.94.22}{Dirichlet-GRH-v930-Volume-XVIII.pdf}}
\expandafter\def\csname r@lem:exchangeroot336\endcsname{{86.33}{26}{}{proposition.86.33}{Dirichlet-GRH-v930-Volume-XVII.pdf}}
\expandafter\def\csname r@lem:lattices247\endcsname{{106.1}{52}{}{proposition.106.1}{Dirichlet-GRH-v930-Volume-XX.pdf}}
\expandafter\def\csname r@lem:openflags353\endcsname{{120.7}{18}{}{proposition.120.7}{Dirichlet-GRH-v930-Volume-XXIII.pdf}}
\expandafter\def\csname r@lem:orbvol263\endcsname{{111.18}{73}{}{proposition.111.18}{Dirichlet-GRH-v930-Volume-XXI.pdf}}
\expandafter\def\csname r@lem:oscillatorcartan710\endcsname{{142.12}{45}{}{proposition.142.12}{Dirichlet-GRH-v930-Volume-XXVII.pdf}}
\expandafter\def\csname r@lem:rank\endcsname{{19.3}{4}{}{proposition.19.3}{Dirichlet-GRH-v930-Volume-IV.pdf}}
\expandafter\def\csname r@lem:sharpplaces709\endcsname{{142.7}{41}{}{proposition.142.7}{Dirichlet-GRH-v930-Volume-XXVII.pdf}}
\expandafter\def\csname r@lem:size97\endcsname{{87.1}{33}{}{proposition.87.1}{Dirichlet-GRH-v930-Volume-XVII.pdf}}
\expandafter\def\csname r@lem:tametraceheight734\endcsname{{13.12}{72}{}{proposition.13.12}{Dirichlet-GRH-v930-Volume-II.pdf}}
\expandafter\def\csname r@lem:terminalclass259\endcsname{{111.2}{58}{}{proposition.111.2}{Dirichlet-GRH-v930-Volume-XXI.pdf}}
\expandafter\def\csname r@lem:top\endcsname{{19.2}{3}{}{proposition.19.2}{Dirichlet-GRH-v930-Volume-IV.pdf}}
\expandafter\def\csname r@lem:translate163\endcsname{{86.19}{16}{}{proposition.86.19}{Dirichlet-GRH-v930-Volume-XVII.pdf}}
\expandafter\def\csname r@lem:vacuumtail709\endcsname{{142.8}{42}{}{proposition.142.8}{Dirichlet-GRH-v930-Volume-XXVII.pdf}}
\expandafter\def\csname r@lem:weightedhermite260\endcsname{{106.10}{64}{}{proposition.106.10}{Dirichlet-GRH-v930-Volume-XX.pdf}}
\expandafter\def\csname r@nav:I-28\endcsname{{6.6}{32}{}{subsection.6.6}{Dirichlet-GRH-v930-Volume-I.pdf}}
\expandafter\def\csname r@nav:VII-12\endcsname{{86.8}{19}{}{subsection.86.8}{Dirichlet-GRH-v930-Volume-XVII.pdf}}
\expandafter\def\csname r@nav:VII-13\endcsname{{93.2}{23}{}{subsection.93.2}{Dirichlet-GRH-v930-Volume-XVIII.pdf}}
\expandafter\def\csname r@prop:besorbit265\endcsname{{88.4}{56}{}{proposition.88.4}{Dirichlet-GRH-v930-Volume-XVII.pdf}}
\expandafter\def\csname r@prop:centeredflag341\endcsname{{99.18}{58}{}{proposition.99.18}{Dirichlet-GRH-v930-Volume-XIX.pdf}}
\expandafter\def\csname r@prop:depth245\endcsname{{95.3}{67}{}{proposition.95.3}{Dirichlet-GRH-v930-Volume-XVIII.pdf}}
\expandafter\def\csname r@prop:graphframe340\endcsname{{83.24}{71}{}{proposition.83.24}{Dirichlet-GRH-v930-Volume-XVI.pdf}}
\expandafter\def\csname r@prop:moment249\endcsname{{106.11}{65}{}{proposition.106.11}{Dirichlet-GRH-v930-Volume-XX.pdf}}
\expandafter\def\csname r@prop:oppline256\endcsname{{102.8}{96}{}{proposition.102.8}{Dirichlet-GRH-v930-Volume-XIX.pdf}}
\expandafter\def\csname r@prop:poincarenorm261\endcsname{{106.3}{55}{}{proposition.106.3}{Dirichlet-GRH-v930-Volume-XX.pdf}}
\expandafter\def\csname r@prop:polarcoefficient697\endcsname{{132.7}{56}{}{proposition.132.7}{Dirichlet-GRH-v930-Volume-XXV.pdf}}
\expandafter\def\csname r@prop:positions264\endcsname{{121.4}{33}{}{proposition.121.4}{Dirichlet-GRH-v930-Volume-XXIII.pdf}}
\expandafter\def\csname r@prop:rank11\endcsname{{6.9}{30}{}{proposition.6.9}{Dirichlet-GRH-v930-Volume-I.pdf}}
\expandafter\def\csname r@prop:ranks47\endcsname{{92.1}{3}{}{proposition.92.1}{Dirichlet-GRH-v930-Volume-XVIII.pdf}}
\expandafter\def\csname r@prop:rowfinite262\endcsname{{88.15}{70}{}{proposition.88.15}{Dirichlet-GRH-v930-Volume-XVII.pdf}}
\expandafter\def\csname r@prop:splitvalue257\endcsname{{25.26}{45}{}{proposition.25.26}{Dirichlet-GRH-v930-Volume-V.pdf}}
\expandafter\def\csname r@prop:terminalcore255\endcsname{{95.4}{68}{}{proposition.95.4}{Dirichlet-GRH-v930-Volume-XVIII.pdf}}
\expandafter\def\csname r@prop:tracelat253\endcsname{{122.7}{63}{}{proposition.122.7}{Dirichlet-GRH-v930-Volume-XXIII.pdf}}
\expandafter\def\csname r@prop:unfold263\endcsname{{88.9}{63}{}{proposition.88.9}{Dirichlet-GRH-v930-Volume-XVII.pdf}}
\expandafter\def\csname r@prop:vacuumcore254\endcsname{{122.8}{64}{}{proposition.122.8}{Dirichlet-GRH-v930-Volume-XXIII.pdf}}
\expandafter\def\csname r@sec:depth163\endcsname{{86.8}{19}{}{subsection.86.8}{Dirichlet-GRH-v930-Volume-XVII.pdf}}
\expandafter\def\csname r@sec:levelmass678\endcsname{{130.1}{3}{}{subsection.130.1}{Dirichlet-GRH-v930-Volume-XXV.pdf}}
\expandafter\def\csname r@sec:missing\endcsname{{6.6}{32}{}{subsection.6.6}{Dirichlet-GRH-v930-Volume-I.pdf}}
\expandafter\def\csname r@sec:pathnorm352\endcsname{{119.7}{11}{}{subsection.119.7}{Dirichlet-GRH-v930-Volume-XXIII.pdf}}
\expandafter\def\csname r@sec:samecritical731\endcsname{{13.1}{60}{}{subsection.13.1}{Dirichlet-GRH-v930-Volume-II.pdf}}
\expandafter\def\csname r@sec:whit125\endcsname{{92.6}{13}{}{subsection.92.6}{Dirichlet-GRH-v930-Volume-XVIII.pdf}}
\expandafter\def\csname r@seriesguide\endcsname{{1.1}{4}{}{subsection.1.1}{Dirichlet-GRH-v930-Volume-I.pdf}}
\expandafter\def\csname r@thm:actualminors804\endcsname{{52.24}{82}{}{proposition.52.24}{Dirichlet-GRH-v930-Volume-X.pdf}}
\expandafter\def\csname r@thm:actualselmer803\endcsname{{52.20}{79}{}{proposition.52.20}{Dirichlet-GRH-v930-Volume-X.pdf}}
\expandafter\def\csname r@thm:auxiliaryflag339\endcsname{{99.17}{57}{}{proposition.99.17}{Dirichlet-GRH-v930-Volume-XIX.pdf}}
\expandafter\def\csname r@thm:auxoverlap341\endcsname{{99.19}{59}{}{proposition.99.19}{Dirichlet-GRH-v930-Volume-XIX.pdf}}
\expandafter\def\csname r@thm:common111\endcsname{{9.13}{118}{}{proposition.9.13}{Dirichlet-GRH-v930-Volume-I.pdf}}
\expandafter\def\csname r@thm:criticalheights734\endcsname{{13.13}{72}{}{proposition.13.13}{Dirichlet-GRH-v930-Volume-II.pdf}}
\expandafter\def\csname r@thm:depth164\endcsname{{110.3}{39}{}{proposition.110.3}{Dirichlet-GRH-v930-Volume-XXI.pdf}}
\expandafter\def\csname r@thm:end14\endcsname{{7.3}{90}{}{proposition.7.3}{Dirichlet-GRH-v930-Volume-I.pdf}}
\expandafter\def\csname r@thm:exchangefourier336\endcsname{{86.34}{26}{}{proposition.86.34}{Dirichlet-GRH-v930-Volume-XVII.pdf}}
\expandafter\def\csname r@thm:finitegram709\endcsname{{142.9}{43}{}{proposition.142.9}{Dirichlet-GRH-v930-Volume-XXVII.pdf}}
\expandafter\def\csname r@thm:fixed\endcsname{{19.6}{6}{}{proposition.19.6}{Dirichlet-GRH-v930-Volume-IV.pdf}}
\expandafter\def\csname r@thm:flagfinite264\endcsname{{90.6}{88}{}{proposition.90.6}{Dirichlet-GRH-v930-Volume-XVII.pdf}}
\expandafter\def\csname r@thm:flagindex336\endcsname{{99.14}{55}{}{proposition.99.14}{Dirichlet-GRH-v930-Volume-XIX.pdf}}
\expandafter\def\csname r@thm:fourierslopes354\endcsname{{120.8}{20}{}{proposition.120.8}{Dirichlet-GRH-v930-Volume-XXIII.pdf}}
\expandafter\def\csname r@thm:globalgram711\endcsname{{142.15}{48}{}{proposition.142.15}{Dirichlet-GRH-v930-Volume-XXVII.pdf}}
\expandafter\def\csname r@thm:graph274\endcsname{{94.27}{59}{}{proposition.94.27}{Dirichlet-GRH-v930-Volume-XVIII.pdf}}
\expandafter\def\csname r@thm:harmonic188\endcsname{{104.28}{24}{}{proposition.104.28}{Dirichlet-GRH-v930-Volume-XX.pdf}}
\expandafter\def\csname r@thm:incidence338\endcsname{{86.37}{29}{}{proposition.86.37}{Dirichlet-GRH-v930-Volume-XVII.pdf}}
\expandafter\def\csname r@thm:initial247\endcsname{{89.2}{74}{}{proposition.89.2}{Dirichlet-GRH-v930-Volume-XVII.pdf}}
\expandafter\def\csname r@thm:invar148\endcsname{{20.11}{31}{}{proposition.20.11}{Dirichlet-GRH-v930-Volume-IV.pdf}}
\expandafter\def\csname r@thm:long232\endcsname{{5.6}{23}{}{proposition.5.6}{Dirichlet-GRH-v930-Volume-I.pdf}}
\expandafter\def\csname r@thm:markedcore253\endcsname{{102.7}{95}{}{proposition.102.7}{Dirichlet-GRH-v930-Volume-XIX.pdf}}
\expandafter\def\csname r@thm:markedstrict337\endcsname{{22.11}{70}{}{proposition.22.11}{Dirichlet-GRH-v930-Volume-IV.pdf}}
\expandafter\def\csname r@thm:midlevel131\endcsname{{86.1}{3}{}{proposition.86.1}{Dirichlet-GRH-v930-Volume-XVII.pdf}}
\expandafter\def\csname r@thm:moment258\endcsname{{83.14}{63}{}{proposition.83.14}{Dirichlet-GRH-v930-Volume-XVI.pdf}}
\expandafter\def\csname r@thm:norm91\endcsname{{94.9}{44}{}{proposition.94.9}{Dirichlet-GRH-v930-Volume-XVIII.pdf}}
\expandafter\def\csname r@thm:ordinary163\endcsname{{86.22}{18}{}{proposition.86.22}{Dirichlet-GRH-v930-Volume-XVII.pdf}}
\expandafter\def\csname r@thm:ordinary164\endcsname{{110.2}{39}{}{proposition.110.2}{Dirichlet-GRH-v930-Volume-XXI.pdf}}
\expandafter\def\csname r@thm:pintable340\endcsname{{83.25}{72}{}{proposition.83.25}{Dirichlet-GRH-v930-Volume-XVI.pdf}}
\expandafter\def\csname r@thm:pointlattice353\endcsname{{99.20}{60}{}{proposition.99.20}{Dirichlet-GRH-v930-Volume-XIX.pdf}}
\expandafter\def\csname r@thm:range239\endcsname{{88.1}{52}{}{proposition.88.1}{Dirichlet-GRH-v930-Volume-XVII.pdf}}
\expandafter\def\csname r@thm:rangenew50\endcsname{{6.29}{62}{}{proposition.6.29}{Dirichlet-GRH-v930-Volume-I.pdf}}
\expandafter\def\csname r@thm:ratio248\endcsname{{89.3}{76}{}{proposition.89.3}{Dirichlet-GRH-v930-Volume-XVII.pdf}}
\expandafter\def\csname r@thm:realgram710\endcsname{{142.13}{46}{}{proposition.142.13}{Dirichlet-GRH-v930-Volume-XXVII.pdf}}
\expandafter\def\csname r@thm:restriction246\endcsname{{89.1}{73}{}{proposition.89.1}{Dirichlet-GRH-v930-Volume-XVII.pdf}}
\expandafter\def\csname r@thm:secondbound260\endcsname{{89.4}{77}{}{proposition.89.4}{Dirichlet-GRH-v930-Volume-XVII.pdf}}
\expandafter\def\csname r@thm:short246\endcsname{{111.1}{57}{}{proposition.111.1}{Dirichlet-GRH-v930-Volume-XXI.pdf}}
\expandafter\def\csname r@thm:tensor249\endcsname{{106.9}{62}{}{proposition.106.9}{Dirichlet-GRH-v930-Volume-XX.pdf}}
\expandafter\def\csname r@thm:terminalformula347\endcsname{{127.1}{26}{}{proposition.127.1}{Dirichlet-GRH-v930-Volume-XXIV.pdf}}
\expandafter\def\csname r@thm:tests165\endcsname{{93.5}{25}{}{proposition.93.5}{Dirichlet-GRH-v930-Volume-XVIII.pdf}}
\expandafter\def\csname r@thm:trace166\endcsname{{83.12}{61}{}{proposition.83.12}{Dirichlet-GRH-v930-Volume-XVI.pdf}}
\expandafter\def\csname r@thm:transverse250\endcsname{{106.2}{53}{}{proposition.106.2}{Dirichlet-GRH-v930-Volume-XX.pdf}}
\makeatother
\begin{document}
\setcounter{part}{13}
\fancyhead[L]{\small\color{muted}Dirichlet GRH \textperiodcentered\ Volume XIV}
\fancyhead[R]{\small\color{muted}Branching and first occurrence}
\thispagestyle{plain}\begin{center}
{\LARGE\bfseries The generalized Riemann hypothesis\\[2pt]for Dirichlet \(L\)-functions}\\[6pt]
{\small 27 November 2026 \quad\textperiodcentered\quad Edition 930}\end{center}
\part{Unitary branching and first occurrence in the theta tower}\label{part:XIV}
\begin{abstract}\noindent We follow a prescribed representation through the unitary theta correspondence. Branching, seesaw identities and terminal Hecke operators determine the rows used for the global lift, its first occurrence and the norms on the even core.\end{abstract}
\setcounter{section}{69}
\noindent The \hyperref[seriesguide]{guide in Volume~I} lists all twenty-seven volumes.
\begingroup\small\setlength{\parskip}{0pt}\tableofcontents\endgroup
\clearpage

\section{Unitary branching and the global seesaw}\label{nav:g-VI-0}\label{nav:volume-XIV}

\subsection{Integral projectors for compact unitary branching}\label{nav:I-37}
Let \(E/E^+\) be a CM extension. A Hermitian lattice in this volume is a projective \(\OO_E\)-module of finite rank with a nondegenerate Hermitian form. It is \emph{unimodular} if the form identifies it with its Hermitian dual, and \emph{definite} if the induced forms at all complex embeddings are positive definite. Fix a coefficient field \(J\) containing the images of the embeddings of \(E\) to be used. For an embedding \(\tau:E\hookrightarrow J\) and a lattice \(A\), put \(A_\tau=A\otimes_{\OO_E,\tau}\OO_J\).

We use the irreducible algebraic representation of \(\GL_n\) of highest weight \(\lambda=(\lambda_1\ge\cdots\ge\lambda_n)\in\ZZ^n\). We specify an integral model \(W_\lambda(A_\tau)\). For the partition \(\nu_i=\lambda_i-\lambda_n\), take the image of an integral Young symmetrizer of shape \(\nu\) on \(A_\tau^{\otimes|\nu|}\) and tensor with \((\det A_\tau)^{\lambda_n}\), using the dual line for negative powers. The empty partition gives the trivial representation. Young symmetrizers in characteristic zero give the desired irreducible space after extension of scalars. The image construction commutes with localization, and its image is torsion free and hence projective over the Dedekind ring \(\OO_J\). In a local basis the operators \(e_{ab}\) of the matrix Lie algebra preserve this lattice. Indeed, their action on a tensor power is a sum of integral actions on the individual factors and commutes with permutations. The derivative on the determinant factor is \(\lambda_n\) times the trace. At a complex embedding we use a Hermitian norm invariant under the unitary group of \(A\); the tensor construction supplies such a norm.

\begin{theorem}[A common denominator for the branching projections]\label{thm:proj}
Suppose \(A\) is definite and unimodular of rank \(n\). Let
\[
 0\subset F_1\subset\cdots\subset F_{n-1}\subset F_n=A,\qquad
                 \operatorname{rank}F_j=j
\]
be a flag of unimodular Hermitian sublattices. Write \(h=\lambda_1-\lambda_n\). For restriction along the corresponding compact groups \(U(F_1)\subset\cdots\subset U(A)\), every joint isotypic projection \(P_{F,\boldsymbol\mu}\) for a complete branching pattern is defined over \(J\). A single positive integer clears all these projections, for every flag of this kind:
\begin{equation}
         B_\lambda=\prod_{j=1}^{n-1}\big((h+j-1)!\big)^j .
         \label{eq:den}
\end{equation}
In the indicated integral model \(B_\lambda P_{F,\boldsymbol\mu}\) preserves the lattice. For a fixed flag the nonzero projections are mutually orthogonal, sum to the identity and have operator norm one.
\end{theorem}
\begin{proof}
For the inclusion \(s:F_j\hookrightarrow A\), the adjoint sends \(A\) into \(F_j\) by unimodularity of \(F_j\), and \(s^*s=1\). Thus \(A=sF_j\oplus\ker s^*\) integrally and orthogonally. The same argument for consecutive members splits the full flag into orthogonal projective lines. The group \(U(F_j)\) acts as the identity on its orthogonal complement. Over \(J\) we obtain the corresponding block \(\mathfrak{gl}_j\subset\mathfrak{gl}_n\). At each finite place a basis can be chosen for all the lines simultaneously.

For this block use the column Capelli polynomial
\[
 C_j(z)=\sum_{\pi\in\mathfrak S_j}\operatorname{sgn}(\pi)
       \prod_{b=1}^{j}\big(e_{\pi(b),b}
                    +\mathbf1_{\pi(b)=b}(z+j-b)\big).
\]
Multiply the factors in increasing order of \(b\). The classical Capelli identity makes its coefficients central in \(U(\mathfrak{gl}_j)\). On highest weight \(\gamma=(\gamma_1,\ldots,\gamma_j)\) its scalar is \(\prod_{b=1}^j(z+\gamma_b+j-b)\). This is the normalization of Molev's formulas (2.21)--(2.22) after a translation of the variable by \(j-1\). Centrality also gives invariance under a change of block basis, since \(\GL_j\) in characteristic zero is connected and generated by its torus and root groups. In any integral block basis \(C_j(k)\) preserves \(W_\lambda(A_\tau)\) locally for \(k\in\ZZ\), by the integral formula above. Basis independence identifies these local operators with a single rational operator. Preservation locally everywhere gives preservation of the global lattice.~\cite[Sec.~2]{XIV-XI-IV-molev}

By the branching law for \(\GL_{j+1}\) restricted to \(\GL_j\), the next highest weight interlaces the previous one, each with multiplicity one. This follows, for example, by expanding the Schur character in its last variable: removing all entries equal to the largest letter in a semistandard tableau leaves a smaller tableau and a horizontal strip. Its row lengths are precisely the interlacing possibilities, and each filling of the smaller shape has one extension for each admissible strip. The determinant twist shifts all entries equally. Restricting the resulting character identity to \(U(j)\) and using independence of irreducible characters gives the assertion there. Iteration places every entry between \(\lambda_n\) and \(\lambda_1\).

Fix a desired weight \(\mu\) at level \(j\), and write
\[
 I_j=[\lambda_n,\lambda_1+j-1]\cap\ZZ,\qquad
 S_\mu=\{\mu_b+j-b:1\le b\le j\}.
\]
The \(j\) members of \(S_\mu\) are distinct and in decreasing order, so they determine \(\mu\). Consider
\begin{equation}
 D_{j,\mu}=\prod_{s\in I_j\setminus S_\mu}C_j(-s),\qquad
 a_{j,\mu}=\prod_{s\in I_j\setminus S_\mu}\ \prod_{d\in S_\mu}(d-s).
 \label{eq:sieve}
\end{equation}
The integer \(a_{j,\mu}\) is nonzero, with value one for the empty product. On a summand of type \(\mu\), \(D_{j,\mu}\) is multiplication by this integer. On any other type there is a shifted entry \(s\) outside \(S_\mu\). The corresponding factor kills that entire summand. Thus \(P_{j,\mu}=D_{j,\mu}/a_{j,\mu}\) projects onto the full isotypic part, including any multiplicity. If \(I_j=[u,v]\cap\ZZ\), the absolute value of the product over \(s\notin S_\mu\) for a fixed \(d\) divides \((d-u)!(v-d)!\), which divides \((v-u)!\). Multiplying over \(d\) proves \(|a_{j,\mu}|\mid((h+j-1)!)^j\).

The centers of the nested blocks commute, since \(U(\mathfrak{gl}_j)\) lies in \(U(\mathfrak{gl}_k)\) for \(j<k\). The product of the single-level projectors over the rows of a pattern is therefore its joint projector. For a unitary representation, inequivalent isotypic spaces are orthogonal. The single-level projectors are self-adjoint, and their commuting products are the orthogonal projections for the successive restrictions. The branching law exhausts the original space. Multiplying the individual denominators proves \eqref{eq:den}.
\end{proof}

\begin{corollary}[The induced integral decomposition]\label{cor:sandwich}
For a fixed flag set \(M=W_\lambda(A_\tau)\). For each pattern with projector \(P\), write \(V_P=P(M\otimes J)\) and \(M_P=M\cap V_P\). Then
\[
       M_P\subseteq P(M)\subseteq B_\lambda^{-1}M_P,\qquad
       B_\lambda M\subseteq\bigoplus_P M_P\subseteq M .
\]
If several embeddings have designated weights \(\lambda_\tau\), every tensor product of such projections on \(\bigotimes_\tau W_{\lambda_\tau}(A_\tau)\) has denominator dividing \(B=\prod_\tau B_{\lambda_\tau}\).
\end{corollary}
\begin{proof}
An element in \(M_P\) is fixed by \(P\). For \(v\in M\), \(B_\lambda P v\) is integral and belongs to \(V_P\). Summing these integral elements over the pattern decomposition gives \(B_\lambda v\), and distinct subspaces have direct sum. Tensoring the integral multiples of the projectors gives the last statement.
\end{proof}

\subsection{Projected coefficients of Hermitian lattice theta series}\label{nav:I-38}
For \(1\le j\le n\) let \(\mathcal C^{(j)}\) be the genus of the standard positive definite Hermitian lattice \(\OO_E^j\), taken up to isometry. Its members are unimodular because unimodularity can be checked at finite places. The genus is finite by the double-coset finiteness theorem for definite unitary groups, and \(w(A)=|\Aut(A)|\) is finite since a positive definite trace form is preserved. For each \(A\in\mathcal C^{(n)}\), let \(\mathcal F(A)\) consist of the embedded flags of Theorem~\ref{thm:proj} whose \(F_j\) belongs to \(\mathcal C^{(j)}\). This is a set of flags inside a chosen representative, without dividing by its automorphisms.~\cite{XIV-XI-IV-gross}

On \(\mathcal C^{(n)}\) take equivariant assignments \(v_A\in M(A)\) and \(h_A\in M(A)^\vee=\operatorname{Hom}_{\OO_J}(M(A),\OO_J)\), where \(M(A)=\bigotimes_\tau W_{\lambda_\tau}(A_\tau)\). Equivariance means that the value transports under isometries, so on representatives these are stabilizer-invariant vectors and covectors. Let \(f\) be any \(\OO_J\)-valued function on \(\mathcal C^{(1)}\). For a pattern at each embedding use \(P_F\) for the resulting tensor projector in the given flag.

\begin{theorem}[A projected theta coefficient on the common family]\label{thm:theta}
The period
\begin{equation}
 {\mathcal P}(f,v,h)=
   \sum_{A\in\mathcal C^{(n)}}\frac{1}{w(A)}
     \sum_{F\in\mathcal F(A)} f([F_1])\,h_A(P_Fv_A)
 \label{eq:period}
\end{equation}
is a mass-weighted sum of polynomially weighted Hermitian theta coefficients. With \(d_E=[E:\QQ]\) and \(D_n=|\GL_{nd_E}(\mathbf F_3)|\) one has
\[
              D_n B\,{\mathcal P}(f,v,h)\in\OO_J.
\]
Summing first from the rank-one end or first from the rank-\(n\) end gives the same period, with the same denominator. On the fields \(E=E_x\) of Theorem \ref{thm:fixed}, take \(n=\lfloor\sqrt x\rfloor+1\) and at most one weight per embedding in a CM type. If \(H\) is the largest weight width \(\lambda_{\tau,1}-\lambda_{\tau,n}\), then
\[
 \log(D_n B)\ll_\psi x^7+
     x^4(H+\sqrt x+1)\log(2+H+\sqrt x).
\]
\end{theorem}
\begin{proof}
Choose representatives \(A_1,\ldots,A_n=A\) from the corresponding genera. For \(j<n\) let \(s_j:A_j\hookrightarrow A_{j+1}\) run over all isometries onto their images. There are finitely many choices because the images of a fixed \(\ZZ\)-basis have fixed lengths in positive definite trace lattices. Put \(S_j=s_{n-1}\cdots s_j\). A tuple of maps determines the flag \(F_j=S_j A_j\). Conversely, for a flag of the specified classes, choose isometries \(A_j\to F_j\) independently for \(j<n\), using the identity at \(j=n\). These choices recover a unique tuple, and their number is \(\prod_{j<n}w(A_j)\). The projector depends only on the images. Consequently \eqref{eq:period} equals
\begin{equation}
 \sum_{[A_1],\ldots,[A_n]}
    \frac{f([A_1])}{\prod_{j=1}^{n}w(A_j)}
      \sum_{s_1,\ldots,s_{n-1}}h_{A_n}(P_{F(s)}v_{A_n}).
 \label{eq:tuple}
\end{equation}
All terms here are finite sums, giving the claimed equality of summation orders. For scalar trivial weights \(B=1\) and \eqref{eq:tuple} is the chain of Theorem \ref{thm:end14}.

We justify the theta coefficient interpretation of the inner sum. Use fixed \(E\)-bases for the rational spaces attached to \(A_j\). In the block for the image of \(S_j\), a matrix unit acts on the ambient space by \(S_j e_{ab} S_j^*\), where \(S_j^*\) is the adjoint for the fixed Hermitian forms. On the isometry locus this agrees with its extension by zero on the orthogonal complement because \(S_j^*S_j=1\). For a general tuple of linear maps \(s_j\) the same expression is polynomial in their entries and their complex conjugates, with constants supplied by the fixed forms. Applying the derived representation is a linear operation on this matrix. Formula \eqref{eq:sieve}, with its constant denominator, thus gives a polynomial insertion on the full product of the spaces of maps whose value on isometries is \(h_{A_n}(P_{F(s)}v_{A_n})\). This description uses rational bases only to identify the polynomial.

In detail, sum this polynomial over \(\prod_{j<n}\operatorname{Hom}_{\OO_E}(A_j,A_{j+1})\), recording in independent variables the pullback Hermitian forms \(s_j^*H_{A_{j+1}}s_j\). This is the polynomially weighted Hermitian theta series meant here. It can either be read coefficientwise or multiplied term by term by Gaussian weights from positive definite trace norms on these spaces of maps. Each coefficient is finite, and the series with these Gaussian weights converges absolutely because a polynomial times a Gaussian is summable on a lattice. Extract the simultaneous coefficient at the tuple of forms \((H_{A_1},\ldots,H_{A_{n-1}})\). Its terms are exactly the isometries in \eqref{eq:tuple}. This identifies the coefficient independently of any automorphic period normalization.

Now use \eqref{eq:period}. The integer \(B\) makes each projected weight integral by Corollary~\ref{cor:sandwich}. Reduction modulo three injects the finite group \(\Aut(A)\) into \(\GL_{nd_E}(\mathbf F_3)\); the prime-order binomial argument is given in Lemma \ref{lem:top}. Hence \(w(A)\mid D_n\), and the claimed denominator follows. The flag count introduces no further denominator.

Finally \(d_E\le4\operatorname{ord}(\psi)2^l\ll_\psi x^3\) for \(0<\eps<1/52\), and \(n\le2\sqrt x\). Thus \(\log D_n\le (nd_E)^2\log3\ll_\psi x^7\). At one embedding \(\log B_\lambda\le n^2(H+n)\log(2+H+n)\) by \eqref{eq:den}. Sum over at most \(d_E/2\) embeddings to obtain the estimate.
\end{proof}

\begin{proposition}[A test with distinct flags]\label{prop:flag}
Take \(A=\OO_E^n\) with its standard form, \(n\ge3\), and the defining representation at one embedding. In each complete integral flag choose the pattern whose projection is onto \(F_1\). If \(P_F\) denotes this projection, then
\[
 \frac{1}{n!}\sum_{F\in\mathcal F(A)}P_F=\frac{I}{n}.
\]
There nevertheless exist distinct flags with \(\operatorname{tr}(P_FP_{F'})=1\).
\end{proposition}
\begin{proof}
A unimodular sublattice determines an orthogonal idempotent \(P\) on \(\OO_E^n\) by integral splitting. At every archimedean embedding each diagonal entry of \(P\) belongs to \([0,1]\). Such an entry is an algebraic integer of \(E^+\). If it is nonzero, its absolute norm is a positive integer at most one. All conjugates must then equal one. Hence every diagonal entry is zero or one. Since \(P=P^*=P^2\), the sum of the squared moduli in any column equals its diagonal entry. This forces all off-diagonal entries to vanish. The unimodular sublattices are therefore coordinate summands. They all lie in the prescribed genera. The \(n!\) complete flags correspond to orderings of the coordinates, and each coordinate is first in \((n-1)!\) of them. This proves the average. The orderings \((1,2,3,\ldots)\) and \((1,3,2,\ldots)\) give different flags and both project onto the first line. Their projectors have product of trace one. The calculation is on the weight space. It also applies at an endpoint with a full marking modulo three, whose stabilizer is trivial.
\end{proof}

This example distinguishes two estimates required in applying a period comparison. Averaging over all flags here gives an exact scalar operator. Translates in the proposed prime count select particular maps. For two rank-one images in an \(n\)-dimensional target, a bound of order \(1/n\) for their squared overlap would already fail for the two flags above. Whether a given arithmetic family of translates has small overlaps depends on its maps and on their distribution.

\subsection{The global seesaw for a definite unitary period}\label{nav:I-39}
Write \(k=E^+\), let \(\mathbb A\) be its adeles, and put \([G]=G(k)\backslash G(\mathbb A)\). Fix a nontrivial additive character of \(\mathbb A/k\). Let \(V=V'\perp D\) be a definite Hermitian space, where \(D\) is a line, and let \(W\) be a nonzero skew-Hermitian space such that \(U(W)\) is compact at all real places. For example, multiply a positive definite Hermitian form by a nonzero scalar of \(E\) negated by conjugation. Set \(G'=U(V')\), \(H=U(W)\), and view \(G'\) inside \(U(V)\) by extending every isometry by the identity on \(D\).

We recall the normalization that matters for the identity. The unitary dual pair \(U(V)\times H\) acts on the symplectic space over \(k\) obtained from \(V\otimes_E W\) with its alternating trace form. Use a rational polarization compatible with the orthogonal sum, with position spaces \(Y=Y'\oplus Y_D\). For each of the three pairs use the global Weil representation for the fixed additive character and an automorphic unitary splitting. Such splittings are supplied by the standard unitary theta construction with splitting characters. They agree with the rational splitting on rational points. If \(\phi\) is an adelic Schwartz function on the position space, the kernel is
\[
 \Theta_V(g,h;\phi)=\sum_{y\in Y(k)}(\omega_V(g,h)\phi)(y).
\]
Use corresponding notation for \(V'\) and \(D\). Haar measures are fixed throughout; in a given integral the same measure is used on both sides. These constructions and the rational invariance of the theta functional are the usual global Weil representation and Poisson summation.~\cite{XIV-XI-IV-gan}

\begin{theorem}[Restriction of a theta lift against a period]\label{thm:seesaw}
There are unitary automorphic characters \(\delta_1\) of \(G'(\mathbb A)\) and \(\delta_2\) of \(H(\mathbb A)\), determined by the chosen splittings, with the following property. Take \(\phi=\phi'\otimes\eta\) and continuous automorphic functions \(a\) on \([H]\), \(b\) on \([G']\). Define
\[
 L_\phi(a)(g)=\int_{[H]}\Theta_V(g,h;\phi)a(h)\,dh .
\]
The codimension-one period of this theta lift is given by
\begin{equation}
 \begin{split}
 &\int_{[G']} L_\phi(a)(g'\oplus1)\,\overline{b(g')}\,dg'\\
 &=\int_{[H]}\delta_2(h)a(h)\Theta_D(1,h;\eta)
       \left(\int_{[G']}\delta_1(g')\Theta_{V'}(g',h;\phi')
                           \overline{b(g')}\,dg'\right)\,dh .
 \end{split}\label{eq:seesaw}
\end{equation}
For compatible splittings \(\delta_1=\delta_2=1\). The equality remains valid after an archimedean isotypic projection on the restricted form, by applying the same projection to \(b\) on the right.
\end{theorem}
\begin{proof}
In the direct sum of symplectic spaces the Heisenberg representation in the chosen polarization is the tensor product on the two position spaces. Under the standard homomorphism of the two metaplectic covers into the cover of the sum, the oscillator action is the tensor product. This can also be seen on the Schr\"odinger generators: changes of coordinates and quadratic characters separate, and the Fourier transform of a product in the direct sum is the product of the two Fourier transforms. On \(G'(\mathbb A)\times H(\mathbb A)\) the two splittings to compare are thus lifts of the same symplectic action. One is the restriction of the splitting for \(V\); the other is the product of the smaller lifts, evaluated at \((g',h)\) and \((1,h)\). Their ratio lies in the central circle and is a character because both lifts are homomorphisms. The rational splittings agree under the direct sum, so the ratio is trivial on rational points. A character on a direct product is the product of its restrictions to the two factors. Call these \(\delta_1,\delta_2\). We have thereby fixed the direction of the ratio by the equality
\[
 \Theta_V(g'\oplus1,h;\phi'\otimes\eta)
  =\delta_1(g')\delta_2(h)\,
           \Theta_{V'}(g',h;\phi')\Theta_D(1,h;\eta).
\]
Indeed the rational points in the chosen position space are \(Y'(k)\times Y_D(k)\).

The quotients in \eqref{eq:seesaw} are compact: a definite unitary group is anisotropic over \(k\), including its norm-one center. Equivalently, finite class number at a finite open compact level and compact archimedean groups give compact representatives. On compact subsets of the adelic group the oscillator translates of a Schwartz function give uniformly bounded Schwartz seminorms at infinity and a common compact open support bound at the finite places. For the latter assertion use smoothness at finite places and a finite covering of the compact set by cosets fixing the test functions. Thus the rational-point series are dominated by the sum of a sufficiently rapid inverse power of the archimedean norm over a fixed lattice. These sums converge uniformly. Both automorphic weights are bounded on their compact quotients. Substitute the kernel identity in the left integral and apply Fubini to obtain \eqref{eq:seesaw}. Finite sums of factorizable test functions follow by linearity.

Finally, on the compact archimedean group \(G'(k\otimes\RR)\), the orthogonal projector onto an irreducible type with character \(\chi_\mu\) is right convolution by \(\dim(\mu)\overline{\chi_\mu}\) for probability Haar measure, in the convention \(P b(g)=\dim(\mu)\int\overline{\chi_\mu(t)}b(gt)\,dt\). Schur orthogonality gives this formula and self-adjointness. Right invariance of quotient measure moves the projector across the pairing on \([G']\). Since restriction to \(G'\) commutes with right translation by \(G'\), the assertion also applies to the restriction of the lifted form. The integral formula moves the projector on the whole restricted form, with the splitting characters kept as in \eqref{eq:seesaw}.
\end{proof}

Theorem~\ref{thm:seesaw} applies to every \(E_x\) on the intervals of Volume~IV. Several orthogonal line summands can be treated by factoring the Schwartz function over all the summands and repeating the same proof. This gives the complex period identity for a fixed \(W\). The integral lattice in Theorem~\ref{thm:theta} is specified by polynomial insertions on definite Hermitian embedding coefficients. To identify it with the lattice of periods of an \emph{iterated} lift with changing \(W\) and prescribed Hecke characters, the map on test vectors at each finite place must carry the specified lattices with controlled indices. The seesaw integral gives equality of complex values and has no index term in its formula.

In the notation of Volume~I, the target is its inequality \eqref{eq:gap}, \(\Delta_u\le x^{5\eps}\mathcal B\) for every translate \(u\). It requires a lower estimate for finite divisibility that represents the Frobenius-weighted count and an upper estimate from archimedean norms and covolumes on the same period lattice. For example, the associated off-diagonal Gram estimates concern distinct Hecke translates; Theorem~\ref{thm:proj} gives orthogonality within one flag, and Proposition~\ref{prop:flag} computes what happens for two distinct flags. To deduce \eqref{eq:gap} one still needs the actual translate pairings and the integral identifications just described. Their uniform comparison on \([V,V^{3/2}]\) remains unresolved here.

\lead{Source of the comparison}Items 7, 8, 12, 15 and 16 on page 4 of Alp\"oge's \emph{Remark} describe the branching, theta lifts, common Bessel period and an estimate between twists. In the supplied source, \nolinkurl{Formalism/SharpRestriction/GlobalTheoremInputs.lean}, lines 65--89, specifies the rationality, common isometry and simultaneous height requirements. The assertions proved in this volume use the mathematical constructions above.~\cite{XIV-XI-IV-remark}

\subsubsection{Finite-order determinant twists through a chain}\label{nav:I-40}
\label{sec:chain}

Let \(T=\operatorname{Res}^{1}_{E/k}\mathbf G_m\), the norm-one torus, and write \([T]=T(k)\backslash T(\mathbb A)\). Take \(G_i=U(V_i)\), \(0\le i\le m\), definite at every archimedean place. Write \(n_i=\dim_E V_i\). Use Hermitian and skew-Hermitian forms in alternating order so that consecutive spaces give unitary dual pairs. Once splittings and Schwartz functions are chosen, let \(\Theta_i(g_{i-1},g_i)\) be the corresponding global theta kernel for \(1\le i\le m\). All groups and the kernels in this subsection are over the same \(E/k\). The Haar measures on their compact quotients are fixed.

Let \(\chi_i\) be finite-order automorphic characters of \(T(\mathbb A)\), \(1\le i<m\), and set \(\xi_i=\chi_i\circ\det\) on \(G_i(\mathbb A)\). For a continuous automorphic function \(f\) on \([G_0]\), form the actual iterated lift with these multipliers inserted at the intermediate groups:
\begin{equation}
 (A_{\boldsymbol\chi}f)(g_m)=
 \int_{\prod_{i=0}^{m-1}[G_i]}
 f(g_0)\prod_{i=1}^m\Theta_i(g_{i-1},g_i)
                 \prod_{i=1}^{m-1}\xi_i(g_i)\ \prod_{i=0}^{m-1}dg_i .
 \label{eq:iter}
\end{equation}
Continuity and absolute convergence follow as in Theorem~\ref{thm:seesaw}, using the finite product of compact quotients. Computing the integral successively gives the individual theta lifts. Changing the splitting at an edge multiplies its kernel by a character of either group on that edge, or by their product, by the central-ratio argument of Theorem~\ref{thm:seesaw}. Characters from neighboring edges multiply at their common group. Formula \eqref{eq:iter} records a finite-order determinant product remaining at an intermediate group. If the products are trivial, only the two endpoint multipliers remain.

\begin{lemma}[Integral endpoint multipliers]\label{lem:multi}
Realize the integral algebraic modular forms on \([G_i]\) at a finite open compact level by their values in prescribed equivariant \(\OO_J\)-lattices at finite-adelic representatives. If \(\chi\) has finite order, is trivial on the determinant of the level, and \(J\) contains its values, multiplication by \(\chi\circ\det\) is an integral automorphism and a Petersson isometry. If all intermediate products from a splitting change are trivial and its two endpoint multipliers are of this kind, their action preserves the finite content of a rational endpoint pairing, as well as its norm at every archimedean embedding.
\end{lemma}
\begin{proof}
The character is trivial on rational points, and its archimedean restriction is trivial since each compact real unitary group is connected. It is constant on each finite level coset and its values are roots of unity, hence units at every finite place of \(J\). Multiplication and division both preserve each value lattice. At every archimedean embedding the scalar has modulus one, and integration in the Petersson norm gives the isometry. For finitely many characters one may take the intersection of their kernels in the original finite level. With the stated cancellation, expansion of the iterated integral leaves just the endpoint multipliers, which can be absorbed into the two test functions. As these are bijections of the integral test lattices, the ideals generated by all paired values agree. The same bijections are unitary for the archimedean operator norms.
\end{proof}

\begin{theorem}[Central characters of the images]\label{thm:center}
Put \(s_i=(-1)^i\). For each edge there is a unitary automorphic character \(\lambda_i\) of \(T(\mathbb A)\), determined by its splitting, such that the following holds. Suppose \(f\) has central character \(\nu_0\), with the convention \(f(zg)=\nu_0(z)f(g)\). The image in \eqref{eq:iter} belongs to the central character \(\nu_{\boldsymbol\chi}\) determined by
\begin{equation}
 \nu_{\boldsymbol\chi}(z^{s_m})
    =\nu_0(z)\prod_{i=1}^{m}\lambda_i(z^{s_{i-1}})
                    \prod_{i=1}^{m-1}\chi_i(z)^{s_i n_i}.
 \label{eq:center}
\end{equation}
In particular two choices with different products \(\prod_{i=1}^{m-1}\chi_i^{s_i n_i}\) have orthogonal images on the entire \(\nu_0\)-space for the Petersson inner product.
\end{theorem}
\begin{proof}
Use the convention that the symplectic action at an edge is the action on the tensor product of the two spaces over \(E\). Then \((zI_{V_{i-1}},z^{-1}I_{V_i})\) acts trivially on that symplectic space. Its lift by the splitting lies in the central circle of the metaplectic group and acts on the Weil representation by a scalar. These scalars define \(\lambda_i(z)\). Multiplicativity of the splitting makes it a character; rational compatibility makes it trivial on \(T(k)\). This proves, for the same Schwartz function,
\[
 \Theta_i(zg_{i-1},z^{-1}g_i)=\lambda_i(z)\Theta_i(g_{i-1},g_i).
\]
Evaluate \eqref{eq:iter} at \(z^{s_m}g_m\) and substitute \(z^{s_i}g_i\) for each previous integration variable. Haar measures are preserved. At consecutive groups the exponents have opposite sign, so each kernel contributes exactly its indicated scalar. The determinant of the scalar action \(z^{s_i}\) on \(V_i\) is \(z^{s_i n_i}\), giving the remaining factors. This is \eqref{eq:center}. For two different output characters choose \(z\) where they differ. Applying central translation by \(z\) to their Petersson pairing multiplies it by a unit complex number different from one while preserving its value, so the pairing is zero.
\end{proof}

\subsubsection{The number of central characters on the common interval}\label{nav:I-41}
For a positive rational integer \(M\), write \(T(\widehat{\OO}_k)\) for the finite-adelic norm-one integral units and \(T(\widehat{\OO}_k;M)\) for the subgroup congruent to one modulo \(M\widehat{\OO}_E\). Define
\[
 C(M)=T(\mathbb A)\big/\big(T(k)T(k\otimes_\QQ\RR)T(\widehat{\OO}_k;M)\big).
\]
Every character of this group is a finite-order determinant multiplier after composition with determinant. Write \(d_E=[E:\QQ]\), \(w_E=|\mu(E)|\), and assume \(E\subseteq\QQ(\zeta_N)\).

\begin{proposition}\label{prop:torus}
The group \(C(M)\) is finite and satisfies \(\log|C(M)|\ll d_E\log(2NM)\), with an absolute constant. For any odd rational prime \(P\nmid N\),
\begin{equation}
           |C(P)|\ \ge\ (P-1)^{[k:\QQ]}/w_E,\qquad
           w_E\le 2d_E^2 .
 \label{eq:ray}
\end{equation}
\end{proposition}
\begin{proof}
At level one, valuations identify \(C(1)\) with the group of fractional ideals \(I\) of \(E\) satisfying \(I\bar I=\OO_E\), modulo the ideals \((a)\) with \(a\bar a=1\). At split places a norm-one idele has opposite valuations in the two factors, and any desired opposite pair is realized. At the other finite places its valuation is zero. The kernel of this description consists of the integral norm-one units and the archimedean factors in the definition.

The image of the map to the ideal class group has order at most \(h_E\). For a principal ideal in its kernel, \(I=(a)\), its norm \(a\bar a\) is a totally positive unit of \(k\). Changing the generator changes this unit by the norm of a unit of \(E\); equality modulo such norms gives the same class of norm-one principal ideals. Thus the kernel injects into \(\OO_k^{\times,+}/N_{E/k}\OO_E^\times\). The norm group contains all squares of units of \(k\), so the unit theorem bounds this index by \(2^{[k:\QQ]}\). Passing from \(C(1)\) to \(C(M)\) adds at most \(|(\OO_E/M\OO_E)^\times|\le M^{d_E}\) by reduction of integral norm-one units. Minkowski's ideal bound and the inequality \(\zeta_E(2)\le\zeta(2)^{d_E}\) give \(\log h_E\ll d_E+\log D_E\), as in Lemma \ref{lem:rank}. The conductor--discriminant formula gives \(D_E\le N^{d_E}\), proving the upper bound.

For the lower bound consider the image of the integral norm-one ideles in \(C(P)\). Its kernel modulo \(T(\widehat{\OO}_k;P)\) consists of reductions of global norm-one units. These units belong to \(\mu(E)\) by Kronecker's theorem, since all their archimedean absolute values are one. As \(P\nmid N\), every local extension at \(P\) is unramified or split. At a split place of \(k\) with residue field of size \(Q\), reduction of the norm-one local units has order \(Q-1\). At an inert place it has order \(Q+1\), the kernel of the residue norm. To see surjectivity in this case, lift any residue root of unity of order dividing \(Q+1\) to its Teichm\"uller root of the same prime-to-\(P\) order. Conjugation raises this root to the \(Q\)-th power, so its norm is one. Independent local choices give at least \(\prod_{v\mid P}(P^{f_v}-1)\ge(P-1)^{[k:\QQ]}\) residue classes, since \(\sum_{v\mid P} f_v=[k:\QQ]\). Divide by the bound \(w_E\) for the global kernel. Finally \(\QQ(\zeta_{w_E})\subseteq E\), and \(\varphi(m)^2\ge m/2\) for every \(m\), as one sees prime power by prime power. The only factor \(\varphi(p^a)^2/p^a\) less than one occurs at \(2^1\), where it is \(1/2\). Thus \(w_E\le2d_E^2\).
\end{proof}

\begin{corollary}[Pairings for explicit twists of the chain]\label{cor:orth}
On the common intervals take \(E=E_x\), \(n_i=i+1\), \(m=\lfloor\sqrt x\rfloor\ge3\). Fix an odd prime \(P\ge5\) prime to \(2F\). For all sufficiently large intervals \(P\nmid N=2^{l+2}Fqd\). For each \(\omega\in\widehat{C(P)}\), put \(\chi_1=\chi_2=\omega\) and take all other interior characters trivial in \eqref{eq:iter}. On any one starting central character the corresponding images in \(L^2([G_m])\) are pairwise orthogonal for distinct \(\omega\). The number of twist choices is at least \(\exp(c_P x^{2+8\eps})\) for all large \(x\), where \(c_P>0\) is independent of \(x,q,d\).
\end{corollary}
\begin{proof}
Definite spaces of these ranks exist over each \(E_x\), using standard forms with a purely imaginary scalar in every skew-Hermitian space. The two auxiliary primes eventually exceed \(P\). In \eqref{eq:center}, the product contributed by the multipliers is \(\omega^{-2+3}=\omega\); hence all output characters are distinct. Apply Theorem~\ref{thm:center} and \eqref{eq:ray}. We have \([k:\QQ]\ge 2^{l-1}\ge x^{2+8\eps}/2\) and \(\log w_E=O_\psi(\log(2x))\). These estimates give the claimed number of choices.
\end{proof}

The conclusion concerns orthogonality, including the possibility that a lift vanishes. Projecting all these images to a \emph{fixed} output central character retains at most one of them. In particular the family in Corollary~\ref{cor:orth} alone cannot prove that many projected maps of positive rank contribute in one designated Bessel component.

\subsubsection{Integral comparison across determinant classes}\label{nav:I-42}
There is also a simultaneous comparison for the finite Fourier transforms that occur when the internal products of characters survive. It specifies the possible denominator in passing between all twists and all determinant classes. Fix a level \(M\) and write \(D_i\) for the image of determinant from \([G_i]\) in \(C(M)\), \(1\le i<m\). This is a finite group: it is the image of \(G_i(\mathbb A)\) as rational determinants map to one. Put \(D=\prod_{i=1}^{m-1} D_i\) and enlarge \(J\) to contain the values of its characters. Pair the output \eqref{eq:iter} with an endpoint test function \(\overline h\). For \(a=(a_i)\in D\), let \(P_a(f,h)\) be the untwisted endpoint integral restricted at intermediate group \(i\) to the fiber \(\det g_i=a_i\). For \(\xi\in\widehat D\) let \(P_\xi^{\rm tw}\) be the full integral with multiplier \(\xi((\det g_i)_i)\).

\begin{proposition}\label{prop:four}
For fixed endpoint test functions,
\[
 P_\xi^{\rm tw}=\sum_{a\in D}\xi(a)P_a,\qquad
 |D| P_a=\sum_{\xi\in\widehat D}\overline{\xi(a)}P_\xi^{\rm tw},\qquad
 \sum_{a\in D}|P_a|^2=|D|^{-1}\sum_{\xi\in\widehat D}|P_\xi^{\rm tw}|^2.
\]
If the fiber values on chosen integral test lattices lie in \(J\) and generate a fractional ideal \(I_{\rm fib}\), their twisted values generate an ideal satisfying
\[
             |D|I_{\rm fib}\subseteq I_{\rm tw}\subseteq I_{\rm fib}.
\]
In particular these ideals agree at every finite place away from \(|D|\) and, for nonzero ideals, their local valuations differ by at most the valuation of \(|D|\). One has \(\log |D|\ll m d_E\log(2NM)\).
\end{proposition}
\begin{proof}
The determinant fibers are open and closed, so partitioning the absolutely convergent integral gives the first formula. For \(a,b\in D\), \(\sum_\xi\overline{\xi(a)}\xi(b)\) equals \(|D|\) if \(a=b\) and vanishes otherwise. Indeed translating the character sum by a character nontrivial on \(ba^{-1}\) multiplies it by a number different from one. Substitution gives the inverse and the stated sum of squares. Character values and their conjugates are algebraic units, so both integral linear combinations imply the ideal inclusions on all test values. The size bound follows from \(D_i\subseteq C(M)\) and Proposition~\ref{prop:torus}.
\end{proof}

For the lattice theta coefficients of Theorem~\ref{thm:theta}, the rationality condition in Proposition~\ref{prop:four} can be realized directly. Use the genera with finite level restricted to the determinant kernel modulo \(M\) at each intermediate group. On each level class the determinant in \(C(M)\) is well-defined because rational points and the level have trivial image. Equivalently, decorate each lattice by a coset of its level inside its original stabilizer. In the finite embedding count of \eqref{eq:tuple}, fix such decorated representatives, count embeddings respecting their induced decorations whenever these are prescribed, and use their actual stabilizers. For definiteness one may allow all isometric embeddings between the underlying lattices while summing over every decoration. Changing an intermediate representative by a decorated isometry permutes the tuples. Each underlying flag with fixed intermediate decoration choices inside the endpoint has \(\prod_{j<n} w(A_j)\) descriptions by decorated representatives, with \(w(A_j)\) now the corresponding stabilizer orders. This follows by transporting each decoration through \(S_j:A_j\to F_j\), just as in the proof of Theorem~\ref{thm:theta}. Restricting the finite sum to any prescribed determinant fibers therefore leaves every value in \((D_nB)^{-1}\OO_J\). For the explicit integral model the two ideals of Proposition~\ref{prop:four} exist, since they lie in that fractional lattice over a Noetherian ring. At each large \(x\), taking \(M=P\) and the family of Corollary~\ref{cor:orth} gives \(\log|D|\ll_{\psi,P} x^{7/2}\log(2N)\).

For the periods prescribed in Section \ref{nav:I-28}, one still needs the lower divisibility that records its prime count and the upper slope estimate on its prescribed common target. Exact orthogonality before projection in \eqref{eq:center} and an identity of complex theta integrals do not identify the integral fibers of that target. In particular, the terms pairing distinct twists with the \emph{same} output central character, the sizes of their diagonal blocks, and the index of the actual finite-place test lattices require estimates. Section~\ref{sec:fixed25} computes the combined Gram expressions in one such character.

\subsubsection{Twists with the same output central character}\label{nav:I-43}
\label{sec:fixed25}

Keep the ranks \(n_i=i+1\), \(m\ge3\), and the finite torus group \(C=C(P)\) of Corollary~\ref{cor:orth}. Every determinant map \(G_i(\mathbb A)\to T(\mathbb A)\) is onto. Indeed, a nondegenerate Hermitian space has an orthogonal line defined over \(E\), by diagonalization over the base field. The same holds for a skew-Hermitian space after scaling its form. Any norm-one idele can act on that line and fix the orthogonal complement. We may therefore take \(D=C^{m-1}\) in Proposition~\ref{prop:four}. Write
\[
  B=\{(z^{s_1 n_1},\ldots,z^{s_{m-1}n_{m-1}}):z\in C\},\qquad H=D/B .
\]
The first two coordinates \(z^{-2},z^3\) have product \(z\), so \(B\simeq C\) and \(h=|H|=|C|^{m-2}\).

Fix a finite-dimensional subspace \(V\) of continuous input forms with one central character and a finite-dimensional Hermitian target \(W\). At the output apply any fixed bounded linear map into \(W\), for example the Petersson projection to a prescribed subspace. Evaluation against a fixed continuous Bessel test on a compact subgroup is also allowed on continuous outputs, since it is bounded in the supremum norm there. Use the same output map for every twist. Fix \(\xi_0\in\widehat D\). The twists having the same output central character as \(\xi_0\) are precisely \(\xi_0\eta\), with \(\eta\in\widehat H\) inflated to \(D\), where the character means the one prescribed by Theorem~\ref{thm:center}, even for a zero lift. Indeed, a ratio preserves that character exactly when it is trivial on \(B\). Denote the resulting operators \(V\to W\) by \(A_\eta\). For \(a\in D\) let \(L_a:V\to W\) be the untwisted operator restricted to determinant fiber \(a\), followed by the same output map. For \(y\in H\) define
\[
        F_y=\sum_{a\bmod B=y}\xi_0(a)L_a,\qquad
        T=\sum_{y\in H} F_yF_y^*,\qquad
        C_g=\sum_{y\in H}F_y^*F_{yg}\quad(g\in H).
\]
These are finite sums of operators given by the actual compact integrals of \eqref{eq:iter}. In the definition of \(F_y\), the notation for the summation means all \(a\in D\) whose image is \(y\).

\begin{theorem}[Second moments after a common projection]\label{thm:gram25}
With adjoints for the chosen Hermitian inner products and \(I_V\) the identity,
\begin{align}
 \frac1h\sum_{\eta\in\widehat H}\|A_\eta^*A_\eta-I_V\|_{\rm HS}^2
    &=\|C_1-I_V\|_{\rm HS}^2+
                     \sum_{g\ne1}\|C_g\|_{\rm HS}^2,\label{eq:gram25}\\
 \sum_{\eta\in\widehat H}A_\eta A_\eta^*&=hT.\label{eq:cov25}
\end{align}
In addition, suppose the entries of the \(F_y\) paired on prescribed integral test lattices generate a fractional ideal \(I_F\) in a coefficient field containing the character values. The corresponding ideal \(I_A\) for all twists in this one central character satisfies \(hI_F\subseteq I_A\subseteq I_F\).
\end{theorem}
\begin{proof}
Splitting the integral by its determinant fibers gives
\[
 A_\eta=\sum_{y\in H}\eta(y)F_y,\qquad
 hF_y=\sum_\eta\overline{\eta(y)}A_\eta.
\]
This identity is preserved by restriction to \(V\) and by the linear map into \(W\). In \(A_\eta^*A_\eta\), group pairs of indices \((y,v)\) according to \(g=y^{-1}v\). This yields \(A_\eta^*A_\eta=\sum_g\eta(g)C_g\). Subtract \(I_V\) from the constant coefficient and apply orthogonality of characters entry by entry in an orthonormal basis of \(V\). Summing the resulting squares gives \eqref{eq:gram25}. In the sum of the products \(A_\eta A_\eta^*\), orthogonality retains only \(y=v\), which proves \eqref{eq:cov25}. Finally both Fourier relations involve algebraic integers, since character values are roots of unity. Apply them to each lattice test pair and then take the ideals generated by all such tests.
\end{proof}

In particular let \(d=\dim W>0\), \(b=\dim V>0\), and index the maps in Proposition \ref{prop:rank11}, by this common family, using one copy of \(V\) for each twist. Its first hypothesis is exactly that the right side of \eqref{eq:gram25} is at most \(eb\). Its second hypothesis is
\[
       d\,\operatorname{tr}(T^2)\le(1+\delta)(\operatorname{tr}T)^2 ,
\]
where \(\delta\) is the nonnegative parameter of that hypothesis. For \(T\ne0\) this asks how evenly the eigenvalues of this particular positive operator fill the common target. Both assertions now concern the same projected maps and need estimates on their actual fiber values.

\section{The first core and the prescribed row}\label{nav:g-VI-1}

\subsection{The cuspidal realization of the prescribed first core}\label{nav:III-40}
Let \([G]=G(F)\backslash G(\mathbb A_F)\), with its quotient measure,
and keep the notation of Proposition~\ref{prop:corecusp251}. Right
translation is \((R(a)F)(s)=F(sa)\).

\begin{theorem}[The first restricted kernel belongs to one copy of the core]
\label{thm:corecopy251}
For the chosen row, lowest-type vector and point in
Theorem~\ref{thm:first250}, the kernel \(P_g\) there belongs to a single
closed isometric copy of the prescribed real discrete series \(D\) in
\(L^2_{\rm cusp}([G])^{K_f}\). It is a smooth vector in that copy.
\end{theorem}
\begin{proof}
The real coefficient used for this kernel is
\(\overline{\langle\pi(s)v',v'\rangle}=\langle D(s)b,b\rangle\)
with \(b=\overline{v'}\) in the conjugate Hilbert space. By the
construction of Theorems~\ref{thm:compact248} and~\ref{thm:direction249},
\(v'\) is obtained by translating a lowest-type tensor and then
integrating its translates over the compact group of the hyperplane,
against a compact matrix coefficient. All intervening translates,
including a compact translate used to select a non-pure prescribed row,
are real. Thus \(b=\int D(a)b_0\,d\nu(a)\) for a compact-finite vector
\(b_0\) and a finite complex measure of compact support. The vector
\(b\) is smooth, since on this compact set every derivative of the
translate of the smooth \(b_0\) is uniformly bounded. The kernel is
exactly \(P_g=Y_b(b)\) with its original finite characteristic function.

On the Harish-Chandra module \(D_K\) define \(Y(u)=Y_b(u)\).
Proposition~\ref{prop:corecusp251} puts its values in cuspidal \(L^2\).
Right translation by the real compact subgroup replaces \(u\) by its
translate. Termwise real differentiation gives \(R_XY(u)=Y(dD(X)u)\).
This is a Hilbert-space derivative as well. Both functions are smooth
on the quotient and square integrable. Integration of the pointwise
derivative along \(\exp(tX)\) gives the same identity in \(L^2\),
since right translation is strongly continuous on that Hilbert space
and the integral of the translated derivative is a Bochner integral.
For example, testing against compactly supported smooth functions
first proves the identity weakly, hence in \(L^2\). Iteration gives
all higher derivatives. These facts make \(Y\) a homomorphism of
\((\mathfrak g,K_\infty)\)-modules.

For a translate of \(b_0\) the series satisfies the elementary
identity \(Y_b(D(a)b_0)=R(a)Y(b_0)\) term by term. Consequently
\[
            P_g=\int R(a)Y(b_0)\,d\nu(a).
\]
This equality holds pointwise by the compact-uniform convergence in
Proposition~\ref{prop:corecusp251}, and in \(L^2\) by strong continuity
and testing on compact sets. Theorem~\ref{thm:first250} gives a nonzero
restricted integral of \(P_g\). It is continuous, so it is nonzero
in \(L^2\), whose measure has full support. In particular \(Y(b_0)\ne0\).

The form \(\langle Y(u),Y(t)\rangle_{[G]}\) on \(D_K\) is invariant,
nonnegative, and nonzero. Its radical is an invariant submodule and
\(D_K\) is irreducible. It is therefore positive definite. By
admissibility any invariant Hermitian form gives an intertwiner to
the compact-finite Hermitian dual. The original unitary form
identifies that dual with \(D_K\), using finite multiplicity on each
compact isotypic space. Schur's lemma now makes the two forms positive
scalar multiples. It follows that \(Y\) extends to a scalar multiple
of an isometry from the whole Hilbert space of \(D\).

This extension intertwines the real group. Indeed compact-finite
vectors of an irreducible unitary Harish-Chandra module are analytic
in its unitary globalization. For each real Lie algebra generator
\(X\), the norms of the successive derivatives of \(Y(u)\) are those
of the corresponding derivatives of \(u\) times the same scalar.
Thus both one-parameter orbit maps are analytic and their Taylor
coefficients correspond at the identity. They agree there on a
neighborhood and hence on the entire real parameter line by real
analyticity and translation. Density extends the equality to all
vectors. The real unitary groups are connected and generated by their
one-parameter subgroups. Hence the closed image is an isometric copy
of \(D\) upon normalization. The compact integral above places \(P_g\)
in that very image and identifies it with the image of the smooth
vector \(b\) under the unnormalized extension. The finite level is
preserved throughout.
\end{proof}

\subsection{The unramified factors on the full even core chain}\label{nav:III-41}
Keep the fields, intervals, \(a_l,g_*,m,n\) and core parameters of
Theorem~\ref{thm:first250}. The ranks of the even core chain are
\(n_i=n+i d\), \(0\le i\le m\), with
\(d=(r+6)-1=n+4\). At the step from \(i\) to \(i+1\), number the
step starting with \(i=0\) and use its unitary splitting character
\(\kappa_i\) on both sides in the upward parameter convention. Its
angular exponent at \(\tau\) is \(-2c_{m-i,\tau}\), where
\(c_{l,\tau}=\lfloor(a_{l-1}+a_l)/2\rfloor\). These are the inverses
of the characters indexed by \(m-i\) in the original grid. They are
trivial on \(\mathbb A_F^\times\), as prescribed for the two even
ranks, and unramified away from the marked finite places. The definitions
of the characters, their order and their real exponents can be found in
\srcloc{Development/AbelianVarietyTheory/Proofs/TwistedOperators/TestGridCommonDataPropertiesHold}{31--59}.
Use these given characters throughout. Let \(S\) contain the exceptional
finite places in \(F\), and write \(S_E\) for all places over \(S\).
If \(w\notin S_E\), write \(Q_w=Nw\) and evaluate an unramified
character at a uniformizer of \(E_w\).

\begin{theorem}[Nonvanishing at every unramified doubling value of the core]
\label{thm:coreeuler252}
Start with the bottom multiset \(B_{0,w}\) of a unitary spherical
Hecke character as in Proposition~\ref{prop:heckecore252}, in the
upward convention, taking its conjugate when the conjugate of the
form enters the kernel. Put
\[
 \begin{gathered}
 A_w=\{Q_w^{(d-1)/2-l}:0\le l<d\},\qquad
 B_{i+1,w}=B_{i,w}\sqcup\kappa_i(w) A_w,\\
 z_*=(d+1)/2,\quad \rho_{hi}=\kappa_h\kappa_i^{-1},\qquad
 C_i=\prod_{h<i}\prod_{a=1}^d L_E^{S_E}(a,\rho_{hi}).
 \end{gathered}
\]
For every \(0\le i<m\), the partial standard Euler product of
\(B_i\otimes\kappa_i^{-1}\), defined initially sufficiently far to the
right, continues holomorphically to a neighborhood of \(z_*\). Its
value there is nonzero, \(C_i>0\), and
\[
 e^{-2fn\log\zeta(3)}\ \le
 \left|\frac{L^{S}(z_*,B_i\otimes\kappa_i^{-1})}{C_i}\right|
 \ \le e^{2fn\log\zeta(3)} .
\]
The constants hold on every original interval, including for the
Hecke character detected by the first period. The transferred good
parameters at any stage still distinguish all distinct systems of
bottom good parameters.
\end{theorem}
\begin{proof}
The recurrence is the unitary unramified theta parameter formula with
equal splitting characters. The old block is twisted by their ratio,
which is one, and the new block is a character times the \(d\)-term
segment. For the split and inert formulas see the unramified transfer
calculation accompanying Theorem~\ref{thm:inert175}. The statement
about these Euler products does not require a choice of ramified test
vectors. Induction gives one copy of \(B_{0,w}\) and exactly the segments
\(\kappa_h(w)A_w\) for \(h<i\). On dividing the multiset by \(\kappa_i(w)\)
and multiplying its Euler factors on a common half-plane of absolute
convergence one therefore obtains
\[
 L^S(z,B_i\otimes\kappa_i^{-1})
 =L^S(z,B_0\otimes\kappa_i^{-1})
   \prod_{h<i}\prod_{l=0}^{d-1}
        L_E^{S_E}\bigl(z-(d-1)/2+l,\rho_{hi}\bigr) .
\]
At \(z_*\) the shifted arguments run once through \(1,\ldots,d\).
For \(h<i\), the arithmetic progression \(a_l=a_{l-1}-g_*\) gives
\(c_{m-i,\tau}-c_{m-h,\tau}=(i-h)g_*\) because adding an integer commutes
with taking a floor. Thus \(\rho_{hi}\) has nonzero angular exponent
\(2(i-h)g_*\) at each selected complex embedding. By
Lemma~\ref{lem:hecke171} its partial Hecke function is holomorphic
near these integers and is nonzero there, also at one. It takes
positive real values at the integers because the ratio character is
unitary and trivial on \(\mathbb A_F^\times\), and the omitted set is
conjugation stable. In particular all \(C_i\) are positive and depend
only on the splitting data and omitted set.

By Lemma~\ref{lem:sphericalcore252}, at \(z_*\) each local logarithmic
series for the initial factor is majorized, per root, by
\(\sum_{a\ge1} Q_w^{-3a}/a\). These series converge locally uniformly
for \(\Re z>(n+1)/2\), by the same estimate with exponent
\(\Re z-(n-1)/2>1\). In particular their exponential is holomorphic
and nonzero on this half-plane. There are at most \(2f\) primes of
\(E\) above each rational prime \(p\), all of norm at least \(p\).
The sum of the absolute logarithmic series at \(z_*\) is consequently
at most \(2fn\log\zeta(3)\), which proves both bounds. This continuation
is for the full partial product with the omitted set held fixed.
Finally each added segment is the same for all initial Hecke characters.
Subtracting its multiplicity at every value of the multiset recovers
the preceding multiset, even when roots coincide. Repetition recovers
\(B_{0,w}\) at each good prime. At a split place this multiset determines
the normalized Satake orbit. At an inert place it is the union of pairs
\(z_j,z_j^{-1}\). It determines the tuple up to permutation and inversion
of individual entries, precisely the relative Weyl orbit. At a
self-inverse root the number of pairs is half its multiplicity. This proves
the asserted separation.
\end{proof}

For the normalized unramified doubling quotient at \(s=d/2\) one must
also divide by \(d_{n_i}^S(s)\). In the normalization following
Theorem~\ref{thm:inert175} it is the product for \(1\le j\le n_i\) of
\(L_F^S(d+j,\omega_{E/F}^{d+j})\) there. All Euler factors at these
integers are positive. Each logarithmic absolute value is bounded by
\(f\log\zeta(d+j)\). The sum over \(j\ge1\) is \(O(f2^{-d})\)
by the absolutely convergent zeta series. Thus the product of the
absolute normalized good-place quotients over all \(m\) steps, divided
by \(\prod_i C_i\), lies between \(e^{-C m f n}\) and \(e^{C m f n}\)
with an absolute \(C\). The remaining local integrals at the exceptional
places and at infinity are needed to infer nonvanishing of a global
lift with the prescribed test.

\subsection{One core form and line test for a whole bottom row}\label{nav:III-42}
Use the bottom hyperplane \(Y\subset X=Y\perp L\) and the next
larger core \(S\). Thus \(\dim Y=r,\ \dim X=u=r+1,\ \dim S=N=u+d\).
Use their original lattices with
\(\Lambda_X=\Lambda_Y\oplus\mathcal O_E z\) on the line \(L=Ez\).
Fix the cuspidal realization \(T\) of one bottom row from
Theorem~\ref{thm:first250}. It has trivial character at its old
finite level \(J_Y\). Denote the core congruence level on \(S\)
by \(J_S\). Use the first seesaw integral \(J(p,q)\) and the
splitting ratios of Theorem~\ref{thm:seesaw245}.

\begin{theorem}[A nonzero first projected seesaw integral at the original depth]\leavevmode
\label{thm:firstsquare255}
For this entire row one can fix one terminal coefficient of
Proposition~\ref{prop:terminalcore255}, all descending tests
from that terminal group down to \(S\), and hence one cusp
form \(A\) on \(S\). One can also fix the finite tests on \(Y,L\)
and the real line test \(q_\infty\). For every nonzero smooth
row vector \(Tw\) some product Hermite test \(p_\infty\) on the
\(Y\)-factor then gives \(J(p,q)\ne0\).
The same period survives averaging the large test over
\(J_S\times J_Y\) in the original splitting. If \(\phi_f\) is
the large finite test, and \(e\) denotes that projector, then
\[
 \|e\phi_f\|^2=\|\phi_f\|^2/D,\qquad
 D\, e\phi_f\ \text{has integral values},\qquad
 1\le D\le (N\mathfrak m)^{r^2}.
\]
Here \(D\) is the order of the image of \(J_Y\) on the chosen
marked halves of \(\Lambda_Y/\mathfrak m\Lambda_Y\).
\end{theorem}
\begin{proof}
At each marked place use a basis of \(X\) ending in \(z\).
The rectangular indicator of Theorem~\ref{thm:markedcore253}
then has its first \(r\) rows centered on the first \(r\)
coordinate row vectors in \(\mathcal O_v^N\), and its last
row centered on the \(u\)-th coordinate row vector. Every
entry has the same displacement modulus \(\varpi_v^{a_v}\).
This is the product of the two indicated coset indicators.
At an unmarked split place the integral rectangular lattice
also decomposes by rows, with the same fractional ideal on
the two summands. At an inert place use the integral
hyperbolic polarization of the larger core \(S\) in
Proposition~\ref{prop:vacuumcore254}. Its position variable
consists of \(N/2\) copies of \(X\). The position lattice
decomposes in each copy as the direct sum of the fractional
lattices from \(Y\) and \(L\). Thus in the direct-sum
polarization the original finite test is
\(\phi_f=p_f\otimes\overline{q_f}\) with both factors integral.
Conjugation changes the form on the line in the seesaw formula.
All these bases may be chosen before applying
Theorem~\ref{thm:upcore254}, whose proof uses only full row rank
of the specified residues.

For one nonzero lowest-type row vector, choose the terminal
coefficient and a successful descent by
Proposition~\ref{prop:terminalcore255}. Keep all its tests
above \(S\), giving \(A\). Its first restriction period is now
the integral \(I(\phi)\) of Theorem~\ref{thm:seesaw245}, by
substituting the last descending kernel into
\(\mathcal B(S_{0,p_0} A,Tw)\).
At infinity take the positive structures and position
polarization respecting the orthogonal sum \(Y\perp L\).
Product Hermite functions on these real spaces are dense in
Schwartz space. With the finite test already a product,
continuity of the seesaw integral gives one real product
\(p_\infty^0\otimes\overline{q_\infty}\) on which it is
nonzero. These Hermite functions are polynomial Fock vectors
for the compatible compact subgroups. Call the resulting form
on \(X\) by descent \(F^0\). It is a compact-finite cusp form
by the reverse occurrence in Theorem~\ref{thm:downcore254}
and is bounded.

The functional \(w\mapsto\mathcal B(F^0,Tw)\) is continuous
in the row Hilbert norm, since \([U(Y)]\) has finite measure
and the restriction of \(F^0\) is bounded. Its value at the
chosen lowest-type vector is nonzero. Thus the orbit of
any nonzero \(w\) in the irreducible row contains a translate
having nonzero value. Changing variables transfers this real
translation to the \(Y\)-action on the large test. In the
direct-sum splitting only \(p_\infty^0\) changes, up to the
nonzero scalar given by the splitting ratio on \(Y\).
All the other data stay fixed. For this fixed smooth \(Tw\),
the resulting functional on the \(Y\)-Schwartz variable is
continuous. Indeed \(A\) and \(Tw\) decrease rapidly on their
Siegel sets and the theta bound of Theorem~\ref{thm:seesaw245}
uses a Schwartz seminorm. The smooth-vector decay at fixed
level follows from
\href{https://arxiv.org/abs/1106.2149}{Miller and Schmid,
\emph{On the rapid decay of cuspidal automorphic forms}}.
Expanding just the translated \(Y\)-test in product Hermite
functions now gives the asserted nonzero \(J(p,q)\).

It remains to compute the projector on this particular test.
The full \(J_S\) already fixes \(\phi_f\), and it commutes with
\(J_Y\). At every unmarked place \(J_Y\) preserves it by
Proposition~\ref{prop:vacuumcore254}. At a marked place
\(k\in J_Y\) operates on the \(u\) rows by \(k\oplus1\)
or its inverse. The leading \(u\)-by-\(u\) square is invertible
modulo the marking, so the translated support intersects the
initial support exactly when \(k\) is the identity modulo
\(\varpi_v^{a_v}\). In that case its determinant character is
trivial by the conductor bounds in Theorem~\ref{thm:markedcore253}.
In general the modulus factor on the compact group is one
and the remaining scalar is a root of unity. The translates
in distinct residue classes thus have disjoint support and
identical norm. The projector is their average with coefficient
\(1/D\), so its squared norm and its denominator satisfy the
stated formulas. The image size is bounded by the number of
all \(r\)-square matrices over the marked residue rings, namely
\(\prod_v(Nv)^{a_v r^2}\). Both input forms are invariant at
their indicated levels. Changing integration variables therefore
leaves the period unchanged under this averaging, as asserted.
\end{proof}

\subsection{Realizing the bottom projection in the first theta test}\label{nav:III-44}
Keep \(v=\bigotimes_\tau v_\tau,H(v),T_v,P_v,T^\dagger\) of
Theorem~\ref{thm:individual259}. Let \(\phi=\bigotimes_\tau\phi_\tau\)
denote the first real test in the downward kernel, with
\(\|\phi\|_2=\prod_\tau\|\phi_\tau\|_2>0\). We take its real
monomials using the positive structures of
Lemma~\ref{lem:weightedhermite260}. Such structures may be fixed
before choosing the monomials, since all monomials in them span
the polynomial Fock model, as in Theorem~\ref{thm:upcore254}.
Let \(L_\tau\) be their energies. Define
\[
 C(v,\phi)=\prod_\tau
    \sqrt{\zeta(2)}\,A(D,L_\tau)M_\tau(v_\tau).
\]
For a terminal coefficient \(e\), write \(F_e(\Psi)\) for its
descent using \(\Psi\) in place of the first real test. In this
notation \(F_e(\phi)=T^\dagger e\). Use linear Schwartz variables
in the downward theta kernel and probability Haar measure on the
terminal quotient.

\begin{proposition}[Approximation of the projected restrictions]\leavevmode
\label{prop:projecttests260}
Fix a smooth bottom row \(u\), and let \(\mathcal E\) be any
finite list of terminal forms in the prescribed compact type.
There are finite Hermite polynomials \(\Phi_j=\sum_\alpha
c_{j,\alpha}h_\alpha\) in the first real test space such that
\[
 \begin{gathered}
 \limsup_{j\to\infty}\|\Phi_j\|_2\le\|\phi\|_2,\qquad
 \sum_\alpha|c_{j,\alpha}|\le C(v,\phi)\|\phi\|_2,\\
 \lim_{j\to\infty}\mathcal B(F_e(\Phi_j),u)
       =\mathcal B(P_vT^\dagger e,u)\qquad(e\in\mathcal E).
 \end{gathered}
\]
The same first test \(\Phi_j\) applies to every \(e\) in the
list. Each \(h_\alpha\) is an orthonormal product on \(Y\) and
the line, and across the real places. All finite tests are the
original ones.
\end{proposition}
\begin{proof}
Use an inner product linear in its second variable. On
\(G_\tau=U(C_0)_\tau\) put
\[
       q_\tau(g)=d_{\pi_\tau}
                 \overline{\langle v_\tau,\pi_\tau(g)v_\tau\rangle}.
\]
This belongs to \(L^1\) by Lemma~\ref{lem:weightedhermite260}.
Schur orthogonality for the discrete series gives
\(\pi_\tau(q_\tau)w=v_\tau\langle v_\tau,w\rangle\).
The coefficient identity at two translated vectors gives
\(q_\tau*q_\tau=q_\tau\). Indeed, in the convolution integral
the middle pair is integrated by the orthogonality formula
\[
 d_{\pi_\tau}\int
 \langle a,\pi_\tau(g)v_\tau\rangle
 \langle\pi_\tau(g)v_\tau,b\rangle\,dg=\langle a,b\rangle .
\]
Also \(q_\tau^*=q_\tau\). Thus in every unitary realization
the integrated operator is an orthogonal projection of norm at
most one. This applies to the actual real oscillator \(\omega_\tau\).

Let \(K_\tau\) be the maximal compact in the signed coordinates.
Sum its central type projectors for the finitely many types in
\(v_\tau\), and denote the resulting character measure by \(b_\tau\).
In any unitary representation it gives an orthogonal projection.
For instance the individual measures on Haar probability are
\((\dim\gamma)\overline{\operatorname{tr}\gamma(k)}\,dk\) by
compact Schur orthogonality. In the oscillator \(K_\tau\) preserves
the chosen positive structure. Its linear action on the positive
creation space preserves total degree. Its splitting character is
a scalar, so \(\omega_\tau(b_\tau)\) commutes with \(\mathcal N\).

Put \(q_{\tau,j}=\mathbf 1_{\{W(g)\le j\}}q_\tau\) for integers
\(j\ge2\). The sets are compact in the matrix group and exhaust
it. Define the Schwartz test
\[
 \Psi_j=\bigotimes_\tau
       \omega_\tau(b_\tau)\omega_\tau(q_{\tau,j})\phi_\tau .
\]
An integral of translates over a compact set defines a Schwartz
vector here, since the Weil action is continuous on Schwartz
space and every Schwartz seminorm of these translates is bounded
on the integration set. The \(L^1\) convergence of \(q_{\tau,j}\)
to \(q_\tau\) implies convergence in operator norm in any unitary
representation. Both limiting factors are contractions. Hence
\(\limsup\|\Psi_j\|_2\le\|\phi\|_2\). Because the compact
projector contracts the weighted number-operator norm, the proof
of Lemma~\ref{lem:weightedhermite260} gives
\[
 \sum_\alpha|\langle h_\alpha,\Psi_j\rangle|
 \le\prod_\tau\sqrt{\zeta(2)}\,A(D,L_\tau)\|\phi_\tau\|_2
       \int_{G_\tau}|q_\tau(g)|W(g)^{2s}\,dg
 = C(v,\phi)\|\phi\|_2 .
\]
Here the Hermite coefficients of a pure tensor multiply.

Right translation of \(F_e(\phi)\) on the bottom core equals
the descent with the corresponding bottom Weil translate of
\(\phi\), by the definition of the theta kernel. This equality
can be integrated against the compactly supported measures just
used. The remaining descent is cuspidal and the theta kernel
has polynomial growth uniformly on compact sets of Schwartz
tests, as in Theorem~\ref{thm:seesaw245}. Thus \(F_e(\Psi_j)\)
is the product of the same compact and truncated convolution
operators in the automorphic representation applied to
\(T^\dagger e\). The reverse real occurrence in
Theorem~\ref{thm:downcore254} puts this vector in the bottom
core isotypic space. The convolution tends in Hilbert norm to
\(P_vT^\dagger e\), since each \(b_\tau\) fixes \(v_\tau\).
All outputs belong to the same finite-dimensional subspace.
Indeed the indicated real representation has finite cuspidal
multiplicity at the fixed level, and its retained compact types
have finite multiplicity in that representation. The
restriction against \(u\) is defined on this subspace of smooth
cusp forms and is continuous there by finite dimensionality.

Finally truncate the Hermite expansion of each \(\Psi_j\).
For a Schwartz vector, applying powers of \(1+\mathcal N\)
gives coefficient decay faster than any power of the degree.
The Schwartz seminorms of a normalized Hermite function grow
at most polynomially in that degree, by the raising and lowering
relations and Euclidean Sobolev embedding. Since the number
of indices of a degree is polynomially bounded at a fixed field,
these expansions converge in Schwartz space. For each \(e\)
the restriction integral as a function of the first test is
continuous in that topology. In the seesaw expression both
the form on \(C_1\) and the row \(u\) have rapid decay, and a
Schwartz seminorm bounds the theta kernel on their Siegel sets
by Theorem~\ref{thm:seesaw245}. Choose the truncation so its
error in all the indicated periods is at most \(1/j\). Orthogonal
coordinate truncation increases neither of the two coefficient
norms. This gives the assertions for \(\Phi_j\).
\end{proof}

\subsection{The exact lowest-type projection on the first core test}\label{nav:III-49}
Take the original even core descent of
Theorem~\ref{thm:downcore254}. Write \(T^\dagger e=F_e(\phi)\)
for its bottom output, where \(\phi=\bigotimes_\tau\phi_\tau\)
is its first real test. The output is in the bottom
discrete-series isotypic space by reverse first occurrence.
At a real place write \(\pi_\tau\) for that representation,
\(K_\tau\) for its maximal compact and \(\kappa_\tau\) for
its lowest compact type. Choose a pure unit vector
\(v_0=\bigotimes_\tau v_\tau\) in these types, as in
Proposition~\ref{prop:poincarenorm261}. On the bottom isotypic
space \(P_{v_0}\) projects to that line in every copy.

\begin{lemma}[A single first test realizing the projected descent]\leavevmode
\label{lem:project271}
There is a first real test \(\phi^0=\bigotimes_\tau\phi^0_\tau\)
such that, for every terminal coefficient \(e\),
\[
 F_e(\phi^0)=P_{v_0}T^\dagger e,\qquad
 F_e(\omega(a)\phi^0)=R(a)P_{v_0}T^\dagger e
 \quad (a\in G_\infty).
\]
Each \(\phi^0_\tau\) is a polynomial Fock vector of energy
at most that of \(\phi_\tau\) and has norm at most
\(\|\phi_\tau\|_2\).
\end{lemma}
\begin{proof}
In the local compact group, with Haar probability and inner
product linear in the second variable, use the measure
\[
 dQ_\tau(k)=
 (\dim\kappa_\tau)\overline{\langle v_\tau,\kappa_\tau(k)v_\tau\rangle}\,dk .
\]
On an irreducible compact representation inequivalent to
\(\kappa_\tau\), its integrated action is zero. On a copy of
\(\kappa_\tau\) it is the rank-one projection with image the
line of \(v_\tau\). Both facts follow by integrating pairs of
matrix entries in an orthonormal basis with first vector
\(v_\tau\). The lowest compact type of a discrete series
occurs with multiplicity one. This standard lowest-type theorem
applies to \(U(p,q)\) by the semisimple theorem on \(SU(p,q)\)
and its compact central character, or directly by the Blattner
formula. In the compact case the whole representation is its
unique compact type.
\href{https://doi.org/10.1007/BF01404112}{Hecht and Schmid,
\emph{A proof of Blattner's conjecture}}.
Thus on every copy of the prescribed discrete series the
measure acts by precisely the desired line projection.

Define \(\phi^0_\tau=\omega_\tau(Q_\tau)\phi_\tau\).
The oscillator restriction to \(K_\tau\) is unitary, so
the same compact orthogonality makes this operator an
orthogonal projection there as well. The compact action
preserves the number operator in the compatible signed
coordinates of Lemma~\ref{lem:weightedhermite260}, giving
the asserted energy and norm. Equivariance of the first
descending kernel moves the right compact translate to
its Schwartz test. The compact integrals may pass through
the descent by the absolute convergence in
Proposition~\ref{prop:projecttests260}. Their product acts
on \(T^\dagger e\) by \(P_{v_0}\) because every real
constituent of this descent has the prescribed type.
Finally the same equivariance for one group element \(a\)
gives the second equality. All these operations are at
infinity.
\end{proof}

\subsection{From a whole discrete summand to the specified vector}\label{nav:IV-31}
At one real place let \(H\subset G\) be the stabilizer of the unit
coordinate \(z\), in ranks \(n-1,n\). Use \(W\) for the defining
operator norms on both groups. Let \(j:\sigma\to\pi|_H\) be an
isometric intertwiner and \(P=j^*\). For the two discrete series take
the parameter bounds of Theorem~\ref{thm:tensor249} with the same
\(n,B,T_0\). For the present local argument \(\mathcal T\) denotes
the quantity in Theorem~\ref{thm:compact248} at this single place.

\begin{theorem}[A quantitative transfer to any lowest-type vector]
\label{thm:direction249}
Let \(v_0\) be a lowest-type unit vector of \(\pi\). Suppose \(A\ge1\)
and there is a \(g_0\in G\) with
\[
      W(g_0)\le T_0^{A n},\qquad
          \|P\pi(g_0)v_0\|\ge T_0^{-A n^3} .
\]
For every unit vector \(w\) in the lowest compact type of \(\sigma\),
there is an \(h\in H\) such that \(v=\pi(hg_0)v_0\) satisfies
\[
 \begin{split}
 W(hg_0)&\le T_0^{C_A n^3},\qquad
          |\langle v,jw\rangle|\ge T_0^{-C_A n^5},\\
 d_\sigma\mathcal T(v)&\le T_0^{C_A n^5} .
 \end{split}
\]
Here \(C_A\) depends only on \(A\). In the product of the real groups,
if these hypotheses hold at every place, any prescribed line of the
product lowest type satisfies the initial-period inequality of
Theorem~\ref{thm:ratio248} on every sufficiently distant original
interval, after choosing the test vector there. In this assertion
assume multiplicity one, the row hypotheses of that theorem and
\(\log T_0=O(\log b)\) uniformly.
\end{theorem}
\begin{proof}
Put \(a=P\pi(g_0)v_0\), \(\eta=\|a\|\), and \(R=W(g_0)\). The
vector \(a\) is smooth under \(H\), since the bounded map \(P\)
intertwines the group and the original vector is smooth. For a word
of length \(j\) in the matrix units of \(H\), move its derivatives
through \(P\pi(g_0)\) by equivariance. One obtains a word on \(v_0\)
with each of its matrix arguments conjugated by \(g_0^{-1}\), hence
of Hilbert--Schmidt norm at most \(R^2\). The word bound in the proof
of Theorem~\ref{thm:tensor249} gives \(J_{n-1}(a)\le(T_0R)^{2n}\).
For \(w\) it gives \(J_{n-1}(w)\le T_0^{2n}\). Proposition~\ref{prop:moment249}
therefore bounds the second moment of \(\langle\sigma(h)a,w\rangle\)
by \((T_0R)^{C n}/d_\sigma\). Its unweighted squared integral is
\(\eta^2/d_\sigma\).

Take \(R_H=\sqrt{2(T_0R)^{C n}}/\eta\) for the constant in that
moment bound. At least \(\eta^2/(2d_\sigma)\) of the squared integral
lies in \(W(h)\le R_H\), since \(W(h)\le\|h\|_{\rm HS}\). The volume
estimate in Theorem~\ref{thm:tensor249} bounds \(d_\sigma\) times the
volume of this compact ball by \(T_0^{C n^2}R_H^{4(n-1)^2}\).
The continuous coefficient attains its maximum on the ball. At a
maximizer its absolute value is at least
\(\eta/(\sqrt{2}T_0^{C n^2}R_H^{2(n-1)^2})\). The starting bounds give
\(R_H\le T_0^{C_A n^3}\), so \(W(hg_0)\le R_H R\) and the claimed
lower bound holds for \(\langle\sigma(h)a,w\rangle=\langle v,jw\rangle\).

Write \(g_1=hg_0\). The first inequality of Theorem~\ref{thm:tensor249}
on \(G\) yields
\[
 |F_v(t)|\le T_0^{C n^3} W(g_1)^{2L}W(t)^{-L},\qquad L=s+2D .
\]
Use the representative \(b_1\) of \(tH\) constructed after
Theorem~\ref{thm:ratio248}. For \(y\in H\), \(W(b_1y)\ge|tz|\) and
\(W(b_1 y)\ge W(y)/(2|tz|)\). Thus multiplication by \(|tz|^s\) in
the definition of \(\mathcal T\) leaves the integrable majorant
\(2^D W(y)^{-D}\) for this last power. Theorem~\ref{thm:tensor249} on
\(H\) bounds its integral times \(d_\sigma\) by \(2^D T_0^{C n^2}\).
Since \(L=O(n^2)\), this proves the asserted bound for \(\mathcal T(v)\).

In the product take the tensor of the local vectors found for any
fixed pure unit vector in the product lowest type \(S\). Its overlap
with that vector has absolute value at least \(T_0^{-C_A f n^5}\),
while \(d_\sigma\mathcal T(v)\le T_0^{C_A f n^5}\) by product
integration. To select any prescribed unit vector \(w\in S\), let
\(e\) be the unit vector in the direction of the orthogonal projection
of \(Pv\) to \(S\). Its overlap is at least the pure-vector bound.
Over the compact group use the matrix unit
\(d_S\overline{\langle\sigma(k)e,w\rangle}\) in place of the diagonal
one in Theorem~\ref{thm:compact248}. Its action on \(S\) is
\(y\mapsto\langle y,e\rangle w\), with zero action on inequivalent
types. Its squared \(L^1\) norm is at most \(d_S\) by orthogonality.
The same calculation gives the identity action on \(w\) with scalar
\(\|\operatorname{pr}_S Pv\|^2/d_\sigma\), and the same tail cost.
Since \(\log d_S=O(f n^2\log T_0)\), the ratio for this real test
has logarithm at most \(C_A f n^5\log T_0\). The depth bound
\(\log N\mathfrak m\ge C_{\rm dep} f^2n^6\log b\) and
\(\log C_n<6n+9\) make it smaller than
\((n+2)\log N\mathfrak m-f\log C_n-\log4\) for large \(x\).
The proof of Theorem~\ref{thm:ratio248} now gives the positive
lower bound at the original level, with this scalar. For the precise ratio of that theorem one can also use
a compact translate of \(v\). Indeed the compact average of
\(|\langle\pi(k)v,jw\rangle|^2\) is \(\|\operatorname{pr}_S Pv\|^2/d_S\),
and \(\mathcal T(\pi(k)v)=\mathcal T(v)\). The extra factor \(d_S\)
obeys the same logarithmic bound. All changes of vectors take place
at the real places.
\end{proof}

\subsection{The omitted coordinate of a tensor projector}\label{nav:IV-32}
Let \(G=U(p,q)\), \(n=p+q\), and let \(H\) fix a signed coordinate
vector \(z\) of Euclidean length one. Order the basis with \(z\) last.
Let \(\pi,\sigma\) be discrete series of \(G,H\) with Harish-Chandra
parameters \(\lambda,\mu\), expressed for their respective diagonal
Cartans with the usual half-sum shift. Suppose \(j:\sigma\to\pi|_H\)
is an isometric embedding and put \(P=j^*\). Use the regularity and
size assumptions of Theorem~\ref{thm:tensor249}. On the smooth tensor
of \(\pi\) with the dual defining representation, write \(P_i\) for
the projector to parameter \(\lambda-e_i\), and \(I_z a=a\otimes e^z\).
If \(y\) is a tensor, \(y_z\) denotes its last component.

\begin{theorem}[A scalar for the projection to the stabilizer]
\label{thm:capelli250}
For every smooth vector \(a\) of \(G\),
\[
 P\bigl((P_i I_z a)_z\bigr)=\gamma_i Pa,\qquad
 \gamma_i=\frac{\prod_{k=1}^{n-1}(\mu_k-\lambda_i+1/2)}
                   {\prod_{j\ne i}(\lambda_j-\lambda_i)} .
\]
In particular \(|\gamma_i|\le(2B+1)^{n-1}\) if every entry of
both parameters has absolute value at most \(B\).
\end{theorem}
\begin{proof}
Work first in the enveloping algebra of \(\mathfrak{gl}_m\). For its
matrix units put \(L_{ip}(u)=u\delta_{ip}+E_{ip}\). Antisymmetrize
in the rows, taking products in order from left to right, and put
\[
 C_m(u)=\sum_{\tau\in S_m}\operatorname{sgn}(\tau)
       L_{\tau(1),1}(u+m-1)\cdots L_{\tau(m),m}(u).
\]
The row antisymmetrization with an arbitrary ordered list of columns
is also alternating in the column indices, keeping the shifts at their
positions. For adjacent columns \(p,q\) the two shifts are \(t,t-1\).
Pairing the two orders of rows \(i,l\), the sum obtained on adding the
original and exchanged column expressions consists of the two commutators
\([E_{ip},E_{lq}]+[E_{iq},E_{lp}]\) and the linear terms
\(\delta_{ip}E_{lq}+\delta_{iq}E_{lp}-\delta_{lp}E_{iq}-\delta_{lq}E_{ip}\).
The quadratic scalar terms cancel. The bracket rule
\([E_{ip},E_{lq}]=\delta_{pl}E_{iq}-\delta_{iq}E_{lp}\) cancels the
remaining terms. Factors on either side of the two positions are
identical in each paired expression. This proves the assertion and in
particular gives zero for any repeated column.

In the commutator with \(E_{ab}\), the Leibniz rule acts on the row
indices by replacing row \(b\) by row \(a\) and on the columns by
subtracting the expression with column \(a\) replaced by column \(b\).
This follows from the same bracket rule, which holds for the shifted
entries in this expression because its two additional scalars cancel.
For \(a\ne b\) both determinants vanish by alternation. For \(a=b\)
they cancel. Thus \(C_m(u)\) is central. For an upper triangular choice
of positive roots apply it to a highest vector of weight \(\Lambda\).
In a nonidentity permutation take its last moved position \(l\). Then
\(\tau(l)<l\), and all factors to its right are diagonal. The raising
operator at \(l\) kills the highest vector. The identity permutation
has eigenvalue \(\prod_l(u+m-l+\Lambda_l)\). In terms of the shifted
parameter \(\nu_l=\Lambda_l+(m+1-2l)/2\) this is
\(\prod_l(u+\nu_l+(m-1)/2)\). By the Harish-Chandra isomorphism the
same symmetric polynomial gives the action at every infinitesimal
character. In particular the formula applies to the discrete series
here, including those without a highest vector.

In rank \(n\) take the first \(n-1\) columns with the shifts in
\(C_n(u)\) and use column \(k\) for the last one. Expansion in that
last position gives polynomials \(c_l(u)\), independent of \(k\), with
\[
 \sum_l c_l(u)L_{lk}(u)=\delta_{kn}C_n(u),\qquad
 c_n(u)= C_{n-1}(u+1).
\]
The second identity uses the upper-left subalgebra. The first holds
for \(k<n\) by alternation. Right multiplication of this row identity
by the inverse of \(L(u)=uI+E\) in formal Laurent series gives
\(C_n(u)(L(u)^{-1})_{zz}=C_{n-1}(u+1)\). The order of multiplication
here follows from the expansion in the last position.

On the smooth tensor in question \(L(u)=uI-A\), where the split
Casimir has entries \(A_{ip}=-d\pi(E_{ip})\). Its resolution in
Theorem~\ref{thm:tensor249} is \((uI-A)^{-1}=\sum_i P_i/(u-\xi_i)\),
with \(\xi_i=-\lambda_i-(n-1)/2\). Apply the last-coordinate
identity to \(a\), then apply \(P\). This map intertwines the
\(H\)-derivatives on smooth vectors, so the rank \(n-1\) central
polynomial acts with parameter \(\mu\). The result is
\[
 \sum_i\frac{P((P_i I_z a)_z)}{u-\xi_i}
 =\frac{\prod_{k=1}^{n-1}(u+\mu_k+n/2)}
         {\prod_{j=1}^n(u+\lambda_j+(n-1)/2)}\,Pa .
\]
Multiplication by the denominator gives a polynomial identity. Taking
the residue at each distinct pole proves the formula. Each denominator
factor in \(\gamma_i\) has absolute value at least one and each numerator
factor has absolute value at most \(2B+1\). This proves the bound.
\end{proof}

\subsection{Cuspidality of the first core kernel}\label{nav:IV-33}
Use the adelic groups \(H\subset G\), the finite level \(K_f=K(\mathfrak m)\)
and the real parameter of Theorem~\ref{thm:first250}. Fix the point \(g\) of
that theorem and write \(c=\iota(g)\). Let \(D\) be the unitary discrete
series of \(G_\infty\) prescribed for the core. Thus \(D\) is the conjugate of
the representation \(\pi\) used to construct the conjugate coefficient
in Theorem~\ref{thm:ratio248}. Write \(K_\infty\) for its standard maximal
compact subgroup. In the next arguments a compact integral of translates
of a vector means \(\int D(a)v\,d\nu(a)\), where \(\nu\) is a complex
measure of finite total variation with compact support in \(G_\infty\).

\begin{proposition}[The periodized core coefficient is a smooth cusp form]
\label{prop:corecusp251}
Let \(u,v\) each be a finite sum of compact integrals of translates of
\(K_\infty\)-finite vectors of \(D\). With inner product linear in the
first variable, put
\[
 Y_v(u)(s)=\sum_{\gamma\in G(F)}
 {\bf1}_{K_f}(c_f^{-1}\gamma_f s_f)
 \langle D(c_\infty^{-1}\gamma_\infty s_\infty)u,v\rangle .
\]
The series and all its right real derivatives converge absolutely and
locally uniformly. The resulting function is bounded and cuspidal in
\(L^2(G(F)\backslash G(\mathbb A_F))\), fixed on the right by \(K_f\).
\end{proposition}
\begin{proof}
Set \(w(s)=\prod_{\tau}\|s_\tau\|_{\rm op}\), so that \(w\) is
submultiplicative, takes values at least one and satisfies
\(w(s^{-1})=w(s)\). Put \(L=2(4n^2+2)+2n+5\) as in
Theorem~\ref{thm:tensor249}. The gap verification of
Theorem~\ref{thm:first250} allows Lemma~\ref{lem:translate163} with
this exponent on every factor. Thus a coefficient between any two
compact-finite vectors is bounded by \(Cw(s)^{-L}\). Indeed a
product compact-finite vector is a finite sum of pure tensors of
factor compact-finite vectors, by admissibility of the discrete
series. The constant in this proof may depend on the fixed vectors
and field. Right differentiation replaces the first vector by an
enveloping-algebra derivative, which is still compact-finite. For
a translated first vector, commute the derivative across the translate
by the adjoint action. On a compact set the resulting derivatives of
a fixed order lie in the span of finitely many derivatives with
uniformly bounded coefficients. Left and right translation of the
group argument over compact sets costs a bounded factor in the weight
estimate by submultiplicativity. These observations justify the same
exponent for both compact integrals in the statement and for all
right derivatives, with constants depending on their order.

Only the finite double coset \(G(F)c_fK_f\) contributes. At the
representative \(s_f=c_f\) the sum is over the arithmetic subgroup
\(\Gamma=G(F)\cap c_fK_fc_f^{-1}\), embedded in \(G_\infty\).
It is discrete because rational points are discrete in the adeles
and the finite subgroup is compact. More explicitly, a compact set
at infinity times this finite subgroup contains only finitely many
rational points. Choose a relatively compact identity neighborhood
\(U\) with the translates \(\gamma U\) disjoint. On such a translate
\(w(t)^{-L}\ge (\sup_U w)^{-L} w(\gamma)^{-L}\). Since
\(\int_{G_\infty} w(t)^{-L}\,dt<\infty\) by Theorem~\ref{thm:tensor249}
and product integration, \(\sum_{\Gamma}w(\gamma)^{-L}<\infty\).
Therefore the absolute sum of every differentiated series is bounded by
\(C w(s_\infty)^L\) on this component, uniformly on real compact sets.
The other representatives in the double coset are obtained by left
rational translation and right translation by \(K_f\). This proves
smoothness and uniform moderate growth. The center of the enveloping
algebra acts through the infinitesimal character of \(D\), by termwise
differentiation.

Consider a proper rational parabolic of \(G/F\) with unipotent radical
\(N\), in this component. The group \(\Gamma\cap N(F)\) is a lattice
in \(N_\infty\) with compact quotient, by the lattice theorem for
unipotent groups. Integrate the series on that quotient after replacing
\(s_\infty\) by \(n s_\infty\). Absolute uniform convergence on a compact
fundamental set permits summing over right cosets of \(\Gamma\cap N(F)\)
in \(\Gamma\). The contribution of each coset is the integral over
\(N_\infty\) of a translated matrix coefficient. For compact-finite
vectors it vanishes by Harish-Chandra's local cusp-form theorem for the
discrete series. Its precise assertion here is that a compact-finite
coefficient belongs to the cuspidal Harish-Chandra Schwartz space,
so its integral along the radical of a proper real parabolic is zero
after any left and right translations. A proper \(F\)-parabolic stays
proper under real base change. For a product, the vanishing follows
factorwise. The finite central covering by the compact real center
and the special unitary group gives the same assertion for our unitary
group since the unipotent radicals lift isomorphically. This local
argument in a Poincar\'e series is also given by
\href{https://www.math.utah.edu/~savin/langlands30.pdf}{Khare, Larsen and
Savin, Section~4.1, pp.~16--17}.

The compact integrals in the statement retain this vanishing. For
absolute interchange, the sum of integrals of the absolute translated
seeds over the compact fundamental set is bounded by the first part
of the proof, uniformly in the two compact translation parameters.
Unfolding this positive majorant dominates each integral on \(N_\infty\).
Finally \(N(\mathbb A_{F,f})=N(F)(N(\mathbb A_{F,f})\cap c_fK_fc_f^{-1})\).
Indeed the exponential coordinates of the unipotent group in
characteristic zero reduce this density assertion to density of \(F\) in \(\mathbb A_{F,f}\). Thus the adelic constant
term on the component is a finite Haar factor times the classical
constant term just calculated. Conjugating the parabolic gives the
same assertion at every representative. On another double coset
one uses its own finite representative and the same density in \(N\)
to express the adelic constant term there by a classical one of the
zero function. This proves cuspidality.

The rapid decrease theorem for smooth cuspidal automorphic forms applies
to this smooth function of moderate growth with a fixed infinitesimal
character. Reduction theory covers an arithmetic component by finitely
many Siegel sets up to rational translates and compact factors. The
norm-one central torus has compact adelic quotient. The function is
therefore bounded, and finite covolume gives square integrability.
See \href{https://arxiv.org/abs/1106.2149}{Miller and Schmid, Theorem~A}.
\end{proof}

\subsection{Cuspidal lifting on the even core chain}\label{nav:IV-34}
Fix an original \(E/F\) and interval as in Theorem~\ref{thm:first250}.
Let \(G_i\) be the unitary group of the \(i\)-th core and \(K_{i,f}\)
its given congruence level. Write \(\mathcal A_i\) for the
\(K_{i,f}\)-fixed cuspidal forms, finite under a real maximal
compact group \(K_{i,\infty}\), of the real type \(P_i\) at all
places. Use the actual lattices of the chain. Fix for the finite
test \(\phi_{i,f}\) the indicators from
Proposition~\ref{prop:vacuumcore254} and
Theorem~\ref{thm:markedcore253}. Their tensor is a finite adelic
Schwartz vector. A rational polarization may be used for the
global theta kernel because the changes from the indicated local
polarizations identify the usual unramified vacuum lines at
almost every place. Normalize the intertwiners there on these
lines. In the complex calculation below any of the finitely
many remaining nonzero scalar normalizations gives the same
nonvanishing assertion.

Write \([G_i]=G_i(F)\backslash G_i(\mathbb A_F)\).
For a polynomial Fock test \(p\) at the real places, denote
the downward theta kernel, written with the small group argument
first, by \(\theta_{i,p}(g,h)\). Use the linear map on cusp forms
\[
   T_{i,p}f(h)=\int_{[G_i]}\overline{\theta_{i,p}(g,h)}\,f(g)\,dg .
\]
The good Hecke parameters below use the quotient convention of
Theorem~\ref{thm:coreeuler252}. Thus at each step the conjugate
of the physical source is the input to the local correspondence.

\begin{theorem}[A nonzero ascent from every bottom core vector]\leavevmode
\label{thm:upcore254}
For every \(0\ne f\in\mathcal A_0\) there are polynomial tests
\(p_0,\ldots,p_{m-1}\) for which \(T_{j-1,p_{j-1}}\cdots T_{0,p_0}f\)
is nonzero in \(\mathcal A_j\) for every \(1\le j\le m\).
The finite tests are the same for all these input vectors.
For an input in one system of good Hecke eigenvalues, every
nonzero output at step \(j\) has good parameters \(B_j\) of
Theorem~\ref{thm:coreeuler252}. One can take a single Fock
monomial at each real place of each step.
\end{theorem}
\begin{proof}
Suppose at step \(i\) that an irreducible cuspidal constituent
has the specified real type and good parameters \(B_i\) arising
from a bottom system. Write \(u=n_i,\ N=u+d,\ s_0=d/2\).
Every smaller target in the rational Witt tower has zero lift.
Indeed it removes at least one hyperbolic plane at a real place,
and a nonzero global lift would give a nonzero local pairing
there after the other variables are fixed. This contradicts
Theorem~\ref{thm:realcore253}. By the constant-term tower property
the lift to this target is cuspidal whenever nonzero. If \(R\)
is the rational Witt index of the target, its smaller real
signature bounds it by \(\lfloor(m-i)/2\rfloor\). Thus
\(N-R>u\), since \(m<d\). This verifies the convergence
condition for the doubled Siegel--Weil integral. The center
is anisotropic over \(F\), so the Petersson products are taken
on the indicated quotients.

Use the mixed Rallis inner product formula for two tests and two
cuspidal vectors in the constituent. The Petersson product of
the lifts is the value at \(s_0\) of the global doubling integral
times a nonzero measure constant. For \(N>2u\) one uses the
convergent Siegel--Weil formula. For \(u<N\le2u\) the constant
Laurent coefficient gives the value as soon as the unfolded
product is holomorphic. This is the formula used in the proof
of Theorem~\ref{thm:global175}.
\href{https://arxiv.org/pdf/1207.4709}{Gan, Qiu and Takeda,
Propositions 11.1(ii) and 11.2 and Section 11.5}.

Unfold first far to the right. Cuspidality removes the non-open
orbits, and the open orbit gives the integral of a global
matrix coefficient times the restricted section. For factorized
data this gives the product of local integrals. At places outside
\(S_0\) the value is the partial standard function divided by
the doubling denominator, by Proposition~\ref{prop:vacuumcore254}.
Theorem~\ref{thm:coreeuler252} proves that this quotient is
holomorphic and nonzero near its required value (shrinking the
neighborhood if necessary). At a marked place the flat section
restricts to a compactly supported coefficient multiplied by
an Iwasawa height to the power \(s-s_0\). On that compact set
the height and its inverse are bounded. The local integral is
therefore entire and at \(s_0\) is \(D_{i,v}\) times the
source inner product. At infinity the discrete series is
tempered, so the flat section integrals are holomorphic for
\(\Re s>-1/2\), by absolute convergence locally uniformly
there. Theorem~\ref{thm:realcore253} gives a nonzero pairing
on the real polynomial tests. These factors give the desired
continuation and its actual value.

We spell out why this argument applies to a prescribed vector.
Let \(D\) be the tensor of the real Harish-Chandra modules of
the source in quotient coordinates. Denote the product of real
pairings by \(Z_\infty(p,p';v,w)\), extending it to finite sums
of tensors. The subspace of \(v\in D\) annihilated for all
\(p,p',w\) is stable under the real maximal compact group and
the complexified Lie algebra. For a compact translate the
invariance of Haar measure moves its action on \(v\) onto the
first oscillator test. Differentiating the same covariance
identity moves a Lie algebra action onto that test. The
identity is one of smooth invariant pairings, and the derived
oscillator action preserves Fock polynomials. The real pairing
is nonzero since tests nonzero in each local pairing can be
tensored. By irreducibility of the exterior tensor product \(D\)
the annihilated subspace is zero.

Choose an orthonormal basis \(e_j\) of the constituent's
finite fixed space. At the good places the spherical lines are
one-dimensional. At the marked places the finite pairing is
scalar on the entire fixed space by
Theorem~\ref{thm:markedcore253}. By multilinearity the mixed
global product on \(f=\sum_j v_j\otimes e_j\) against the lift
of \(w\otimes e_j\) is consequently a nonzero fixed scalar
times \(Z_\infty(p,p';v_j,w)\), with complex conjugates of the
vectors in the quotient convention. If \(v_j\ne0\), some such
product is nonzero. This gives a nonzero lift of \(f\).
For orthogonal invariant cuspidal copies, even copies of the
same representation, the mixed product between their lifts
vanishes. Indeed the open-orbit unfolding pairs a translate
in one copy with a vector in the other, giving zero identically
in the absolute convergence region and hence on continuation.
At fixed level and real infinitesimal character the cuspidal
multiplicity is finite. Resolving a nonzero vector into the
orthogonal constituents and testing against a companion in one
nonzero constituent therefore proves the same assertion for
their sum.

The output is compact-finite at infinity and fixed by the
required finite groups. Local Howe duality and the factorization
of a nonzero cuspidal theta lift give its real type by
Theorem~\ref{thm:realcore253}. At the good primes its parameters
are those of Theorem~\ref{thm:inert175} and
Proposition~\ref{prop:split175}. Equal splitting characters leave
the old multiset in place and add \(\kappa_i(w)A_w\). At a split
place the larger space was written in dual coordinates. Pulling
back by inverse transpose and then conjugating the unitary
representation for the next quotient input returns its transferred
multiset. At an inert place the multiset is already reciprocal.
This proves the required recursion. Each bottom system satisfies
the initial spherical bound of Lemma~\ref{lem:sphericalcore252},
so induction applies through all \(m\) steps. At any step
expanding the successful polynomial into products of monomials
leaves at least one product with nonzero output. Subsequent tests
can be chosen starting with that output.
\end{proof}

\subsection{The first restriction on the original fields and intervals}\label{nav:V-32}
Take \(E=B_J(\sqrt{-q},\sqrt{-d_0})\) and the intervals of
Theorem~\ref{thm:invar148}, with \(F=E^+\), \(f=[F:\QQ]\) and
\(b=\exp(x^\eps)\). Let \(r=k+m-1\), \(n=r+1\), and let
\(Y\subset X\) be the bottom hyperplane and its first core, of
respective dimensions \(r,n\). Write \(H=U(Y)\), \(G=U(X)\), with
inclusion \(\iota\). Keep the level \(K(\mathfrak m)\) on \(X\),
\(J=\iota^{-1}K(\mathfrak m)\) at the finite places, and the old
row level \(K_H(m_0)\), where \(m_0=g_1g_2\). A row here means a
nonzero realization \(T\) of the specified bottom real series
\(\sigma\) in cuspidal \(L^2\) at that old level with trivial level
character. Normalize the realization to be isometric. For a unit
vector \(w\) in its lowest compact type set \(u_w=Tw\). Use the
smooth representative and the Haar measures of
Theorem~\ref{thm:restriction246}.

\begin{theorem}[A uniform lower bound for the first restricted period]
\label{thm:first250}
On every sufficiently distant original interval, at each \(x\), for
each such row and each prescribed \(w\), there is a real matrix
coefficient for the prescribed core, in the convention of
Theorem~\ref{thm:ratio248}, such that at every \(g\) satisfying
\(|u_w(g)|>\|u_w\|_\infty/2\) one has
\[
   |\mathcal P_g|\ \ge
    \frac{\operatorname{vol}(J)}{d_\sigma}\,\|u_w\|_\infty\,
           \exp\{-C f n^5\log b\} .
\]
Here \(\mathcal P_g\) is the restriction integral of the original
finite-level Poincar\'e kernel against \(u_w\), as defined in
Theorem~\ref{thm:restriction246}. The constant is uniform in the
field, point of the interval, row and vector for the fixed family.
\end{theorem}
\begin{proof}
We identify the two real parameters and check the hypotheses of the
local estimates. In the notation of \eqref{eq:letters111} the integers
\(v_i=4\Delta(k-1-i)\) decrease from \(v_1\) to \(v_{k-1}=0\).
Put \(s_*=r+6\), \(N_*=k+m(t+s_*+1)\), where \(t=t_{\rm r}\)
is the integer of Theorem~\ref{thm:rangenew50}. The gap used in the
first core of this grid is
\[
 g_*=\max\bigl\{\lceil N_*^3+n^7+n^3\sqrt{v_1-v_{k-1}}\rceil,
             \ 2N_*^3+2s_*+3,\ \lceil48n^2\sqrt m\rceil\bigr\} .
\]
The spread in this formula is the maximum over the real places. It
has the value displayed since all the \(v_i\) are the same at each
place. This is the gap with scale constant one in
\srcloc{Development/AbelianVarietyTheory/Definitions/Divisors/LeastSplitAuxiliaryPrimeThetaChainSpec}{121--142}.
It is polynomially bounded in \(x\) and at least \(n^7\).

At a real place \(\tau\) write \(b_j=B_\tau+2(j-1)H_0\) and
\(a_l=v_1+1/2+(m+1-l)g_*\). In decreasing order the core parameter
and the dual of the row parameter are respectively
\[
 \begin{aligned}
 \lambda&=(a_0,\ldots,a_m, v_1-1/2,\ldots,v_{k-1}-1/2),\\
 \mu&=(v_1,\ldots,v_{k-1},-b_1,\ldots,-b_m).
 \end{aligned}
\]
These formulas follow by taking the empty added blocks in the grid
parameter lists in
\srcloc{Development/AbelianVarietyTheory/Definitions/Foundations/StrictInterlaceIffDescEnum}{774--814}.
All cloud entries have sign \(+\). For the \(a_l\) the sign is positive
at even \(l\), since \(m\) is even. On the \(-b_j\) it is positive
at odd \(j\). Thus the larger group has one more positive coordinate
and the same number of negative coordinates. On merging the two lists,
the first block consists of \(a_0,\ldots,a_m\), with alternating
signs ending in \(+\). It is followed by \(v_1,v_1-1/2,\ldots,
v_{k-1},v_{k-1}-1/2\), alternating between lists, all with positive
sign. The final block is \(-b_1,\ldots,-b_m\), beginning with \(+\)
and alternating in sign. Adjacent entries from one list have opposite
signs, while those from different lists have the same sign. The corank-one
branching theorem therefore embeds the discrete series with parameter
\(\mu\) into the restriction of the one with parameter \(\lambda\),
with multiplicity one. See \href{https://arxiv.org/abs/1508.02032}{H.~He,
Definition~1 and Theorem~0.1}. Taking the complex conjugate of this
isometric embedding gives the embedding of the row into the dual of the
core for the conjugate coefficient in Theorem~\ref{thm:ratio248}.
Simultaneous duality preserves every absolute gap and absolute size.
The parameter convention can also be located in
\srcloc{Development/AbelianVarietyTheory/Proofs/TwistedOperators/BranchingHypothesesGridParametersLargeScale}{80--98}.

Within the core the successive gaps are drawn from
\(g_*,g_*+1,4\Delta\). Within the row they are drawn from
\(4\Delta,B_\tau,2H_0\). In the spacing used in
Theorem~\ref{thm:ordinary164}, \(\Delta>8(f+1)n^3\),
\(H_0\ge r^{19}f^5x^\eps\), and the least \(B_\tau\) exceeds
every fixed power of \(x\) eventually. Since \(n\) tends to infinity
and is polynomially bounded, all the gaps exceed \(24n^2+2\) uniformly.
The formulas for the weights and the polynomial bounds on \(g_*,k,m\)
give a common bound \(B\le\exp(C x^\eps)\) for their absolute values.
Consequently \(T_0=10^4 n^4(B+1)\) satisfies \(\log T_0=O(\log b)\).

The two finite invariances in Theorem~\ref{thm:ratio248} hold for the
same row. Its old level is the principal congruence on the lattice of
\(Y\) modulo \(m_0\), by the definition in
\srcloc{Development/AbelianVarietyTheory/Definitions/TwistedInnerProducts/OffQOfSizeRows}{25--33}.
Both \(g_i\) split completely in \(E\), as in Section~\ref{sec:whit125}.
For each place of \(F\) over either marked prime the modulus selects one
of the two primes of \(E\) to positive depth. Its conjugate selects the
other. Thus \(\mathfrak m\bar{\mathfrak m}\subset m_0\mathcal O_E\).
The self-dual lattice calculation of Theorem~\ref{thm:short246} identifies
the level on \(X\) with the congruence subgroup for this product ideal.
Intersecting with \(H\), using the orthogonal lattice decomposition
\(\Lambda_X=\Lambda_Y\oplus\mathcal O_E z\) in finite coordinates,
gives \(J\subset K_H(m_0)\). The fixed-level compact-finite cusp form
\(u_w\) is bounded by rapid decrease of cusp forms on Siegel sets. There
are finitely many such sets at the fixed level, and the norm-one center
has compact adelic quotient. See \href{https://arxiv.org/abs/1106.2149}{Miller
and Schmid, Theorem~A}. Its supremum is positive by isometry. No
uniform constant in this qualitative boundedness assertion is used.

At each real place Theorem~\ref{thm:transverse250} now gives the
translate required by Theorem~\ref{thm:direction249} with an absolute
\(A\). Take the product of the local witnesses, including the compact
translate for the specified vector in the product lowest type. By
that theorem the resulting vector \(v\) has
\[
 |\langle v,jw\rangle|\ge e^{-C f n^5\log b},\qquad
 d_\sigma\mathcal T(v)\le e^{C f n^5\log b} .
\]
This use of compact averaging also covers a unit vector in the product
type that is a sum of pure tensors. The Weyl dimension formula gives
\(\log d_S=O(f n^2\log T_0)\). Hence the logarithm of the ratio in
Theorem~\ref{thm:ratio248} is at most \(C f n^5\log b\). At the
original auxiliary depth, \(\log N\mathfrak m>C_{\rm dep} f^2 n^6\log b\),
as calculated in Theorem~\ref{thm:initial247}, and
\(\log C_n<6n+9\). It follows uniformly for large \(x\) that the ratio
is at most \((N\mathfrak m)^{n+2}/(4 C_n^f)\). The compact projection in
Theorem~\ref{thm:compact248} supplies the real coefficient. The last
inequality of Theorem~\ref{thm:ratio248} gives the claimed lower bound
after increasing \(C\) to absorb the square and the factor four.
\end{proof}

This period pairs the bottom row with the first core kernel. The ensuing
successive lifts, the prescribed integral projection and its normalization,
and the unrestricted strict cohomology sum in \eqref{eq:region198} still
enter the comparison leading to \eqref{eq:window}.

\subsection{Every smooth bottom vector has a compact-finite core partner}\label{nav:V-33}
In the notation of Theorem~\ref{thm:first250} use \([H]=H(F)\backslash
H(\mathbb A_F)\). The first restriction is the bilinear integral
\[
        \mathcal B(F_1,Tw)=\int_{[H]} F_1(\iota h)\,(Tw)(h)\,dh .
\]
It converges absolutely when \(F_1\) is bounded on \([G]\) and
\(Tw\in L^2([H])\), since \([H]\) has finite measure. The symbol
\(T\) denotes the isometric realization of the nonzero old cuspidal row.

\begin{theorem}[Nonvanishing of the first period on every smooth row vector]
\label{thm:allrows251}
For each such \(T\) and \(x\) as in Theorem~\ref{thm:first250} is
one cuspidal copy of the prescribed core at \(K(\mathfrak m)\) with the
following property. For every nonzero smooth vector \(w\) of the bottom
row, some \(K_\infty\)-finite vector \(F_1\) in that same copy satisfies
\(\mathcal B(F_1,Tw)\ne0\).
\end{theorem}
\begin{proof}
Choose one unit lowest-type vector \(w_0\), then a point above half the
supremum in Theorem~\ref{thm:first250}. That theorem gives a nonzero
pairing \(\mathcal B(P_g,Tw_0)\), and Theorem~\ref{thm:corecopy251}
places \(P_g\) in a copy \(Y(D)\) of the required core, fixed for the
rest of the proof. By Proposition~\ref{prop:corecusp251} the function
\(P_g\) is continuous and bounded on \([G]\). Its restriction to
\([H]\) belongs to \(L^2\), so \(e\mapsto\mathcal B(P_g,Te)\) is a
continuous linear functional on the row Hilbert space. This functional
is nonzero at \(w_0\). For any nonzero smooth \(w\), irreducibility
implies that the linear span of its \(H_\infty\)-orbit is dense.
Hence \(\mathcal B(P_g,T\sigma(a)w)\ne0\) for some \(a\in H_\infty\).
Right translation invariance of quotient measure and equivariance of
\(T\) give
\[
       \mathcal B\big(R(\iota(a^{-1}))P_g,Tw\big)\ne0 .
\]
Write \(F_2=R(\iota(a^{-1}))P_g\). This smooth bounded function
lies in the same closed core representation and has the same level.

Put \(A(k)=\mathcal B(R(k)F_2,Tw)\) on \(K_\infty\). It is continuous
by dominated convergence with majorant \(\|F_2\|_\infty |Tw(h)|\),
and \(A(1)\ne0\). By Peter--Weyl orthogonality and density, some
matrix coefficient \(p\) of an irreducible compact-group representation
has \(\int_{K_\infty}p(k)A(k)\,dk\ne0\). For example, if every such
integral vanished, uniform approximation of \(\overline{A}\) by their
finite linear span would give \(\int |A|^2=0\), contradicting
continuity at the identity. The compact integral
\(F_1=\int p(k)R(k)F_2\,dk\) belongs to the same closed copy.
Its compact translates span a finite-dimensional space, since the
left translates of \(p\) do. Its level is unchanged, and Fubini's
theorem identifies \(\mathcal B(F_1,Tw)\) with the nonzero integral
just obtained. The preimage of \(F_1\) in \(D\) is compact-finite.
\end{proof}

The core partner here may depend on the bottom vector. For lowest-type
vectors Theorem~\ref{thm:first250} gives a quantitative estimate for
its explicitly constructed kernel. The compact integration in the last
proof gives qualitative nonvanishing for general smooth vectors.
To use this period in the full theta argument one must still place a
partner in the span of the columns reached by the subsequent lifts.
The corresponding two assertions are distinguished in
\srcloc{Development/AbelianVarietyTheory/Definitions/TwistedInnerProducts/OffQOfSizeRows}{119--128}.

\subsection{Selecting the unramified Hecke character of the first core}\label{nav:V-34}
At the level \(K_f=K(\mathfrak m)\) of Theorem~\ref{thm:first250}, let
\(S\) contain the exceptional finite places of \(F\). Away from \(S\) use
the spherical Hecke algebras of the bottom core, with each hyperspecial
group of volume one. Let \(M_D\) be the multiplicity space of the
prescribed real series \(D\) in \(L^2_{\rm cusp}([G])^{K_f}\). Give it
the inner product determined by composition of Hilbert intertwiners and
the unitary form of \(D\).

\begin{proposition}[A first core partner with one system of spherical eigenvalues]
\label{prop:heckecore252}
Let \(h_D\) be the number of distinct simultaneous Hecke characters on
\(M_D\). For the lowest-type vector \(w_0\) and the kernel \(P_g\) of
Theorem~\ref{thm:first250}, one of its projections \(P_{g,\eta}\) to a
single Hecke character satisfies
\[
 |\mathcal B(P_{g,\eta},Tw_0)|\ge h_D^{-1}|\mathcal B(P_g,Tw_0)|>0 .
\]
That character and one of its cuspidal copies of \(D\) can be fixed for
the entire bottom row. Every nonzero smooth vector of the row has a
compact-finite partner with nonzero first period in that copy.
\end{proposition}
\begin{proof}
Harish-Chandra finiteness for cuspidal forms with a fixed infinitesimal
character, finite level and compact type gives finite dimension for
\(M_D\). Indeed evaluation at any fixed nonzero vector of the lowest
compact type is injective on this space of intertwiners by cyclicity.
Its image lies in the specified finite-dimensional space of forms.
The finiteness theorem on the arithmetic components applies since the
adelic center is compact modulo rational points, as in
Theorem~\ref{thm:range239}. The fixed-level \(D\)-isotypic space is
therefore \(D\otimes M_D\). The good Hecke operators act on its second
factor, since they commute with the real group. They preserve the
finite level and the cuspidal subspace. Each spherical algebra is
commutative by Satake's theorem and closed under Hilbert adjoints,
and different primes commute. These operators are simultaneous normal
matrices on \(M_D\).

Their joint eigenspaces are mutually orthogonal. The projection to an
eigenspace is a polynomial in finitely many of these operators. Indeed
for each other eigencharacter choose an operator separating the two,
subtract the value at the other character and divide by the difference
of the two values. The product of these finitely many factors is the
projection. Its action on the original fixed-level function is thus a
finite linear combination of right finite-adelic translates because each
compact double coset is a finite union of sets \(aK_f\) on which
the integrand in the convolution is constant. This realizes the projected functions by
bounded smooth cusp forms, with the same level. Their sum is \(P_g\)
pointwise by continuity and equality in \(L^2\).

By Theorem~\ref{thm:corecopy251}, \(P_g=Y(b)\) for a single
intertwiner \(Y\) and a smooth \(b\in D\). If \(Y_\eta\) is the
projection of \(Y\) in \(M_D\), the corresponding function is
\(P_{g,\eta}=Y_\eta(b)\). The positive pairing of
Theorem~\ref{thm:first250} in absolute value, linearity and the
triangle inequality give the asserted lower bound for at least one
\(\eta\). In particular \(h_D\ge1\) and \(Y_\eta\ne0\). The image
of this intertwiner is a closed copy of \(D\), by Schur's lemma.
Apply the translation and compact averaging argument of
Theorem~\ref{thm:allrows251} with the bounded continuous starting
function \(P_{g,\eta}\). Real translations preserve that same image
and hence its finite-place Hecke character. This gives all the claimed
partners. The numerical loss \(h_D\) is retained explicitly.
\end{proof}

\subsection{Descent to the period of a prescribed bottom row}\label{nav:V-35}
Use the forms \(\mathcal A_i\), tests and kernels of
Theorem~\ref{thm:upcore254}. In the opposite direction set
\[
 S_{i,p}b(g)=\int_{[G_{i+1}]}\theta_{i,p}(g,h)b(h)\,dh .
\]
At \(i=m\) the group \(G_m\) is compact at every real place.
Its rational unitary group is anisotropic and its automorphic
quotient is compact. Denote by \(\mathcal B(F,u)\) the first
restriction period of Theorem~\ref{thm:allrows251}, with the
original old row level and the given level on the core.

\begin{theorem}[Complex span of the descending core periods]\leavevmode
\label{thm:downcore254}
On each of the same fields the space \(\mathcal A_0\) is spanned
over \(\CC\) by the forms
\(S_{0,p_0}\cdots S_{m-1,p_{m-1}}b\), \(b\in\mathcal A_m\).
All finite tests may be fixed as in Theorem~\ref{thm:upcore254}.
For each nonzero smooth vector \(u\) in a bottom row of
Theorem~\ref{thm:allrows251}, there is a single such descent
whose first restriction period against \(u\) is nonzero.
Each real test in it can be a monomial.
\end{theorem}
\begin{proof}
We first compute the real first occurrence in the reverse
direction, since cuspidality of the input alone does not give
cuspidality of its theta lift. Work in the quotient coordinates
of Theorem~\ref{thm:realcore253}, with \(t=m-i,\ A=k+i(d+1)\),
\(y=\lfloor t/2\rfloor,\ z=\lceil t/2\rceil\). Interchange
the factors, moving the scalar \(\sqrt{-1}\) from the smaller form to
the larger at the real place. This gives the same oscillator
with the groups interchanged. Equivalently one may conjugate
the physical downward pairing to use exactly the quotient
coordinates of the upward pairing. The two characters both
have half exponent \(-c_t\).
In the centered parameter of the larger representation there
is now a complete plus string
\((d-1)/2,(d-3)/2,\ldots,-(d-1)/2\).
Every other entry has absolute value greater than \((d+3)/2\).
The string is maximal, so \(k_{\lambda'}=d\) in the reverse
criterion. After removing it, all \(A\) positive entries have
plus sign, and among the negative entries \(z\) have plus sign
and \(y\) have minus sign. Hence \(r_{\lambda'}=A+y\) and
\(s_{\lambda'}=z\). This is exactly the signature of \(G_i\).
In the equal-offset case of the discrete-series criterion the
target \((r_{\lambda'}+l,s_{\lambda'}+l)\) gives a nonzero
lift exactly for \(l\ge0\). Removing a hyperbolic plane would
give \(l=-1\), and removing more makes \(l\) smaller still.
Local Howe reciprocity identifies the lift at \(l=0\) with
\(\check P_i\), since that smaller representation paired with
\(\check P_{i+1}\) in Theorem~\ref{thm:realcore253}. Returning
to the physical downward kernel gives \(P_i\).
\href{https://arxiv.org/pdf/2002.09148}{Ichino, Sections 2.2
and 4.3, Theorem 4.3(ii)}.

It follows that every smaller rational lift in this target tower
vanishes by localization at infinity. The tower property gives
cuspidality at \(G_i\). The given finite test is invariant under
\(K_{i,f}\), and the real test is compact-finite, so
\(S_{i,p}\mathcal A_{i+1}\subseteq\mathcal A_i\).
This holds for arbitrary constituents in the indicated spaces.
For cusp forms the two operators are adjoint for the Petersson
inner products. With the first variable linear one has
\[
 \langle S_{i,p} b,f\rangle_{[G_i]}
       =\langle b,T_{i,p} f\rangle_{[G_{i+1}]} .
\]
Indeed both sides integrate
\(\theta_{i,p}(g,h)b(h)\overline{f(g)}\) on the product of the
two quotients. On Siegel sets the theta kernel for a Schwartz
test has polynomial growth in the two heights, and each
compact-finite cusp form decreases faster than every power
in its own height. This gives absolute integrability and
permits the change of order, as in
Theorem~\ref{thm:seesaw245}.

Let \(R_0\subseteq\mathcal A_0\) be the algebraic span in the
statement. For any nonzero \(f\in\mathcal A_0\), choose the
ascending tests of Theorem~\ref{thm:upcore254} and take
\(b=T_{m-1,p_{m-1}}\cdots T_{0,p_0} f\). Repeated adjunction
gives a member of \(R_0\) paired with \(f\) by the strictly
positive number \(\|b\|^2\). Hence \(R_0\) has zero orthogonal
complement inside \(\mathcal A_0\). To obtain algebraic equality,
fix a \(K_{0,\infty}\)-type \(\tau\).
The group \(K_{0,\infty}\) preserves \(R_0\) because right
translation at the bottom can be absorbed by the first polynomial
test. Its type projection also preserves \(R_0\), since for each
individual compact-finite form the projection takes place in the
finite-dimensional span of its compact translates. The
\(\tau\)-space of \(\mathcal A_0\) is finite-dimensional by finite
cuspidal multiplicity at this level and admissibility at infinity.
If its projection of \(R_0\) were proper, a nonzero vector in
its orthogonal complement inside that type would be orthogonal
to all of \(R_0\), a contradiction. Every form in \(\mathcal A_0\)
is a finite sum of its type components. This proves the asserted
span.

Finally Theorem~\ref{thm:allrows251} provides a compact-finite
\(F\in\mathcal A_0\) with \(\mathcal B(F,u)\ne0\). Express \(F\)
as a finite linear combination of the displayed descents. At
least one term has nonzero period. Expanding its polynomial
tests in monomials and using the same argument gives the last
assertion. The existence statement allows its terminal form
and real monomials to depend on the prescribed row vector.
\end{proof}

The positive norm in the proof is for an ascending complex
theta lift. The period used in the slope comparison of
Section~\ref{sec:missing} includes an integral lattice and
a specified Fourier projection. A bound for that projected
period still requires quantitative control of the terminal
form and of the coefficients of descent in the integral model.

\subsection{Fixed forms in the initial seesaw period}\label{nav:VIII-36}
Use \(X=Y\perp L\) and \(S\) from Proposition~\ref{prop:depth245}
on the given interval. Write \([S]=U(S)(F)\backslash U(S)(\mathbb A_F)\),
and use the same convention for \(Y\). For \(B=X,Y,-L\), write
\(\Theta_B(g,h;f)\) for the theta kernel for \(S,B\), with \(g\)
acting on \(S\). At finite places use the depth spaces of that
proposition. At infinity use Schwartz functions in the compatible
real position polarizations. For automorphic functions \(A,B_0\) on
\([S],[Y]\), respectively, set
\[
 I(f)=\int_{[Y]}\int_{[S]} A(g) B_0(h)
                \Theta_X(g,h\oplus1;f)\,dg\,dh .
\]
For the bottom of a descending column, \(A\) is its first core form
above the bottom and \(B_0\) is the initial row form, with conjugation
on both if the conjugate lifting convention is used. The character
and the field in this definition are fixed.

\begin{theorem}[Transfer of nonvanishing with both forms fixed]
\label{thm:seesaw245}
Suppose \(A,B_0\) are smooth and rapidly decreasing on Siegel sets.
Let \(\delta_S,\delta_Y\) be the ratios of the splitting for \(X\)
to the product splitting for \(Y\) and the conjugate of \(-L\),
restricted to the two groups. For all product tests put
\[
 \begin{split}
 J(p,q)=\int_{[S]} &\delta_S(g)A(g)\,
                \overline{\Theta_{-L}(g,1;q)}\\[-2pt]
        &{}\cdot\left(\int_{[Y]}\delta_Y(h)B_0(h)
                           \Theta_Y(g,h;p)\,dh\right)dg .
 \end{split}
\]
Then \(I(p\otimes\overline q)=J(p,q)\). If \(I(f)\ne0\) for
any Schwartz test of the given finite depth, some such \(J(p,q)\)
is nonzero with the finite components of \(p,q\) individual integral
basis functions of Proposition~\ref{prop:depth245}. Their real
components may be chosen as individual Hermite functions in the
two position spaces.

More precisely, let \(A(gj)=\chi_S(j)A(g)\) and
\(B_0(hk)=\chi_Y(k)B_0(h)\) at their original compact open levels.
On the large test take the orthogonal projector \(e_\zeta\) to the
character \(\zeta(j,k)=(\chi_S(j)\chi_Y(k))^{-1}\), using the
actual splitting on \(X\). Then
\[
                 I(e_\zeta f)=I(f),\qquad
     I(e_\zeta(p\otimes\overline q))=J(p,q).
\]
In particular an integral of the original forms already known to be
nonzero survives that joint projection.
\end{theorem}
\begin{proof}
For a product test the Weil action for the sum is the tensor of the
actions on the summands up to the two splitting characters. In
Schr\"odinger coordinates this follows from the tensor product of the
two Heisenberg representations and the direct-sum action of the
metaplectic group, as in Theorem~\ref{thm:seesaw}. Reversing the
alternating form conjugates the additive character, the Weyl phases
and the theta sum. The ratio of splittings is a character on a direct
product and hence the product of its two restrictions. Both are
trivial on rational points. The position lattice of rational points
is the product of the two rational position lattices. Factoring its
sum therefore gives precisely the integrand defining \(J\), with
both displayed characters. In the Fourier--Jacobi period at this
step the isotropic flag has length zero. Its unipotent group is
trivial and its Haar probability integral is evaluation at the identity.

We give the convergence argument because the groups at the bottom
can be noncompact at infinity. Take finite representatives and Siegel
sets for the two groups by reduction theory. For a fixed finite
Schwartz test the translates by the remaining compact finite groups
take finitely many values. For each such value the rational positions
in its support lie in a finite union of translates of a full real
lattice, by clearing finite denominators in a rational basis.
The absolute theta sum at the identity is consequently bounded by
a Schwartz seminorm on the real test. One may use a finite sum of
\(L^2\)-norms of words in coordinate multiplications and coordinate
derivatives. These control weighted supremum norms by Euclidean
Sobolev embedding. Conjugation by the real oscillator action carries
each such first-order operator to a linear combination of the same
operators, with coefficients given by the symplectic matrix or its
inverse. Unitarity now bounds the theta sum at \((g,h\oplus1)\) by
a fixed seminorm times a polynomial in the real matrix norms of
\(g,h\) and their inverses. The entries of the symplectic action
are polynomial in these matrices and their inverses. On the finitely
many Siegel sets rapid decrease absorbs this polynomial. The quotient
measures are finite, since the centers of the unitary groups are
anisotropic. Equivalently one can integrate the resulting bound in
Siegel coordinates. This proves absolute convergence, justifies
interchanging the integrals, and proves continuity on every fixed
finite test space.

Expand the finite variables in the integral product basis of
Proposition~\ref{prop:depth245}. The finitely many coefficient
functions at infinity are Schwartz, since each is evaluation in a
finite coordinate. Expand each in product Hermite functions. The
expansion converges in the Schwartz topology: applying powers of the
positive harmonic oscillator gives arbitrarily rapid decay of the
Hermite coefficients by its eigenvalue equation and Cauchy--Schwarz.
The creation and annihilation formulas bound every fixed word in
coordinates and derivatives on a normalized Hermite vector by a
polynomial in its degree. Multiplicities grow polynomially for the
fixed real dimension, so the rapidly decreasing coefficients give
absolute convergence in every such seminorm. Completeness of the
Hermite functions in \(L^2\) identifies the limit. Real Hermite
functions may also be used in the conjugate factor. Applying the
continuous functional \(I\) proves the nonvanishing assertion.

Finally, right translating both integration variables shows directly
that \(I(\omega_X(j,k)f)=\zeta(j,k)I(f)\). Integrating with weight
\(\overline{\zeta(j,k)}\) and Haar probability gives the first
projector identity. The product formula gives the second one.
The two automorphic forms in the argument have been held fixed
throughout.
\end{proof}

\subsection{Selecting a prescribed real row vector}\label{nav:VIII-37}
Write \(G_\infty,H_\infty\) for the two real groups in the first
restriction, and suppress the embedding of \(H_\infty\) in
\(G_\infty\). Let \(K\) be the product of the standard maximal
compact subgroups of \(H_\infty\). Its elements preserve the
Euclidean norms and fix \(z\). Give \(K\) probability Haar measure.
Use inner products linear in the first argument. For a vector
\(v\) in a unitary representation \(\pi\) of \(G_\infty\), set
\[
 F_v(t)=\overline{\langle\pi(t)v,v\rangle},\qquad
 \mathcal T(v)=\sup_{t\in G_\infty}
    \left(\prod_\tau|t_\tau z|_\tau^{2n+5}\right)
                         \int_{H_\infty}|F_v(th)|\,dh .
\]
Let \(\sigma\) be an irreducible discrete series of \(H_\infty\),
with formal degree \(d_\sigma\) in this measure. Assume its
compact-finite coefficients are integrable. Lemma~\ref{lem:translate163}
gives this for the spaced row parameters by taking a power larger
than \(4(\dim_EY)^2\) at each place.

\begin{theorem}[A matrix coefficient acting on the prescribed line]
\leavevmode\label{thm:compact248}
Suppose \(\sigma\) occurs with multiplicity one in the Hilbert
restriction of \(\pi\), with isometric embedding \(j\). Let \(w\)
be a unit vector in a single multiplicity-one \(K\)-type \(S\) of
\(\sigma\), of dimension \(d_S\). If \(\mathcal T(v)<\infty\) and
\(\alpha=\langle v,jw\rangle\ne0\), there is a vector \(v'\) in the
same representation \(\pi\) such that
\[
 \sigma(F_{v'}|_{H_\infty})w
      = \frac{|\alpha|^2}{d_\sigma}w,\qquad
                    \mathcal T(v')\le d_S\mathcal T(v).
\]
Here \(\sigma(f)=\int f(h)\sigma(h)\,dh\) for \(f\in L^1\).
One can use every line in the lowest \(K\)-type, including a
prescribed linear combination of weight vectors there.
\end{theorem}
\begin{proof}
Put \(P=j^*\). First consider an arbitrary \(a\) with
\(F_a|_{H_\infty}\in L^1\). We claim that on \(\sigma\)
\[
     \sigma(F_a|_{H_\infty})e
          =d_\sigma^{-1}\langle e,Pa\rangle Pa .
\]
For a unit compact-finite \(e\) of \(\sigma\), put
\(p_e(h)=d_\sigma\overline{\langle\sigma(h)e,e\rangle}\).
It is integrable. Formal degree orthogonality gives
\(p_e=p_e^*=p_e*p_e\) and
\(p_e*\delta_h*p_e=\langle\sigma(h)e,e\rangle p_e\), by substituting
the two matrix coefficients in each convolution. Consequently
\(Q_e=\pi|_{H_\infty}(p_e)\) is an orthogonal projection fixing \(je\).
If \(Q_e x=x\ne0\), its matrix coefficient on \(H_\infty\) is
\(\|x\|^2\langle\sigma(h)e,e\rangle\) by the second convolution
identity. The assignment \(\sigma(h)e\mapsto\pi(h)x/\|x\|\)
therefore preserves the Gram matrix of every finite set of orbit
vectors and extends to an isometric intertwiner. Multiplicity one
places \(x\) on the line \(\mathbb C je\). Hence
\(\langle Q_e a,a\rangle=|\langle a,je\rangle|^2\).
The diagonal coefficient of \(\sigma(F_a|_{H_\infty})\) at \(e\) is
the conjugate of \(d_\sigma^{-1}\langle Q_e a,a\rangle\). It is real
and equals \(d_\sigma^{-1}|\langle e,Pa\rangle|^2\). Homogeneity,
polarization on the compact-finite vectors and density give the
operator identity, since the integral is bounded by the \(L^1\) norm.

On \(K\) form the compact matrix unit
\(p_w^K(k)=d_S\overline{\langle\sigma(k)w,w\rangle}\) and take
\(v'=\pi(p_w^K)v\). By compact Schur orthogonality and multiplicity
one of \(S\), \(Pv'=\sigma(p_w^K)Pv=\alpha w\).
The double integral for \(F_{v'}(th)\) uses the values of
\(F_v(k_1^{-1}thk_2)\) at \(k_1,k_2\in K\), with absolute weights
\(|p_w^K(k_1)p_w^K(k_2)|\). Right translation by \(k_2\) preserves
Haar measure on \(H_\infty\), and
\(|k_{1,\tau}^{-1}t_\tau z|_\tau=|t_\tau z|_\tau\). The triangle
inequality and Tonelli therefore give
\(\mathcal T(v')\le\|p_w^K\|_1^2\mathcal T(v)\).
Schur orthogonality on \(K\) gives \(\|p_w^K\|_2^2=d_S\), so
\(\|p_w^K\|_1^2\le d_S\). In particular the restriction to
\(H_\infty\) is integrable. Applying the operator identity to \(v'\)
and \(w\) gives the assertion. The lowest compact type of a discrete
series has multiplicity one at each factor and also in the product.
\end{proof}

For the corank-one real unitary pair the Hilbert multiplicity
assumption has a direct parameter test. Order the two signed
Harish-Chandra lists together, in coordinates where the omitted
line is positive. For a pair of adjacent entries from one list the
signs must differ, and for adjacent entries from different lists
the signs must agree. For regular parameters with no common entry
this is exactly the interlacing in \href{https://arxiv.org/abs/1508.02032}{He,
Definition 1 and Theorem 0.1}. His branching theorem gives a single
isometric copy in precisely that case. The simultaneous change of
sign of the Hermitian forms gives the same groups for an omitted
negative line, with the parameters expressed in the new signature.
For a product, one takes the tensor product of the embeddings.
For uniqueness take a matrix entry of any
intertwiner against vectors in all the other Hilbert spaces. The resulting
operator is a scalar multiple of the factor embedding. Thus its image lies in that
factor's embedded space. Applying this at every place puts the image
in the tensor product copy, where Schur's lemma applies.

\subsection{First occurrence at the real places of the even core}\label{nav:VIII-38}
Fix a field and point of the interval in Theorem~\ref{thm:first250}.
Use its \(a_\ell,v_j,g_*\) and the notation \(d=n+4,\ n_i=n+i d\)
from Theorem~\ref{thm:coreeuler252}. Suppress a selected complex
embedding \(\tau\). Write \(c_\ell=\lfloor(a_{\ell-1}+a_\ell)/2\rfloor\).
In the signed parameter \(P_i\) of the \(i\)-th core, the entry
\(a_\ell\), \(0\le\ell\le m\), is in the plus block if
\(\ell\ge m-i\) or \(m-i-\ell\) is even. All \(v_j-1/2\),
\(1\le j\le k-1\), are in the plus block. The other plus entries are
\[
  c_\ell+(d+1)/2-h,\qquad
  m-i+1\le\ell\le m,\quad 1\le h\le d .
\]
Write \(\check P_i\) for the parameter obtained by negating the
entries in each block and putting the block back in decreasing order.

\begin{theorem}[The real theta lift on the even core chain]\leavevmode
\label{thm:realcore253}
At every real place and every \(0\le i<m\), use the quotient
convention for the upward kernel, with both half angular exponents
equal to \(-c_{m-i}\). The discrete series with parameter \(\check P_i\)
has first occurrence at the specified target of rank \(n_{i+1}\) in
its Witt tower. Its local theta lift has parameter \(\check P_{i+1}\)
in dual target coordinates. The local doubling pairing at \(s=d/2\)
is nonzero on polynomial Fock data and compact-finite vectors.
\end{theorem}
\begin{proof}
Here is the choice of coordinates. In the downward kernel a map
from the larger space to the smaller belongs to the tensor of the
smaller space with the dual of the larger. The larger Hermitian
form occurs with a minus sign. Complex conjugation of the kernel
changes the sign of the symplectic pairing for the fixed additive
character and inverts the splitting characters. The two minus signs
cancel. At the chosen embedding the purely imaginary scalar
\(-1/(N\mathfrak c\sqrt{-d_0})\) is a positive multiple of \(i\).
Rescaling by a positive real square root gives the skew form
\(iH_i\) on the smaller factor and the Hermitian dual form on the
larger factor. The real additive character is \(u\mapsto e^{-2\pi i u}\).
A form with parameter \(P_i\) enters the quotient convention through
its conjugate, of parameter \(\check P_i\). Pullback from the dual
target changes the signs of the weights within each compact block
because a diagonal matrix acts there by its inverse. These choices
explain the two dual parameters in the statement. The kernel and
character conventions are specified in
\srcloc{Development/AbelianVarietyTheory/Definitions/Divisors/BasePlaceNotMemChiGridBadPlaces}{401--449,535--546}.

Put \(t=m-i,\ c=c_t,\ A=k+i(d+1),\ y=\lfloor t/2\rfloor,\
z=\lceil t/2\rceil\). The centered entries of \(\check P_i\) are
\(c-u\) for entries \(u\) of \(P_i\). Precisely \(A\) are positive.
They consist of the \(i+1\) top entries from \(a_t\) through \(a_m\),
the \(k-1\) cloud entries and the \(id\) entries of the old
strings. All have plus sign. The other centered entries, in order
outward from zero, correspond to \(a_{t-1},\ldots,a_0\), whose
signs start with minus and alternate. Thus \(y\) of them have plus
sign and \(z\) have minus sign. The nearest top entries have
distance at least \(g_*/2-1\) from zero. Also
\(c_\ell=c_t-(\ell-t)g_*\), so every old string stays at distance
at least \(g_*-(d-1)/2\) from zero. The cloud lies below \(a_m\).
The gap in Theorem~\ref{thm:first250} makes every distance
greater than \((d+3)/2\). It also separates distinct blocks, and
entries within each block are distinct half-integers.

In Atobe's criterion the absence of the two central half-integers
and equality of the dimension parities give \(k_\lambda=0\).
Counting positive plus and negative minus entries gives
\(r_\lambda=A+z\), and the other two counts give \(s_\lambda=y\).
The target signature is \((A+z+d,y)\), since its minus block has
\(\lceil(t-1)/2\rceil=y\) entries. In the large-side formula write it
as \((r_\lambda+l+2T,s_\lambda+l)\) with \(l=0,\ T=d/2\).
The required inequalities are \(l\ge k_\lambda\) and
\(\#C^\pm_\lambda(l+T)\le l\). A member of either counting set
would have absolute value at most \(T-1/2\), even before the
deletions in the definition of the sets. Both are empty.
Removing \(h>0\) hyperbolic planes gives \(l=-h\), which fails
the necessary inequality. This proves first occurrence also when
\(y=0\), in which case the tower starts with the compact target.
\href{https://arxiv.org/pdf/2002.09148}{Ichino, Theorem 4.3(ii),
stating Atobe's discrete-series criterion}.

For the parameter of this lift use Ichino's Theorem 4.1(i).
The positive centered entries keep their compact signs, and the
negative centered entries exchange signs. The middle block has
signature \((d,0)\). Its Levi factor \(U(d)\) and the remaining
\(n_i\) one-dimensional Levi factors are compact. To see the
infinitesimal character explicitly, let \(N=n_i+d\). Put the
middle block between the positive and negative entries. For an
old centered entry \(\beta\) in position \(b\), the inducing
weight is \(\beta-((N+1)/2-b)-c\). In each middle position it
is \(A-n_i/2-c\). These are the inducing weights calculated in
the proof of Theorem~\ref{thm:tower169}, with the two half
exponents both equal to \(-c\). Adding the ordered half sum of
the roots restores \(\beta-c\) at an old position and gives
\(-c+(d+1)/2-h\) at middle position \(A+h\). The resulting
list is strictly decreasing, which verifies the good-range
inequalities between the Levi blocks. Induction in stages from
the compact Cartan therefore gives a discrete series. Exchanging
the signs on indices \(\ell<t\) replaces the old threshold \(t\)
by \(t-1\), and the middle adds exactly the reflected string
centered at \(-c_t\). This gives \(\check P_{i+1}\).
Finally the source is discrete series, the target has larger
rank, and \(d/2>-1/2\). The real local doubling integral is
absolutely convergent there, and its pairing on polynomial
Fock tests and compact-finite vectors detects this nonzero
lift. \href{https://arxiv.org/pdf/2002.09148}{Ichino,
Sections 4.2 and 7.2, Lemma 7.1 and Proposition 7.2}.
\end{proof}

\subsection{Cuspidal base change at the bottom of the two towers}\label{nav:VIII-39}
The elementary spherical bound in Lemma~\ref{lem:sphericalcore252}
allows the trivial representation at a good place. For the global
cuspidal spaces of Theorem~\ref{thm:first250}, the widely spaced
infinitesimal character imposes a stronger bound at all those
places simultaneously.

\begin{proposition}[The good parameters of the actual bottom cusp forms]\leavevmode
\label{prop:basechange258}
Let \(\pi\) be an irreducible cuspidal constituent in the indicated
bottom core or bottom row on any sufficiently distant common
interval. Use the prescribed real type in each case, or its dual.
At every finite place \(v\) where \(\pi\) and the unitary group are
unramified, every member \(\alpha\) of its standard base-change
Satake multiset at \(w\mid v\) satisfies
\[
               (Nw)^{-1/2}\le|\alpha|\le(Nw)^{1/2}.
\]
\end{proposition}
\begin{proof}
Write \(N\) for the rank. The field \(F=E^+\) has at least two
real places, since \(E\) contains the prescribed growing real
cyclotomic layer. Both groups are unitary groups of Hermitian
spaces over \(E/F\). The finite places of \(E/F\) are split or
unramified, and the given unimodular lattices give hyperspecial
groups at the unmarked places. At a real place the unsigned
infinitesimal characters in decreasing order are, up to duality,
\[
 (a_0,\ldots,a_m,v_1-1/2,\ldots,v_{k-1}-1/2),\qquad
 (v_1,\ldots,v_{k-1},-b_1,\ldots,-b_m).
\]
Their entries have parity \((N-1)/2\bmod\ZZ\). Every two distinct
entries differ by more than one by the gap estimates of
Theorem~\ref{thm:first250}. Thus these are real integral strongly
regular infinitesimal characters. Indeed after subtracting the
ordered half sum of roots the successive gaps are still positive.
The discrete series with these parameters are cohomological for
the algebraic coefficient with the corresponding infinitesimal
character, using its dual when pairing in cohomology.

We check the regularity hypothesis for Labesse's cohomological
base change. An irreducible unitary representation with this
infinitesimal character is cohomologically induced from a unitary
character on a theta-stable Levi, by Salamanca-Riba's strongly
regular theorem. The Levi of such a parabolic in \(U(p,q)\) is
a product of unitary blocks. Within a block of size \(b\) its
character has constant differential, and the half sum from the
nilradical is also constant on this block. The half sum of the
Levi roots therefore contributes \(b\) consecutive entries of
step one to the infinitesimal character. This follows directly
from the list \(((b+1)/2-j)_{1\le j\le b}\) on its Cartan.
Consequently each block has size one. The Levi is a compact
Cartan and the corresponding cohomological modules are discrete
series in the packet of the given infinitesimal character.
Labesse's condition \((*)\) follows. A discrete-series
pseudo-coefficient has trace zero on the other tempered members
and the prescribed sign on the chosen member. This is the
sufficient condition preceding his Corollary~5.3.
\href{https://doi.org/10.1215/S0012-7094-99-09616-3}{S.~Salamanca-Riba,
the strongly regular unitary dual} and
\href{https://webusers.imj-prg.fr/~michael.harris/CBCM.pdf}
{J.-P.~Labesse, Section 5.1}.

Apply that corollary. It supplies an isobaric sum of discrete
automorphic representations of general linear groups over \(E\).
Its infinitesimal character is the base change of the displayed
one. It agrees with the local base change at the unramified
places and at split places. In particular this applies at each
spherical place under consideration. By the discrete spectrum
classification for general linear groups, a noncuspidal discrete
summand has cuspidal support a string of length \(b\ge2\) on one
unitary cuspidal representation, with norm twists
\((b-1)/2,(b-3)/2,\ldots,-(b-1)/2\). At a complex place
the factor \(|z|_{\CC}^{a}=(z\bar z)^a\) adds \(a\) to the
holomorphic infinitesimal character as well as to its
antiholomorphic counterpart. Even one holomorphic entry of
the cuspidal factor then produces two entries differing by
one in the holomorphic list of the isobaric sum. Stable base
change of the real infinitesimal character has the displayed
list on one factor and its dual on the other. Each excludes
such a pair. Hence every discrete summand is unitary cuspidal.
\href{https://webusers.imj-prg.fr/~michael.harris/CBCM.pdf}
{Labesse, Corollary 5.3} and
\href{https://www.numdam.org/articles/10.24033/asens.1595/}
{C.~M\oe glin and J.-L.~Waldspurger, the residual spectrum of
\(\mathrm{GL}_N\)}.

At a good place the constituents of the isobaric sum have
unramified local parameters. For example, their Weil--Deligne
parameters form the direct sum for the isobaric sum, so
trivial inertia and zero monodromy for that sum impose the
same conditions on each summand. The Satake multiset is their
multiset union. For a unitary cuspidal representation of
\(\mathrm{GL}_b\) over a number field, the bound of Luo, Rudnick
and Sarnak at unramified places gives exponent
\(1/2-1/(b^2+1)\) for \(b\ge2\). A unitary character in rank
one has modulus one. Apply the same upper bound to the dual
to obtain the lower bound. This proves the assertion also
in the conjugate quotient convention.
\href{https://doi.org/10.1090/pspum/066.2/1703764}
{W.~Luo, Z.~Rudnick and P.~Sarnak, bounds for cuspidal
Satake parameters}.
\end{proof}

\subsection{An everywhere unramified finite parity character}\label{nav:VIII-40}
\begin{proposition}[A single parity character for a whole interval]\leavevmode
\label{prop:parity262}
On every sufficiently distant interval of Theorem~\ref{thm:first250} is a unitary Hecke character on the fixed biquadratic subfield
\(M=\QQ(\sqrt{-q},\sqrt{-d_0})\) whose pullback by \(N_{E/M}\)
has the following properties for every field \(E\) on the interval.
It is unramified at every finite place, has angular character
\(z/|z|\) at every embedding in the CM type above
\(K=\QQ(\sqrt{-d_0})\), and restricts to \(\omega_{E/F}\) on
\(\mathbb A_F^\times\). In particular one can use it as the parity
character \(\eta\) in both theta directions. On each of these
fields one can simultaneously choose all horizontal characters
\(\gamma_i\) with half angular exponents \(c_{m-i,\tau}\),
trivial restriction to \(F\), conductor exponent two at the
four level and type primes, and exponent one at the single
marked auxiliary place. They can be taken trivial on residue
roots of unity at both type primes and unramified elsewhere.
\end{proposition}
\begin{proof}
Put \(F_0=M^+=\QQ(\sqrt{q d_0})\). The two imaginary quadratic
fields have coprime fundamental discriminants \(-q,-d_0\).
Either one generates \(M\) over \(F_0\), so \(M/F_0\) is
unramified at every finite place by unramified base change.
Write \(\omega_f\) for the finite part of its quadratic idele
class character, viewed as a character of fractional ideals.
Let \(\Phi\) be the product of \(z/|z|\) over the two complex
embeddings of \(M\) extending the chosen embedding of \(K\).
On \(F_{0,\infty}^\times\) it is the product of the real signs.

Let \(\mathfrak p\) be the unique prime of \(F_0\) above \(d_0\).
Its residue field is \(\mathbb F_{d_0}\). The primes in
Theorem~\ref{thm:long232} are both \(3\bmod4\) and
\(d_0\equiv-1\bmod q\). Quadratic reciprocity gives
\[
 \left(\frac{d_0}{q}\right)=-1,\qquad
 \left(\frac{q}{d_0}\right)=1,\qquad
 \left(\frac{-q}{d_0}\right)=-1 .
\]
Thus \(\mathfrak p\) is inert in \(F_0(\sqrt{-q})=M\).
For \(\delta=\sqrt{-d_0}\) we have
\((\delta)=\mathfrak p\mathcal O_M\), since both ideals have
square \(d_0\mathcal O_M\). Also
\(\Phi(\delta)=-1=\omega_f(\mathfrak p)\).

We first prescribe the finite character on the subgroup of
fractional ideals of \(M\) generated by principal ideals and
ideals extended from \(F_0\). On \((\alpha)\) prescribe
\(\Phi(\alpha)^{-1}\) and on
\(\mathfrak a\mathcal O_M\) prescribe \(\omega_f(\mathfrak a)\).
To check all relations, suppose \((\alpha)=\mathfrak a\mathcal O_M\).
Then \(\alpha/\bar\alpha\) is a unit of absolute value one under
every embedding, hence a root of unity by Kronecker's theorem.
The only roots of unity in \(M\) are \(\pm1\). Indeed its Galois
group has exponent two. A primitive root of prime order \(p\)
would force \(p-1\mid2\), and its three quadratic subfields exclude
both \(\QQ(i)\) and \(\QQ(\sqrt{-3})\) since \(q,d_0>3\).
It follows that \(\alpha=t\) or \(\alpha=t\delta\) with
\(t\in F_0^\times\). In the first case the required identity is
the principal idele relation for \(\omega_{M/F_0}\). In the second
it follows by multiplying that relation by the one for \(\delta\).
Extension of ideals is injective, as is seen from valuations.
This verifies well-definedness, including the relation for units.

The unit circle is divisible, so the resulting ideal character
extends to the group of all fractional ideals. Multiplying this
extension on the finite idele ideal by \(\Phi\) at infinity gives
a continuous unitary character trivial on \(M^\times\). It is
trivial on all finite local units and has the asserted restriction
on \(\mathbb A_{F_0}^\times\) by construction.
Finally \(M\cap F=F_0\) and \(E=FM\), since \(M\) supplies an
imaginary generator of the quadratic extension \(E/F\). Thus
\(N_{E/M}\) on \(\mathbb A_F^\times\) is \(N_{F/F_0}\).
Functoriality of quadratic reciprocity gives
\(\omega_{M/F_0}\circ N_{F/F_0}=\omega_{E/F}\). Local norms
preserve units, and the infinite norm multiplies exactly the
coordinates above each chosen embedding. This proves the
assertion about \(\eta\) with a single choice on \(M\).

For \(\gamma_i\) work on \(A=\ker N_{E/F}\) and then pull
back along \(z\mapsto z/\bar z\). At infinity prescribe powers
\(c_{m-i,\tau}\) on the unit circles. At every place over the
four marked rational primes use the chosen split component to
identify \(A(\mathcal O_v)\) with \(\ZZ_p^\times\). Prescribe
a nontrivial character on \((1+p\ZZ_p)/(1+p^2\ZZ_p)\), extended
trivially across the residue roots of unity. At the single
auxiliary place prescribe a nontrivial residue character.
The auxiliary rational prime is distinct from the two type primes
by its defining congruence, and from \(g_1,g_2\) by its upper
bound \(b\) and their lower bounds \((kb)^8\).
At all remaining places prescribe the trivial character on
units. The rational points of \(A\) integral everywhere finite
are exactly \(\mu(E)\), by Kronecker's theorem. At one place
over \(g_1\), reduction is injective on \(\mu(E)\), since
\(1+g_1\ZZ_{g_1}\) is torsion free. Multiply the unit character
there by a residue character canceling the product of the
prescriptions on \(\mu(E)\). Such a character extends from
the subgroup of the cyclic residue group, and leaves the
nontrivial conductor-two part intact. The type primes are
distinct from \(g_1\). The prescriptions now descend to the
image of \(A(F_\infty)\prod_{v<\infty}A(\mathcal O_v)\) in
\(A(F)\backslash A(\mathbb A_F)\). This image is open. Extend
its character by divisibility of the circle as above. Continuity
on the open subgroup gives continuity of the extension.
Pullback doubles the angular exponents and preserves the local
conductors because at a split place it is the ratio of the two
coordinates. The construction works independently for every
\(i\) with the same marked places and levels.
\end{proof}

\subsection{The terminal exchange on the prescribed row}\label{sec:terminalexchange336}
Fix the field, word and terminal prime \(T=\varpi\bar\varpi\)
of Theorem~\ref{thm:common111}. At \(v\mid T\) use the split
half where \(\varpi\) has valuation one, and put \(k=F_{+,v}\),
\(O=\mathcal O_k\), \(q=|O/\pi O|\). Use the terminal frame
of Theorem~\ref{thm:midlevel131}, with \(z=\varpi e\).
In the simple character of Lemma~\ref{lem:exchangeroot336} the
annihilator of \(O\) is \(O\). Indeed \(T\) splits completely
and the standard finite character of \(\mathbb Q_T\) has this
conductor. In the absolute-trace normalization of
Theorem~\ref{thm:graph274} the multiplier is a unit at \(T\).
Write \(m_i=\mu_i(O^a)\) for exchange Haar volumes. Assume \(s\ge1\).
For \(s=0\) the period already has its symmetric value.

\begin{theorem}[The terminal sum at the original lattice level]\label{thm:terminalexchange336}
Suppose the row \(F\) is fixed by the original standard lattice
stabilizer at \(v\). Let \(L'\) be the symmetric period, possibly
with the prescribed exchanges already performed at other places.
Put \(b(x)=L'(R(X_s(x))F)\). The period after all exchanges at
\(v\) is
\[
 \left(\prod_{i=1}^s m_i\right) q^{-a}
                      \sum_{\bar x\in (O/\pi O)^a} b(x).
\]
If \(\phi\) has the prescribed covector level
\(P_\ell=\{h\in\GL_a(O):\ell h\equiv\ell\bmod\pi\}\)
for a primitive linear form \(\ell\), the residue sum equals
\[
 b(0)+(q^{a-1}-1)b(x_*)+
             q^{a-1}\sum_{t\in (O/\pi O)^\times} b(x_t).
\]
Here \(\bar x_*\) is any nonzero vector in \(\ker\bar\ell\), and
\(\bar\ell(\bar x_t)=t\). The middle term is omitted for \(a=1\).
\end{theorem}
\begin{proof}
The original middle lattice is \(O^a\oplus Oe\), which in the
coordinates \(A,z\) is \(O^a\oplus\pi^{-1}Oz\). Thus
\(Y_i(O^{a*})\) fixes \(F\) for every \(i\). For \(i<s\) the
group \(X_i(O^a)\) also fixes it. At \(i=s\) the group of
integral \(X_s(x)\) has \(x\in\pi O^a\), since it sends \(z\)
to \(z+x\). For \(i<s\), apply
Theorem~\ref{thm:exchangefourier336} to \(L_{>i}\) with
\(U=O^a=V^\perp\). Removing each of these integrations leaves
the factor \(m_i\) and the original \(F\).
The last exchange has \(V^\perp=O^a\), \(U=\pi O^a\) and
\(\mu_s(U)=m_s q^{-a}\). This proves the first formula.

For \(h_0\in P_\ell\), conjugation by \(h_0\) sends
\(X_s(x)\) to \(X_s(h_0x)\), and \(R(h_0)F=F\).
In \(L'\), moving \(h_0\) onto the column amounts to right
translation of \(\phi\) by \(h_0^{-1}\), by invariance of the
quotient measure on \([H]\). This fixes \(\phi\).
Thus \(b\) is constant on the orbits of the stabilizer of
\(\bar\ell\). Every matrix in the residue stabilizer lifts to
\(P_\ell\). In coordinates \(\bar\ell(w,t)=t\) the action is
\((w,t)\mapsto(Bw+tv,t)\) with \(B\in\GL_{a-1}(O/\pi O)\)
and \(v\) arbitrary. This gives precisely the orbits and sizes
in the second formula, including the case \(a=1\).
\end{proof}

Applying the first formula at every terminal place contributes
at most \(T^{a[F_+:\mathbb Q]}\) summands on this lattice level,
since each residue field has \(T\) elements. The other marked
places can keep their original exchanges. All formulas are linear
in the row, hence apply after any prescribed finite linear
combination for which the indicated invariances hold.

\begin{proposition}[The congruence lattice at the transported level]\label{prop:graphexchange336}
At the terminal level \(K^0\) of Theorem~\ref{thm:midlevel131}
put \(J=\{w\in O^a:\ell(w)\in\pi O\}\). For a \(K^0\)-fixed row,
one may take \(V=O^{a*}\) in every step of
Theorem~\ref{thm:exchangefourier336}. On that row one may take
\(U=J\) for \(i<s\) and \(U=\pi J\) for \(i=s\). The respective
numbers of terms in the single-step transforms are \(q\) and
\(q^{a+1}\). The joint sum on the transported row satisfies
Theorem~\ref{thm:flagindex336} with \(d=2,\ c=0\).
\end{proposition}
\begin{proof}
In coordinates \((w,b)\) for \(w+bz\), the middle lattice
is given by \(w\in O^a\) and
\(b-\ell(w)/\varpi\in O\). It meets \(A\) in \(J\).
Its projections to \(A\) and the \(b\)-line are \(O^a\)
and \(\pi^{-1}O\), respectively. The last equality uses
primitivity of \(\ell\). Adding \(y(w)e_i\) consequently
preserves the full lattice for \(y\in O^{a*}\).
For \(i<s\), adding \(x\) to \(d_i=e_{i+1}\) preserves it
exactly when \(x\in J\). For the last root the displacement
is \(b x\). This belongs to \(J\) for every possible \(b\)
exactly when \(x\in\pi J\). All are shears, so checking
the negative displacement also verifies preservation by the
inverse. As \([O^a:J]=q\), the two indices follow. These
single-step indices concern evaluations on \(F\). In a nested
sum the later arguments can be translates of \(F\).

Finally the graph lattice is \(C\Lambda\) for the coordinate
lattice \(\Lambda\) in \(A,z\) and the flag, with
\(C(w,b)=(w,b+\ell(w)/\varpi)\). Both \(C\) and \(C^{-1}\)
have entries in \(\pi^{-1}O\). For \(g\in K(2)\) the entries
of \(C^{-1}(g-1)C\) and the corresponding entries for \(g^{-1}\)
are integral. Hence \(K(2)\subset K^0_v\), which gives the
joint bound. For a further specified translate its depth is
given by the elementary-divisor calculation following
Theorem~\ref{thm:flagindex336}.
\end{proof}

The formulas retain the values of the actual automorphic row
at the listed translates. Bounding their signed sum after
the specified cuspidal projection is still required for the
comparison in Section~\ref{sec:missing}.


\subsection{Transporting the terminal graph to the column}\label{sec:graphtransport337}
Use the terminal prime and the flag of Section~\ref{sec:terminalexchange336}.
Write \(H(k)=\GL(A_k)\), with its given embedding fixing the flag
and the marked line. In the middle coordinates \(w+t e\) let
\(\Lambda_*\) have standard middle lattice \(O^a\oplus Oe\).
For a nonzero residue covector \(\bar\ell\), take an integral lift
\(\ell\). Give \(\Lambda_\ell\) middle lattice
\[
 \{w+t e:w\in O^a,\ t\in O,\ t\equiv\ell(w)\pmod\pi\}.
\]
Both lattices have the same integral flags. Put \(J_\ell=\{w\in
O^a:\ell(w)\in\pi O\}\). Choose \(h_\ell\in H(k)\) taking \(O^a\)
to \(J_\ell\), and write \(a_\ell\) for a linear transformation
taking \(\Lambda_*\) to \(\Lambda_\ell\). One such choice on the
middle block is \((w,t)\mapsto(w,\ell(w)+\pi t)\).

A unitary character \(\kappa\) of the compact determinant extends
to \(k^\times\) by assigning the value one to \(\pi\). Work with
\(\widetilde R(g)=\kappa(\det g)^{-1}R(g)\) on the row and its
restriction to \(H(k)\) on the column. Thus rows of the given
compact type are fixed in this action. Put
\(\mathcal J_\phi(F)=\int L_\phi(R(\prod_i X_i(x_i))F)\prod_i
d\mu_i(x_i)\). Exchanges at other marked places may already be
included in \(L_\phi\). Theorem~\ref{thm:exchangefourier336} gives
absolute convergence, and \(\mathcal J\) is conjugate linear in \(\phi\).
\begin{theorem}[Exact transport at the first Bessel position]\label{thm:graphtransport337}
Let \(F_*\) be fixed under \(\widetilde R\) by the stabilizer of
\(\Lambda_*\), and let \(s\ge1\). Then
\[
 \mathcal J_\phi(\widetilde R(a_\ell)F_*) =
        q^{-s}\mathcal J_{\widetilde R(h_\ell^{-1})\phi}(F_*).
\]
Every row of the prescribed type at \(\Lambda_\ell\) has the
displayed form. The right side is unchanged on replacing \(\phi\)
by \(\widetilde R(v)\phi\) for any \(v\in H(k)\) preserving
\(J_\ell\). This includes all scalars in \(O^\times\).
\end{theorem}
\begin{proof}
If \(u\in O^a\), then \(X_s(\varpi u)\), with \(z=\varpi e\),
takes \((w,t)\) to \((w',t)=(w+tu,t)\). Its image of the graph
satisfies \((1+\ell(u))t\equiv\ell(w')\pmod\pi\).
The converse follows by setting \(w=w'-tu\). Primitivity gives
\(u\) with \(1+\ell(u)\equiv0\). For that \(u\) the image lattice
is \(J_\ell\oplus Oe\) with the original flags, namely
\(h_\ell\Lambda_*\). Thus
\(h_\ell^{-1}X_s(\varpi u)a_\ell\) stabilizes \(\Lambda_*\).
Translation in the final exchange variable leaves the Haar integral
unchanged, since the \(X_i\) commute. The left side therefore equals
\(\mathcal J_\phi(\widetilde R(h_\ell)F_*)\).

For \(h\in H(k)\), moving \(h\) to the left of the exchange factors
changes \(x_i\) to \(h^{-1}x_i\). This gives \(|\det h|^s\).
In \(L_\phi(R(h)F)\) put \(h'=h_0h\) in the integral over \([H]\).
Right invariance of its measure moves \(R(h^{-1})\) onto the column.
Unitarity of \(\kappa\) and conjugate linearity there give exactly
\[
 \mathcal J_\phi(\widetilde R(h)F)=
 |\det h|^s\mathcal J_{\widetilde R(h^{-1})\phi}(F).
\]
Now \(|\det h_\ell|=q^{-1}\). Given a graph row \(F\), take
\(F_*=\widetilde R(a_\ell^{-1})F\). Finally, for \(vJ_\ell=J_\ell\),
the element \(h_\ell^{-1}vh_\ell\) belongs to \(\GL_a(O)\), fixes
\(F_*\) in the twisted action and has absolute determinant one.
Applying the last formula with this element proves the column
invariance.
\end{proof}

\begin{proposition}[The integral cost of averaging the column]\label{prop:columntrace337}
Suppose \(\phi\) is fixed in the twisted action by the covector
stabilizer \(P_\ell\) of Theorem~\ref{thm:terminalexchange336}.
In the transport formula its translated column can be replaced by
the mean of \(q^a-1\) translates, fixed under \(\GL_a(O)\).
For a column fixed by the line stabilizer the corresponding number
is \((q^a-1)/(q-1)\). At the \(f\) terminal places the total
transport factor on the evaluated periods is \(T^{-sf}\).
\end{proposition}
\begin{proof}
In a basis with \(\bar\ell\) the last coordinate, take
\(h_\ell=\operatorname{diag}(1,\ldots,1,\pi)\).
Conjugating \(P_\ell\) by \(h_\ell^{-1}\) gives the inverse
image in \(\GL_a(O)\) of the stabilizer of the last basis vector.
Indeed the last row of an element of \(P_\ell\) reduces to
\((0,\ldots,0,1)\). Conjugation divides the off-diagonal part
of this row by \(\pi\) and multiplies the entries above the last
diagonal entry in the last column by \(\pi\). The inverse
calculation gives all elements with the specified last column
modulo \(\pi\). The residue group is transitive on nonzero vectors.
With the line in place of the covector, the last diagonal entry
may reduce to any nonzero element and the orbit consists of lines.
These two orbit sizes give the asserted finite averages. The
spherical row pairs equally with every translate in the mean,
by the transformation formula for \(h\in\GL_a(O)\). Thus averaging
creates no denominator in the evaluated equality. The independent
terminal actions commute and \(q=T\) at every one of the \(f\)
places, giving the stated product.
\end{proof}

If one keeps the average as a vector in a lattice of integral
class functions stable under the indicated translations and
character values, the integer \(H_x=T^{sf}(T^a-1)^f\) clears
all the rational coefficients of this operation. An \(n\)-row
minor uses at most \(H_x^n\) by multilinearity. Lemma~\ref{lem:size97}
bounds \(s,a,f,\log T\) by fixed powers of \(x\), so
\(\log H_x\le f(s+a)\log T=o(\exp(x^\eps)/x^B)\) for every fixed
\(B>0\). This estimate concerns the explicitly specified transport
and the averaging of its tests. Comparing the integral values of
the resulting global period with a theta column is a further step.


\subsection{Following the refined level through the whole terminal flag}\label{sec:terminalwalk338}
We retain the chosen split half and simple character of
Section~\ref{sec:terminalexchange336}, and the column operator \(D\)
of Theorem~\ref{thm:incidence338}. The graph at the symmetric
endpoint is the prescribed \(\Lambda_{\ell_0}\). We specify one
translate at the first position that compares the two endpoint
pairings on the separately zero-sum columns. Write \(R_G^\sim\)
for the compatible twisted action on rows. A lattice on the chosen
half also specifies its dual on the opposite half.

Use the following coordinates on the string of length \(s+1\):
\[
          b_i=\pi^{-1}e_i\quad(1\le i\le s),\qquad
          b_{s+1}=\pi^{-1}z=e .
\]
For \(1\le i\le s+1\), give the \(j\)-th string line lattice
\(\pi O b_j\) if \(j<i\) and \(O b_j\) if \(j>i\).
On the sum of the column space and \(kb_i\) use the graph
\[
     w+t b_i,\qquad w\in O^a,\quad t\in O,\quad
                     t\equiv\ell_0(w)\pmod\pi.
\]
All dual flag lines \(Of_j\) on the chosen half stay fixed. A
common row fixed at the standard lattice in the twisted action
defines a row \(F_i\) at this lattice by translation. Any two
matrices taking the standard lattice to it give the same row.
Write \(F_i^0\) for the row with zero in place of \(\ell_0\).
Thus \(F_{s+1}\) is an arbitrary row of the given graph type.
The row \(F_1\) is \(R_G^\sim(g)F_{s+1}\), where \(g\) first interchanges
\(e\) and \(e_1\) and then multiplies each \(e_i\) by \(\pi^{-1}\),
fixing \(A_k,e,f_1,\ldots,f_s\) in this second operation.

Denote the Bessel pairing with integration in \(X_j\) for
\(j\ge i\) by \(L_{i,\phi}\). In particular \(L_{s+1,\phi}\)
has no exchange at this place. Other places have the same data and
the same exchange variables on the two sides of each identity.
\begin{lemma}[One step with the original Haar measure]\label{lem:flagstep338}
For each \(i\le s\) and every \(P\)-fixed column \(\phi\) with
zero scalar sum at this place,
\[
 L_{i,\phi}(F_i)=\mu_i(\pi O^a)\,q^{s-i}
                                      L_{i+1,D\phi}(F_{i+1}).
\]
\end{lemma}
\begin{proof}
Choose an integral lift of \(\ell_0\). Then
\(F_i=R_G^\sim(Y_i(\pi^{-1}\ell_0))F_i^0\).
The row \(F_i^0\) is fixed by \(Y_i(O^{a*})\), since its
\(i\)-th line is \(\pi O b_i=O e_i\). It is fixed by
\(X_i(\pi O^a)\), since the next line has lattice
\(O b_{i+1}=\pi^{-1}O d_i\). These determinant-one shears
give the two lattices in Theorem~\ref{thm:exchangefourier336}.
That theorem, at \(y=\pi^{-1}\ell_0\), gives the finite sum
\[
 \mu_i(\pi O^a)\sum_{v\in V}
         \overline{\psi_0(\ell_0(v))}
             L_{i+1,\phi}(R_G^\sim(X_i(\widetilde v))F_i^0),
\]
where \(\widetilde v\in O^a\) lifts \(v\).
The zero term pairs with a \(K\)-fixed row. Moving a scalar
unit to the column in the remaining exchanges has Jacobian one,
so this term vanishes by the zero scalar sum of \(\phi\).

For \(v\ne0\), take \(h_v=k_v h\) as in Theorem~\ref{thm:incidence338}.
We claim
\[
          R_G^\sim(X_i(\widetilde v))F_i^0=R_G^\sim(h_v)F_{i+1}.
\]
Indeed the left lattice in \(A_k\oplus k b_{i+1}\) is generated
by \(O^a\) and \(b_{i+1}+\pi^{-1}\widetilde v\). Its projection
to \(A_k\) is \(A'=O^a+O\pi^{-1}\widetilde v=h_vO^a\).
The \(b_{i+1}\)-coordinate modulo \(\pi\) is the unique residue
functional on \(A'\) vanishing on \(O^a\) and taking value one on
\(\pi^{-1}\widetilde v\). This functional is
\(\ell_0 h_v^{-1}\bmod\pi\), since \(k_v v_0\equiv\widetilde v\).
The allowed coefficient of \(b_{i+1}\) for each given \(w\in A'\)
is exactly that residue class. This describes the lattice on the
right, and its other summands are identical. Twisted compact
invariance proves the claim. Moving \(h_v\) onto the column as in
Theorem~\ref{thm:graphtransport337} gives the factor
\(|\det h_v|^{s-i}=q^{s-i}\).

It remains to identify the resulting column in this pairing.
The row \(F_{i+1}\) has \(P\)-invariance, so \(L_{i+1,\xi}(F_{i+1})\)
is unchanged on averaging \(\xi\) over \(P\). With Haar probability
\(dk\) on \(K\), the column sum after this averaging is
\[
 (q^a-1)\rho(h^{-1})
       \int_K\psi_0(\ell_0(kv_0))\rho(k^{-1})\phi\,dk .
\]
To check it, partition \(K\) into the cosets \(k_vP'\).
Conjugation by \(h^{-1}\) transports the normalized average of
\(P'\) to that of \(P\). The coefficient here is conjugated in
the evaluated pairing, which is conjugate linear on columns.
Now replace \(k\) by \(t^{-1}\), which preserves Haar probability,
and partition by \(\eta=\ell_0 t^{-1}\). The integrand is constant
on each \(t_\eta P\) since \(\phi\) is \(P\)-fixed. Each coset has
measure \(1/(q^a-1)\). This gives exactly \(D\phi\), proving the identity.
\end{proof}

\begin{theorem}[Comparison at every terminal place on the original refined columns]\label{thm:terminalwalk338}
Use all \(f\) places above the same \(T\) of Theorem~\ref{thm:common111}.
Let \(A\) be any row of the prescribed graph type at these places.
Write \(g_T=\prod_{v\mid T}g_v\) for the displayed translate,
\(D_T=\prod_{v\mid T}D_v\) on the column, and set
\[
 C_T= T^{fs(s-1)/2}
                 \prod_{v\mid T}\prod_{i=1}^s\mu_{i,v}(\pi_v O_v^a).
\]
Let \(L^T\) have no exchange above \(T\) and let \(J^T\) have
all the first-position exchanges there. On the refined columns
with zero scalar sum separately at each terminal place one has
\[
 J^T_\phi(R_G^\sim(g_T)A)= C_T L^T_{D_T^s\phi}(A),
       \qquad \|D_T^s\phi\|=T^{afs/2}\|\phi\|.
\]
This holds with common arbitrary data at all other places. On
these column spaces the inverse of \(D_T^s\) is a finite integral
combination of translations divided by \(T^{afs}\), with
coefficients extended by the finite character values.
\end{theorem}
\begin{proof}
The lemma can be iterated because \(D\) preserves the level and the
zero scalar sum. The powers of \(q\) multiply to \(q^{s(s-1)/2}\).
Exchanges and translations at distinct places commute and are
absolutely integrable by Theorem~\ref{thm:exchangefourier336}.
Applying the one-place identity at each place proves the comparison.
The norm identity and the formula for an integral multiple of the
inverse follow by multiplying those in Theorem~\ref{thm:incidence338}.
\end{proof}

After removal of the positive Haar factor \(C_T T^{afs/2}\), the
two terminal pairings therefore have identical row Gram forms on
these columns. This conclusion covers a specified translation
at \(T\), with the original symmetric endpoint. Theorem~\ref{thm:markedglobal339} below compares the remaining
exchanges at \(g_1,g_2\) on a common auxiliary translate. The value of a
symmetric projected period on the signed theta column also remains
to be estimated.


\subsection{The two Bessel positions at one pair of finite translates}\label{sec:markedglobal339}
At a point of an interval of Theorem~\ref{thm:markedstrict337}
keep the original word, degree, flag, signed auxiliary row
levels, terminal graph rows and definite refined columns.
At the marked halves over \(g_1,g_2\) take the product \(t_G\)
of the diagonals in Theorem~\ref{thm:auxiliaryflag339}.
At \(T\) take the translate \(g_T\) of
Theorem~\ref{thm:terminalwalk338}. Write \(\mathcal B\) for the
ordinary symmetric Bessel pairing and \(\mathcal J\) for the
full first-position pairing, which integrates the exchanges at
all these halves. At \(T\) use the twisted action on rows and
columns, and write \(R(t_G)\) for ordinary translation at the
auxiliary primes.

\begin{theorem}[Comparison of all the marked exchanges]\label{thm:markedglobal339}
On columns having zero scalar sum separately at each place
over \(T\), for every row \(A\) of the stated types,
\[
 \mathcal J_\phi\big(R(t_G)\widetilde R(g_T)A\big)
     = C_T\!\prod_{v\mid g_1g_2} C_v\,
                        \mathcal B_{D_T^s\phi}\big(R(t_G)A\big).
\]
Here \(v\) runs over the marked half at each base place.
The scalar \(C_T\) uses the original Haar measures of
Theorem~\ref{thm:terminalwalk338}, and each \(C_v\) is the
volume in Theorem~\ref{thm:auxiliaryflag339}. After dividing
by the additional factor \(T^{afs/2}\), the column operator
\(D_T^s\) is unitary. Thus for the two specified translates
the row Gram forms on these full refined column spaces differ
by the square of \(C_T T^{afs/2}\prod_v C_v\).
\end{theorem}
\begin{proof}
Translations at the auxiliary places preserve the terminal graph
type since the local groups commute. First apply
Theorem~\ref{thm:terminalwalk338} with every auxiliary exchange
still present. The column \(D_T^s\phi\) has its prescribed levels
at the auxiliary places, which commute with all its summands.
At each of those places apply Theorem~\ref{thm:auxiliaryflag339}
to remove all remaining exchanges. The product contains exactly
the places in the definition of the first position and no exchange
is left on the right. Absolute integrability in
Theorem~\ref{thm:exchangefourier336} justifies these changes of
integration order. For the Gram assertion take an orthonormal
basis of the separately zero-sum columns. The normalized \(D_T^s\)
maps it to another such basis, so summing products of the two
pairings over this basis gives the claimed scalar square.
\end{proof}

Both translates in this theorem are fixed before varying the
row or column. The source comparison in
\srcloc{Formalism/SharpRestriction/Divisibility/ZetaProductPairs}{199}
uses two translates belonging to further theta and period data.
For the pair above the equality of Bessel positions is exact.
Its common symmetric Gram form, and the identification and
integral size of its value on the prescribed signed theta image,
still enter Section~\ref{sec:missing}. The compact local witness
at an auxiliary place gives no lower bound for that global value.


\subsection{Transport of the projected theta chain at the terminal prime}\label{sec:thetatransport340}
Fix the fields, word, real vectors and measures in
Section~\ref{sec:missing}. On a sequence of theta edges let $R(u)$
denote right translation of automorphic forms. A Schwartz kernel on
an edge is denoted by $\Phi$. At any vertex one may apply a specified
linear projector $P$, including the projector to a compact type with
its prescribed character. In the following identities the splitting
of the Weil representation on each edge is fixed.

\begin{proposition}[Translation of the actual projected image]\leavevmode\label{prop:thetaimage340}
Choose finite adelic elements $g_i$ at the vertices. Replace each test
on an edge with ends $i,j$ by $\omega(g_i,g_j)\Phi$, each projector
at $i$ by $R(g_i)P_iR(g_i)^{-1}$ and the input at its vertex by its
right translate. Then the output is the right translate at the last
vertex of the original projected theta form. This holds in either
direction along the chain and preserves all Hilbert inner products
between outputs, when the same translates are used for each input.
\end{proposition}
\begin{proof}
For an edge directed from $i$ to $j$ use the theta series
$\Theta(h,k;\Phi)$. Its defining sum gives
$\Theta(h,k;\omega(g_i,g_j)\Phi)=\Theta(hg_i,kg_j;\Phi)$.
The input factor at $h$ is now the original factor at $hg_i$,
also when the convention conjugates this factor. Change the variable
from $h$ to $hg_i$ in the quotient integral defining the lift.
Right invariance gives the original output at $kg_j$. This argument
applies to the convergent cuspidal theta integrals used here. It also
applies after reversing the two groups, with the same variable change
on the integrated group. The relation
$(R(g_j)P_jR(g_j)^{-1})R(g_j)=R(g_j)P_j$ permits iteration.
Right translation preserves the automorphic inner product by the
same change of variable. Compact-level character projectors are
transported with the conjugate group and the transported Haar
measure. Any scalar in an additional determinant twist of the right
action is retained separately on the resulting translated form.
\end{proof}

At a terminal place take $g_0=g$ of Theorem~\ref{thm:terminalwalk338}
and $g_i=\operatorname{diag}(g,I_{n_i-r})$ in compatible local frames
for the increasing ranks $n_i=r+it_{\rm r}$. Write
$\mathcal M_i=M_{n_{i+1},n_i}(O)$, and denote the additive lattice
$g_{i+1}\mathcal M_i g_i^{-1}$ by $\mathcal M'_i$. Put
$\mathcal D_i=\mathcal M_i\cap\mathcal M'_i$.

\begin{theorem}[An integral common subdivision on every theta edge]\leavevmode\label{thm:thetacells340}
Take the graph frame at the first vertex, with $h=s+1$, and set
\[
   E_i=(s+2)(n_i+n_{i+1})/2-(h^2+1),\qquad
   A_i=s(n_{i+1}-n_i)/2 .
\]
For any integer $c$, on subdividing by translates of $\pi^c\mathcal D_i$,
one original matrix cell of lattice $\pi^c\mathcal M_i$ gives
$q^{E_i+A_i}$ cells, and one transported cell gives $q^{E_i-A_i}$.
These counts are exact. For the standard inclusion
$Y_i=\binom{I_{n_i}}{0}$, the normalized absolute matrix coefficient
of the indicator of $Y_i+\pi^c\mathcal M_i$ under
$\omega(g_i,g_{i+1})$ is $q^{-E_i}$.

Use this subdivision at every terminal place in the projected chain
of Proposition~\ref{prop:thetaimage340}. For the fine first test of
Theorem~\ref{thm:tests165} and the other matrix indicators of
Proposition~\ref{prop:split173}, the expansion of the transported
Schwartz tensor on these cells has coefficient ideal generated, at
one terminal place, by
\[
               q^{s(n_m-r)/2}\prod_{i=0}^{m-1}\eta_i(g_i,g_{i+1}).
\]
Here $\eta_i$ is the splitting character relative to the matrix model.
The formula concerns the tensor of tests before the integrations.
The logarithm of the number of subdivisions per original tensor cell
over all the terminal places is bounded by a fixed power of $x$ on
the intervals of Section~\ref{sec:missing}, independently of the
chosen congruence depths. The same holds per transported tensor cell.
\end{theorem}
\begin{proof}
In Smith frames let $t_l$ and $u_j$ be the exponents of $g_i$ and
$g_{i+1}$. The relative exponents of $\mathcal M'_i$ are $u_j-t_l$,
one per matrix entry. Changing the two frames by units preserves the
reference matrix lattice. By Proposition~\ref{prop:graphframe340}
each list contains $h$ copies of $-1$, one copy of $1$, and zeros
in the other coordinates. In the square block of size $n_i$,
half the sum of the absolute differences is
$(n_i-h-1)(h+1)+2h=n_i(h+1)-(h^2+1)$. Each of the additional
$n_{i+1}-n_i$ rows contributes $(h+1)/2$ to that half sum.
The result is $E_i$. The sum of the signed differences is
$s(n_{i+1}-n_i)=2A_i$. The positive differences therefore sum
to $E_i+A_i$ and the absolute values of the negative ones to
$E_i-A_i$. These sums give the two indices of $\mathcal D_i$.
Multiplying all lattices by $\pi^c$ preserves their indices. The
intersection is their largest common translation lattice, so each
subdivision is the partition by all cosets it contains.

The center stays fixed since $g_{i+1}Y_i=Y_i g_i$. The absolute
inner product of the two unit vectors supported on these cells is
the measure of their intersection divided by the geometric mean of
their measures. The indices just computed give $q^{-E_i}$.
For the fine support at the first edge, choose $c$ at least one,
the depth, and the conductor depth of the compact character.
Then the original first test is constant on the ordinary cells because its
invertible top block stays in its coset, with the same character,
and the lower block stays divisible by $\pi$. At each other edge
one can take the original indicator depth. The transported support
is a union of transported cells with values multiplied by
\[
 \eta_i(g_i,g_{i+1})|\det g_i|^{n_{i+1}/2}
                    |\det g_{i+1}|^{-n_i/2}
               =\eta_i(g_i,g_{i+1})q^{A_i}.
\]
On the common subdivision evaluation on one cell is one coefficient.
Every nonzero original value is a root of unity and each support is
nonempty. This proves the ideal formula on multiplying the edge
factors, since $\sum A_i=s(n_m-r)/2$. A theta integral or a period
with the other inputs fixed can now be evaluated as the finite sum
of its values on these cell indicators, with complex conjugation
of coefficients on a conjugate-linear input. This is an equality
for each actual form and each transported projector.

Finally the ranks $n_i$, the length $m$, the number $f$ of terminal
places, and $\log T$ are each bounded by a fixed power of $x$ by
Theorem~\ref{thm:common111} and the formula for $n_i$. Hence both
$f\sum_i(E_i\pm A_i)\log T$ obey the claimed bound. The fine first
edge uses the whole coset matrix of Theorem~\ref{thm:pintable340}
when it is paired back against its original test space.
\end{proof}

These formulas follow a given projected automorphic form under a
specified change of tests and projectors. The independent signed
column in \srcloc{Formalism/SharpRestriction/Divisibility/ZetaProductPairs}{107}
also prescribes the value at the original endpoint. To obtain its
comparison in Section~\ref{sec:missing}, one must evaluate the
resulting signed sum of cell periods with that endpoint projector.
The cell count bounds the number of terms. Their phases and values
on the automorphic forms determine the cancellation in the sum.


\subsection{All the marked tests in the same theta chain}\label{sec:markedtests341}
Use the centered diagonal at every auxiliary half, the translate of
Theorem~\ref{thm:terminalwalk338} at the terminal halves, and the
same fields and interval as in Section~\ref{sec:missing}. Take
all the test functions at the marked places of the common theta
operator in Section~\ref{sec:common177}. Keep its first-edge
terminal depth and allow all residue cosets at that depth. Each
later terminal edge and each auxiliary edge has its specified
single matrix indicator. Write $M_v$ for the normalized matrix of
Theorem~\ref{thm:pintable340} at a terminal place. Every factor
uses its original local norm to normalize inner products.

\begin{corollary}[The coefficient ideal and rank on the marked tensor of tests]\leavevmode\label{cor:markedtests341}
Up to the product of the scalar splitting characters on the singleton
edges, the mixed Gram matrix of this tensor and its translate is
\[
 \left(\prod_{v\mid g_1g_2}q_v^{-\sum_{i=0}^{m-1}L_i}\right)
 T^{-f\sum_{i=1}^{m-1} E_i}\ \bigotimes_{v\mid T}M_v .
\]
Here $L_i$ is the auxiliary cost of Theorem~\ref{thm:auxoverlap341},
$E_i$ is the terminal cost of Theorem~\ref{thm:thetacells340}, and
one marked half is used per base place. In the coefficient field
with these scalars adjoined, the ideal of its entries is the product
of the displayed scalar ideals, the splitting character ideals of
the singleton edges, and the ideals $(m_v)$ in
Theorem~\ref{thm:pintable340}. If $b$ terminal places have positive
first-edge depth, its rank fraction is at most
\[
             \left(\frac{{s+1\brack s}_T}{{r\brack s}_T}\right)^b .
\]
These formulas use the same pair of local translations as the
comparison of Bessel positions in Proposition~\ref{prop:centeredflag341}.
\end{corollary}
\begin{proof}
Choose the extension by identity on the partner coordinates at each
successive vertex, independently at every marked place. For a
singleton indicator its center stays fixed. Its normalized matrix
coefficient is its splitting character times the positive overlap
computed in the respective theorem. In the Bessel comparison the
symmetric endpoint also carries the auxiliary translate. The matrix
here measures the change from the original Schwartz tensor to the
translate at the first position. Its auxiliary factors also measure
the common auxiliary change at the symmetric endpoint. The fine
first-edge computation at $T$ retains all residue cosets in $M_v$.
An inner product of simple tensors is the product of their inner
products, which proves the identity entry by entry. Each local
matrix has a nonzero entry, and all nonzero entries in each $M_v$
are equal. Taking one such entry in each factor gives a generator
of the asserted ideal. Ranks multiply under tensor products of
finite-dimensional linear maps, giving the bound at the positive
depths. The auxiliary splitting characters are units by
Theorem~\ref{thm:auxoverlap341}.
\end{proof}

The rank here refers to Schwartz tests before integrating a theta
kernel. Proposition~\ref{prop:thetaimage340} transports the full
automorphic projector. An assessment of the period on the fixed
original projector still needs the signed global sum of the
subdivided tests, even when all the local matrix entries above are
known. This is the endpoint requirement of Section~\ref{sec:missing}.

\section{Global theta lifting through the chain}\label{nav:g-VI-2}

\subsection{First occurrence on the real row}\label{nav:VIII-00}
\label{sec:first167}

Fix a field \(E/F\), \(F=E^+\), and a translate on the power
intervals of Section \ref{nav:I-28}. Write \(t=t_{\rm r}\),
\(m=2s+2\), \(r=k+2s+1\), and \(w=r+t\).
Write \(G=U(Y_0)\) and \(H_1=U(Y_1)\). Their signatures at
every real place are \((k+s,s+1)\) and \((w-s,s)\) by
Proposition~\ref{prop:ranks47} and Theorem \ref{thm:rangenew50}.
We work with the forward theta lift having
the row representation as a quotient of the local oscillator.
For global functions this convention pairs the theta kernel
with the complex conjugate of the row cusp form.

At a chosen complex embedding above each real place we use
Atobe's Hermitian source, skew-Hermitian target and signature
convention. The local unitary splitting character attached to
the target has angular exponent \(-2e_{1,\tau}\), so it sends
\(z\) to \((z/|z|)^{-2e_{1,\tau}}\). Here an infinity exponent
of \(a\) denotes \((z/\bar z)^a\), interpreted as angular
exponent \(2a\). The first target has even dimension. The
source character has odd angular exponent, and the choice
\(-2e_{1,\tau}-1\) gives the first pair of the archive.
The half exponents are given at
\srcloc{Formalism/LocalRingCardinality/DataTheta}{610}
and in the alternating prescription on line 727 there.
The scalar multiplying the target Hermitian form to make it
skew-Hermitian may be chosen with the indicated imaginary sign
at every selected complex embedding. Indeed, multiply a fixed
nonzero trace-zero element of \(E\) by an element of \(F^\times\)
with prescribed real signs, using weak approximation. Positive
real changes of scale give isometric real forms. More
explicitly, the real additive character in Ichino's
coordinates is \(y\mapsto e^{-2\pi i y}\), the infinite
component of the usual adelic trace character with this
sign. Ichino identifies this normalization with the one in
Atobe's criterion in Section 4.3 of his paper.

\begin{theorem}[The first real lift in the specified tower]
\leavevmode\label{thm:occ167}
Let \(\sigma_\tau\) have the signed row of Volume~I, equation
\textup{\eqref{eq:src109}}, as specified below. At every real
place its theta lift to \(U(w-s,s)\) is nonzero. Its lift to
\(U(w-s-h,s-h)\) is zero for each \(1\le h\le s\).
These assertions hold simultaneously for every word.
\end{theorem}
\begin{proof}
Drop \(\tau\) from the notation. In decreasing order the row
consists of \(v_1,\ldots,v_{k-1}\), all marked \(+\), followed
by \(-b_1,\ldots,-b_m\) marked \(+,-,+,-,\ldots,+,-\).
Here \(v_{k-1}=0\), all \(v_i\ge0\),
\[
 b_j=B+2(j-1)H_0,\qquad e_1=B+H_0,\qquad B>0,\qquad
 H_0>(t+3)/2 .
\]
The signed entries specify the two compact blocks of the
Harish-Chandra parameter \(\lambda\).

Atobe's criterion first subtracts half the target angular
exponent from each entry. Thus the entries here are
\(v_i+e_1\) and \((3-2j)H_0\). They are integers of absolute
value at least \(H_0\). In the plus block there are \(k\)
positive and \(s\) negative entries. In the minus block
there are zero positive and \(s+1\) negative entries.
In particular there is no zero in either block and no
centered string of positive length. Definition 1.6 of Atobe
therefore gives
\[
     k_\lambda=-1,\qquad r_\lambda=k+s+1,\qquad s_\lambda=s .
\]
The second count adds the positive plus entries to the negative
minus entries. The third adds the other two counts.

For \(k_\lambda=-1\) the nonvanishing criterion at signature
\((r_\lambda+2a+l+1,s_\lambda+l)\), \(a\ge0\), requires
\(l\ge0\) and at most \(l\) surviving entries of each sign
\(\epsilon\) with \(0\le\epsilon\alpha-1<a+l\), where
\(\alpha\) is the shifted entry. The deletion procedure in the
criterion only removes entries. Choose \(a=(t-1)/2\), an
integer because \(t\) is odd. For \(l=0\) even the original
list contains no entry in the indicated range. Indeed
\(\epsilon\alpha\le-H_0\) or \(\epsilon\alpha\ge H_0>a+1\).
The resulting signature is exactly \((w-s,s)\).
Removing \(h\) hyperbolic planes gives the same \(a\) and
\(l=-h\), which violates the necessary lower bound on \(l\).
The parameters used here are common to the words at this
place. We have used the discrete-series case of
\href{https://arxiv.org/pdf/1610.07794}{Atobe, Definition 1.6
and Theorem 1.7}, which is unconditional.
\end{proof}

At a place \(v\) let \(\omega_v\) be the oscillator restricted
to \(G_v\), with invariant inner product linear in the first
variable. For a unitary representation \(\pi_v\) of \(G_v\)
and smooth vectors \(u,u'\), write
\[
 Z_v(\phi,\phi';u,u')=\int_{G_v}
 \langle\omega_v(g)\phi,\phi'\rangle\,
 \overline{\langle\pi_v(g)u,u'\rangle}\,dg .
\]
Write \(Z_v(\phi;u)\) for the diagonal value. We specify
convergence for the data at which we use this integral.

\begin{proposition}[Positive integrals on the chosen row vectors]
\label{prop:local167}
At each real place, for every unit vector \(u\) in the lowest
compact type of \(\sigma_\tau\), there exists a polynomial
Fock test \(\phi\) for which \(Z_\tau(\phi;u)>0\).
At a terminal split place take the test \(\phi_{1,j}\) of
Theorem~\ref{thm:tests165}, writing \(\phi_{1,0}=\phi_*\),
\(L_0=K\), and \(I_0=1\). For every local unitary row one has
\[
 Z_v(\phi_{1,j};u)
   =\|\phi_{1,j}\|^2\operatorname{vol}(L_j)
       \|e_{L_j,\kappa}u\|^2 .
\]
In particular this is positive on every nonzero
\((L_j,\kappa)\)-vector. On an original \((K,\kappa)\)-vector
it is exactly \(I_j^{-2}Z_v(\phi_*;u)\).
\end{proposition}
\begin{proof}
At infinity the rank inequality \(w\ge r\) and
Theorem~\ref{thm:occ167} meet the hypotheses of the local
inner product theorem for discrete series.
\href{https://arxiv.org/pdf/2002.09148}{Ichino,
Proposition 7.2} gives a nonzero integral between two tests
and two compact-finite vectors of \(\sigma_\tau\).
All the row coefficients here belong to \(L^1\) by
Lemma~\ref{lem:translate163}. Multiplying the conjugate
coefficients by the formal degree \(d_\sigma\) gives
integrable matrix units. Polarization expresses a mixed
coefficient as a linear combination of diagonal ones.
Consequently \(\omega_\tau(p_z)\ne0\) for some unit
compact-finite \(z\), where
\(p_z(g)=d_\sigma\overline{\langle\sigma_\tau(g)z,z\rangle}\).
The projection and orbit-vector argument in
Theorem~\ref{thm:ordinary163} embeds a Hilbert copy of
\(\sigma_\tau\) into the oscillator. On that copy
\(\omega_\tau(p_u)\) acts as the projection onto the line
of \(u\). This bounded orthogonal projection on the full
oscillator is therefore nonzero. Fock polynomials are dense
in its Hilbert space. In fact some coordinate monomial \(\phi\) has
\(\omega_\tau(p_u)\phi\ne0\), and
\[
 Z_\tau(\phi;u)=d_\sigma^{-1}\|\omega_\tau(p_u)\phi\|^2>0 .
\]
At \(v\mid T\) the oscillator coefficient equals
\(\|\phi_{1,j}\|^2\kappa(g)\mathbf1_{L_j}(g)\) by
Theorem~\ref{thm:tests165}, including the same calculation
with the entire support at \(j=0\). Integration gives the
orthogonal character projector. The norm of the test squared
and the volume of the subgroup each decrease by \(I_j\)
from their values at \(j=0\).
\end{proof}

\subsection{A nonzero global period on the first theta pair}\label{nav:VIII-01}
\label{sec:global167}

Keep the splitting at infinity from Section~\ref{sec:first167}.
Fix compatible global unitary splitting characters. Their
restrictions to the ideles of \(F\) are the required powers
of the quadratic character of \(E/F\). For completeness the
real specifications above admit such characters. In the
norm-one idele class group of \(E\), take the images of the
norm-one class group of \(F\) and of \(\prod_{\tau\mid\infty}S^1\).
Both images are compact. Their intersection consists of real
sign tuples. Indeed, if their representatives differ by a
principal idele \(a\in E^\times\), its finite components lie
in the local copies of \(F\). This forces \(a=\bar a\).
The specified characters agree on the intersection by parity.
They thus give a continuous character of the product of the
two compact subgroups. Character extension for closed
subgroups of locally compact abelian groups extends it to the
norm-one class group. Extend trivially along the positive
idele norm direction chosen in \(F\). This argument establishes
existence. The statements below apply to any choice with these
real specifications, including preassigned compatible finite
characters.

Write \([G]=G(F)\backslash G(\mathbb A_F)\) and use analogous
notation for other groups over \(F\). For a Schwartz test
\(\phi\) and a cuspidal row \(f_0\), set
\[
 \Theta_\phi f_0(h)=\int_{[G]}\theta_\phi(g,h)
                                \overline{f_0(g)}\,dg .
\]
At every terminal place \(v\mid T\) prescribe the test
\(\phi_{1,j_v}\) of Proposition~\ref{prop:local167}, where
\(j_v=0\) is allowed. At other places we shall choose the
tests. We call a vector \emph{of the prescribed terminal
type} if \(L_{j_v}\) acts on it through \(\kappa_v\) for every
\(v\mid T\).

\begin{lemma}[Finite places away from the terminal prime]
\label{lem:finite167}
For the ranks of Theorem \ref{thm:rangenew50}, and all sufficiently
large \(x\), the following hold.
At each nonterminal finite place and any sufficiently small
principal subgroup \(J\) of the row group there is a Schwartz
test whose row coefficient is a positive constant times
\(\mathbf1_J\). At almost every finite place let \(Q_v\)
be the residue cardinality of \(F_v\) and give the hyperspecial
row subgroup volume one. For the unit lattice test \(\phi_v^0\)
and any unit spherical row vector \(u\),
\[
       |Z_v(\phi_v^0;u)-1|
          \le \sum_{h\ge1} Q_v^{-hD},\qquad
       D=\lfloor(w-2)/2\rfloor-4r^2\ge3 .
\]
These unramified integrals are real and their product outside
any fixed sufficiently large finite set is positive.
\end{lemma}
\begin{proof}
The ranks satisfy \(w=8[F:\QQ]r+4+\lceil x^{377}\rceil n_W\)
and \(r=k+m-1\) with \(k=O(x^{9/2+10\eps})\),
\(m=O(\sqrt x)\) in the range of \(\eps\) of Volume~I.
In particular \(w>8r^2+10\) for large \(x\).
At a nonsplit finite place a Hermitian space has anisotropic
dimension at most two. One can also see the bound by taking
the underlying quadratic form over \(F_v\). A Hermitian
anisotropic subspace of dimension \(d\) gives an anisotropic
quadratic form of dimension \(2d\), and a nonarchimedean
local field of characteristic zero has quadratic \(u\)-invariant
four. This is also the dimension bound in the
\href{https://doi.org/10.2307/2372982}{local classification of
Hermitian forms by Jacobowitz}. Witt decomposition gives at least \(r\)
hyperbolic pairs in the target, including at dyadic places.
In a mixed Schr\"odinger model each such pair supplies one
row-space variable. In coordinates for \(r\) pairs they form
a square matrix \(X\), acted on by change of row-space frame.
This follows by polarizing the hyperbolic pairs into their two
isotropic halves. The isometry \(g\) acts on the position
half by its linear action, with the square root of the
Jacobian and a smooth splitting character. The orthogonal
remainder supplies a separate oscillator factor.

Use \(\mathbf1_{1+\mathfrak p_E^aM_r(\OO_E)}(X)\) and a
fixed nonzero smooth test on the remainder. Over a split
place use the square top block of the matrix model
\eqref{eq:weil161} over \(F_v\) instead, and a smooth
test on its other rows. At both types of places the overlap
of the two square-block supports implies \(g\in J\), the
subgroup congruent to the identity at depth \(a\) in the
chosen frame. Indeed both \(X\) and its transform must lie
in the same matrix subgroup. In dual coordinates this gives
the same assertion after taking the dual frame. For large
\(a\), \(J\) fixes the remainder test and the splitting
character is trivial on \(J\). The Jacobian there is one.
Hence the coefficient is precisely the squared norm of the
test on \(J\), and vanishes outside it.

Outside finitely many places the extensions and the data are
unramified, the additive character has integral conductor,
the lattices are self-dual, and all the hyperspecial groups
fix the unit lattice test. At an unramified nonsplit place
of odd residue characteristic one can split off at least
\(\lfloor(w-2)/2\rfloor\) integral hyperbolic pairs from
the unimodular lattice. To see this, an isotropic vector of
the reduced nondegenerate Hermitian space lifts to a primitive
isotropic vector by Hensel's lemma, since a nonzero residual
vector has nonzero linear derivative in some direction.
Such residual vectors exist in dimension at least two by
surjectivity of the finite-field norm. Pair the lift with a
vector of pairing one and subtract half its norm times the
isotropic vector to give a hyperbolic pair. Its unimodular
complement permits iteration. The lattice test in this
polarization is a product of normalized lattice indicators
and a unit vector for the remainder.

For a right coset \(gK_v\) put \(\Lambda'=g\Lambda\) in
the row space and let \(h\) be least with
\(\varpi^h\Lambda\subseteq\Lambda'\subseteq\varpi^{-h}\Lambda\).
Take one component at a split place. At these good places
\(K_v\) is the full stabilizer, so the lattice determines
the coset. The quotient modulo \(\varpi^h\Lambda\) is
generated by the images of \(r\) basis elements of \(\Lambda'\).
There are at most \(Q_v^{4hr^2}\) ordered choices since
the coefficient residue degree is at most two. In the
nonsplit case the two lattice covolumes are equal. The
coefficient of one row variable then has absolute value
\([\Lambda:\Lambda\cap\Lambda']^{-1}\le Q_v^{-h}\).
Indeed the elementary-divisor exponents sum to zero and the
largest absolute exponent is \(h\). At a split place the
coefficient in absolute value for each of the \(w\) row
variables is
\[
  \bigl([\Lambda:\Lambda\cap\Lambda']
          [\Lambda':\Lambda\cap\Lambda']\bigr)^{-1/2}
                      \le Q_v^{-h/2}.
\]
This follows by dividing the intersection volume by the
geometric mean of the two volumes. The residual oscillator
coefficient in the nonsplit case and the unit row coefficient
in both cases have absolute value at most one. The integral
over \(K_v\) is one. Summing over the shells \(h\ge1\) gives
the bound, also for the integral of the absolute value off
\(K_v\). Inversion makes each diagonal integral real, and
the displayed error is at most \(1/7\). Finally
\(\sum_v Q_v^{-3}<\infty\) for a fixed number field, by
comparison with \([F:\QQ]\sum_{n\ge2}n^{-3}\).
The product of the positive factors therefore converges to
a positive number.
\end{proof}

\begin{theorem}[Detection of a cuspidal row on the first pair]
\leavevmode\label{thm:global167}
Fix a unit pure tensor \(u_\infty=\bigotimes_\tau u_\tau\)
of lowest compact types of the original signed row.
Let \(A=S_{\boldsymbol j}(u_\infty\otimes\ell)\), with an
arbitrary finite class vector \(\ell\) at the chosen levels.
Whenever \(A\ne0\), there is a factorizable complex Schwartz
test \(\phi\), with the prescribed terminal factors, such that
\(\Theta_\phi A\) is a nonzero cusp form on \([H_1]\).
Equivalently for this existence assertion, there is a cusp
form \(D_1\) on \([H_1]\) with
\[
               \langle \Theta_\phi A,D_1\rangle_{[H_1]}\ne0 .
\]
The assertion includes \(\boldsymbol j=0\) whenever the
original sum \(A\) is nonzero.
\end{theorem}
\begin{proof}
By Theorem~\ref{thm:ordinary163} and its finite-level
restriction in Section~\ref{sec:depth163}, the square-integrable
vector \(A\) is fixed by the real matrix unit \(p_{u_\infty}\).
It generates under the real group a copy of the discrete
series \(\sigma=\bigotimes_\tau\sigma_\tau\). At a fixed
finite level the adelic quotient is a finite union of
arithmetic real quotients. Restrict the real representation
to each such component by the corresponding projection.
Its nonzero image is again a Hilbert copy of \(\sigma\),
as one sees from the matrix-unit argument. Its maximal
compact-finite vectors are smooth automorphic eigenfunctions
(the Casimir together with the compact Laplacian gives
elliptic regularity). The discrete series is tempered.
Wallach's constant-term theorem places its full square-integrable
multiplicity on each arithmetic quotient in the cuspidal
spectrum. This proves adelic cuspidality. Indeed the finite
adelic unipotent group is the product of its rational points
and any given compact open subgroup by strong approximation
for unipotent groups. Thus the adelic constant term at a
finite class representative is a positive multiple of the
arithmetic real constant term for that representative.
\href{https://mathweb.ucsd.edu/~nwallach/tempered-cuspidal.pdf}{Wallach,
Theorem 4.3}.

Theorem~\ref{thm:occ167} forces all theta lifts of these
cuspidal vectors to smaller spaces in the same rational Witt
tower to vanish. A smaller space is obtained by removing
hyperbolic planes at every place. A nonzero global lift
would give a nonzero equivariant pairing at a real place by
fixing the other tensors and evaluating the lift. In
particular the local oscillator there would have the row
as a quotient. The tower property of the global theta
correspondence expresses the constant term of the lift of
a cusp form along a target isotropic parabolic in terms of
the lifts to its smaller Witt space. Consequently every
lift under consideration to \(H_1\) is cuspidal.
\href{https://arxiv.org/pdf/1207.4709}{Gan, Qiu and Takeda,
introduction and Section 11}.

We give the norm argument since the vectors at \(T\) are
fixed. The global Witt index \(d_1\) of \(Y_1\) is at most
\(s\), as seen at a real place. We have \(w-d_1>2r\)
and \(w>2r+2\). Hence the ordinary doubled Siegel--Weil
theta integral is in Weil's convergent range, and its Siegel
Eisenstein series at \((w-r)/2\) converges absolutely.
For cusp forms \(f_1,f_2\) the convergent inner product
formula at a common test gives, with a fixed \(c>0\),
\begin{align}
 \langle\Theta_\phi f_1,\Theta_\phi f_2\rangle_{[H_1]}
  = c\int_{G(\mathbb A_F)}
     \langle\omega(g)\phi,\phi\rangle\,
       \overline{\langle R(g)f_1,f_2\rangle_{[G]}}\,dg .
 \label{eq:rallis167}
\end{align}
Here \(R\) denotes right translation. To recall the unfolding,
the product of a theta kernel with the conjugate of a second
one is the restriction of the doubled kernel. Under the
polarization by the diagonal and antidiagonal, evaluation of
the transformed tensor at zero gives the oscillator inner
product. The determinant twist on the second factor cancels
the corresponding twist in the doubling embedding. The
convergent Siegel--Weil formula identifies the target integral
with the Siegel Eisenstein series, up to the Haar normalization.
Its rational terms are indexed by maximal isotropic planes
in \(Y_0\oplus(-Y_0)\). An open-orbit plane is a graph of
an isometry, with diagonal stabilizer. If a plane meets the
first summand in \(D\ne0\), its projection to that summand
lies in \(D^\perp\). The unipotent radical stabilizing \(D\)
acts trivially on \(D^\perp/D\), preserves the plane and acts
there with determinant one. Its integral on the first cusp
form is zero. The same argument applies to the second summand,
and all remaining planes are graphs. Unfolding the graph
orbit gives exactly the integral in \eqref{eq:rallis167}.
Weil convergence and rapid decrease of cusp forms justify
the kernel integrals. Absolute convergence of the open-orbit
integral for the tests below also follows directly from
Lemma~\ref{lem:finite167}, the compact support of the bad
finite coefficients and Lemma~\ref{lem:translate163} at
infinity. These are the classical convergent doubling formula
and its local evaluation at zero.
\href{https://arxiv.org/pdf/1207.4709}{Gan, Qiu and Takeda,
Sections 1 and 3} and
\href{https://arxiv.org/pdf/2002.09148}{Ichino, Sections 6.3
and 7.2}.

Project \(A\) onto an irreducible Hilbert summand \(\Pi\) of
the cuspidal spectrum on which it has nonzero projection
\(A'\). Orthogonal projection commutes with \(R\). The
projected vector has the same terminal type and is still
fixed by \(p_{u_\infty}\). At infinity it therefore factors
as \(u_\infty\) times its finite part in the restricted
tensor product realization of \(\Pi\). At almost all finite
places the finite part has the spherical line. At the
remaining finite places it lies in a tensor product of
finite-dimensional spaces of vectors fixed by compact open
subgroups, by smoothness and admissibility. These facts follow
also by applying the finite-level averaging projectors to
the restricted tensor product. They allow the finite part
to be a sum of pure tensors.

Choose the real tests by Proposition~\ref{prop:local167}
at \(u_\tau\), and the good finite lattice tests by
Lemma~\ref{lem:finite167}. At every remaining nonterminal
finite place choose the small subgroup in that lemma to
fix the entire finite-dimensional space of components of
\(A'\). The integrated local form is then a positive scalar
times its Hilbert inner product on those components. The
same holds on the terminal type by Proposition~\ref{prop:local167}
and polarization. At the good places it holds on the
one-dimensional spherical line, with positive scalar whose
product converges. At the real places only the line \(u_\tau\)
occurs in \(A'\), and its scalar is positive. Expanding the
finite part into pure tensors and using the absolutely
convergent product of the local integrals, we obtain
\[
      \|\Theta_\phi A'\|^2
          =c\|A'\|^2\prod_v c_v>0 ,
\]
where \(c_v\) is the value of the local form on a unit vector
in the indicated subspace. This argument permits arbitrary
unitary finite components of \(\Pi\). Finally
\(\langle R(g)(A-A'),A'\rangle=0\) for every \(g\).
The mixed identity \eqref{eq:rallis167} yields
\(\langle\Theta_\phi A,\Theta_\phi A'\rangle
=\|\Theta_\phi A'\|^2>0\).
Taking \(D_1=\Theta_\phi A'\) gives the asserted period.
\end{proof}

\begin{corollary}[One test away from the terminal prime]
\label{cor:uniform167}
Fix the field, translate, splitting, pure real row \(u_\infty\)
and levels away from \(T\). There is a single choice of all
theta factors away from \(T\) and a number \(C_x>0\) such
that at every terminal depth \(\boldsymbol j\) and for every
cuspidal row \(A\) at these levels fixed by \(p_{u_\infty}\)
and of the prescribed terminal type,
\begin{equation}
  e^{-\delta_x} C_x I^{-2}\|A\|^2
    \ \le\ \|\Theta_{\phi_{\boldsymbol j}} A\|^2
    \ \le\ e^{\delta_x} C_x I^{-2}\|A\|^2,\qquad
  \delta_x=4[F:\QQ]\,2^{-D},\quad I=\prod_{v\mid T} I_{j_v}.
  \label{eq:family167}
\end{equation}
The estimate holds simultaneously for the Bessel columns of
all words with these data. The chosen real factors can be
Fock monomials. At finitely many nonterminal finite places
the factors can be characteristic functions in the mixed
coordinates of Lemma~\ref{lem:finite167}.
\end{corollary}
\begin{proof}
Choose a finite set of finite places containing the terminal
places, all ramification and all deviations from the
hyperspecial levels and lattice tests. At each bad nonterminal
place the kernel of the specified level character contains
an open neighborhood of the identity. Take \(J\) of
Lemma~\ref{lem:finite167} in this kernel, using any fixed
lattice indicator for its residual test and shrinking \(J\)
to fix that test. The two local factors in the coefficient
calculation are then independent of the cuspidal representation,
the column, and the terminal depth. For a finite list of level
characters one can take the intersection of their kernels.
Proposition~\ref{prop:local167} supplies the same monomial at
each real place for the whole fixed real row.

In the product formula for an irreducible projection, let
\(C_x\) be \(c\) times the product of the real scalars and
the nonterminal bad finite scalars, times
\(\prod_{v\mid T}\|\phi_{*,v}\|^2\operatorname{vol}(K_v)\).
Each scalar is positive. The terminal contribution for a unit
vector of the fine type differs from this contribution at
depth zero by precisely \(I^{-2}\), by the norm and volume
calculation in Proposition~\ref{prop:local167}. Each good
factor \(z_v\) satisfies
\[
   |\log z_v|\le 2Q_v^{-D},\qquad
   \sum_v Q_v^{-D}\le [F:\QQ]
       \left(2^{-D}+\int_2^\infty y^{-D}\,dy\right)
                         \le 2[F:\QQ]\,2^{-D}.
\]
The first bound follows from the error at most \(1/7\) and
the geometric series in Lemma~\ref{lem:finite167}. Each
rational prime has at most \([F:\QQ]\) places above it.
These estimates place the good product in
\([e^{-\delta_x},e^{\delta_x}]\) independently of the
irreducible summand. Distinct orthogonal summands have zero
mixed integral in \eqref{eq:rallis167}. At any fixed depth
the space of cusp forms of the given levels, compact type and
infinitesimal character is finite-dimensional by the finiteness
theorem for automorphic forms. Adding the squared norms in
these summands proves \eqref{eq:family167} for \(A\).
\end{proof}

For a nonzero Bessel column at depth \(\boldsymbol j\),
Corollary~\ref{cor:uniform167} now gives a nonzero global
first lift with common tests away from \(T\). At the original
level it detects precisely the columns surviving periodization.
The factor \(C_x\) contains the archimedean matrix coefficients
of the chosen monomials and the volumes of the small subgroups.
The estimate gives no arithmetic height or valuation for these
factors. The period in Section \ref{nav:I-28}, also uses the
remaining steps of the specified theta chain.

\subsection{The signed real theta chain to the compact group}\label{nav:VIII-02}
\label{sec:realtower169}

Fix a real place and suppress its subscript on the weights.
Write \(t=t_{\rm r}\), \(d=(t-1)/2\),
\(n_j=r+jt\), \(s_j=\lfloor(m-j)/2\rfloor\), and
\[
          h_j=-\tfrac12\mathbf1_{j\ {\rm odd}},\qquad
          G_j=U(n_j-s_j,s_j) \quad(0\le j\le m).
\]
The target \(G_m\) is compact. On the ordered pair
\(G_j,G_{j+1}\) use the real additive character and the
signature convention of Section~\ref{sec:first167}. In the
normalization of Ichino, the half angular exponents of the
characters attached to target and source are, respectively,
\[
            a_T=-e_{j+1}+h_j,\qquad
            a_S=-e_{j+1}+h_{j+1}.
\]
The parity conditions hold since the ranks differ by the odd
integer \(t\). These exponents occur in
\srcloc{Development/AbelianVarietyTheory/Definitions/TwistedInnerProducts/IsStepWeil}{25}.

Define a list \(P_j\) of signed numbers as follows, sorting
within each compact block after forming the list.
\begin{enumerate}
\item Put \(v_1,\ldots,v_{k-1}\) in the plus block.
\item Include \(-b_i\) for \(1\le i\le m\). The entries with
 \(i\le j+1\) have sign plus. If \(i>j+1\), give the entry
 sign plus when \(i-j-1\) is even and sign minus otherwise.
\item For each \(1\le i\le j\) put the entire string
 \(-e_i+d,-e_i+d-1,\ldots,-e_i-d\) in the plus block.
\item Add \(h_j\) to every entry.
\end{enumerate}
This is the signed parameter at
\srcloc{Formalism/LocalRingCardinality/DataTheta}{538}.
It has \(n_j\) entries. All entries before the final shift are
integers. Each string centered at \(-e_i\) is strictly below
\(-b_i\) and, for \(i<m\), strictly above \(-b_{i+1}\), since
\(e_i-b_i=b_{i+1}-e_i=H_0>d\) whenever the indicated
indices exist. The strings and the other entries are therefore
pairwise disjoint, and the final shift gives the parity
\((n_j-1)/2\bmod\ZZ\). Exactly \(s_j\) entries have sign
minus. The Harish-Chandra classification gives a discrete
series \(\sigma_j\) with parameter \(P_j\) on \(G_j\), including
the usual irreducible finite-dimensional representation when
\(s_j=0\). The parameter at \(j=0\) is the signed row in
Theorem~\ref{thm:occ167}.

\begin{theorem}[Every real lift and its first occurrence]
\label{thm:tower169}
For \(0\le j<m\) the real theta lift of \(\sigma_j\) to
\(G_{j+1}\) is \(\sigma_{j+1}\). Removing any positive
number of hyperbolic planes from this target gives zero
lift in the same Witt tower. Thus the forward real lifts
of the original row give \(P_m\) at the compact endpoint.
\end{theorem}
\begin{proof}
Subtract \(a_T\) from \(P_j\). The entries arising from
\(v_i\) become \(v_i+e_{j+1}>0\). Those from \(-b_i\)
become
\[
                   e_{j+1}-b_i=(2j+3-2i)H_0 .
\]
This is positive exactly for \(i\le j+1\). Each old string
becomes \(2(j+1-i)H_0+d,\ldots,2(j+1-i)H_0-d\),
hence is positive in its entirety. All the positive entries
belong to the plus block, and their number is
\(A=k+j(t+1)\). Among the negative entries the initial sign
is minus and the signs alternate. There are
\(y=s_{j+1}\) negative plus entries and \(z=s_j\)
negative minus entries. In particular \(n_j=A+y+z\).
All centered entries are integers of absolute value at least
\(H_0\).

In Atobe's notation this gives \(k_\lambda=-1\),
\(r_\lambda=A+z\) and \(s_\lambda=y\). As in the proof of
Theorem~\ref{thm:occ167}, the count for each sign at \(d\)
is zero because \(H_0>d+1\). With \(l=0\) the target is
\((A+z+2d+1,y)=(n_{j+1}-s_{j+1},s_{j+1})\).
Removing any \(h\ge1\) hyperbolic planes changes \(l\) to
\(-h\). The necessary condition \(l\ge0\) then fails.
This proves nonvanishing and the first occurrence for this
row at every step.
\href{https://arxiv.org/pdf/1610.07794}{Atobe, Theorem 1.7}.

To identify the nonzero lift, use the discrete-series
correspondence of \href{https://arxiv.org/pdf/2002.09148}{Ichino,
Theorem 4.1(i) and the compact-Levi calculation in Section 5}.
We give the calculation
because the signs of the negative entries change in this
correspondence. In the parabolic defining the cohomologically
induced target, positive entries stay in their original
blocks. Negative entries exchange their blocks. The new
middle block has signature
\[
 ((n_{j+1}-s_{j+1})-(A+z),\,s_{j+1}-y)=(t,0).
\]
All other Levi blocks are single compact lines because the
old entries are distinct. Thus the Levi is compact. Order the
old centered entries decreasingly and place the \(t\) slots
of the middle block between the positive and negative entries.
Put \(N=n_{j+1}\).
For an old centered entry \(\beta\) at combined position \(i\),
Ichino's inducing weight, inclusive of its scalar twist, is
\[
                  \beta-\bigl((N+1)/2-i\bigr)+a_S .
\]
For a positive entry this follows by counting the larger
positive entries ahead of it. For a negative one there are
\(A+t\) earlier slots before the negative list and then all
larger negative entries, in either block. Equivalently there
are \(N-i\) smaller negative entries. Substitution into the
two negative-entry formulas in Theorem 4.1(i) gives the
displayed weight. For every middle slot the inducing weight
is \(A-n_j/2+a_S\).

Refine the parabolic to the compact Cartan by putting the
middle slots in decreasing order. For a compact Levi the
cohomological induction of its character, in the regular
good range, is the discrete series with parameter equal to
the inducing weight plus the half sum of the positive roots
of this ordering. Indeed one may first induce from the
compact Cartan inside the compact Levi, giving its character
by the compact highest-weight formula, and then use induction
in stages. For the complex root system of \(\mathrm{GL}_N\) the \(i\)-th
entry of that half sum is \((N+1)/2-i\). Adding it back gives
\(\beta+a_S\) at each old entry and, at the middle positions
\(A+i\) for \(1\le i\le t\), gives
\[
     A-n_j/2+a_S+(N+1)/2-(A+i)
                     =a_S+(t+1)/2-i .
\]
The resulting entries strictly decrease in the combined
ordering. This verifies regularity and the good-range
inequalities just used, including across both ends of the
middle block. The old entries have now undergone the full
translation \(a_S-a_T=h_{j+1}-h_j\). Their negative
entries before translation were exactly those coming from
\(-b_i\) with \(i\ge j+2\). Flipping those signs turns the
old alternating tail into the prescribed tail for \(j+1\).
All old strings retain their plus signs. The middle slots
give precisely the new string centered at
\(-e_{j+1}+h_{j+1}\), with plus signs. We have obtained
\(P_{j+1}\), proving the assertion.
\end{proof}

\begin{corollary}[Real test vectors through the compact endpoint]
\label{cor:vectors169}
Fix a nonzero compact-finite vector at the first row and a
nonzero linear functional on the final compact representation.
At this real place there are polynomial Fock tests, one at
each step, whose successive local theta maps give a nonzero
value for that functional. The tests can be chosen to be
monomials in fixed Fock coordinates if only nonvanishing of
the final value is prescribed. The product over all the
real places is nonzero for pure initial vectors and terminal
functionals of this kind.
\end{corollary}
\begin{proof}
Let \(V_j\) be the compact-finite module of \(\sigma_j\).
By local Howe duality and Theorem~\ref{thm:tower169} the
oscillator has the irreducible quotient
\(\sigma_j\boxtimes\sigma_{j+1}\) in the smooth local
correspondence. On polynomial Fock vectors its image is the
algebraic tensor \(V_j\otimes V_{j+1}\).
Indeed the polynomial vectors (in the Schr\"odinger model,
Hermite polynomials times the Gaussian) are dense in the
smooth Schwartz space and are stable under the complexified
Lie algebras and compact groups of the pair. Their image
is nonzero, consists of compact-finite vectors, and is an
invariant submodule of an irreducible Harish-Chandra module.
It must be the entire module.

Given a nonzero input \(u_j\in V_j\) and a desired
\(u_{j+1}\in V_{j+1}\), choose a polynomial mapping to
\((u_j/\|u_j\|^2)\otimes u_{j+1}\). Contracting the first
factor against \(u_j\) in the invariant Hermitian pairing
gives \(u_{j+1}\). Choose all intermediate vectors nonzero
and choose the final one with nonzero value under the
given functional. This constructs the successive tests.
Expand each polynomial into its finitely many monomials.
Successive contraction is additive in each test, with complex
conjugation on the scalar in alternating positions. If every
tuple of monomials gave value zero, the polynomial tuple
would give zero as well. Hence one monomial tuple suffices.
Multiplying nonzero scalars over the finite set of real
places proves the last assertion.
\end{proof}

Here \(P_j\) labels the forward quotient of the oscillator.
The contragredient model negates its Harish-Chandra entries,
as in \srcloc{Formalism/LocalToGlobal/IndefiniteReal/ScalarOnPartsFull}{15}
and line 34 there. Complex conjugation of the entire
intertwiner passes to the dual unitary representations and
the conjugate oscillator. Individual monomial tests in the
corollary have coefficient one in the chosen Fock model.
Their images under the quotient maps depend on the scalars
normalizing those maps. The integral comparison along the
global chain requires those scalars and the prescribed finite
data.

\subsection{The finite standard \texorpdfstring{\(L\)}{L}-values in the theta chain}\label{nav:VIII-03}
\label{sec:values171}

Keep the fields and spacings of Section \ref{nav:I-28}. Write
\(F=E^+\), \(f=[F:\QQ]\), \(t=t_{\rm r}\), \(d=(t-1)/2\),
\(z_0=d+1\) and \(n_j=r+jt\). Write \(S\) for a finite set of
finite places of \(F\), and \(S_E\) for all the places of \(E\)
above \(S\). An \(L\)-function with superscript \(S\) or
\(S_E\) omits the indicated Euler factors and all gamma factors.
We study the scalar at \(z_0\) that occurs before the ramified
local integrals in the doubling formula.

The splitting characters can be written
\[
 D_j=\eta^{-\mathbf1_{j\ {\rm odd}}},\qquad
 \chi_{S,i}=\mu_iD_i,\qquad \chi_{T,i}=\mu_iD_{i-1}.
\]
Here all the characters are unitary Hecke characters of \(E\).
The \(\mu_i\) are trivial on \(\mathbb A_F^\times\) and at a
selected complex embedding \(\tau\) have character
\((z/|z|)^{-2e_{i,\tau}}\). The exponent for \(\eta\) is \(1\).
We take \(S\) large enough to contain the finite ramification
of the data. Put \(\rho_{ik}=\mu_i\mu_k^{-1}\) for \(i<k\).
The exponent of \(\rho_{ik}\) at \(\tau\) is
\(4(k-i)H_0\). In particular the common translation of all
the \(e_{i,\tau}\), which can be very large, disappears.
The splitting convention is recorded in
\srcloc{Formalism/LocalToGlobal/IndefiniteReal/HasInfExpF}{17}.

\begin{lemma}[A bound for the initial spherical parameters]
\label{lem:sat171}
Let \(B_{0,w}\) be the standard unramified base-change parameters
of a unitary spherical representation of the row group at a
good finite place \(v\), for \(w\mid v\). Then
\[
     \# B_{0,w}=r,\qquad
     |\alpha|\le (Nw)^{c_r}\quad(\alpha\in B_{0,w}),\qquad
     c_r=10(r+1)^3 .
\]
\end{lemma}
\begin{proof}
Write \(q\) for the residue cardinality of \(F_v\), let \(K\)
be the hyperspecial subgroup, and give \(K\) volume one.
At a split place take \(G=\mathrm{GL}_r(F_v)\). At an inert
place the unimodular odd-dimensional Hermitian space has the
usual unramified group with relative Weyl group consisting
of signed permutations on \((r-1)/2\) coordinates.
In either cocharacter lattice put
\(\lambda_i=(1,\ldots,1,0,\ldots,0)\) with \(i\) ones.
At an inert place the diagonal in \(G\) also contains the
reversed negatives of these coordinates. Write \(M_\lambda\)
for the sum of the distinct Laurent monomials in the Weyl
orbit of \(\lambda\).

We use the normalized Satake isomorphism. The transform of
the indicator of \(K\lambda_i(\varpi)K\) is Weyl invariant,
has nonnegative coefficients, is supported on dominant
weights \(\nu\le\lambda_i\) and their orbits, and has
coefficient \(q^{\langle\varrho,\lambda_i\rangle}\ge1\)
on \(M_{\lambda_i}\). Here \(2\varrho\) sums the positive
relative roots with their multiplicities. These are the
triangularity and leading coefficient in Satake's transform.
Nonnegativity also follows directly by taking the unipotent
constant term of an indicator and multiplying by the square
root of the modulus. The sum of the coefficients is its
spherical eigenvalue in the normalized principal series
induced from the trivial unitary character. Unitarity bounds
that eigenvalue, and the absolute value of the eigenvalue in
the given representation, by the volume of the double coset.
\href{https://www.numdam.org/item/PMIHES_1963__18__5_0.pdf}{Satake,
\emph{Theory of spherical functions on reductive algebraic
groups over \(p\)-adic fields}, Sections 5--7}.

The volume is at most \(V=q^{4r^2}\). Indeed, the right cosets
give distinct lattices between \(\varpi\Lambda\) and
\(\varpi^{-1}\Lambda\). Their images modulo \(\varpi\Lambda\)
are determined by \(r\) generators in that quotient, each
with at most \(q^{4r}\) possibilities. At a split place one
component suffices. In the split chamber, dominance bounds the
largest coordinate by one and the smallest by zero from below,
and fixes the sum \(i\). At an inert place the coordinates are
decreasing nonnegative integers, the largest is at most one
and the sum at most \(i\). Thus the only possible dominant
weights are \(\lambda_i\) in the split case and
\(\lambda_0,\ldots,\lambda_i\) in the inert case. At the
Satake point of our representation, triangularity consequently
gives \(|M_{\lambda_i}|\le(2V)^i\). To see the induction,
subtract the lower terms from the eigenvalue, use the bound
\(V\) on the sum of their nonnegative coefficients, and divide
by a leading coefficient at least one. Start with \(M_0=1\).

For the split standard representation the elementary symmetric
functions of the parameters are exactly these orbit sums.
For the inert one the parameters are
\(z_1,\ldots,z_a,1,z_a^{-1},\ldots,z_1^{-1}\),
where \(a=(r-1)/2\). Indeed, a diagonal dual Frobenius
\(s\vartheta\) restricts over \(E_w\) to \(s\vartheta(s)\),
where \(\vartheta\) reverses and inverts the diagonal.
The \(z_i\) are the values on individual relative coordinates.
Each elementary symmetric function is now a nonnegative
integral combination of the same orbit sums, with coefficient
sum at most \(2^r\), by expanding the subsets of this list.
The elementary root bound for a monic polynomial gives
\(|\alpha|\le1+2^r(2V)^r\le q^{c_r}\). At the other component
of a split place apply the same argument to the contragredient,
which is again unitary. Since \(q\le Nw\), this suffices.
\end{proof}

\begin{lemma}[Hecke characters at the endpoint of convergence]
\label{lem:hecke171}
Let \(K\) be a CM field of degree \(2f\). Let \(\rho\) be a
unitary Hecke character with pure angular exponents
\(\nu_\tau\in\ZZ\), at least one nonzero. Let \(U\) contain its
finite ramification and put
\[
 M(\rho,U)=1+\log(2|D_K|N\mathfrak f_\rho)
       +f\log(3+\max_\tau|\nu_\tau|)+\#U .
\]
The value \(L_K^U(1,\rho)\) is nonzero and
\[
 \big|\log|L_K^U(1,\rho)|\big|\ll M(\rho,U),\qquad
 \sum_{a=2}^{\infty}\big|\log|L_K^U(a,\rho)|\big|\ll f .
\]
Both implied constants are absolute. If \(K/E_0\) is CM
quadratic, \(\rho\) is trivial on \(\mathbb A_{E_0}^\times\)
and \(U\) is stable under conjugation, all these values are
positive real numbers.
\end{lemma}
\begin{proof}
Neither \(\rho\) nor \(\rho^2\) is an imaginary power of
the idelic norm, by the nonzero angular exponent. Hecke's
continuation and functional equation give entire \(L\)-functions
with complex factors \(\Gamma_{\CC}(s+|\nu_\tau|/2)\) and
conductor \(|D_K|N\mathfrak f_\rho\). The completed functions
are bounded on closed vertical strips away from their poles
by the global zeta integral, and the incomplete functions have
finite order on strips by Stirling's formula. The trivial
character gives the Dedekind zeta function with a simple pole
at \(1\).
\href{https://sites.math.rutgers.edu/~alexk/2023S572/Tate1950.pdf}
{Tate, \emph{Fourier analysis in number fields and Hecke's
zeta-functions}, Chapters II and IV}.

Here are uniform constants sufficient near \(1\).
On \(\Re s=2\) the Euler product is bounded by
\(\zeta(2)^{2f}\). On \(\Re s=-1\), the functional equation
costs \((|D_K|N\mathfrak f_\rho)^{3/2}\) and a factor at most
\(C(3+|\nu_\tau|+|\Im s|)^3\) at each complex place.
Indeed the absolute gamma quotient there reduces by conjugation
and three applications of \(\Gamma(u+1)=u\Gamma(u)\) to
the absolute value of a cubic times an absolute constant.
This expression extends across poles of the denominator.
Apply the strip maximum principle of Phragm\'en--Lindel\"of
after dividing by \((s+3)^{3f}\). For \(\zeta_K\) first
multiply by \((s-1)/(s+3)\). The same bounds apply to
\(\rho^2\) since its conductor divides \(\mathfrak f_\rho\)
and its angular exponents double. Removing the local factors
in \(U\) costs at most \(2^{\#U}\) on \(|s-1|\le1/2\).
Consequently on that disk
\[
 |(s-1)\zeta_K^U(s)|,\quad |L_K^U(s,\rho)|,\quad
 |L_K^U(s,\rho^2)|\ \le A=\exp(C_0 M(\rho,U))
\]
for an absolute \(C_0\) with \(A\ge2\).
Cauchy's inequality on the smaller disk of radius \(1/4\)
gives \(|(L_K^U)'(s,\rho)|\le4A\).

For real \(\sigma>1\), take Euler logarithms and apply
\(3+4\cos u+\cos(2u)=2(1+\cos u)^2\ge0\) to each prime
power. This gives
\[
  \zeta_K^U(\sigma)^3\,|L_K^U(\sigma,\rho)|^4\,
                |L_K^U(\sigma,\rho^2)|\ge1 .
\]
Hence \(|L_K^U(1+h,\rho)|\ge h^{3/4}/A\) for \(0<h\le1/4\).
Choose \(h=(16A^2)^{-4}\). Moving back to \(1\) costs at most
\(4Ah\le h^{3/4}/(4A)\), giving the claimed lower bound after
taking logarithms. For \(a\ge2\), absolute Euler logarithms
are bounded by \(2f\log\zeta(a)\), whose sum is \(O(f)\) by
the geometric series over \(a\).

For positivity, conjugation interchanges a split pair of Euler
factors with their complex conjugates since
\(\rho(z)\rho(\bar z)=\rho(Nz)=1\). At a nonsplit place a
real value \(+1\) or \(-1\) on an unramified uniformizer
follows from the same identity and conjugation invariance of
its valuation. Each unramified local contribution at real
\(\sigma>1\) is positive. Pass to the nonzero limit at one.
\end{proof}

\begin{theorem}[The successive partial standard values]
\label{thm:values171}
For each \(w\notin S_E\) set \(Q_w=Nw\),
\(\mathcal A_w=\{Q_w^a:-d\le a\le d\}\) as a multiset,
and evaluate an unramified character at a uniformizer of \(w\).
Starting with any initial spherical row as in Lemma~\ref{lem:sat171},
form the successive parameter lists
\[
 B_{i,w}=(\chi_{S,i}/\chi_{T,i})(w)B_{i-1,w}
                       \ \sqcup\ \chi_{S,i}(w)\mathcal A_w .
\]
For \(0\le j<m\) and \(k=j+1\), the partial standard product of
\(B_j\otimes\chi_{T,k}^{-1}\) continues holomorphically to a
neighborhood of \(z_0\). Its value is
\begin{equation}
 L_j^S(z_0)=L^S(z_0,B_0\otimes\mu_k^{-1})
                  \prod_{i=1}^j\prod_{a=1}^t L_E^{S_E}(a,\rho_{ik})
                       \ne0 .                              \label{eq:L171}
\end{equation}
Suppose \(U_0=\max_{i<k} \log(2|D_E|N\mathfrak f_{\rho_{ik}})\),
with \(\log(2|D_E|)\) used for an empty set, and put
\(M_*=1+U_0+f\log(3+4mH_0)+\#S_E\). For large \(x\),
\[
 \sum_{j=0}^{m-1}\big|\log|L_j^S(z_0)|\big|
       \ll m^2 M_*+mfr\,2^{-(z_0-c_r)} .
\]
After dividing each value by its segment product on the right
of \eqref{eq:L171}, its logarithmic absolute value is
\(O(fr\,2^{-(z_0-c_r)})\), uniformly in the initial row and
in the splitting characters.
\end{theorem}
\begin{proof}
For \(\Re z\) large, the standard product of a multiset \(B\)
is \(\prod_{w\notin S_E}\prod_{\alpha\in B_w}
(1-\alpha Q_w^{-z})^{-1}\). Induction on \(j\) gives
\[
 B_{j,w}=D_j(w)\Big(B_{0,w}\sqcup
                     \bigsqcup_{i=1}^j\mu_i(w)\mathcal A_w\Big).
\]
Indeed the old terms gain \(D_j/D_{j-1}\), and the new
segment has scalar \(D_j\mu_j\). Division by \(\chi_{T,k}\)
cancels \(D_j\). Disjoint union multiplies Euler products,
and a monomial \(Q_w^a\) shifts the Hecke argument from \(z\)
to \(z-a\). This gives \eqref{eq:L171} by continuation, since
\(z_0-a\) runs exactly once through \(1,\ldots,t\).

For the initial factor put \(u=\Re z-c_r\).
For \(u\ge2\) the sum of the absolute logarithms of its local
factors is at most \(2fr\log\zeta(u)\ll fr\,2^{-u}\),
by Lemma~\ref{lem:sat171}. Comparison uses at most \(2f\)
primes over each rational prime and \(Nw\ge p\).
The choice \(t\ge x^{377}\) and the row rank of
Lemma~\ref{lem:finite167} give \(z_0-c_r>3\) eventually.
The initial product is thus holomorphic and nonzero there.
All shifted factors are entire, and Lemma~\ref{lem:hecke171}
applies at their positive arguments. Summing that lemma over
the \(m(m-1)/2\) pairs proves the stated bounds. The
recursion and location of the special value can also be
read in \srcloc{Development/AbelianVarietyTheory/Definitions/LatticeFrames/ChainParSucc}{121}
and line 128 there.
\end{proof}

\begin{corollary}[Characters of controlled conductor on the common fields]
\label{cor:chars171}
The \(\mu_i\) with the required real exponents can be chosen
unramified away from the existing split auxiliary primes,
with conductor exponent at most one there. With this choice
the logarithm of each ratio conductor is at most
\(2f\log(g_1g_2T)\). Multiplying all the \(\mu_i\) by the
same finite-order character preserves all segment factors
for a fixed omitted set. For omitted primes above
\(2Ng_1g_2T\), where \(E\subset\QQ(\zeta_N)\), the sum
of logarithmic absolute values in Theorem~\ref{thm:values171}
is bounded by a polynomial in \(x\) on the common intervals.
\end{corollary}
\begin{proof}
Let \(A=\ker(N_{E/F}:\mathrm{Res}_{E/F}\mathbb G_m
\to\mathbb G_m)\). In \(A(\mathbb A)\) prescribe the character
of \(A(F_\infty)\) with weights \(-e_{i,\tau}\).
The elements of \(A(F)\) integral at all finite places form
the roots of unity in \(E\). They are units of absolute value
one at all embeddings, so Kronecker's theorem applies.
At any chosen \(v\mid g_1\), reduction on these roots of
unity is injective. Indeed \(E\) splits completely at the
odd prime \(g_1\), and \(1+g_1\ZZ_{g_1}\) is torsion free.
Choose a character of the cyclic residue-unit group at \(v\)
canceling the prescribed character on these roots. Such a
character exists by extending from the cyclic subgroup.
One may first prescribe other residue characters at split
auxiliary places and cancel their product in the same way.
Take the trivial character on all remaining finite unit
groups. The resulting character on
\(A(F_\infty)\prod_{v<\infty}A(\OO_v)\) kills its intersection
with \(A(F)\). It descends to an open subgroup of
\(A(F)\backslash A(\mathbb A)\) and extends to a unitary
character there. For example, the circle is divisible, so
an abelian character extends algebraically, and continuity
on the open subgroup gives continuity throughout. Pull back
under \(z\mapsto z/\bar z\). This doubles the angular
exponents, kills \(\mathbb A_F^\times\), and gives the
asserted finite conductors. The same conductor bound applies
to each quotient.

For completeness one can choose \(\eta\), whose conductor
cancels from the formula, on the same fields. On
\(\QQ(\sqrt{-d_0})\) take the angular character \(z/|z|\)
and the quadratic residue character on the units at
\(\sqrt{-d_0}\), trivial on the other finite unit groups.
The roots of unity are \(\{\pm1\}\) and \(d_0\equiv3\bmod4\),
so the product on them is one. Extension from the subgroup
of units and the infinite component, as above, gives a
unitary Hecke character. Its restriction to the rational
ideles is the quadratic character of this imaginary quadratic
field. They agree at infinity and on all finite units, and
\(\QQ\) has class number one. Composing with the idele norm
from \(E\) gives \(\eta\) with the desired exponents on the
selected CM type. On \(\mathbb A_F^\times\) it restricts to
\(\omega_{E/F}\) by quadratic base change. Its possible finite
ramification lies over \(d_0\), which divides \(N\).

Finally \(\#S_E\ll f(1+\log(2Ng_1g_2T))\). The
conductor--discriminant formula gives
\(\log|D_E|\le 2f\log N\). The field degrees, \(m,H_0\)
and the logarithms of all the indicated rational primes and
of \(N\) are polynomially bounded in the construction of
Volume~I. Substitution proves the assertion.
\end{proof}

The theorem concerns the standard parameters propagated by
the displayed recursion. At a step where a spherical theta
lift has these parameters, it computes its partial scalar.
The nonzero local integrals for a prescribed Schwartz vector
and the integral normalization of a global theta lift are
additional requirements. They occur separately in
\srcloc{Development/AbelianVarietyTheory/Definitions/LatticeFrames/ChainParSucc}{197}
and line 215 there, which impose separate conditions on the
local integrals and the resulting global lift.

\subsection{A common trace pairing and the finite theta tests}\label{nav:VIII-04}
\label{sec:lattices173}

Consider the actual fields \(F=E^+\) and
\(E=B_J(\sqrt{-q},\sqrt{-d_0})\) of Theorem \ref{thm:invar148}.
Here \(d_0\equiv7\bmod8\), \(C_J=\QQ(\zeta_{2^{J+2}})^+
\subset F\), and \(J\ge2\) on sufficiently distant intervals.
Choose the CM embeddings that are positive on
\(\sqrt{-d_0}/i\). In this section \(\delta_0=\sqrt{-d_0}\),
\(\mathfrak d_F\) denotes the absolute different of \(F\),
and \(\omega=\omega_{E/F}\) denotes the quadratic character.
At finite places \(E/F\) is unramified, including the split
case, by Theorem \ref{thm:norm91}. In particular \(\omega\)
is a character of the narrow ideal class group of \(F\).
All dyadic places split since \(2\) splits in \(\QQ(\sqrt{-d_0})\).

\begin{lemma}[A narrow norm class]
\leavevmode\label{lem:norm173}
There is an integral ideal \(\mathfrak a_0\) of \(F\) with
\(\mathfrak a_0\OO_E=(\delta_0)\) and
\(\omega(\mathfrak d_F\mathfrak a_0)=1\).
An integral ideal \(\mathfrak b\) of \(E\) and a totally
positive \(\alpha\in F^\times\) can be chosen with
\begin{equation}
   (\alpha)\,\mathfrak d_F\mathfrak a_0
                 = N_{E/F}\mathfrak b,\qquad
   \log N_{E/\QQ}\mathfrak b \le \tfrac12\log|D_E|+
                                  f\log(4/\pi).             \label{eq:norm173}
\end{equation}
\end{lemma}
\begin{proof}
The valuations of \(\delta_0\) agree at conjugate primes and
the relative ramification index is one everywhere. They
therefore define an integral ideal downstairs, supported
above \(d_0\). All its valuations at primes above \(d_0\)
are equal since \(E/\QQ\) is abelian and every conjugate
of \(\delta_0\) is \(\pm\delta_0\). The primes of \(F\)
over \(d_0\) also have the same value under \(\omega\).
There are evenly many of them: \(d_0\equiv7\bmod8\) splits
in \(\QQ(\sqrt2)\subset F\), and the two primes there have
the same number of prolongations by Galois symmetry. Hence
\(\omega(\mathfrak a_0)=1\).

The global root number of \(\omega\) is one, by the
factorization of the completed Dedekind zeta function of \(E\)
as the completed zeta function of \(F\) times the completed
\(L\)-function of \(\omega\). Both zeta functional equations
have sign one. For the global trace additive character, the
local root number of an unramified quadratic character at
\(v\) is \(\omega_v(\varpi_v)^{v(\mathfrak d_F)}\).
Indeed the local constant for unit conductor additive character
is one, and scaling that character by a uniformizer multiplies
the root number by its character value. Reversing the
conductor convention inverts this sign and gives the same
number. At each real place the character is the sign
character, with root number \(i\) or \(-i\) according to the
single global sign convention. The number of real places is
divisible by four since \(C_J\subset F\) with \(J\ge2\).
Multiplying the local constants gives \(\omega(\mathfrak d_F)=1\).
These local constants follow from the unramified and real
local functional equations in
\href{https://sites.math.rutgers.edu/~alexk/2023S572/Tate1950.pdf}
{Tate, \emph{Fourier analysis in number fields and Hecke's
zeta-functions}, Chapters II and IV}.

Let \(H^+\) be the narrow Hilbert class field of \(F\).
The unramified-at-finite-primes abelian extension \(E/F\)
is contained in \(H^+\). By Artin reciprocity and Chebotarev
there is a prime in the narrow class of
\(\mathfrak d_F\mathfrak a_0\). Its Frobenius in \(E/F\)
is trivial by the first two paragraphs, so either prime of
\(E\) above it has ideal norm that prime. Replace one such
prime upstairs by an integral ideal \(\mathfrak b\) in its
ordinary ideal class satisfying Minkowski's bound in degree
\(2f\):
\[
 N_{E/\QQ}\mathfrak b\ \le\
  (4/\pi)^f\,\frac{(2f)!}{(2f)^{2f}}\,\sqrt{|D_E|}.
\]
The norm of a principal ideal from \(E\) has a totally
positive generator in \(F\), so this replacement preserves
the narrow class of the norm. Equality of narrow classes
now gives \(\alpha\) and the asserted size estimate.
\end{proof}

\begin{theorem}[Simultaneous self-dual lattices at the inert places]
\label{thm:lattices173}
Put \(\delta=\alpha\delta_0\) using Lemma~\ref{lem:norm173}.
Use this one trace-zero scalar to make the Hermitian tensor
pairing alternating at every step, with the global additive
character obtained by trace from \(\QQ\). At each inert finite
place \(v\), let \(L_{j,v}\) be the tensor of the two standard
unimodular local Hermitian lattices at step \(j\to j+1\).
There are homothetic lattices \(L_{j,v}\subseteq L'_{j,v}\)
which are self-dual for the additive pairing and are preserved
by both of the original local compact stabilizers. They are
equal to \(L_{j,v}\) away from finitely many \(v\), and
\begin{equation}
 \sum_{j=0}^{m-1}\ \sum_{v\ {\rm inert}}
        \log[L'_{j,v}:L_{j,v}]
   \ \le\ \Big(\sum_{j=0}^{m-1} n_jn_{j+1}\Big)
                     \log N_{E/\QQ}\mathfrak b.              \label{eq:index173}
\end{equation}
The right side is polynomially bounded in \(x\) throughout
the intervals of Section~\ref{sec:values171}. All the real
signatures and the oriented real additive pairings of
Section~\ref{sec:realtower169} can be used at the same time.
\end{theorem}
\begin{proof}
Fix \(v\) with prime \(w\) above it, and put \(a=v_w(\mathfrak b)\).
Equation~\eqref{eq:norm173} gives
\(v_w(\delta)+v(\mathfrak d_F)=2a\) since the residue degree
of \(E_w/F_v\) is two. Write \(\Lambda=L_{j,v}\).
The pairing has the form
\[
 (x,y)\longmapsto
 \psi_{\QQ,p}\bigl(\operatorname{Tr}_{F_v/\QQ_p}
      \operatorname{Tr}_{E_w/F_v}(\delta\,h(x,y))\bigr),
\]
up to a sign or a factor \(1/2\), both units here. The
Hermitian tensor form \(h\) is unimodular. The annihilator
of \(\OO_{E_w}\) for the absolute trace character is its
inverse different, namely
\(\mathfrak d_F^{-1}\OO_{E_w}\), because \(E_w/F_v\)
is unramified. In a lattice basis inversion of the
unimodular Gram matrix thus gives
\(\Lambda^\#=\varpi_v^{-2a}\Lambda\). Taking
\(L'_{j,v}=\varpi_v^{-a}\Lambda\) makes it self-dual.
Both compact stabilizers preserve every such homothety.
The tensor rank over \(E_w\) is \(n_jn_{j+1}\), and
\(Nw=(Nv)^2\). Its index is \((Nw)^{a n_jn_{j+1}}\).
Multiplication over inert primes and then the steps gives
\eqref{eq:index173}, because retaining all primes in the
norm of the integral ideal \(\mathfrak b\) only increases
the bound. Polynomial growth follows from
\(\log|D_E|\le2f\log N\) and the rank bounds.
At each chosen embedding \(\alpha\) is positive. The new
trace-zero scalar consequently has the prescribed imaginary
sign there. Any change of magnitude is absorbed by a positive
real rescaling in the complex vector space. This verifies
the assertion about the real pairs.
\end{proof}

At an inert place the self-dual lattice also determines a
compact vector in the lattice model. To see the local
normalization explicitly, take the Heisenberg group of the
symplectic tensor space and induce the character of the
subgroup above \(L'_{j,v}\) that is trivial on the lattice
and is the given additive character on the center.
Integrality of the pairing and odd residue characteristic make
this a character with the usual half-pairing in the group
law. Self-duality gives the irreducible Heisenberg model.
Every isometry preserving the lattice acts on this model by
pullback, and fixes the function supported on the subgroup and
equal there to its inducing character. This describes the
vector for the lattice splitting over the compact stabilizer.
Its use with a globally specified Weil splitting requires
the agreement of the two compact characters.

\begin{proposition}[Removing the finite inert ramification of the splitting characters]
\label{prop:splittings173}
The characters \(\mu_i,\eta\) of Section~\ref{sec:values171}
can be chosen with the stated archimedean types and restrictions
to \(\mathbb A_F^\times\), unramified at every nonsplit finite
place. In fact \(\eta\) can have finite ramification confined
to the existing split prime \(g_1\) with conductor exponent at
most one.
\end{proposition}
\begin{proof}
Use the \(\mu_i\) from Corollary~\ref{cor:chars171}.
Start with the \(\eta\) constructed at the end of its proof.
Its finite unit characters have conductor exponent at most one
at places above \(d_0\): local norms carry principal units
into principal units, and the original character at the
quadratic prime is a residue character. At all other finite
places it is unramified. At a place \(v\mid d_0\) its product
on \(\OO_{E,v}^{\times}\) is trivial on the diagonal
\(\OO_{F,v}^{\times}\), since \(\omega_{E/F}\) is unramified.
Over the residue fields the quotient by these scalar units
identifies via \(z\mapsto z/\bar z\) with the group of
norm-one elements. For an inert place this follows by
counting in the cyclic group of the quadratic residue
field. At a split place it is the ratio of the two coordinates.
There is therefore a character on the local norm-one residue
group whose pullback cancels the finite unit character.
Prescribe these characters on the norm-one torus at all
places over \(d_0\) and the trivial character at infinity.
Use one place above \(g_1\) to cancel their product on the
global roots of unity, exactly as in
Corollary~\ref{cor:chars171}. Extend and pull back along
\(z\mapsto z/\bar z\). Multiplying \(\eta\) by this Hecke
character preserves its infinite type and restriction to \(F\),
removes the ramification above \(d_0\), and gives the
claimed support and conductor.
\end{proof}

\begin{proposition}[A common depth for the split local tests]
\label{prop:split173}
At a split finite place put \(q_v=Nv\) and let
\(K_n(a)=1+\varpi_v^a M_n(\OO_v)\), \(a\ge1\).
For any consecutive ranks \(n\le N'\), use the matrix
Schr\"odinger model there and suppose both determinant
splitting characters are trivial on \(1+\varpi_v^a\OO_v\).
There is an integral indicator fixed by \(K_{N'}(a)\) such
that its diagonal local doubling integral at the theta point,
on every unit \(K_n(a)\)-fixed source vector, equals
\begin{equation}
  q_v^{-a(N'n+n^2)}
                 \prod_{\ell=1}^n(1-q_v^{-\ell})^{-1}.     \label{eq:split173}
\end{equation}
Here the integral matrix lattice has additive volume one and
\(\mathrm{GL}_n(\OO_v)\) has multiplicative volume one.
\end{proposition}
\begin{proof}
For \(X\in M_{N',n}(F_v)\) take the indicator of
\[
       X\equiv\begin{pmatrix}1_n\\0\end{pmatrix}
                         \pmod{\varpi_v^a M_{N',n}(\OO_v)} .
\]
Both congruence groups preserve it in the split model
\eqref{eq:weil161}. The determinant modulus on each is one,
and the character there is trivial by hypothesis. For the
inner product with a source translate to be nonzero, the
two top square blocks must lie in \(K_n(a)\); their quotient
then gives \(g\in K_n(a)\). Conversely inside \(K_n(a)\) the
two indicators coincide. The oscillator coefficient is thus
\(q_v^{-aN'n}\mathbf1_{K_n(a)}\). The source matrix
coefficient is one on that subgroup. Finally reduction
modulo \(\varpi_v^a\) has image of cardinality
\(q_v^{an^2}\prod_{\ell=1}^n(1-q_v^{-\ell})\).
Integration gives \eqref{eq:split173}. Since the output
test is fixed by \(K_{N'}(a)\), this construction can use
the same depth at the following step. It applies at dyadic
split places as well.
\end{proof}

The lattice index in \eqref{eq:index173} and the local
integral in \eqref{eq:split173} measure different parts of
a global comparison. The first bounds the finite cost of
making the trace lattice self-dual for a permissible global
scalar. The second evaluates a particular compactly supported
zeta integral. The action of global projection on the
designated arithmetic tensors and their periods requires a
separate estimate.

\subsection{Spherical parameters of a nonzero local lift}\label{nav:VIII-05}
\label{sec:transfer175}

Let \(E_v/F_v\) be unramified quadratic with odd residue
characteristic, \(Q=\#\kappa_{E_v}\), and take the two
unramified unitary groups of ranks \(n,N\). The groups have
their unimodular lattices and are quasi-split. Write
\(\chi_T,\chi_S\) for the splitting characters attached to
the spaces of dimensions \(N,n\), respectively. At an inert
place their respective values on the common uniformizer are
\((-1)^N,(-1)^n\). This follows by restriction to \(F_v^\times\).
For a spherical representation \(\pi\) let \(B(\pi)\) be
its standard base-change Satake multiset. It has a reciprocal
pair for each coordinate of the relative split torus and,
in odd dimension, one additional entry \(1\). For \(b\ge1\)
put \(A_b=\{Q^{(b-1)/2-i}:0\le i<b\}\) and put \(A_0=\varnothing\).

\begin{theorem}[Spherical transfer at an inert place]
\leavevmode\label{thm:inert175}
Suppose that a spherical \(\pi\) in rank \(n\) and a
spherical \(\sigma\) in rank \(N\) occur together as an
irreducible quotient of the local oscillator. If \(N\ge n\),
then
\begin{equation}
    \chi_S(\varpi)^{-1} B(\sigma)
       =\chi_T(\varpi)^{-1}B(\pi)\ \sqcup\ A_{N-n}.
                                                       \label{eq:inert175}
\end{equation}
This assertion allows non-tempered spherical representations.
\end{theorem}
\begin{proof}
A spherical representation embeds in a normalized principal
series with unramified torus character in its Satake orbit.
Indeed it is a subquotient of an unramified principal series,
has unramified cuspidal support, and the subrepresentation
theorem embeds it in a principal series with that support.
Composing the assumed quotient with such embeddings on both
sides and using normalized Frobenius reciprocity gives a
joint torus-character quotient of the normalized Jacquet
module of the oscillator for the two minimal parabolics.
We prove the corresponding multiset identity for any such
unramified character quotient. We work in the free abelian
group on \(\CC^\times\), extending \(A_b\) by
\(A_{-b}=-A_b\) for \(b>0\). Interchanging the two groups
then negates both sides of the desired difference formula.
The unramified splitting values and their inverses agree
at this place.

When both ranks are at most one there are no reciprocal
pairs. The middle entry in odd rank is \(1\). With ranks
\((0,1)\), the splitting value for the zero-dimensional
space is \(1\); with ranks \((1,1)\) both splitting values
are \(-1\). These facts verify the assertion in all base
cases. Now take \(N\ge n\) with \(N\ge2\) and take first
the Jacquet module on the \(N\)-side for the parabolic of
an isotropic line. Write \(u\) for the unramified character
on its \(\mathrm{GL}_1(E_v)\) factor, evaluated at \(\varpi\).
The two possible strata of Kudla's filtration have ranks
zero and one. Normalized Jacquet functors are exact and
transitive. A nonzero map from a finite filtration to a
character restricts nontrivially to some successive quotient,
by taking the last filtration subspace with nonzero image.
We spell out both possibilities.
\href{https://arxiv.org/pdf/1602.01299}{Atobe and Gan, Lemma 5.1,
with \(k=1\) and \(\epsilon_0=0\)}.

In rank zero the line carries
\(\chi_S|\cdot|_{E_v}^{(n-N+1)/2}\), and the remaining
oscillator has ranks \(N-2,n\). Set \(b=N-n\). After
removing \(\chi_S(\varpi)\) the two entries supplied by
the line are \(Q^{(b-1)/2},Q^{-(b-1)/2}\). In the free
abelian group one has
\[
 A_b-A_{b-2}=[Q^{(b-1)/2}]+[Q^{-(b-1)/2}]
 \qquad(b\ge0).
\]
This identity also covers \(b=0,1\) by the signed definition.
The inductive formula for the remaining ranks, interchanged
if necessary, now proves the assertion in rank zero.

In rank one the other group also has an isotropic line.
The quotient is normalized induction from its line
parabolic. Its inducing representation is the oscillator
in ranks \(N-2,n-2\) tensored with the functions of compact
support on the isomorphisms between the two conjugate
lines. For the two general linear factors this last space
is a copy of \(C_c^\infty(E_v^\times)\), with left and right
translation and the two splitting characters. The element
\((\varpi,\varpi)\) acts there by the scalar
\(\chi_S(\varpi)\chi_T(\varpi)\). The conjugate-linear
isomorphisms cause no change because \(\varpi\in F_v\).
Taking the further minimal Jacquet module on the \(n\)-side
uses the geometric lemma for parabolic induction. Its finite
Bruhat filtration places the line character in one Weyl
coordinate or its conjugate inverse. The other coordinates
belong to the minimal Jacquet module of the smaller group.
This description follows by taking the unipotent coinvariants
on the Bruhat cells of the line-parabolic flag variety.
Normalized induction and Jacquet functors cancel the modulus
on each cell. It applies to the smooth compactly supported
function representation on the line as well.

If \(v'\) is the value of the coordinate so placed (after
this possible inversion), then
\(uv'=\chi_S(\varpi)\chi_T(\varpi)\). Hence the two
normalized reciprocal pairs on opposite sides coincide
as multisets. Fixing any line function on which the character
map is nonzero for some remaining tensor gives a nonzero
character map for the two smaller minimal Jacquet modules.
The rank difference is still \(b\), so induction proves
the formula in rank one also. The parity of each group
stays the same whenever two dimensions are removed, and
the splitting values in every invocation of induction
are consequently the stated ones. This completes the proof.
\end{proof}

\begin{proposition}[The split pair and its two components]
\label{prop:split175}
At a split place with unramified data the same formula holds,
with \(Q=Nw\) on each component \(w\mid v\), in the unitary
tensor convention. Thus at either type of good place a nonzero
spherical theta pair has exactly the parameter update of
Theorem~\ref{thm:values171}.
\end{proposition}
\begin{proof}
Write \(w^c\) for the other component. The position
\(V_w\otimes W_w\) of the symplectic tensor can be viewed
as \(\operatorname{Hom}(W_{w^c},V_w)\), since the Hermitian
form identifies \(W_{w^c}\) with the dual of \(W_w\).
In the matrix splitting the right group is therefore the
source in its \(w^c\) coordinates, and the left group is
the target in its \(w\) coordinates. The normalized
matrix action has the square roots of the two determinant
Jacobians. The unitary splitting twists the two actions
in their \(w\) coordinates by
\(\chi_{T,w}\circ\det\) and \(\chi_{S,w}\circ\det\).
Equivalently, the change from the trivial split characters
is through the isomorphism
\(E_v^\times/F_v^\times\to E_v^1,\ z\mapsto z/\bar z\).
The determinant with coordinates \((a,a^{-1})\) lifts
to \((a,1)\). This gives precisely the indicated
component character. The same split unitary convention and
its reduction to the matrix correspondence are used in
\href{https://ems.press/content/serial-article-files/51452}
{Chen and Zou, \emph{Arthur's multiplicity formula for even
orthogonal and unitary groups}, Remark 6.4(2)}.

Remove both twists. A spherical source with list \(B\) in
its \(w\) coordinates then has list
\((B/\chi_{T,w}(\varpi))^{-1}\) in its \(w^c\) coordinates,
where inversion is termwise. For the normalized type II
matrix oscillator the larger representation has cuspidal
support given by the contragredient of the smaller one
together with the support of the trivial representation of
\(\mathrm{GL}_{N-n}\).
\href{https://arxiv.org/abs/0709.4384}{M\'inguez,
\emph{Correspondance de Howe explicite}, Theorem 1(2)}.
The second list is \(A_{N-n}\) by normalized induction
from the diagonal. The contragredient reverses the inversion
from \(w^c\). Restoring the target twist proves the formula
at \(w\). At \(w^c\) both lists and both component characters
invert, while \(A_{N-n}\) is invariant under inversion.
\end{proof}

\subsection{A nonzero global theta chain with a choice of tests}\label{nav:VIII-06}
\label{sec:global175}

There is another use of the narrow-class calculation of
Lemma~\ref{lem:norm173}. For an existence assertion over
\(\CC\) one can arrange that the additive character and the
trace-zero scalar are separately unramified at every inert
place. This choice may introduce split places at which the
Schwartz function is allowed to vary.

\begin{lemma}[Simultaneous unramified data off split places]
\label{lem:good175}
On each field of Section~\ref{sec:lattices173} is a global
additive character and a trace-zero scalar of the signs of
Section~\ref{sec:realtower169} such that at every inert finite
place the additive character has conductor \(\OO_v\) and
the scalar is a unit. The choices in this statement use at
most two additional prime ideals of \(F\) as representatives
of narrow ideal classes, both split in \(E\).
\end{lemma}
\begin{proof}
The proof of Lemma~\ref{lem:norm173} shows separately that
\(\omega(\mathfrak d_F)=\omega(\mathfrak a_0)=1\).
Choose distinct prime representatives \(\mathfrak p,\mathfrak q\)
of these two narrow classes outside the existing finite levels,
by Chebotarev in the narrow Hilbert class field. Both split
in \(E\). There exist totally positive \(\beta,\gamma\) with
\[
        (\beta)\mathfrak d_F=\mathfrak p,\qquad
        (\gamma)\mathfrak a_0=\mathfrak q.
\]
Use the negative trace character evaluated at \(\beta x\)
and take the skew scalar \(\gamma\sqrt{-d_0}\) with the sign
specified for the pair. The annihilator of \(\OO_v\) for
the additive character is
\(\beta^{-1}\mathfrak d_{F,v}^{-1}\), which is \(\OO_v\)
off \(\mathfrak p\). The principal ideal of the skew scalar
is \(\mathfrak q\OO_E\). Positivity of \(\beta,\gamma\)
retains the real orientation. Positive real changes of
magnitude are absorbed by an equivariant scalar isometry
of the local tensor spaces. Together with
Proposition~\ref{prop:splittings173} this gives unramified
data at every inert finite place. The size estimate for
the original trace character is the separate ideal bound
\eqref{eq:index173}.
\end{proof}

\begin{theorem}[Cuspidal lifts to the compact endpoint]
\label{thm:global175}
On a field and interval of Theorem \ref{thm:invar148}, let
\(\Pi_0\) be an irreducible cuspidal automorphic representation
of the row group with the real parameters \(P_0\) and a
hyperspecial fixed vector at every inert finite place.
Use the splitting characters of Proposition~\ref{prop:splittings173}
and the global data of Lemma~\ref{lem:good175}. With
Schwartz functions free at a finite set of split places and
at infinity, the successive global theta lifts through
\(U(Y_m)\) are all nonzero and cuspidal. Their real parameters
are \(P_j\), and their standard parameters at the good finite
places are precisely \(B_j\) of Theorem~\ref{thm:values171}.
In particular this applies to the cuspidal constituents of a
nonzero row \(S_{\boldsymbol j}\) at the levels of
Section~\ref{sec:global167}. It also gives a nonzero complex
iteration on any nonzero compact-finite vector in such an
irreducible constituent fixed by the good hyperspecial groups.
\end{theorem}
\begin{proof}
In the quotient convention for theta lifting, the kernel
is integrated against the complex conjugate of the input
cusp form. A local nonzero pairing then gives the source
representation and its forward lift as a quotient of the
oscillator. Choose a finite set \(S\) of split places
including the two primes of Lemma~\ref{lem:good175}, the
dyadic places, the auxiliary levels and any split places
where \(\Pi_0\) is ramified. Outside \(S\), all spaces have
unimodular lattices, the finite characters are unramified
and the additive character has integral conductor. The
standard unramified pair therefore has its lattice vacuum,
fixed by both hyperspecial subgroups. In a self-dual lattice
model this is the vector supported on the lattice. The
unramified Weil splitting on the hyperspecial groups is the
lattice splitting, equivalently the usual spherical data in
the local doubling integral. Fix these vacua outside \(S\).

Suppose \(\Pi_j\) has been obtained with the indicated real
components and good finite parameters. Write \(n=n_j\),
\(N=n+t\), \(s_0=t/2\). Theorem~\ref{thm:tower169}
excludes a nonzero local lift to any smaller member of
the target Witt tower at every real place. A nonzero
global theta integral to a smaller member would give such
a local pairing by fixing all other test vectors. Thus all
the smaller global lifts vanish. The constant-term tower
property then makes the lift to \(Y_{j+1}\) cuspidal if it
is nonzero. Moreover the rational Witt index \(R\) of
\(Y_{j+1}\) is at most \(s_{j+1}\), the smaller real
signature, so \(N-R\ge n+t-s_{j+1}>n\). This is Weil's
convergent range for the doubled unitary pair at every step.

We apply the convergent-case Rallis inner product formula.
For a pure cusp vector and a pure Schwartz function the
squared Petersson norm is a nonzero Haar constant times
the value at \(s_0\) of the global doubling integral for
the associated Siegel--Weil section. For \(N>2n\) this
is the convergent formula used in Section~\ref{sec:global167}.
For \(n<N\le2n\) it is
\href{https://arxiv.org/pdf/1207.4709}{Gan, Qiu and Takeda,
Propositions 11.1(ii) and 11.2}. The latter formula uses
meromorphic continuation from large \(\Re s\). We verify
below that the unfolded product is holomorphic and nonzero
at the required point, which also identifies its constant
Laurent coefficient there.

At a good finite place take the spherical input and the fixed
vacuum with norms one. The partial product of the local
integrals is
\[
   \frac{L^S(s+1/2,B_j\otimes\chi_{T,j+1}^{-1})}
        {d_n^S(s)} .
\]
Here \(d_n^S(s)=\prod_{i=1}^{n}
L_F^S(2s+i,\omega_{E/F}^{\,t+i})\) is the unramified
doubling denominator. This evaluation is the
unramified doubling calculation of
\href{https://doi.org/10.1515/crll.1992.428.177}
{J.-S. Li, \emph{Non-vanishing theorems for the cohomology of
certain arithmetic quotients}, Theorem 3.1}. The unitary
normalization is displayed in
\href{https://arxiv.org/pdf/2002.09148}{Ichino, Section 6.3},
and the meromorphic integral is treated in
\href{https://arxiv.org/pdf/1207.4709}{Gan, Qiu and Takeda,
Section 11.5}. In the norm pairing the source coefficient
is conjugated. One may equivalently use the dual list
twisted by \(\chi_T\).
Conjugation on \(E/F\) exchanges this with the list displayed
above. Indeed \(B_{j,w^c}=B_{j,w}^{-1}\) at split places,
the inert list is reciprocal, and
\(\chi_T(z)\chi_T(\bar z)=1\) because its restriction to
\(\mathbb A_F^\times\) is trivial on norms. The two partial
products agree since all places above \(S\) are omitted.
The calculation initially holds for \(\Re s\) large.
Theorem~\ref{thm:values171} continues the numerator to
\(s_0\) with nonzero value. The denominator has an
absolutely convergent Euler logarithm near \(s_0\), since
\(t+1>1\). Its logarithmic absolute value there at \(s_0\)
is \(O(f\,2^{-t})\) by summing
\(f\log\zeta(t+i)\) over \(1\le i\le n\).

At \(v\in S\) take a unit smooth vector fixed by
\(K_n(a_v)\), with \(a_v\ge1\) also large enough for the
splitting characters. Use Proposition~\ref{prop:split173}.
The resulting local integral is strictly positive. It is
holomorphic in a neighborhood of \(s_0\). Indeed the
restriction of a flat extension of the Siegel--Weil section
at \((g,1)\) equals the section at \(s_0\) times its
positive Iwasawa height to the power \(s-s_0\).
The local oscillator coefficient has compact support there.
The multiplying height is bounded above and below on that
compact set, so integration gives an entire function for
this local test. Each place in \(S\) is split, where the
matrix argument works for every rank \(N\ge n\).

At a real place the input is the discrete series of
Theorem~\ref{thm:tower169}, including a compact representation
when the source has compact real points. Its local theta
lift to the specified target is nonzero. For such a
discrete series and \(N\ge n\), the local Rallis integral
is a nonzero sesquilinear pairing, and the flat section
integrals converge in an open neighborhood of \(s_0>0\).
\href{https://arxiv.org/pdf/2002.09148}{Ichino, Lemma 7.1
and Proposition 7.2}. By polarization in each of the two
vectors, some diagonal choice of a Schwartz vector and a
source vector has nonzero integral. One may take compact-finite
vectors by density and continuity, since the convergent local
zeta integral is a continuous functional on smooth vectors.
For a compact source the same holomorphy follows immediately
by compact integration.

For these factorizable data the Eulerian unfolding of the
doubling integral holds first in its half-plane of absolute
convergence. Cuspidality annuls the non-open orbits there.
The finite product of the integrals at \(S\) and infinity,
times the partial product already computed, is its meromorphic
continuation. Every factor is holomorphic and nonzero at
\(s_0\). The norm formula therefore gives a nonzero cusp
form. The nonzero cuspidal theta lift of an irreducible
cuspidal representation is irreducible and is the restricted
tensor product of the local small lifts. In this setting this
is the cuspidal global Howe correspondence as recalled in
\href{https://arxiv.org/pdf/2002.09148}{Ichino, Section 6.3}.
Our nonzero vector is fixed by the good target hyperspecial
subgroups. By Theorem~\ref{thm:tower169} its archimedean
type is \(P_{j+1}\), and Theorem~\ref{thm:inert175} and
Proposition~\ref{prop:split175} give exactly the required
parameter recursion at every good finite place. This closes
the induction from \(\Pi_0\). The hypotheses on a surviving
row constituent follow from Section~\ref{sec:global167}
and its hyperspecial levels at the inert places.

Finally fix a nonzero compact-finite vector in an obtained
\(\Pi_j\), spherical outside \(S\). The common kernel of
all lifts using the fixed vacua outside \(S\) and variable
tests inside \(S\) and at infinity is stable under the
finite groups in \(S\), the real maximal compact groups
and the complexified real Lie algebras. This follows by
translating or differentiating the Schwartz function, using
cuspidal decay in the theta integral. On the subspace
spherical outside \(S\) the resulting smooth, compact-finite
module is irreducible: it is the tensor of its irreducible
local factors at the indicated places with one-dimensional
spherical lines elsewhere. The common kernel is proper by
the nonzero norm just proved and hence is zero. Some test
therefore gives a nonzero output on the fixed vector. Its
output is compact-finite for compact-finite tests and again
spherical outside \(S\). Repetition proves the assertion
about a nonzero composition.
\end{proof}

The source distinguishes the standard \(L\)-value from
the nonzero upward integral in
\srcloc{Development/AbelianVarietyTheory/Definitions/LatticeFrames/ChainParSucc}{215}.
Theorem~\ref{thm:global175} obtains an actual nonzero complex
iteration with its stated choices. The estimate for the
arithmetic periods calls for particular rational tensors,
a prescribed projection and quantitative norms on the
integral lattices. It must also handle each required word
at the original finite level.

\subsection{One iterated theta operator on the prescribed row}\label{nav:VIII-07}
\label{sec:common177}

Fix a field, a translate, a signed real row and the splitting
data of Theorem~\ref{thm:global175}. Fix one unit compact-finite
vector \(u_{0,\tau}\) at each real place, and set
\(u_0=\bigotimes_\tau u_{0,\tau}\). Choose the finite exceptional
set \(S\) there to include all places \(v\mid T\). Let
\(\mathbf j=(j_v)_{v\mid T}\) be the terminal depths of the
initial Bessel periodization in Sections \ref{nav:VII-12}--\ref{nav:VII-13}, with
the first-edge functions \(\phi_{1,j_v}\) and the indices
\(I_{j_v}\) of Section \ref{nav:VII-13}. Put \(I=\prod_{v\mid T}I_{j_v}\).
The character \(\kappa_v\) on the terminal row compact
subgroup is fixed. On each \(v\in S\) choose a depth
\(a_v\ge1\) on which all determinant characters at that
place are trivial. At \(v\mid T\) also require \(a_v\ge j_v\)
and triviality of the top-block character
\(\kappa_v/\zeta_v\) on \(K_r(a_v)\), in the matrix frame
of Section \ref{nav:VII-13}. Here \(\zeta_v\) is the scalar splitting
character in that frame. At a nonterminal place one may increase \(a_v\) so that
its principal subgroup is contained in the prescribed source level.

Let \(\mathcal H\) be the space of cusp vectors on the row
group with real factor in the specified line, spherical
outside \(S\), of the prescribed fine type at \(v\mid T\),
and fixed by \(K_r(a_v)\) at the other places of \(S\).
The real line is taken in each copy of the fixed irreducible
archimedean representation in the cuspidal spectrum.
Define
\[
 \mathcal J_v=
 \begin{cases}
    \{1,\dots,m-1\} & v\mid T,\\
    \{0,\dots,m-1\} & v\in S,\ v\nmid T ,
 \end{cases}
 \qquad
 P(\mathbf a)=\prod_{v\in S}(Nv)^{
       -a_v\sum_{j\in\mathcal J_v}n_j(n_{j+1}+n_j)} .
\]
The missing first edge in \(\mathcal J_v\) at \(T\) has its
own normalization \(I_{j_v}^{-2}\).

\begin{lemma}[Real monomials on prescribed successive lines]
\label{lem:monomials177}
Suppose that at least one cuspidal constituent of the
specified row is present. There are monomial Fock tests
\(\phi_{j,\tau}\) and unit compact-finite vectors
\(u_{j+1,\tau}\) at all real places and steps such that
the following holds. Contraction of \(\phi_{j,\tau}\)
against \(u_{j,\tau}\) in the local theta correspondence
spans the line of \(u_{j+1,\tau}\). The diagonal real Rallis
integral for these data is nonzero. Both assertions use
only the fixed representations and real tests.
\end{lemma}
\begin{proof}
Consider a single step and place, with source vector \(u\).
Denote the local Rallis pairing by
\(Z(\phi,u;\phi',u')\), in the convention in which its
diagonal gives the local factor of the theta norm.
Ichino's nonvanishing criterion, applied as in
Theorem~\ref{thm:global175}, makes this four-vector pairing
nonzero. If \(Z(\phi,u;\phi,u)=0\) for every \(\phi\), then
polarization in the test would give zero for
\(Z(\phi,u;\phi',u)\) for all \(\phi,\phi'\).
Left and right invariance of Haar measure in the local
integral allow a translate of \(u\) in either slot to be
moved inversely to that slot's Schwartz test. Consequently
all pairings would vanish on the two spans of translates of
\(u\), and then on their smooth closures. The compact-finite
module generated by a nonzero \(u\) is irreducible. One
may also apply Lie derivatives in this argument, by
continuity of the convergent local integral. This would
contradict the nonzero pairing.

For this fixed \(u\) the diagonal in \(\phi\) has a constant
sign after multiplication by a fixed nonzero scalar. To see
this, use one reference cuspidal constituent from
Theorem~\ref{thm:global175} at this step. Fix its other local
data with nonzero diagonal factors, including an arbitrary
nonzero diagonal choice at the other real places. Inserting
the variable test at the present place expresses a global
squared norm as a fixed nonzero scalar times
\(Z(\phi,u;\phi,u)\). Thus the resulting local sesquilinear
form is positive semidefinite. If it vanished on each
monomial diagonal, Cauchy--Schwarz would kill every pair
of monomials. Polynomial Fock vectors are dense in the
smooth oscillator, so the whole form would vanish, a
contradiction. Choose a monomial with nonzero diagonal.
The same global norm computation shows that its contraction
in the local small lift is nonzero. Local Howe multiplicity
one makes the line of that contraction independent of the
choice of a nonzero quotient map. Choose a unit vector on it
and repeat. The reference constituent has every successive
lift by Theorem~\ref{thm:global175}, so the argument applies
at every step.
\end{proof}

\begin{theorem}[Norm of the full complex iteration]
\label{thm:norm177}
There are common choices of the real monomials, independent
of \(\mathbf j,\mathbf a\), and a positive constant \(C_x\)
depending on the fixed field, row, set \(S\), splittings,
real lines and measure normalizations. Use normal lattice
vacua outside \(S\), the fine test \(\phi_{1,j_v}\) on the
first edge at each \(v\mid T\), and on every other edge at
each exceptional place the matrix indicator of
Proposition~\ref{prop:split173} at depth \(a_v\). Write
\(\mathcal T_{\mathbf a,\mathbf j}\) for their composed
global theta operator. For every \(A\in\mathcal H\),
\begin{equation}
 e^{-\Delta_x} C_x I^{-2}P(\mathbf a)\|A\|^2
 \ \le\ \|\mathcal T_{\mathbf a,\mathbf j}A\|^2
 \ \le\ e^{\Delta_x} C_x I^{-2}P(\mathbf a)\|A\|^2,
 \qquad
 \Delta_x=Cmfr\,2^{-(z_0-c_r)}.                       \label{eq:common177}
\end{equation}
Here \(c_r=10(r+1)^3,\ z_0=(t+1)/2\), and \(C\) is
absolute. In particular the same operator detects every
nonzero vector in the fixed row space. After the displayed
level factors have been extracted \(C_x\) is independent
of both depth vectors and of the finite cuspidal constituent.
\end{theorem}
\begin{proof}
If there is no cuspidal representation of the specified
real row and good spherical type at any of the levels, the
assertion is vacuous. Otherwise fix a reference representation
of this type. Choose the real data by Lemma~\ref{lem:monomials177}.
At a terminal place the support and compact character of
the first test are preserved by the target principal group
\(K_{n_1}(a_v)\). Indeed the change of the integral position
matrix under left multiplication is zero modulo
\(\varpi_v^{a_v}\); its invertible top block thus stays in
the same coset modulo \(\varpi_v^{j_v}\) and in the same
top-block character. The lower block stays divisible by
\(\varpi_v\), and the target splitting character is trivial.
Away from these first terminal edges, the congruence test
of Proposition~\ref{prop:split173} preserves the target
depth \(a_v\) for the next lift. All good places preserve
their spherical lines.

First take the input in one irreducible cuspidal summand.
The archimedean line at the next step is the one constructed
in the lemma for every finite part of this summand. Indeed
the global theta map is equivariant in each pair of local
variables and factors locally through the unique irreducible
theta pairing. For a pure tensor this fixes the target
archimedean line, and finite sums give the same conclusion
for an arbitrary finite part. With unit vectors on the
lines, the real factors of the Rallis norm formula at
each step are therefore fixed scalars.

At \(v\mid T\) on the first edge the oscillator coefficient is
\(\|\phi_{1,j_v}\|^2\kappa_v\mathbf1_{L_{j_v}}\),
where \(L_{j_v}=(B_v\cap K_v)K_r(j_v)\) for positive
depth and \(L_0=K_v\). Integration against the source
coefficient projects onto the prescribed character there.
On that space the local factor is
\(\|\phi_{*,v}\|^2\operatorname{vol}(K_v)/I_{j_v}^2\).
One index comes from the test volume and one from the
group volume, as in Section \ref{nav:VII-13}. On the other exceptional
edges the same calculation with a principal subgroup acts
as a scalar on its fixed vectors. By
\eqref{eq:split173} their product is \(P(\mathbf a)\)
times a positive constant independent of all the depths.
These projection identities hold for arbitrary tensors of
fixed local type. They can be checked first for pure tensors,
then by taking the corresponding mixed pairings.

At step \(j\to j+1\), write \(k=j+1\). In the unramified
part of the continued norm formula from
Theorem~\ref{thm:global175} the denominator is fixed, as are
all factors \(L_E^{S_E}(a,\mu_i\mu_k^{-1})\) of
\eqref{eq:L171}. The only factor varying with the
constituent is \(L^S(z_0,B_0\otimes\mu_k^{-1})\),
whose logarithmic absolute value has size
\(O(fr\,2^{-(z_0-c_r)})\). This bound uses the initial
spherical parameters throughout. In particular its exponent
does not contain the growing rank \(n_j\). The Haar
constant in the global norm formula is fixed for each pair
of groups. Expressing it using the input Petersson norm
removes any scalar used to identify an automorphic copy
with its abstract unitary tensor product. Take the product
of the absolute values of all the common factors at all
steps, including the real factors and the terminal
depth-zero constants, to define \(C_x>0\).
Taking absolute values is legitimate since each complete
identity is a squared norm. The logarithmic error is
additive over the \(m\) applications, giving \eqref{eq:common177}
in one summand.

For two orthogonal invariant cuspidal summands the lifts
are orthogonal, also for independently chosen tests. In the
half-plane of absolute convergence the mixed doubling
integral unfolds against the matrix coefficient
\(\langle R(g)A_1,A_2\rangle=0\). Its continuation is
identically zero. The convergent-range Siegel--Weil
formula used in Theorem~\ref{thm:global175} consequently
makes the mixed theta inner product zero. This argument
also applies to equivalent representations in orthogonal
multiplicity spaces. Varying both tests gives orthogonality
of the invariant lifted spans, so it persists through the
iteration. At any fixed level the real type and infinitesimal
character in question occur in a finite-dimensional space
of cusp vectors of that level and compact type. Summing
the squared norms proves \eqref{eq:common177} on \(\mathcal H\).
\end{proof}

\begin{corollary}[The full iteration on the Bessel columns]
\label{cor:gram177}
Restrict the Bessel periodization \(S_{\mathbf j}\) to
the unit input line \(u_0\) and its finite column domain,
with depths as above. Write
\(B=\mathcal T_{\mathbf a,\mathbf j}S_{\mathbf j}\) on that
domain. Its kernel equals \(\ker S_{\mathbf j}\), including
at the original depth. For the exponent \(\eta=L-4fr^2>0\)
of Section \ref{nav:VII-12} and its field-dependent constant \(C'_x\),
\[
 \left\|\frac{I^3}{C_x P(\mathbf a)} B^*B-1\right\|_{\rm op}
       \le (e^{\Delta_x}-1)
           +e^{\Delta_x} C'_x T^{-\eta\sum_{v\mid T}j_v}.
\]
At sufficiently large total terminal depth for each fixed
large \(x\), this gives a nonzero full complex iteration
on every nonzero finite column.
\end{corollary}
\begin{proof}
The hypotheses of the theorem on the nonterminal source
groups hold after choosing \(a_v\) inside their given
levels. Apply \eqref{eq:common177} to \(S_{\mathbf j}\ell\).
Positivity gives the kernel assertion. The estimate
\(\|I S_{\mathbf j}^*S_{\mathbf j}-1\|_{\rm op}
\le C'_x T^{-\eta\sum_v j_v}\) from Sections \ref{nav:VII-12}--\ref{nav:VII-13}
continues to hold on the fixed unit real line by compression.
Use its upper and lower quadratic-form bounds in the preceding
inequality. If its lower bound becomes negative, positivity
of \(B^*B\) still gives the stated absolute error. Taking
the supremum over unit columns gives the formula. At
sufficient depth the Bessel bound alone gives injectivity,
and the kernel identity transfers it to the iteration.
\end{proof}

\begin{proposition}[Cost of the two auxiliary split primes]
\label{prop:cost177}
The prime representatives in Lemma~\ref{lem:good175} can be
chosen above distinct rational primes outside \(2Ng_1g_2T\)
so that, on the intervals of Section \ref{nav:I-28},
\[
      \log N\mathfrak p+\log N\mathfrak q\ll_{\psi,\epsilon} x^c
\]
for a constant \(c=c(\psi,\epsilon)\). For the prescribed row
with its original exceptional places and these two primes,
put \(A=\max(1,\max_{v\in S}a_v)\). Then
\[
 0\le -\log P(\mathbf a)+2\log I\ll_{\psi,\epsilon} A x^{c'}
\]
for a constant \(c'\) depending only on the fixed family.
In particular, for depths bounded by a fixed power of \(x\),
this explicit logarithmic loss is \(o(b/x^B)\) for every fixed
\(B>0\), where \(b=\exp(x^\epsilon)\).
\end{proposition}
\begin{proof}
Let \(H^+\) denote the narrow Hilbert class field of \(F\).
Write \(D_F\) for the absolute discriminant and \(f=[F:\QQ]\).
We first find a prime representing any given narrow class
while excluding an arbitrary finite set \(U\) of rational
primes. For \(p\in U\) take \(k_p\) least positive with
\(\varphi(p^{k_p})>f\), and put \(M=\prod_{p\in U}p^{k_p}\).
If \(k_p>1\) then \(p^{k_p-1}(1-1/p)\le f\), hence
\(p^{k_p}\le 2fp\). This bound holds also for \(k_p=1\).
Consequently
\[
       \log M\le\sum_{p\in U}\log p+|U|\log(2f).
\]
In the abelian extension \(L=H^+F(\zeta_M)\) of \(F\), every
place above \(p\in U\) ramifies. In fact its completion
contains \(\QQ_p(\zeta_{p^{k_p}})\), whose absolute ramification
index \(\varphi(p^{k_p})\) exceeds the absolute ramification
index of any completion of \(F\) at \(p\).

The norm of the finite conductor of every character of
\(\operatorname{Gal}(L/F)\) is at most \(M^f\). To see this,
embed this Galois group by restriction in the product of
\(\operatorname{Gal}(H^+/F)\) and
\(\operatorname{Gal}(\QQ(\zeta_M)/\QQ)\). A character of a
subgroup of a finite abelian group extends to the whole
group. Adjoining one generator at a time requires only
taking a root of a complex number of absolute value one.
The first component is unramified at finite places. By
functoriality of Artin reciprocity the second component on
\(F\) is the pullback of a Dirichlet character by the norm.
For \(v\mid p\) and \(y\in\OO_v\) the norm of
\(1+p^{k_p}y\) is \(1\) modulo \(p^{k_p}\), as follows from
the determinant of multiplication on the free
\(\ZZ_p\)-module \(\OO_v\). Away from \(M\) this character is
unramified. Its finite conductor thus divides \(M\OO_F\).

Lift the chosen narrow class to an element of
\(\operatorname{Gal}(L/F)\) via the surjective restriction
to \(H^+\), identifying classes with their Frobenius elements.
The unconditional least-prime theorem of
\href{https://arxiv.org/pdf/1604.01750}{Thorner and Zaman,
Theorem 1.1} gives an unramified prime \(v\) of this Frobenius
with rational prime norm. Indeed, for an abelian extension
one takes the full Galois group for the abelian subgroup in
their theorem. The fixed field is then \(F\), and the
quantity called \(\mathcal Q\) is the maximum finite
conductor norm just bounded. Taking logarithms gives
\[
      \log Nv\ll 1+\log |D_F|+f\log M+f\log(2f)
\]
with an absolute implied constant. Such a prime lies outside
\(U\), and its narrow class is the required one. If the
class lies in the kernel of \(\omega_{E/F}\), then it splits
in \(E\), since \(E\subset H^+\).

Start with the primes dividing \(2Ng_1g_2T\) and the class
of \(\mathfrak d_F\). Adjoin the rational prime below the
resulting \(\mathfrak p\) to \(U\) and repeat for the class of
\(\mathfrak a_0\). Both classes have trivial quadratic
character by Lemma~\ref{lem:good175}. The bounds on the
common intervals give \(\log |D_F|\le f\log N\) and bound
\(f,\log(2Ng_1g_2T)\) by powers of \(x\). Also
\(|U|\le \sum_{p\in U}\log p/\log 2\). Inserting the first
prime's logarithmic bound into the second proves the
assertion for both representatives.

For the stated choice of \(S\), the sum of \(\log Nv\)
over its original places is at most \(f\log(2Ng_1g_2T)\),
by the residue degree formula. The two extra terms were just
bounded. Each sum of rank factors defining \(P\) is at most
\(2m n_m^2\). At a terminal place,
\(I_{j_v}\le |\GL_r(\OO_v/\varpi_v^{j_v})|\le (Nv)^{j_vr^2}\)
for \(j_v>0\), while \(I_0=1\). Apply \(j_v\le a_v\) and
the polynomial bounds on the ranks and on \(m\). This proves
the loss estimate. Exponential growth of \(b\) gives the
last assertion.
\end{proof}

The finite tests here are explicit integral-valued indicators
or indicators times finite character values, and each real
test is a Fock monomial in the chosen coordinates. The
Petersson constant \(C_x\) includes the normalization of
the real theta pairings. The Gram estimate controls the norm in the full complex
target. Passage to the prescribed arithmetic projection requires
a bound for the image in its lattice, including the finite
valuations and the absolute archimedean normalization.
The depth needed for injectivity from the estimate in
Corollary~\ref{cor:gram177} also depends on \(C'_x\).

\subsection{The first global theta lift of the Bessel row}\label{nav:VIII-17}
Keep the choices of Theorem~\ref{thm:ordinary164}. For the
signed row \eqref{eq:src109} fix a pure unit vector \(u_\infty\)
in its lowest compact type. At a terminal depth \(j_v\ge0\)
use the residue test of Theorem~\ref{thm:trace166}, with the
coarse test at depth zero. Put \(I=\prod_{v\mid T}I_{j_v}\).
The following period concerns the first forward lift
\(U(Y_0)\to U(Y_1)\) of the indefinite chain.

\begin{theorem}[The first theta period at the prescribed terminal tests]
\leavevmode\label{thm:first168}
Fix \(x\) on a common interval of Theorem~\ref{thm:invar148},
a translate and a compatible unitary splitting with the
prescribed real exponents. A common choice of tests away
from \(T\) and a positive real number \(C_x\) give
\[
 e^{-\delta_x} C_x I^{-2}\|A\|^2
 \ \le\ \|\Theta_{\boldsymbol j} A\|^2
 \ \le\ e^{\delta_x} C_x I^{-2}\|A\|^2,\quad
 \delta_x=4[F:\QQ]\,2^{-\lfloor(w-2)/2\rfloor+4r^2},
\]
where \(F=E^+\), \(w=r+t_{\rm r}\), and
\(A=S_{\boldsymbol j}(u_\infty\otimes\ell)\) for any finite
class column \(\ell\). The same tests and \(C_x\) work for
every depth and every word with the fixed row, splitting and
levels away from \(T\). Every nonzero such lift is cuspidal.
For sufficiently large total depth at each fixed \(x\),
the inequality gives a positive first period for every
nonzero \(\ell\).
\end{theorem}
\begin{proof}
At a real place the first target character has half angular
exponent \(-e_1\). Centering \eqref{eq:src109} gives \(k\)
positive and \(s\) negative entries in the plus block, and
\(s+1\) negative entries in the minus block. All their
absolute values are at least \(H_0>(t_{\rm r}+3)/2\).
Theorem \ref{thm:occ167}, applies the real theta
nonvanishing theorem to these four counts. The first target
in this Witt tower is the one of signature \((w-s,s)\).
The lower lifts vanish, and the constant-term formula puts
the global lift of a cusp form into the cuspidal spectrum.
The square-integrable row \(A\) of Theorem~\ref{thm:ordinary164}
is cuspidal by the tempered multiplicity theorem, as proved
in Theorem \ref{thm:global167}.

In the convergent doubled inner product formula, the local
quadratic form on each terminal type has scalar
\(\|\phi_{*,v}\|^2\operatorname{vol}(K_v) I_{j_v}^{-2}\).
At infinity one can choose a Fock monomial with positive
scalar on the line of \(u_\tau\). At each bad nonterminal
place a matrix neighborhood of an integral frame gives a
positive scalar on the entire specified level. The logarithm
of the product of all other local factors has absolute value
at most \(\delta_x\), uniformly in the cuspidal constituent.
Lemma \ref{lem:finite167} and Corollary \ref{cor:uniform167}, prove these local
estimates and sum the positive forms over the spectrum.
This gives the asserted norm inequality. Lastly
Theorem~\ref{thm:trace166} bounds
\(\|I S_{\boldsymbol j}^*S_{\boldsymbol j}-1\|_{\rm op}\)
by \(C'_xT^{-\eta\sum_v j_v}\), with \(\eta>0\).
For fixed \(x\) this is less than one at sufficient depth.
For a nonzero input take the Petersson pairing of its lift
with itself to obtain the nonzero period.
\end{proof}

\begin{corollary}[The scale of the first lifted norm]
\label{cor:scale168}
Put \(d_{\boldsymbol j}=C'_xT^{-\eta\sum_v j_v}\).
When \(d_{\boldsymbol j}<1\), each nonzero \(\ell\) satisfies
\[
 e^{-\delta_x}(1-d_{\boldsymbol j})
 \ \le\ \frac{I^3\|\Theta_{\boldsymbol j}
        S_{\boldsymbol j}(u_\infty\otimes\ell)\|^2}
        {C_x\|\ell\|^2}
 \ \le\ e^{\delta_x}(1+d_{\boldsymbol j}) .
\]
At the original level a common first theta test detects any
column with \(S_0(u_\infty\otimes\ell)\ne0\).
\end{corollary}
\begin{proof}
The Gram bound in the preceding proof puts \(I\|A\|^2\)
between \((1-d_{\boldsymbol j})\|\ell\|^2\) and
\((1+d_{\boldsymbol j})\|\ell\|^2\). Multiply the two
inequalities of Theorem~\ref{thm:first168} by \(I^3\) and
substitute. Its lower bound at \(I=1\) proves the last
assertion without a depth estimate at zero.
\end{proof}

The choices away from \(T\) give a complex first lift.
The number \(C_x\) includes real integrals of the selected
monomials and finitely many volumes. We still need a surviving
ramified column at the original level for every required word,
and a comparison with the \emph{specified}, arithmetically
normalized periods along the entire chain. The first-lift norm
inequality provides neither their valuations nor the remaining
strict cohomology orders in \eqref{eq:tail114}.

\subsection{The entire theta chain at the real places}\label{nav:VIII-18}
The first-period bound of Theorem~\ref{thm:first168} uses the
prescribed real row and an actual global first lift. We can
also determine all its successive \emph{real} theta lifts.
For \(0\le j\le m\) put
\[
   n_j=r+jt_{\rm r},\quad s_j=\lfloor(m-j)/2\rfloor,\quad
   h_j=-\tfrac12\mathbf1_{j\ {\rm odd}},\quad d=(t_{\rm r}-1)/2 .
\]
At a real place form a signed list as follows. Take the
\(v_i\) with sign plus. Include \(-b_i\) with sign plus for
\(i\le j+1\) and alternating signs starting with minus for
\(i=j+2,\ldots,m\). For each \(1\le i\le j\) add
\(-e_i+d,-e_i+d-1,\ldots,-e_i-d\), all with sign plus.
Finally translate all entries by \(h_j\). Call the result \(P_j\).

\begin{theorem}[The archimedean representations through the compact endpoint]
\label{thm:tower170}
On the same fields, intervals and weights as
Theorem~\ref{thm:invar148}, the real row with parameter
\eqref{eq:src109} lifts at each successive step to the
discrete series with parameter \(P_j\) on \(U(n_j-s_j,s_j)\).
At each such step every smaller target in the same local
Witt tower gives zero. At \(j=m\) one obtains the irreducible
compact representation whose highest weight is
\(\bigl(\lambda_i-(n_m+1)/2+i\bigr)_{1\le i\le n_m}\), where
\(\lambda_1>\cdots>\lambda_{n_m}\) lists \(P_m\).
At every real place, monomial Fock tests exist along the
entire chain with nonzero composed local pairing between any
prescribed nonzero initial compact-finite vector and any
nonzero final covector.
\end{theorem}
\begin{proof}
For the step \(j\to j+1\) its two half angular characters
are \(-e_{j+1}+h_j\) on the target and \(-e_{j+1}+h_{j+1}\)
on the source. Subtract the former from \(P_j\).
All \(k+j(t_{\rm r}+1)\) positive centered entries are in the
plus block. The negative entries consist of \(s_{j+1}\) plus
entries and \(s_j\) minus entries, all of absolute value
at least \(H_0\). Theorem \ref{thm:tower169}, substitutes these
counts into Atobe's first-occurrence criterion. The parabolic
in Ichino's formula for that lift has a new compact block of
size \(t_{\rm r}\). The old negative entries exchange their
blocks, and the new plus string is centered at
\(-e_{j+1}+h_{j+1}\). Adding the half sum of the positive
roots to the inducing weight gives exactly \(P_{j+1}\).
This proves the induction from \(P_0\), and at the compact
endpoint subtracting that half sum gives the stated highest
weight. Corollary \ref{cor:vectors169}, constructs the local tests
by successively contracting the irreducible oscillator
quotients. Expanding the resulting polynomials proves the
monomial assertion.
\end{proof}

For the global period these computations determine all the
real representations. They leave a finite-place and arithmetic
question at each subsequent lift. Already on the second pair,
the source rank is \(r+t_{\rm r}\) and the target rank is
\(r+2t_{\rm r}<2(r+t_{\rm r})\). Thus the large-rank bound
used for the absolute convergence of the first doubled
Eisenstein series in Theorem \ref{thm:global167}, does not extend
by the same argument. A nonzero value for a product of local
monomials also leaves undetermined its normalization against
the specified global integral lattices.

\subsection{The successive standard \texorpdfstring{\(L\)}{L}-values on the common fields}\label{nav:VIII-19}
Use the ranks and weights of Theorem~\ref{thm:tower170}. Put
\(t=t_{\rm r}\), \(d=(t-1)/2\), \(z_0=d+1\), and \(f=[E^+:\QQ]\).
At a finite prime \(w\) of \(E\), away from the omitted set \(S_E\),
write \(Q_w=Nw\). Let \(\mu_i\) be the unitary splitting characters
of half angular exponent \(-e_{i,\tau}\), trivial on the ideles
of \(E^+\). Write \(D_j=\eta^{-\mathbf1_{j\ {\rm odd}}}\).
At step \(i\) the source and target splitting characters are
\(\mu_iD_i\) and \(\mu_iD_{i-1}\), respectively.

\begin{theorem}[The partial standard factor through the last step]
\label{thm:values172}
Start with the standard spherical parameters \(B_0\) of a
unitary constituent of the initial row outside \(S_E\). At
each step \(i\), multiply the existing parameters by
\(D_i/D_{i-1}\) and append the segment with terms
\((\mu_iD_i)(w)Q_w^a\) for \(-d\le a\le d\). Denote the
result by \(B_i\). For \(0\le j<m\), \(k=j+1\), the partial
standard \(L\)-function of \(B_j\otimes(\mu_kD_j)^{-1}\)
has the following nonzero value at \(z_0\):
\[
 L^S(z_0,B_0\otimes\mu_k^{-1})\,
                    \prod_{i=1}^j\prod_{a=1}^t
                    L_E^{S_E}(a,\mu_i\mu_k^{-1}).
\]
The characters can be chosen on the same \(E\) with finite
ramification only at existing auxiliary or conductor primes.
For this choice, and \(S_E\) above the primes of \(2Ng_1g_2T\),
the sum over \(j<m\) of the absolute logarithms of the moduli
of these values is bounded by a polynomial in \(x\) throughout
the common intervals. Here \(E\subset\QQ(\zeta_N)\).
\end{theorem}
\begin{proof}
In the recursion all the entries at step \(j\) have a common
factor \(D_j\). Dividing it out and shifting the variable in
each Euler product gives the displayed formula at first far
to the right. The initial spherical roots satisfy
\(|\alpha|\le Q_w^{10(r+1)^3}\) by Lemma \ref{lem:sat171}.
Hence the initial factor has an absolutely convergent Euler
logarithm of size \(O(fr\,2^{-(z_0-10(r+1)^3)})\) at the
required point. Each ratio character \(\rho=\mu_i\mu_k^{-1}\)
has angular exponent \(4(k-i)H_0\) at every selected embedding.
Hecke continuation applies both to \(\rho\) and to \(\rho^2\).
For the integer \(a=1\), Lemma \ref{lem:hecke171}, combines the
functional equation with positivity of
\(3+4\cos u+\cos(2u)\) to give
\[
  \big|\log|L_E^{S_E}(1,\rho)|\big|
     \ll 1+\log(2|D_E|N\mathfrak f_\rho)
           + f\log(3+4mH_0)+\#S_E .
\]
The Euler logarithms of all the higher integer factors for
one pair contribute \(O(f)\) in total. There are fewer than
\(m^2/2\) pairs. Corollary \ref{cor:chars171}, constructs the
characters with \(N\mathfrak f_\rho\le(g_1g_2T)^{2f}\).
Together with \(\log|D_E|\le2f\log N\) and the common-family
parameter bounds this proves the assertion.
\end{proof}

The partial standard value is the scalar denoted by the
nonvanishing condition at
\srcloc{Development/AbelianVarietyTheory/Definitions/LatticeFrames/ChainParSucc}{128}.
The global theta integral on the designated input is a separate
quantity in line 215 of that file. Establishing its arithmetic
comparison still requires a surviving original-level column
for each word and control of the projected rational periods.
The remaining strict cohomology lengths are those in
\eqref{eq:tail114}.

\subsection{The trace pairing at the finite places}\label{nav:VIII-20}
Write \(F_+=E^+\) in this paragraph and retain the fields
and intervals of Theorem~\ref{thm:invar148}. The integral
tensor of two Hermitian lattices enters the oscillator
through its alternating trace form. Even a unimodular
Hermitian lattice can have a nontrivial duality index for
the trace down to \(\QQ\). This index can be treated
simultaneously at all inert places.

\begin{theorem}[The index for a single global trace pairing]
\label{thm:lattices174}
There exists a trace-zero scalar \(\delta\in E^\times\) of
the required sign at every selected CM embedding, usable at
all \(m\) steps. At any inert prime \(v\) the integral tensor
lattice \(L_{j,v}\) at step \(j\to j+1\) has a self-dual
homothety \(L'_{j,v}\supseteq L_{j,v}\) for this trace pairing.
The original compact stabilizers on both sides preserve
this homothety, and
\[
 \sum_{j=0}^{m-1}\ \sum_{v\ {\rm inert}}\log[L'_{j,v}:L_{j,v}]
 \ \le\ \left(\tfrac12\log|D_E|+
       [F_+:\QQ]\log(4/\pi)\right)\sum_{j=0}^{m-1} n_jn_{j+1}.
\]
The splitting characters can also be chosen unramified at
every inert prime. At any finite collection of split primes,
for source vectors fixed at a principal congruence depth
on which the splitting characters are trivial, one can use
the same depth at each subsequent step. For a split prime
of residue cardinality \(q_v\) and depth \(a\ge1\), the
resulting local integral on a fixed unit vector from rank
\(n\) to rank \(n+t\) is
\[
      q_v^{-a((n+t)n+n^2)}
           \prod_{\ell=1}^{n}(1-q_v^{-\ell})^{-1}.
\]
The measures here give volume one to the full matrix lattice
and to the integral general linear group, respectively.
\end{theorem}
\begin{proof}
Let \(\delta_0=\sqrt{-d_0}\) and
\(\mathfrak a_0\mathcal O_E=(\delta_0)\). Lemma \ref{lem:norm173}, proves that
\(\mathfrak d_{F_+}\mathfrak a_0\) is a narrow ideal norm
class from \(E\). For the different the obstruction is the
finite product of local root numbers of \(\omega_{E/F_+}\),
which is one since \(4\mid[F_+:\QQ]\). For \(\mathfrak a_0\)
it is one because \(d_0\) splits in \(\QQ(\sqrt2)\).
Choose an integral ideal \(\mathfrak b\) in a preimage class
by Minkowski's theorem and write
\((\alpha)\mathfrak d_{F_+}\mathfrak a_0=N_{E/F_+}\mathfrak b\)
with \(\alpha\) totally positive. With \(\delta=\alpha\delta_0\),
at \(w\mid v\) inert the dual of \(L_{j,v}\) is
\(\varpi_v^{-2v_w(\mathfrak b)}L_{j,v}\).
Half this exponent gives the claimed homothety. Its index
is \((Nw)^{v_w(\mathfrak b)n_jn_{j+1}}\).
Multiplication and Minkowski's bound give the displayed
inequality, as in Theorem \ref{thm:lattices173}.

At the finite places above \(d_0\), the original
\(\eta\)-character on residue units factors through
\(z/\bar z\). Proposition \ref{prop:splittings173}, cancels it
with a single character of the norm-one torus, using the
existing split prime \(g_1\) to enforce compatibility on roots
of unity. For the final assertion take the indicator of
matrices congruent to the first \(n\) coordinate columns
modulo \(\varpi_v^a\). Its overlap with a source translate
is supported on the depth-\(a\) subgroup, on which that
overlap is exactly the additive volume of its support. This
gives the integral by the finite general linear group order
and fixes the same target depth. The calculation is in
Proposition \ref{prop:split173}.
\end{proof}

The theorem permits a choice of the global trace scalar
and uses chosen matrix tests at split places. Its
polynomial bound applies to the local trace-lattice indices.
Identification with the specified projected global periods,
including their arithmetic divisibility, remains necessary
in Section~\ref{sec:missing}.

\subsection{A nonzero iteration over the complex numbers}\label{nav:VIII-21}
There is a qualitative global consequence of the preceding
calculations. It applies even at the steps where the doubled
Eisenstein series is outside its half-plane of absolute
convergence. The choices of finite tests and the scope of
the conclusion are part of the statement.

\begin{theorem}[The global theta chain of a surviving row constituent]
\label{thm:global176}
On the fields and intervals of Theorem~\ref{thm:invar148},
let \(\Pi_0\) be an irreducible cuspidal constituent of a nonzero
row \(S_{\boldsymbol j}\) with the levels of
Theorem~\ref{thm:first168}. There are a global additive
character, a choice of skew pairing of the required real
signs, and splitting characters with the prescribed real
exponents, for which the forward theta chain reaches the
compact endpoint with nonzero cuspidal lifts at every step.
One can fix the unramified vacua outside a finite set of
split finite places (and vary the archimedean Schwartz
functions). The finite set may include two additional split
prime ideals. The archimedean representations are those of
Theorem~\ref{thm:tower170}, and the good finite parameters
are those of Theorem~\ref{thm:values172}. On any nonzero
compact-finite input in the constituent fixed at the good
places, a choice of successive tests gives a nonzero
complex cusp form at the endpoint.
\end{theorem}
\begin{proof}
Lemma \ref{lem:good175}, chooses split prime representatives
for the separate narrow ideal classes of the different of
\(E^+\) and the ideal underlying \((\sqrt{-d_0})\).
Rescaling the additive character and skew scalar by totally
positive elements makes both unramified at every inert place.
The splitting characters of Proposition \ref{prop:splittings173},
are also unramified there. At an inert place, taking
the Jacquet module of the oscillator along an isotropic line
gives two terms. One removes the outer two parameters of
the new segment. The other matches a reciprocal pair on
the two sides. This gives the parameter update for every
spherical constituent, even after the first step when
temperedness is unavailable. At a split place the same
update follows from the normalized matrix correspondence
and the two dual components of the Hermitian space. The
calculations are Theorem \ref{thm:inert175} and Proposition \ref{prop:split175}.

For the step from rank \(n_j\) to \(n_j+t\), the global Witt
index of the target is at most its smaller real signature
\(s_{j+1}<t\). The doubled theta integral is therefore
in Weil's convergent range. Every smaller global lift
vanishes by its real first occurrence, so the tower
property gives cuspidality. The Rallis norm formula uses
the continued open-orbit integral at \(s=t/2\). Its
unramified numerator is the nonzero partial standard
value of Theorem~\ref{thm:values172}. Its denominator has
an absolutely convergent Euler product there. All the
exceptional finite places are split, where
Theorem~\ref{thm:lattices174} gives nonzero integrals on
fixed vectors. For a real discrete series its nonzero
local lift gives a nonzero local integral by Ichino's
criterion. Thus the global norm is nonzero and the
cuspidal lift is irreducible, with the local components
just computed. This permits induction through all \(m\)
steps. In an irreducible constituent the common kernel
over the allowed tests is invariant and proper, hence
zero. Theorem \ref{thm:global175}, gives the continuation and
kernel arguments in detail.
\end{proof}

The new conclusion is existence for a choice of Schwartz
functions on a surviving irreducible constituent. The prime
count in \eqref{eq:window} requires the projected values for
the designated rational theta tensors with their prescribed
integral normalization. In addition, the argument requires
a surviving ramified column for every word at the original
level and the strict cohomology estimate in \eqref{eq:tail114}.

\subsection{One theta operator for all cuspidal inputs of a row}\label{nav:VIII-22}
Fix a field and translate as in Theorem~\ref{thm:invar148},
a signed real row, and one unit vector in its lowest compact
type at each real place. Use the splittings and separately
unramified global data of Theorems~\ref{thm:lattices174}
and~\ref{thm:global176}. Let \(S\) consist of the exceptional
split places, including those above \(T\) and the two auxiliary
places. Denote the terminal Bessel depths by \(j_v\ge0\).
Write \(I=\prod_{v\mid T}[K_v:L_{j_v}]\) with the groups of
Theorem~\ref{thm:trace166}. At each \(v\in S\) take a principal
depth \(a_v\ge1\) inside the kernels of the compact determinant
characters. At \(v\mid T\) require also \(a_v\ge j_v\) and
triviality of the top-block character on \(K_r(a_v)\).
At each other exceptional place take this subgroup inside
the source level. The rank at step \(j\) is \(n_j=r+jt\),
where \(t=t_{\rm r}\).

\begin{theorem}[A common comparison through the last lift]
\label{thm:common178}
Let \(\mathcal H\) be the cuspidal space on this real line
with the specified terminal character at \(L_{j_v}\),
principal fixed vectors at the other exceptional places,
and spherical vectors elsewhere. There is a single iterated
theta operator \(\mathcal T\) on this space using monomial
Fock tests at infinity, the prescribed terminal tests on the
first edge, and explicit congruence tests at every other
exceptional finite edge. Put
\[
 \mathcal J_v=
 \begin{cases}\{1,\ldots,m-1\}&v\mid T,\\
               \{0,\ldots,m-1\}&v\in S,\ v\nmid T,
 \end{cases}
 \qquad P=\prod_{v\in S}(Nv)^{-a_v
                  \sum_{j\in\mathcal J_v}n_j(n_j+n_{j+1})}.
\]
For a constant \(C_x>0\) independent of both depth vectors
and of the finite cuspidal constituent, every \(A\in\mathcal H\)
satisfies
\[
 e^{-\Delta_x} C_x I^{-2}P\|A\|^2
 \le \|\mathcal TA\|^2 \le e^{\Delta_x} C_x I^{-2}P\|A\|^2,
 \qquad
 \Delta_x=Cmfr\,2^{-((t+1)/2-10(r+1)^3)}.
\]
Here \(C\) is absolute. For the original row the two auxiliary
primes can be chosen so that
\(-\log P+2\log I\ll_{\psi,\eps} x^c\max_{v\in S}a_v\)
on these same fields and intervals.
\end{theorem}
\begin{proof}
Choose the real tests successively. For the unit source
vector at a real place the local Rallis form is nonzero
by cyclicity and the discrete-series nonvanishing theorem.
One reference global constituent from
Theorem~\ref{thm:global176} realizes a fixed nonzero
multiple of its diagonal as a squared norm. Cauchy--Schwarz and the density of Fock
polynomials then give a monomial with nonzero diagonal.
The unique irreducible local theta pairing determines the
next unit line. This choice uses only the real representations.
When the specified cusp space is zero the estimate is vacuous.

At the first terminal edge the local integral is a scalar
multiple of the projector onto the prescribed character.
Its scalar is the depth-zero scalar times \([K_v:L_{j_v}]^{-2}\).
The top-block support and character survive left multiplication
by \(K_{n_1}(a_v)\). On the other exceptional edges the
integral of Theorem~\ref{thm:lattices174} supplies the
displayed power of \(Nv\) and preserves the target principal
depth. In the unramified formula of Theorem~\ref{thm:values172}
all factors are common except the one from the original
spherical row. Its logarithmic absolute value is
\(O(fr2^{-((t+1)/2-10(r+1)^3)})\) at every edge. Multiply the
\(m\) norm identities. Orthogonal invariant cuspidal summands
have orthogonal lifts because their mixed doubling integral
unfolds to a zero matrix coefficient and hence has identically
zero continuation. This permits summing the norm inequalities
over the whole row.

For the prime bound take the two narrow ideal classes of the
different and of the ideal below \((\sqrt{-d_0})\). To avoid
previous rational primes, adjoin to the narrow Hilbert class
field the roots of unity of orders \(p^{k_p}\) with
\(\varphi(p^{k_p})>[E^+:\QQ]\) at those primes. Every one
of them ramifies, while finite conductor norms of all the
abelian characters are at most \((\prod p^{k_p})^{[E^+:\QQ]}\).
The unconditional least-prime theorem of Thorner and Zaman
now bounds the logarithmic norm of each representative by
a power of \(x\). Both classes split in \(E\). Counting matrix
entries in the finite congruence subgroups gives the asserted
depth cost. The full proofs of the local pairing, the
simultaneous iteration and the least-prime estimate are in
Section \ref{nav:VIII-07}.
\end{proof}

\begin{corollary}[Kernel on the finite Bessel columns]
\label{cor:common178}
On the finite column space at the fixed real line let
\(S_{\mathbf j}\) denote the Bessel periodization and put
\(B=\mathcal T S_{\mathbf j}\). Then
\(\ker B=\ker S_{\mathbf j}\), including at the original
depth. With the exponent \(\eta>0\) and field constant
\(C'_x\) of the terminal-depth estimate, one has
\[
 \left\|\frac{I^3}{C_x P} B^*B-1\right\|_{\rm op}
 \le e^{\Delta_x}-1+e^{\Delta_x}C'_x T^{-\eta\sum_{v\mid T}j_v}.
\]
\end{corollary}
\begin{proof}
The lower bound of the theorem gives the kernel equality.
Insert the upper and lower quadratic-form bounds for
\(I S_{\mathbf j}^*S_{\mathbf j}\) from
Theorem~\ref{thm:depth164} and the simultaneous-depth
calculation of Section \ref{nav:VII-13}. Taking the supremum
over unit columns gives the displayed inequality.
\end{proof}

Thus a column which survives at the original level can be
detected through the compact endpoint by the common tests.
The prime count requires a surviving ramified column for
every word there, with a nonzero pairing against the specified
arithmetic projection and a bound for its integral size. The
remaining strict lengths in \eqref{eq:tail114} also enter the
estimate for \eqref{eq:window}.

\section{Terminal Hecke operators and the first global lift}\label{nav:first-global-lift}

\subsection{Terminal covectors through the definite theta chain}\label{sec:covectorpaths352}

Take the base auxiliary tuple and the definite column used in
Theorem~\ref{thm:terminalformula347}. Keep its chosen period
field \(E_x\), its word, the marks at \(g_1,g_2\), and
\(S_T=\{P:P\mid\pi\}\) for the same terminal generator \(\pi\).
We compute the finite sum on this column before comparing it
with the analytic Bessel functional. At a vertex retain the
lattice \(M_i\), its frames and levels away from \(T\), its
finite type index and all its coefficient data. A terminal
refinement is a tuple \(f_i=(f_{i,P})_{P\in S_T}\) of nonzero
linear forms on \(V_{i,P}=M_i/PM_i\). The unrestricted
coefficient vector space at the vertex is common to these
refinements. Its integral lattice can depend on the kernels of
the \(f_{i,P}\).

An edge term uses a rational matrix \(x_i:M_i\otimes E\to
M_{i+1}\otimes E\). Its terminal factor is one precisely when
\(x_i\) is integral at the indicated \(P\)'s and
\[
       f_{i,P}=f_{i+1,P}\,\overline{x}_{i,P}
                   \qquad(P\in S_T),
\]
and zero otherwise. Here \(\overline{x}_{i,P}:V_{i,P}\to V_{i+1,P}\)
is reduction. This follows directly by applying the two covectors
to every element of the local lattice. In the finite theta sum
the matrix ranges over the quotient of \(c_A^{-1}\mathcal H_i\)
by \(c_B\mathcal H_i\), where \(\mathcal H_i\) is the module of
rational homomorphisms integral between the local lattices at
the indicated finite places. The denominator \(c_A\) is prime
to \(T\), and \(T\mid c_B\). Thus reduction of the matrix at
\(S_T\) is defined on each coset. This quotient is finite:
localizing at the finite set of places gives projective modules
over a semilocal Dedekind ring, and reducing the homomorphism
module modulo \(c_Ac_B\) gives a product of finite residue
modules. This description also follows by choosing bases over
each of the finitely many local rings and applying the Chinese
remainder theorem.

The other factors in an edge term are its coefficient-valued
theta distribution on that matrix coset and its local tests at
the other marked places. They use \(M_i,M_{i+1}\) and the
prescribed frames there. A lift of a frame at \(w\mid g_1g_2\)
is specified by its lattice and its reduction at \(w\), so the
same lift can be used for every \(T\)-refinement. In the final
two coupled edges the additional vector coset is taken from the
intermediate lattice with bounds only at the other marked
places. Its tests on the two edges use this vector coset and
the same frames. In particular these factors have no dependence
on the intermediate terminal covectors. The edge sums and the
three terminal indicators are given in
\srcloc{Formalism/LocalToGlobal/Columns/FiniteSumNonzeroLargeDepth}{636},
lines 1115--1139 and 1162--1170 there.

Fix unrefined initial class representatives \(a\). Fix a small
test \(F\) at the endpoint, possibly after its finite type
projection. Write \(\gamma\) for a choice of the initial and
intermediate unrefined classes, the final small class \(b\), the
finite type indices, the matrix cosets at each edge, and any
coupling coset used by the two adjoining edges. Write
\(W_{\gamma,F}\) for the following linear functional on the
initial coefficient space. Compose the coefficient maps of these
matrix cosets in their original order and apply the terminal
lowering \(L_z\). Contract with \(F(b)\), and multiply by the
product of the signed local tests outside \(T\), including the
weights of the chosen finite linear combinations. Divide by
\(|\Gamma_b|\prod_{i<m}|\Gamma_{a_i}|\), where \(m\) is the
number of edges, \(\Gamma_b\) is the small class stabilizer and
the earlier stabilizers are at the unrefined levels. All changes
of frame use the same class representatives in the two incident
maps. This prescription retains the coefficient maps of the
actual cosets, including their dependence on the matrix residues
at \(T\).

\begin{theorem}[Summing the intermediate terminal refinements]\leavevmode
\label{thm:covectorpaths352}
Let \(e_N^\vee\) read the normal coordinate at the lifted final
class \(i(b)\). For a path starting at \(a\), put
\[
 \eta_{\gamma,P}=(\bar\pi\bmod P)\,\overline{e_N^\vee}\,
       \overline{x}_{m-1,P}\cdots\overline{x}_{0,P}.
\]
Discard a path if any \(\eta_{\gamma,P}=0\). The signed period
on an initial column vector \(\phi\) is exactly
\[
             \mathcal P(\phi,F)
           = \sum_{a}\ \sum_{\gamma\ {\rm starting\ at}\ a}
                  W_{\gamma,F}\big(\phi(a,\eta_\gamma)\big),
\]
using the remaining paths and the stated finite class measure.
In particular every path whose composite maps into the final
marked hyperplane at even one of the places contributes zero.
If the terminal covectors of the lifted class are scaled by
\(c\in A_T\) in this finite recurrence, the same formula reads
the initial column at \(c\eta_\gamma\).
\end{theorem}
\begin{proof}
Above an unrefined representative \(a_i\) the refined classes
are the orbits of the finite group \(\Gamma_{a_i}\) on the
tuples of nonzero covectors. The refined stabilizer is the
stabilizer of the tuple in this action. In a contracted class
sum a change by \(\Gamma_{a_i}\) transports the preceding and
following matrix variables, and their coefficient transports
cancel on the shared value. At the initial vertex the same
statement uses the defining equivariance of \(\phi\).
Reindexing the adjacent coset sums therefore gives a function
constant on each orbit. Orbit--stabilizer replaces the sum over
refined orbits with their inverse stabilizer orders by
\(|\Gamma_{a_i}|^{-1}\) times the sum over all tuples above
this representative. One may make this replacement at every
vertex, or take the product of the finite stabilizers and make
it once on all the vertices. At the final small vertex the
covector is already prescribed by the normal marking.

Now fix the remaining matrix choices and coupling data. Only
the displayed edge indicators and \(\phi\) depend on the
intermediate covectors. At each \(P\) start with the final
covector \((\bar\pi\bmod P)\overline{e_N^\vee}\) and pull
it back through the matrices in reverse order. There is exactly
one possible value at every vertex. If any pullback is zero,
all the pullbacks back to the initial vertex are zero there.
Conversely, nonzero initial pullback implies that all these
values are admissible. The product of all the terminal
indicators sums to one for that initial refinement, and to
zero for every other one. Interchanging the finite sums gives
the formula with precisely the unrefined masses already obtained.
Multiplying the final covector by \(c_P\) multiplies every
one of its pullbacks by \(c_P\) and changes none of the other
factors. This gives the last assertion.
\end{proof}


\subsection{Fourier transformation of the definite matrix recurrence}\label{sec:vectorchain353}

Keep the matrices, signed coefficient maps and class measures of
Theorem~\ref{thm:covectorpaths352}. At a lattice \(M_i\) of rank
\(n_i\) put \(\Omega_i=\prod_{P\mid\pi}M_i/PM_i\) and
\(Q_i=|\Omega_i|=T^{fn_i}\), with \(f=|S_T|\). Use
\(\Omega_i^\vee=\prod_{P\mid\pi}(M_i/PM_i)^\vee\).
Choose the prescribed nontrivial residue additive character
\(\psi_P\) at each place and write
\(\psi(\xi(v))=\prod_P\psi_P(\xi_P(v_P))\).
Define the unnormalized transform and its opposite by
\[
   (\mathcal F_i H)(v)=\sum_{\xi\in\Omega_i^\vee}\psi(\xi(v))H(\xi),
   \qquad
   (\mathcal F_i^-G)(\xi)=\sum_{v\in\Omega_i}\psi(-\xi(v))G(v).
\]
A separately zero-sum table has sum zero on every scalar orbit in
each coordinate. Extend it to the full product by zero wherever
a coordinate vanishes. We use this convention on both sides.

For a residue matrix tuple \(C:\Omega_i\to\Omega_{i+1}\) and
the coefficient map \(W\) of that matrix term, the covector edge
acts by \((K_C H)(\xi')=W H(\xi' C)\). On vector tables use
\[
                (K_C^\to G)(w)=\sum_{v:\,Cv=w} W G(v).
\]
Denote by \(K\) the definite recurrence up to its large terminal
table, before evaluating at the lift of a small test class.
Use \(K^\to\) for the same sum with the vector operations on
its matrix terms. This keeps the order of the coefficient maps,
the unrefined class weights, the common vector coset in the two
coupled edges and all local tests outside \(T\).

\begin{theorem}[Exact vector recurrence with the original theta coefficients]\leavevmode
\label{thm:vectorchain353}
On the separately zero-sum tables one has
\[
 \mathcal F_m K=\frac{Q_m}{Q_0}K^\to\mathcal F_0,\qquad
 K=Q_0^{-1}\mathcal F_m^- K^\to\mathcal F_0 .
\]
In the class Hilbert norms \(Q_i^{-1/2}\mathcal F_i\) is unitary
between the two tables. In particular, when the large terminal
value is read at a lifted small class \(b\), its value before
lowering is the sum of \((K^\to\mathcal F_0 H)(b,v)\) over the
terminal vector tuples with coefficients
\(Q_0^{-1}\psi(-f_b(v))\), where
\(f_{b,P}=(\bar\pi\bmod P)\overline{e_N^\vee}\).
The only Fourier inversion denominator in this evaluation is
\(Q_0=T^{fk}\), whose logarithm is polynomially bounded on the
period intervals.
\end{theorem}
\begin{proof}
For each fixed vector, summing the characters over the full
dual group gives its order at the zero vector and zero otherwise.
For a nonzero vector this follows by translating the summation
by a covector with nontrivial character value. This proves both
\(\mathcal F_i^-\mathcal F_i=Q_i I\) and the norm identity.
At a zero vector coordinate the transform of an extended input
vanishes, because the character there is constant on each scalar
orbit. At a nonzero coordinate, summing the output over scalar
multiples sums the input over the corresponding scalar action,
by a change of variables. The result is zero. The same argument
with the inverse shows that the two zero-sum spaces correspond.
These transforms commute with changes of frame when the vector
and covector frames are dual. At a fixed unrefined class the
refined mass norm is the sum of all point norms divided by the
unrefined stabilizer order, by orbit--stabilizer. This proves
unitarity also on the equivariant class tables.

Insert \(H=Q_i^{-1}\mathcal F_i^-G\) into \(K_C\) and apply
\(\mathcal F_{i+1}\). For each \(v,w\) the inner character sum
over \(\xi'\) is that of \(\psi(\xi'(w-Cv))\). It equals
\(Q_{i+1}\) exactly when \(w=Cv\). This proves the first
identity on one edge, without a rank condition on \(C\).
The covector formula already gives zero if a pulled-back
coordinate vanishes. Thus it uses exactly the terminal indicators
of the original edge. The vector formula preserves the zero-sum
space by Fourier conjugacy. In a product of edges the factors
\(Q_{i+1}/Q_i\) telescope. Apply this term by term with the
coupling data fixed, then sum the signed maps and class weights
as in Theorem~\ref{thm:covectorpaths352}. Inverting the last
transform gives the second formula and evaluation at \(f_b\).
Finally \(f,k\) and \(\log T\) have polynomial bounds by
Theorem~\ref{thm:common111}.
\end{proof}


\subsection{The terminal Hecke operator on the marked initial classes}
\label{sec:neighbormove356}

Fix an ordinary word on the fields of Theorem~\ref{thm:common111}.
Write \(S\) for its full separately zero-sum integral initial
table, summed over the ordinary classes. Its rank is \(d\).
At a chosen coefficient place \(w\mid T\), let \(A_w\) be the
set of \(P\mid\pi\) induced at \(w\) by a high embedding
\(\tau\). Write \(s=|A_w|\), and use \(\tau(P)\) for this
embedding. The initial weights have \(\lambda_{\tau,k}=-b_*\)
at each such place. The operation required on this branch of
the integral construction is \(H_w^m\), where
\[
                  H_w=T^{s b_*}\prod_{P\in A_w} U_P .
\]
Here \(U_P\) sums over upward neighbors with unit coefficient.
We give its local definition below. The scalar, the selected
places and the power \(m\) are specified in
\srcloc{Formalism/LocalToGlobal/Columns/ScaledRamifiedLattice}{219},
and the use of the resulting image at a prime over \(T\) is in
\srcloc{Formalism/LocalToGlobal/Columns/FiniteSumNonzeroLargeDepth}{1336}.
The initial lattice is the Weyl lattice of class functions with
zero scalar sums at every marked prime, as in
\srcloc{Formalism/LocalToGlobal/Columns/ColForms}{108}
and \srcloc{Formalism/LocalToGlobal/Columns/ScaledRamifiedLattice}{159}.

\begin{theorem}[Incidence, integrality and determinant of the terminal Hecke move]\leavevmode
\label{thm:neighbormove356}
The operator \(H_w\) maps \(S_w=S\otimes\mathcal O_{J_w}\)
into itself and is invertible on its fraction space. In the
valuation \(v_w(T)=1\) it satisfies
\[
 \frac1d\,v_w\det H_w
      =s b_*+\frac{s(k-1)}2+
           \frac1k\sum_{P\in A_w}|\lambda_{\tau(P)}|.
\]
Here \(|\lambda|=\sum_i\lambda_i\). At each complex embedding
all its singular values in the ordinary class Hilbert space
are \(T^{s b_*+s(k-1)/2}\). For \(s=0\) the map is the identity.
\end{theorem}
\begin{proof}
At a split pair \(P,\bar P\), write \(M\) for the rank \(k\)
position lattice at \(P\) and \(M^\vee\) for its partner under
the perfect Hermitian pairing. Start with a nonzero marked
covector \(c\) on \(M/TM\). The neighbors read by \(U_P\)
are indexed by \(x\in M/TM\) with \(c(x)=1\). The position
lattice of such a neighbor is \(L_x=M+\mathbb Z_T T^{-1}\tilde x\)
with \(\tilde x\) any lift, and its partner is \(L_x^\vee\).
Its covector \(f_x\) vanishes on \(M\) and takes value one on
\(T^{-1}\tilde x\). Indeed \(f_x(y)=c(Ty\bmod TM)\) on \(L_x\).
Away from this pair the lattices and markings carry over by the
identity. This is the marked neighbor definition in
\srcloc{Formalism/LocalToGlobal/Columns/StdLatUnimodular}{139},
lines 145--163 there. Intersecting the prescribed local
lattices in the rational space produces a global lattice with
those localizations, by the Chinese remainder theorem after
clearing a common denominator. The perfect duals at conjugate
places give self-duality. A basis at the position place and its
dual at the other place identify the two local Hermitian
lattices with the standard split pair, so the new lattice is
in the same genus. The residue vector \(x\) can also be lifted
inside the global lattice by reduction modulo \(P\).

Conversely, given the upward lattice and its covector, taking
the kernel at \(P\) recovers \(M\). The relation
\(f_x(T^{-1}\tilde x)=1\) recovers \(x\) uniquely in \(M/TM\):
any two elements of \(L_x\) with covector value one differ by
\(M\), and multiplication by \(T\) gives the same residue.
Thus the source markings above the common kernel are the nonzero
vectors \(x\), and the target markings are the nonzero covectors
\(c\) on that kernel lattice. Scaling the original \(f_x\)
scales \(x\) inversely. In particular the scalar augmentation
is the same condition after this reindexing. These operations
are equivariant for integral frame changes. The initial
unrefined stabilizers are trivial by Section~\ref{sec:pathnorm352}.
Consequently summation on classes at this level in these
coordinates uses the matrix \(I(c,x)=\mathbf1_{c(x)=1}\)
with coefficient identity and unit mass. Alternatively one can
sum on all markings and divide by the common unrefined
stabilizer, as in Theorem~\ref{thm:covectorpaths352}.
At different split pairs all modifications are independent.
Performing them simultaneously reindexes the composite as
the tensor product of the incidence matrices, with all other
data fixed.

The incidence matrix preserves the two scalar augmentation
spaces. Indeed summing \(I(c,x)\) over the nonzero scalar
multiples of \(c\) gives \(\mathbf1_{c(x)\ne0}\), which is
constant as \(x\) ranges over any scalar line. For two vectors
\(x,y\), the number of covectors taking value one on both is
\(T^{k-1}\) when \(x=y\), zero for unequal proportional
vectors, and \(T^{k-2}\) for independent vectors. On an
augmented table the independent terms vanish by summing over
each of the other lines. Hence \(I^*I=T^{k-1}\) there. The
transpose gives the same argument on covectors. This proves
invertibility and the claimed archimedean norm, including the
product of the places and the positive scalar. Transport by
rational unitary elements in the class representatives is
isometric. Finite level character factors have absolute
value one.

We compute the coefficient frames in this reindexing. Since
\(T\) splits completely, there is at most one coefficient
embedding inducing \(P\) at \(w\). The embeddings form a CM
type. If \(\tau(P)\) induces \(P\), its conjugate induces
\(\bar P\) and is outside that type. No other embedding then
induces either member of the pair. At \(\tau(P)\), the
description by the ordinary dual of the Hermitian dual in
Theorem~\ref{thm:pointlattice353} uses the lattice from the
opposite place. The marking restricts the position place
alone. Thus the target coefficient frame at \(c\) is
\(W_{\lambda_{\tau(P)}}(M^\vee)\), and the source frame
at \(x\) is \(W_{\lambda_{\tau(P)}}(L_x^\vee)\). In the
dual position lattice one has
\[
       L_x^\vee=\{z\in M^\vee:z(\tilde x)=0\bmod T\}.
\]
All other coefficient factors are unchanged for this move.
Each selected source has precisely one coordinate multiplied
by \(T\) in an adapted integral frame. Its least weight is
\(\lambda_{\tau(P),k}=-b_*\), by
Lemma~\ref{lem:openflags353}. Hence the scalar \(T^{s b_*}\)
puts all entries into the target frame. It follows that
\(H_w S_w\subseteq S_w\).

Use the target coefficient frame also as a reference for the
source at every vector marking. Relative to this reference,
the unscaled source exponents are distributed as sums of the
single-coordinate weights at the selected embeddings, with
their Weyl multiplicities. Indeed the source coefficient lattice
is constant on each tuple of scalar lines. Each line contributes
its integral basis of augmentation differences, exactly as
in Theorem~\ref{thm:fourierslopes354}. Other marked places may
first be grouped by their scalar lines. Their coefficient
factors and their augmentation bases occur on both sides.
By symmetry of the diagonal coordinates the mean weight at
\(P\) is \(|\lambda_{\tau(P)}|/k\), as in
Corollary~\ref{cor:fouriermean354}. On the reference lattices
identify vector and covector indices by dual bases. The resulting
single-place incidence endomorphism is symmetric and its square
on the augmented array is \(T^{k-1}I\). Its determinant therefore
has valuation \((k-1)/2\) per dimension. Products add these
means, as does the coefficient scalar. Integral coordinate
changes in the class representatives preserve all of the actual
source and target lattices. This gives the stated local
determinant after summation over the classes.
\end{proof}


\subsection{The auxiliary level factor in the single-lift determinant}
\label{sec:auxdet805}
On a shorter period field put \(s=r_x,\ t=5s,\ n=s+t\).
For a single lift with source \(W_j\), use the block test of
Lemma~\ref{lem:sharpplaces709} and its fixed archimedean vector.
The good source of this lift is the marked space at the original
principal levels. Its dimension \(d_j\) is independent of the
tame twist. Normalize the source integral using the same Haar
measure before and after twisting. On the definite target use
the Hermitian class sum with reciprocal stabilizer weights, as
in Theorem~\ref{thm:actualminors804}. A level character acts by
a root of unity and hence disappears from the squared norm.

For residue cardinality \(q\) and a finite twist set \(\mathcal Q\)
put
\[
     F_{s,t}(q)=\prod_{a=t+1}^{t+s}(1-q^{-a}),\qquad
     C_{\mathcal Q}=\prod_{w\in\mathcal Q}q_w^{st}F_{s,t}(q_w).
\]
In products over marked or twist places we take one member of
each split pair in \(E\); it gives one place of \(E^+\).
In particular \(C_{\varnothing}=1\).

\begin{lemma}[Conversion from maximal-compact measure to the target class sum]
\leavevmode\label{lem:auxmass805}
At an original level place the local factor in the class-sum
norm is \(F_{s,t}(q)/|\mathrm{GL}_s(\mathbb F_q)|\).
At a new twist place it is \(q^{st}F_{s,t}(q)\).
These factors include the local source convolution on its
level-fixed vectors and the change to the target class measure.
\end{lemma}
\begin{proof}
Initially give each maximal compact group and the integral
matrix lattice volume one. Write \(g_s=|\mathrm{GL}_s(\mathbb F_q)|\).
At an original level the marked matrix is one specified
inclusion. By Lemma~\ref{lem:sharpplaces709} its source
coefficient is \(q^{-ns}\) times the characteristic function of
the principal subgroup, which has volume \(1/g_s\). At a twist
place the top \(s\)-by-\(s\) block is invertible modulo the
maximal ideal. Its nonzero values are character values, so its
squared norm is \(g_s/q^{s^2}\), the fraction of invertible
blocks. The source determinant character cancels the change of
this test under the maximal compact, giving that same scalar
on its fixed vectors.

A compact level of index \(b\) in the target maximal compact
has mass \(1/b\) for the initial measure. Its class contributes
\(1/(b|\operatorname{Stab}|)\) under this measure; hence the
class-sum convention multiplies the norm squared by \(b\).
At an original place the target level fixes an ordered
\(s\)-frame. Its index is the number of full rank \(n\)-by-\(s\)
residue matrices, namely
\[
        I_{n,s}(q)=\prod_{k=0}^{s-1}(q^n-q^k)
                      =q^{ns}F_{s,t}(q).
\]
Here each next column is chosen outside the span of the preceding
ones. At a twist place the target level preserves the subspace
of that frame. Its index is \(I_{n,s}(q)/g_s\), because every
such subspace has \(g_s\) ordered bases. Passing to the
parahoric with its prescribed unitary character gives the same
measure on pairings of two forms with that character. The two
products of source masses and target indices are therefore
\(q^{-ns} I_{n,s}(q)/g_s\) and
\((g_s/q^{s^2})I_{n,s}(q)/g_s\), as asserted.
\end{proof}

\begin{theorem}[Exact conductor factor in the auxiliary Gram determinants]
\leavevmode\label{thm:auxdet805}
For a finite twist \(\beta\) let \(\Theta_{j,\beta}\) be the
single lift in the above normalization. Write
\(\mathcal V_j(\beta)=\det(\Theta_{j,\beta}^{*}\Theta_{j,\beta})\)
on the good source, and \(\mathcal R_{j,\beta}^{\,g}\) for the
orthogonal compression to that source of the unmarked finite
operator in Theorem~\ref{thm:finitegram709}. For the simultaneous
twists of Corollary~\ref{cor:actualdenom803}, at every sufficiently
large node and every \(j\) one has
\[
  \frac{\mathcal V_j(\vartheta^j)}{\mathcal V_j(1)}
       = C_{\mathcal Q}^{\,d_j}
            \frac{\det\mathcal R_{j,\vartheta^j}^{\,g}}
                 {\det\mathcal R_{j,1}^{\,g}}.
\]
All these determinants are positive, with empty determinant one.
For the product over \(j\), after division by
\(C_{\mathcal Q}^{\sum_j d_j}\), the absolute logarithm is at
most
\[
          \left(\sum_j d_j\right)
          \log\frac{1+\varepsilon_s}{1-\varepsilon_s},
              \qquad \varepsilon_s=512 [E_x^+:\mathbb Q]\,2^{-s}.
\]
\end{theorem}
\begin{proof}
The unfolded identity for two source forms in the proof of
Theorem~\ref{thm:globalgram711} gives a sesquilinear form equal
to the pairing with the unmarked finite operator times the
product of the real integrals, the positive Haar constant and
the two sorts of level factors. To compare it with the
algebraic-form class pairing, realize the fixed vector-valued
form by its coefficient along the chosen highest line of the
compact real target. For a unit vector \(u\) in its irreducible
compact type, averaging the rank-one operators
\(\rho(k)u\,(\rho(k)u)^*\) over probability Haar gives the
identity divided by the dimension of the type. Indeed the
average commutes with the group and has trace one. Thus passage
between that scalar coefficient norm and the vector class norm
introduces a positive scalar depending only on the real type
and the normalization of its vector. The finite measure
change is exactly the one in the lemma. Equivalently one can
apply the unfolded diagonal formula and polarize on the good
source, since both forms are Hermitian.

The target and source for this \(j\), their real polynomials,
the Haar measures initially giving the maximal compacts mass one,
and the scaling of the theta kernel are fixed as \(\beta\) varies.
The real factors are positive and independent of finite twists
by Theorem~\ref{thm:realgram710}. The original level product is
fixed as well. Consequently, for a positive constant \(A_j\)
independent of \(\beta\), the operator on the good source is
\[
       \Theta_{j,\beta}^*\Theta_{j,\beta}
             = A_j C_{\mathcal Q(\beta)}
                        \mathcal R_{j,\beta}^{\,g}.
\]
This formula uses the unscaled source integral. Multiplying that
integral by any further scalar would multiply the operator by
the squared modulus of that scalar. The estimates of
Theorem~\ref{thm:finitegram709} apply before and after twisting.
The compressed operator is self-adjoint with all eigenvalues in
\([1-\varepsilon_s,1+\varepsilon_s]\), by its Rayleigh quotients.
At the nodes in question \(\varepsilon_s<1/2\). This gives
positivity on every good source, including those on \(W_j\)
with indefinite real signature. Taking determinants and dividing
by the untwisted identity cancels \(A_j^{d_j}\).
Every power \(\vartheta^j,\ 1\le j\le m_x\), ramifies at each
chosen twist place, so \(\mathcal Q(\vartheta^j)=\mathcal Q\).
In each ratio of positive determinants, the logarithm lies
between plus and minus \(d_j\log((1+\varepsilon_s)/(1-\varepsilon_s))\).
Adding proves the assertion for the product.
\end{proof}

\begin{proposition}[Valuation of the common scalar at the coefficient prime]
\leavevmode\label{prop:auxvaluation805}
Suppose \(q\) is the norm of a new split place for a twist
at the coefficient prime \(\ell\); in particular \(\ell\nmid q\).
For \(q\equiv1\bmod\ell\) one has
\[
\begin{split}
 v_\ell\big(q^{st}F_{s,t}(q)\big)
   ={}&s\,v_\ell(q-1)+v_\ell\!\left(\frac{(s+t)!}{t!}\right)\\
   &+\mathbf1_{\ell=2}
       \left(\left\lfloor\frac{s+t}{2}\right\rfloor-
             \left\lfloor\frac{t}{2}\right\rfloor\right)
                          \big(v_2(q+1)-1\big).
\end{split}
\]
If the tame order at \(w\) is \(\ell^{h_w}>1\), then
\(v_\ell(C_{\mathcal Q})\ge s\sum_{w\in\mathcal Q} h_w\)
for a twist of \(\ell\)-power order. In particular a small
valuation of \(1-\vartheta(\tau_w)\) at every new generator
is compatible with a large valuation of the common scalar.
\end{proposition}
\begin{proof}
The factor \(q^{st}\) and all powers of \(q\) in the denominators
are \(\ell\)-adic units. We must sum \(v_\ell(q^a-1)\) for
\(t<a\le s+t\). If \(\ell\) is odd and \(q\equiv1\bmod\ell\),
raising \(q\) to a power prime to \(\ell\) preserves
\(v_\ell(q-1)\) by the geometric series formula. Raising a
number \(1+u\), with \(v_\ell(u)\ge1\), to the \(\ell\)-th
power increases that valuation by one: in the binomial
expansion the term \(\ell u\) has strictly smaller valuation
than all the other terms of positive degree. Thus
\(v_\ell(q^a-1)=v_\ell(q-1)+v_\ell(a)\).
At \(\ell=2\), \(q\) is odd. An odd exponent preserves both
the valuation of \(q-1\) in the difference and the valuation
of \(q+1\) in the sum, by factoring and reducing the odd
quotients modulo two. The first squaring gives
\(v_2(q^2-1)=v_2(q-1)+v_2(q+1)\); after a squaring the new
base is \(1\bmod 8\), so each further squaring adds exactly
one. The correction to \(v_2(q-1)+v_2(a)\) is therefore
\(v_2(q+1)-1\) precisely for even \(a\). Summing proves the
identity.

The tame character at the chosen split place factors through
\(\mathbb F_q^\times\), hence \(\ell^{h_w}\mid q-1\). Each
of the \(s\) differences \(q^a-1\) is divisible by \(q-1\),
giving the stated lower bound after summation over \(w\).
On the other hand
\(v_\ell(1-\vartheta(\tau_w))=1/(\ell^{h_w-1}(\ell-1))\)
by Lemma~\ref{lem:tametraceheight734}. For a nonempty auxiliary
set the latter number can be made arbitrarily small in the
twisting construction, and doing so forces the lower bound for
\(v_\ell(C_{\mathcal Q})\) to grow. This accounts for the power
of \(C_{\mathcal Q}\) in the determinant identity before a
comparison across different lifts.
\end{proof}
The determinant in this subsection concerns one prescribed lift.
Applying it to the composite chain requires an identity comparing
the composite Gram operator with the product of the single-lift
operators, including the level scalar and the source multiplicities.
\clearpage

\subsection{The common finite source of the auxiliary determinants}
\label{sec:commonsource806}
Fix a node of the shorter family of Corollary~\ref{cor:fieldrelations712}.
Write \(F=E_x^+\), \(E=E_x\), \(m=m_x\), and \(s=r_x=k_x+m_x-1\).
Both \(k_x\) and \(m_x\) are even. Consider the rank \(s\) spaces
\[
                         Y_0,\quad W_1,\ldots,W_m .
\]
For each \(V\) in this list let \(H_V\) be its marked automorphic
space at the level of Theorem~\ref{thm:auxdet805}. At infinity
we take the coefficient on the specified line of the discrete
series. At finite places take the principal congruence subgroup
at the chosen places above \(p_0,p_1\) and the lattice stabilizer
elsewhere. Write \(J\) for this compact open subgroup and \(e\)
for the two commuting central projections of
Section~\ref{sec:levelmass678}, taken together as a product.

The finite frames identify all these finite adelic groups, together
with \(J\) and \(e\). To see the identification explicitly, write
\(B_V\) for the rational Hermitian matrix of a space and \(u_{V,v}\)
for an orthonormal frame of its specified local lattice. Then
\(u_{V,v}^{*} B_V u_{V,v}=1\), and
\(u_{V',v}u_{V,v}^{-1}\) carries the lattice of \(V\) isometrically
to that of \(V'\). Conjugation by this isometry identifies their
unitary groups. Such frames exist for the diagonal sign lattices
and for the unimodular row lattice
(Lemma~\ref{lem:augmentedindex687}). At an inert finite place
\(E/F\) is unramified. A nondegenerate Hermitian space over
the quadratic residue field has a vector of nonzero norm.
Indeed, if all diagonal values vanished, variation of a scalar
in the norm of \(a u+v\) would, by nondegeneracy of the residue
trace pairing, force each pairing of two vectors to vanish.
Lift such a vector, use surjectivity of the norm on local units
to give it norm one, and split off its orthogonal complement.
Surjectivity follows from the norm on the residue fields and the
trace on each successive principal-unit quotient. This proves
the frame assertion by induction, also in residue characteristic
two. At a split place choose a basis on the prescribed member
of the pair and its dual on the other member. In particular the
chosen reduction coordinates can be used at the marked places.
In these coordinates \(J\) and the finite residue characters
defining \(e\) are identical for every \(V\).

\begin{lemma}[Two split supercuspidal components]
\leavevmode\label{lem:twoelliptic806}
Every irreducible automorphic constituent of \(H_V\) has a
supercuspidal component at a place above each of \(p_0,p_1\).
It belongs to the cuspidal spectrum.
\end{lemma}
\begin{proof}
The product of finite character projections defining \(e\) selects
tuples with a primitive member above each rational mark. Decompose
the \(J\)-fixed vectors of a constituent into isotypic spaces
for the finite residue groups and take a tuple on which \(e\)
acts as the identity. The two primitive members occur at distinct
split places because \(p_0\ne p_1\). At either place, of residue
cardinality \(q\), the primitive label is one Frobenius orbit of
length \(s\), with partition consisting of one box. In Green's
classification this is a cuspidal representation of
\(\mathrm{GL}_s(\mathbb F_q)\), so it is absent from all finite
inductions from proper parabolics.

Let \(K=\mathrm{GL}_s(\mathcal O_{F_v})\) and \(K^1\) be its
reduction kernel. If the local representation were a subquotient
of a proper parabolic induction, its \(K^1\)-fixed vectors would
be a subquotient of the \(K^1\)-fixed vectors of that induction.
Averaging over \(K^1\) gives exactness over \(\mathbb C\).
For a parabolic \(P=MN\), the Iwasawa decomposition \(G=PK\)
identifies the latter fixed vectors with functions on \(K/K^1\)
induced from \(P(\mathbb F_q)\) with coefficient space the
\((M\cap K^1)\)-invariants of the inducing representation.
Indeed, restriction to \(K\) determines the function, its value
is fixed by \(M\cap K^1\), and the induction modulus is one on
\(P\cap K\). Thus the finite induction cannot contain the
cuspidal label. This proves local supercuspidality at both
places. The noncuspidal terms in the automorphic spectral
decomposition are obtained from proper parabolic induction,
including subquotients at poles of Eisenstein series. A proper
global parabolic remains proper at either split place. The
supercuspidal local factor therefore excludes these terms.
\end{proof}

\begin{theorem}[The finite Hecke modules on the row sources]
\leavevmode\label{thm:commonsource806}
For \(V=Y_0,W_1,\ldots,W_m\), the spaces \(H_V\) are unitarily
equivalent modules for the finite convolution algebra
\(e\mathcal H(G_f,J)e\). Here \(G_f\) is the common finite
adelic group, \(\mathcal H(G_f,J)\) consists of compactly
supported \(J\)-bi-invariant functions, and \(e\) acts with
its idempotent normalization. In particular
\[
                         \dim H_{W_j}=\dim H_{Y_0}=:D .
\]
\end{theorem}
\begin{proof}
First take an automorphic constituent \(\pi\) on one of the
spaces. Lemma~\ref{lem:twoelliptic806} supplies two split
places where the local group is the general linear group and
the representation is supercuspidal. The rank \(s\) is odd.
Theorem~2.2.2 of Harris and Labesse thus gives a cuspidal weak
base change \(\Pi\) on \(\mathrm{GL}_s(\mathbb A_E)\). Its
unramified local parameters at almost all places of \(E\)
are the ordinary base changes of those of \(\pi\).
\cite{XIV-XI-harrislabesse806}

We identify the global parameter to which the multiplicity
formula applies. By Corollary~3.3.2 of Kaletha, M\'inguez,
Shin and White the unramified Hecke data of a discrete
constituent, with the ordinary base-change embedding, are
those of a global parameter for the quasi-split unitary
group. Its image on the general linear group is an isobaric
sum of cuspidal representations and their real determinant
twists. The unramified data just computed agree with those
of \(\Pi\) outside a finite set. Strong multiplicity one
in the isobaric form identifies that entire sum with \(\Pi\).
In particular there is one cuspidal summand with multiplicity
one and the auxiliary \(\mathrm{SL}_2\) is trivial. Indeed, a
block of length \(a\) contributes \(a\) twists with successive
exponents \((a-1)/2,\ldots,(1-a)/2\), and the cuspidal
support of \(\Pi\) consists of its single rank \(s\) summand.
This is a simple generic global parameter. The correct
unitary-group parity is included in the support statement of
Corollary~3.3.2. \cite[\S\S1.3.2, 3.3]{XIV-XI-innerunitary806}

Each group here is a pure inner form of the quasi-split
unitary group. Over an algebraic closure an isometry from a
fixed reference Hermitian space to \(V\) gives the inner
isomorphism on groups, and its Galois coboundary takes values
in the reference unitary group itself. For a simple generic
parameter on pure inner forms the global multiplicity formula
is unconditional. In this case the component group governing
the multiplicity is trivial. The commutant of the single
summand is scalar, and conjugate duality fixes among these
scalars just \(1,-1\). Dividing by the Galois-fixed center
removes both. Equivalently the local packet characters multiply
to one on this center by the global product formula. Hence
every restricted tensor product of local packet members for
the given inner form occurs once. \cite[\S1.7, Theorem~1.7.1
and the paragraph following it]{XIV-XI-innerunitary806}

At every finite place use the isometry of framed Hermitian
spaces above and keep the same local representation. At a
real place the unsigned Harish--Chandra parameters of the
row sources agree. In fact the first \(k_x-1\) entries run
from \(k_x/2-1\) down to \(1-k_x/2\); the last \(m\) entries
are the half-gap minus the successive lift weights. These
decrease by the positive gap, with the first below \(1-k_x/2\)
by the choice of the gap and weights. Thus all entries are
distinct. The prescription for a row assigns the entries to
positive and negative coordinates of its Hermitian space.
For \(U(p,q)\) a regular discrete parameter has precisely the
discrete series obtained by assigning its distinct entries to
the \(p\) positive and \(q\) negative coordinates, modulo
permutation inside each group of coordinates. This is the
compact-Cartan description of the real discrete-series
packet. In particular the prescribed member on the second
space belongs to this same local packet. Using the conjugate
coefficient convention for all rows dualizes the entire list
and gives the same comparison. Replacing the real components
by these chosen members therefore gives an automorphic
constituent on the second space, with multiplicity one. The
reverse replacement gives the inverse correspondence.

The coefficient projection on the fixed real line has rank
one within each such copy by Schur orthogonality, as in
Lemma~\ref{lem:bilinearsource710}. A finite representation
\(\pi_f\) contributes exactly \(e\pi_f^J\). Its nearly
everywhere finite data determine its global parameter by
the same isobaric uniqueness, so counting its contribution
cannot count two distinct global parameters. We have obtained
the same finite direct sum of the spaces \(e\pi_f^J\) on
each row. The finite convolution action on a summand is the
ordinary compressed action in \(\pi_f\). The invariant
Hilbert structures of equivalent irreducible unitary
representations agree up to a positive scalar under an
intertwiner, and rescaling gives a unitary identification.
Taking the direct sum proves the assertion.
\end{proof}

\begin{corollary}[The exponent in the product of auxiliary determinants]
\leavevmode\label{cor:scalarpower806}
Use the characters and the set \(\mathcal Q\) of
Theorem~\ref{thm:auxdet805}. On \(H_{Y_0}\) write
\(R^0_{j,\beta}\) for the compression of the normalized finite
convolution for \((j,\beta)\), transported in the finite frames.
At every sufficiently large node,
\[
 \prod_{j=1}^m\frac{\mathcal V_j(\vartheta^j)}{\mathcal V_j(1)}
    = C_{\mathcal Q}^{\,mD}
       \frac{\det(R^0_{1,\vartheta}\cdots R^0_{m,\vartheta^m})}
            {\det(R^0_{1,1}\cdots R^0_{m,1})}.
\]
The absolute logarithm of the determinant quotient on the
right is at most
\(mD\log((1+\varepsilon_s)/(1-\varepsilon_s))\), where
\(\varepsilon_s=512[F:\mathbb Q]\,2^{-s}\).
\end{corollary}
\begin{proof}
In the normalized finite convolution of
Theorem~\ref{thm:finitegram709}, finite partial sums of double
cosets, averaged on each side by \(J\) and compressed by \(e\),
belong to \(e\mathcal H(G_f,J)e\). The absolute tail bound of
Lemma~\ref{lem:vacuumtail709} gives norm convergence. Therefore
the equivalence of Theorem~\ref{thm:commonsource806} intertwines
the limits as well. Apply Theorem~\ref{thm:auxdet805} for each
\(j\), use \(d_j=D\) in its scalar, and multiply. Determinants
are multiplicative for endomorphisms of one finite-dimensional
space even when the operators do not commute. The interval for
the eigenvalues of each factor gives the asserted logarithmic
bound; with the empty determinant convention the identity also
holds when \(D=0\).
\end{proof}
This identifies the power of \(C_{\mathcal Q}\) contributed by
the auxiliary lifts. To cancel it in the comparison of periods
one must compute the Gram determinant of the composite map
from \(Y_0\) to \(Y_m\) in its own lattice normalization.
\clearpage

\subsection{The terminal flag and the remaining Gram determinant}
\label{sec:terminalflag807}
Keep the short-family node and the notation of
Theorem~\ref{thm:commonsource806}. Set \(n_i=s+it\), \(0\le i\le m\).
Let \(\mathcal M\) be the set of old marked places of \(F\),
\(\mathcal Q\) the set supporting the tame twist, and \(q_w\) the
residue cardinality. Each of these places splits in \(E\); use the
prescribed member of each split pair. Write
\[
\begin{split}
  g_b(q)&=|\mathrm{GL}_b(\mathbb F_q)|,\qquad
  N_i(\vartheta)=\prod_{w\in\mathcal M\cup\mathcal Q}g_{n_i}(q_w),\\
  M(\vartheta)&=\prod_{w\in\mathcal M\cup\mathcal Q}q_w^{s(s+t)},
  \qquad Q_*=\prod_{w\in\mathcal Q}q_w .
\end{split}
\]
Here \(N_i\) is the index of the principal subgroup in the
maximal compact at row \(i\), taken over these places.
\(M(\vartheta)\) counts matrices with \(s+t\) rows and \(s\)
columns over the same product of residue fields.

The point test on the edge from \(Y_{i-1}\) to \(Y_i\) specifies
the standard inclusion on reduction. At a new twist place it
is convenient to write \(\psi_i\) for the tame character on the
determinants of that edge. Use the split matrix action with
target determinant character \(\psi_i\) and source determinant
character \(\psi_i^{-1}\). If the integral changes of coordinates
on target and source have reductions \(h_i,h_{i-1}\) satisfying
\(h_i\iota_i=\iota_i h_{i-1}\), they preserve this point test
up to \(\psi_i(\det h_i/\det h_{i-1})\). Indeed, the matrix
variable is changed by left and right multiplication. The displayed
condition fixes its prescribed residue, the lattice is preserved,
and both determinant absolute values are one. Only the indicated
splitting characters remain. Thus the action on a terminal
coefficient can be followed back to the initial coefficient by
changing the variable on each source.

At \(w\in\mathcal Q\) prescribe the parabolic \(P_m(q_w)\)
preserving the flag of dimensions \(n_0,\ldots,n_{m-1}\)
inside \(\mathbb F_{q_w}^{\,n_m}\). If \(h_m\) preserves the flag,
take \(h_i\) to be its successive restrictions. Arbitrary
invertible lifts of those restrictions to the compact groups
suffice, since the test only sees their reduction. The determinant
quotient at step \(i\) is the determinant on the \(i\)-th
successive quotient of dimension \(t\). The initial compact
change fixes the initial form at this new place. The resulting
terminal character is therefore the product of the \(\psi_i\)
of these \(m\) quotient determinants. At \(w\in\mathcal M\) take
the identity reduction at the terminal level; all restrictions
then have identity reduction and the initial principal level
fixes the form. Denote by \(J_{\mathrm{end}}(\vartheta)\) this
specified compact level, with the full lattice stabilizer at
the remaining finite places, and by \(I_\vartheta\) its index
in the product of the maximal compacts. This description
specifies a level even for a vector that happens to have
additional symmetry.

\begin{proposition}[The metric at the terminal compact level]
\leavevmode\label{prop:terminalmetric807}
Give the terminal group at infinity and its finite maximal
compacts volume one. Let \(d_\infty\) be the dimension of the
fixed irreducible compact representation at infinity. If \(z,z'\)
are terminal vector-valued forms of this weight and the specified
unitary level character, and \(u,u'\) are their scalar
coefficients on the same unit highest-weight vector, then
\[
       \langle z,z'\rangle_{\rm cl}
            =d_\infty I_\vartheta\,\langle u,u'\rangle_{\rm Pet}.
\]
\end{proposition}
\begin{proof}
Over a finite double class of representative \(c\), the quotient
measure divides the compact measure by the order of
\(\Gamma_c\), its rational stabilizer in the real compact
group times \(cJ_{\mathrm{end}}c^{-1}\). The stabilizer is
finite since the rational group is discrete and this product
is compact. The finite volume factor is \(1/I_\vartheta\).
The terminal character has modulus one, so integration over
the finite level introduces no additional factor in the paired
coefficients. Averaging the rank-one projector of the unit
vector over the real compact group gives \(d_\infty^{-1}\)
times the identity by Schur orthogonality. Thus the scalar
Petersson pairing above that class is
\(\langle z(c),z'(c)\rangle/(d_\infty I_\vartheta|\Gamma_c|)\).
Summing over the classes proves the formula with the reciprocal
stabilizer weights of the class-sum pairing.
\end{proof}

\begin{theorem}[The normalized terminal index]
\leavevmode\label{thm:terminalflag807}
Put \(u_m=mst+\binom{m}{2}t^2\), and set
\[
       L_\vartheta =
       I_\vartheta
          \frac{N_0(\vartheta)N_1(\vartheta)^m}
               {N_m(\vartheta)M(\vartheta)^m}.
\]
For \(C_{\mathcal Q}\) as in Theorem~\ref{thm:auxdet805},
\[
                \frac{L_\vartheta}{L_1}
                         =C_{\mathcal Q}^{\,m}\,Q_*^{-u_m}.
\]
In particular at the coefficient prime \(\ell\) of the
simultaneous twists the valuation of this quotient is exactly
\(m\,v_\ell(C_{\mathcal Q})\).
\end{theorem}
\begin{proof}
The successive diagonal blocks of \(P_m(q)\) have sizes
\(s,t,\ldots,t\). Its diagonal choices number \(g_s(q)g_t(q)^m\).
The entries strictly above these blocks can be chosen freely.
There are \(mst\) such entries in the first block row and
\(\binom{m}{2}t^2\) in the rest. Consequently
\[
       [\mathrm{GL}_{s+mt}(\mathbb F_q):P_m(q)]
          =\frac{g_{s+mt}(q)}{g_s(q)g_t(q)^m q^{u_m}} .
\]
Reduction of the split integral group is onto, since a lift of a
matrix invertible on reduction has determinant a unit. At each
old place the terminal index is \(g_{s+mt}(q)\), the same for
every twist. The finite level is a direct product, so all its
old factors cancel from \(L_\vartheta/L_1\).

At one new place the factor remaining in this quotient is
\[
     \frac{g_{s+mt}}{g_sg_t^m q^{u_m}}\,
     \frac{g_sg_{s+t}^m}{g_{s+mt}q^{ms(s+t)}}
       =q^{-u_m}\left(\frac{g_{s+t}}{g_t q^{s(s+t)}}\right)^m .
\]
Counting ordered bases gives
\(g_b(q)=q^{b^2}\prod_{a=1}^b(1-q^{-a})\).
The expression in parentheses is therefore
\(q^{st}\prod_{a=t+1}^{t+s}(1-q^{-a})\), exactly the local
factor of \(C_{\mathcal Q}\). Multiply over the new places.
At a place of tame \(\ell\)-power order \(q_w\) is prime to
\(\ell\), so \(v_\ell(Q_*)=0\). This gives the last assertion.
\end{proof}

\begin{corollary}[Reduction to the Gram form with fixed Haar measures]
\leavevmode\label{cor:remaininggram807}
Take the scalar row integrals in the fixed Haar measures,
using on the \(i\)-th edge the prescribed principal-level
weight \(N_{i-1}(\vartheta)\). Write \(U_\vartheta\) for this
weighted composite, before conversion of its terminal
coefficient to a vector-valued form. On \(H_{Y_0}\) write
\[
 \begin{split}
  K_\vartheta
     &=\frac{N_m(\vartheta)M(\vartheta)^m}
             {N_0(\vartheta)N_1(\vartheta)^m}
                    U_\vartheta^* U_\vartheta,\\
  S_\vartheta
     &=\prod_{i=1}^m \det R^0_{i,\vartheta^i},\qquad
  A_\vartheta=\prod_{i=1}^m\mathcal V_i(\vartheta^i).
 \end{split}
\]
In the first line the adjoint uses the terminal scalar
Petersson pairing. Whenever the row integrals give
square-integrable terminal coefficients on this source, let
\(Z_\vartheta\) be their Gram determinant in the vector-valued
terminal class sum. Then
\[
       \frac{Z_\vartheta}{A_\vartheta}
         =(d_\infty L_1)^D\frac{S_1}{A_1}\,
            Q_*^{-D u_m}\,\frac{\det K_\vartheta}{S_\vartheta}.
\]
The identity permits a zero composite determinant.
\end{corollary}
\begin{proof}
Proposition~\ref{prop:terminalmetric807} identifies the
operator of the terminal class-sum pairing with
\(d_\infty I_\vartheta U_\vartheta^*U_\vartheta
=d_\infty L_\vartheta K_\vartheta\).
Thus \(Z_\vartheta=(d_\infty L_\vartheta)^D\det K_\vartheta\).
By Corollary~\ref{cor:scalarpower806},
\(A_\vartheta/A_1=C_{\mathcal Q}^{mD} S_\vartheta/S_1\).
The factors \(A_1,A_\vartheta,S_1,S_\vartheta\) are positive
by Theorem~\ref{thm:auxdet805}. Substitute the terminal index
formula and divide by \(A_\vartheta\).
\end{proof}
The finite flag calculation cancels the entire explicit factor
\(C_{\mathcal Q}^{mD}\) in this ratio. The remaining expression
contains the actual operator \(K_\vartheta\); an estimate or
identity for its determinant relative to \(S_\vartheta\) still
requires control of the intervening theta integrals.
\clearpage

\subsection{Split doubling factors along the row}
\label{sec:flagdoubling808}
At a place \(w\) of \(\mathcal M\cup\mathcal Q\) put
\(q=q_w\), and write \(P_b=\prod_{a=1}^b(1-q^{-a})\), with
\(P_0=1\). For a row edge of ranks \(n,n+t\) set
\[
      H_{n,t}(q)=\frac{P_{n+t}}{P_t}
                    =\prod_{a=1}^{n}(1-q^{-(t+a)}).
\]
This is the inverse split spherical doubling denominator at
the theta point. In the normalization of
Theorem~\ref{thm:inert175} its factors at parameter \(z\)
are \(1-q^{-(2z+a)}\) for \(1\le a\le n\), since the quadratic
character of \(E/F\) is trivial at a split place. Here
\(z=t/2\). We use \(H_{n,t}\) only as an explicit rational
scalar. The local integral for a ramified test is computed
directly below.

\begin{lemma}[The principal-level weight in the doubled integral]
\leavevmode\label{lem:pointdoubling808}
Use the specified inclusion-point datum on the row edge
of ranks \(n,n+t\). On the principal-level fixed space the
split local doubled pairing at the theta point, including
the two copies of the local edge weight \(g_n(q)\), is the
original pairing times
\[
                    B_{n,t}(q)=g_n(q)q^{-n(n+t)}
                                   =q^{-nt}P_n .
\]
This holds at every stage, also when the source carries a
nontrivial character on the larger flag subgroup.
\end{lemma}
\begin{proof}
In the split matrix model the test is the indicator of
\(X\equiv\binom{1_n}{0}\) in the integral
\((n+t)\)-by-\(n\) matrices. If \(X\) and \(Xg\) both lie in
its support, their top square blocks belong to
\(1+\varpi M_n(\mathcal O_{F_w})=:J_n\).
The quotient of those blocks is \(g\), so \(g\in J_n\).
Conversely \(J_n\) preserves that residue class and its
determinants are in the kernel of the tame character.
Its action leaves the indicator fixed. The squared norm
of the indicator is the additive volume \(q^{-n(n+t)}\)
of one residue class. Thus the matrix coefficient of the
test is \(q^{-n(n+t)}{\bf1}_{J_n}\).

The local doubling section at the theta point restricts
on the open orbit to this oscillator coefficient, as in
Proposition~\ref{prop:split173}. On \(J_n\) the matrix
coefficient of two level-fixed source vectors is precisely
their inner product. With maximal compact volume one,
\(\operatorname{vol}(J_n)=1/g_n(q)\). The scalar for the
unweighted edge is consequently \(q^{-n(n+t)}/g_n(q)\).
Multiplication of each lift by its local edge weight
\(g_n(q)\) multiplies its polarized Gram form by \(g_n(q)^2\).
The output point test is fixed by the principal subgroup
on the next group, by the same reduction argument. Starting
with the initial fixed vectors therefore verifies the level
condition at each successive stage.
\end{proof}

\begin{theorem}[Cancellation along the complete split flag]
\leavevmode\label{thm:flagdoubling808}
Let \(n_i=s+it\), \(0\le i\le m\). Denote by \(I_w\) the local
factor of the specified terminal level index \(I_\vartheta\).
The product of the \(m\) local row factors in
Lemma~\ref{lem:pointdoubling808} obeys
\[
   \frac{I_w\displaystyle\prod_{i=0}^{m-1} B_{n_i,t}(q)}
        {\displaystyle\prod_{i=0}^{m-1} H_{n_i,t}(q)}
    =
    \begin{cases}
       1, & w\in\mathcal Q,\\
       |P_m(q)|, & w\in\mathcal M ,
    \end{cases}
\]
where \(P_m(q)\) in vertical bars denotes the flag subgroup
of Theorem~\ref{thm:terminalflag807}.
\end{theorem}
\begin{proof}
Each quotient in the product before multiplication by \(I_w\)
is \(q^{-n_it}P_{n_i}P_t/P_{n_{i+1}}\). The middle products
cancel, giving
\[
         \prod_{i=0}^{m-1}\frac{B_{n_i,t}}{H_{n_i,t}}
            =q^{-u_m}\frac{P_s P_t^m}{P_{s+mt}},
       \qquad u_m=\sum_{i=0}^{m-1}n_it
                         =mst+\binom{m}{2}t^2 .
\]
At a new place \(I_w=g_{s+mt}/(g_sg_t^m q^{u_m})\).
Substitution of \(g_a=q^{a^2}P_a\) gives exactly the inverse
of the last display, since
\((s+mt)^2-s^2-mt^2=2u_m\). At an old place the prescribed
level is principal, so \(I_w=g_{s+mt}\). This increases the
preceding result by \(|P_m(q)|\).
\end{proof}

\begin{corollary}[Normalized comparison with all single lifts at the marks]
\leavevmode\label{cor:flagunit808}
Use the unscaled auxiliary source integrals and the terminal
class measures of Theorem~\ref{thm:auxdet805}. At each marked
place divide the product of the weighted row factors with its
terminal index by the product of the \(m\) auxiliary local
factors. Then divide by the ratio
\(\prod_{i=0}^{m-1}H_{n_i,t}(q)/H_{s,t}(q)^m\) of inverse
spherical doubling denominators. The resulting scalar is
\[
       q^{-mst}\quad(w\in\mathcal Q),\qquad
       |P_m(q)|\,g_s(q)^m\quad(w\in\mathcal M).
\]
At each node the old factor is identical before and after
twisting, and the product of the new normalized factors has
coefficient-prime valuation zero, including after taking the
\(D\)-th power.
\end{corollary}
\begin{proof}
Lemma~\ref{lem:auxmass805} gives
\(q^{st}H_{s,t}(q)\) for one auxiliary local factor at a
new place and \(H_{s,t}(q)/g_s(q)\) at an old one. The same
set of new places occurs for each of the \(m\) characters.
Divide the formula of Theorem~\ref{thm:flagdoubling808} by
the \(m\)-th power of the corresponding normalized auxiliary
factor. The new factors multiply to \(Q_*^{-mst}\), and
\(\ell\nmid Q_*\) for the simultaneous twists at \(\ell\).
\end{proof}
This calculation evaluates the marked local data in the
doubled pairings. To apply it to \(\det K_\vartheta\) in
Corollary~\ref{cor:remaininggram807} one still needs the
global inner product identities for the successive prescribed
theta vectors, including their archimedean factors.
\clearpage

\subsection{Cuspidality and the real pairings of the prescribed row}
\label{sec:actualrow809}
Fix one node of the shorter intervals of
Corollary~\ref{cor:fieldrelations712}, with its original marked level
and any one of the twists of Corollary~\ref{cor:actualdenom803}.
Use the row \(Y_0,\ldots,Y_m\) and the point tests and measures of
Corollary~\ref{cor:remaininggram807}. Write \(s=k+m-1\), \(t=5s\),
\(n_j=s+jt\), \(q_j=\lfloor(m-j)/2\rfloor\), and \(b=(t-1)/2\).
Both \(k\) and \(m\) are even. The source at step \(j+1\) has real
signature \((n_j-q_j,q_j)\). Write \(H=s^8\) for the half-gap.
At a chosen complex embedding above a real place the weights of the
successive characters are \(-w_a\), where
\(w_a=\kappa+(2a+1)H\) and \(\kappa>0\) is the weight at that
embedding. The same field, point of the interval and embedding are
used throughout the calculation.

\begin{lemma}[First occurrence for the bilinear row kernel]
\leavevmode\label{lem:rowreal809}
Let \(h_j=-\tfrac12{\bf1}_{j\ {\rm odd}}\). Form the real parameter
\(P_j\) by adding \(h_j\) to each of the following entries.
\begin{enumerate}
\item The \(k-1\) consecutive integers from \(k/2-1\) down to
\(1-k/2\) have plus sign.
\item The entries \(H-w_a\), \(1\leq a\leq m\), have plus sign
through \(a=j+1\). The subsequent signs alternate, starting with minus.
\item For \(1\leq a\leq j\), the \(t\) entries
\(-w_a+b,\ldots,-w_a-b\) all have plus sign.
\end{enumerate}
Within each of the two compact blocks put the entries in decreasing
order. For a physical input of parameter \(P_j\), the first possible
nonzero real lift in the target Witt tower is at \(Y_{j+1}\).
The local correspondence there gives parameter \(P_{j+1}\) on the
physical target. This holds for all sufficiently distant nodes, for
every \(0\leq j<m\) and every specified twist.
\end{lemma}
\begin{proof}
The base weights are the centered \(k-1\) consecutive integers and
the \(m\) shifted weights \(H-w_a\). After \(j\) insertions the
underlying list has acquired the \(j\) displayed strings.
On the \(a\)-th shifted weight put \(e=a-1\). The row sign is
\((-1)^{e+\min(e,j)}\), giving exactly the signs in the statement.
The number of minus signs is \(q_j\). Before the parity shift all
entries are integers. Each inserted string lies between \(H-w_a\)
and \(H-w_{a+1}\) when both are present, since the distance from
its center to either neighbor is \(H>b\). The first shifted weight
is below the last consecutive base weight. Thus the parameter is
regular and its parity is \((n_j-1)/2\bmod\ZZ\).

We check how this signed list enters the local correspondence.
On the edge put \(n=n_j\), \(N=n+t\), and
\(\nu=\mu^0_{j+1}\vartheta^{j+1}\). On the norm-one torus the
rank-\(d\) splitting datum is \(\nu\upsilon^{\lfloor d/2\rfloor}\),
where \(\upsilon\) in this expression is restricted to that torus.
Its character on the quadratic field is
\(\chi_d=\nu^{\rm BC}\upsilon^{-(d\bmod 2)}\).
Indeed base change on characters of the torus is pullback by
\(z\mapsto z/\bar z\). Conjugate self-duality takes the restricted
\(\upsilon\) to \(\upsilon^2\), and the space character also has
the factor \(\upsilon^{-d}\). The finite-order twist is trivial
on the connected real norm-one torus. At the chosen embedding
\(\upsilon\) has angular exponent one. The half angular exponents
of \(\chi_N,\chi_n\) are therefore
\[
              a_T=-w_{j+1}+h_j,\qquad
              a_S=-w_{j+1}+h_{j+1}.
\]

The theta integral multiplies the input form by the kernel.
Thus the source quotient of the oscillator is the
contragredient of the physical input. In matrix coordinates the
action is \(X\mapsto hXg^{-1}\). Write its tensor space as
\(Y_{j+1}\otimes Y_j^\vee\). The source acts on the second factor by
\({}^t g^{-1}\), which inverts the compact torus coordinates.
Transporting the contragredient to this dual space therefore
restores \(P_j\), with the compact coordinates sorted in decreasing
order. The target torus has its original coordinates.

One can check the additive-character orientation and both twists
on the Fock vacuum. Put \((p,q)=(n-q_j,q_j)\) and
\((r,u)=(N-q_{j+1},q_{j+1})\). In a target row of sign \(\epsilon\)
the determinant shift is \(\lfloor n/2\rfloor-w_{j+1}\).
Each of the \((n-\epsilon(p-q))/2\) oppositely signed columns
contributes \(-1\) from normal ordering. The vacuum weight there
is \(a_S+\epsilon(p-q)/2\). In a physical source column of sign
\(\epsilon'\) the derivative of right inverse multiplication gives
the opposite weight. Its constant before normal ordering is
\(N-\lfloor N/2\rfloor+w_{j+1}\), and one subtracts the number
\((N+\epsilon'(r-u))/2\) of target rows of the same sign.
In dual source coordinates the result is
\(a_T+\epsilon'(r-u)/2\). A Fock variable has weight \(+1\) on
both corresponding tensor coordinates for equal signs and \(-1\)
on both for opposite signs. These are precisely the compact
action and vacuum characters for Ichino's negative real additive
character. This checks the polarisation on the tensor space.
For a fixed symplectic action two unitary splittings can differ
only by a character. On these connected unitary groups such a
character is determined on the compact tori. See
\href{https://arxiv.org/pdf/2002.09148}{Ichino,
\S4.1 and Lemma~7.8}.

Subtract \(a_T\) from \(P_j\). The consecutive base entries are
positive and greater than \(H\), since \(w_{j+1}\geq 3H\) and
\(H>k\). The \(a\)-th shifted entry becomes
\((2j+3-2a)H\). An old string is centered at \(2(j+1-a)H\)
and has radius \(b\). Hence all centered entries have absolute
value at least \(H>b+1\). There are \(A=k+j(t+1)\) positive
entries, all of plus sign. The negative entries have alternating
signs. Among them \(y=q_{j+1}\) have plus sign and \(z=q_j\)
have minus sign, so \(A+y+z=n\).

Apply Atobe's first-occurrence criterion in the form of
\href{https://arxiv.org/pdf/2002.09148}{Ichino, Theorem~4.3(i)}.
The centered list is integral and misses zero, so its alternating
central string has parameter \(k_\lambda=-1\).
The two signature counts are \(r_\lambda=A+z\) and
\(s_\lambda=y\). Membership of either exceptional counting set
at \(b\) requires a centered entry with absolute value between
\(1\) and \(b\). Thus both sets are empty. The required signature
is \((r_\lambda+2b+1,s_\lambda)=(N-q_{j+1},q_{j+1})\),
with the common displacement \(l=0\) in the criterion.
Deleting \(h\) hyperbolic planes preserves the rank parity and
the splitting character and replaces \(l\) by \(-h\). For every
\(1\leq h\leq q_{j+1}\) this violates \(l\geq0\).

Finally, in the increasing-rank correspondence of Ichino's
Theorem~4.1(i), old positive centered entries keep their signs
and old negative entries change sign. The middle Levi factor
has signature \((t,0)\). Every other factor is a compact line.
For an old centered entry \(\xi\) in combined position \(i\)
the inducing weight is \(\xi+a_S-((N+1)/2-i)\).
The middle weight is \(A-n/2+a_S\). Adding the compact-Cartan
half sum of roots, whose \(i\)-th entry is \((N+1)/2-i\),
restores \(\xi+a_S\) and gives \(a_S+(t+1)/2-h\) in middle
position \(A+h\), \(1\leq h\leq t\). The strict gaps on both
sides of this middle string and the strict order of the old
entries give the regular good range. These computations are also
the compact-Levi calculation of Theorem~\ref{thm:tower169}.
The changed signs belong exactly to the shifted entries with
\(a\geq j+2\). We have obtained the stated list for \(j+1\).
\end{proof}

\begin{theorem}[Cuspidality of the prescribed composite]
\leavevmode\label{thm:rowcuspidal809}
On \(H_{Y_0}\) of Theorem~\ref{thm:commonsource806} all the
successive bilinear row integrals exist. Every intermediate output
is cuspidal, with real types \(P_j\) wherever it is nonzero.
The weighted terminal coefficient defining \(U_\vartheta\) belongs
to the Petersson space of \(U(Y_m)\). The doubled theta integrals
on these edges are in Weil's convergent range.
\end{theorem}
\begin{proof}
Lemma~\ref{lem:twoelliptic806} gives cuspidality at \(Y_0\).
For a smooth compact-finite cusp form, the rapid-decay theorem
on Siegel sets gives decay faster than every power of their
height. A Schwartz theta kernel has at most polynomial growth
in that height, locally uniformly in the other argument. The
center of each unitary group here is the anisotropic norm-one
torus. Consequently the defining integral converges absolutely
for a fixed target argument. The same bounds apply to derivatives
of the Schwartz test, giving a smooth lift. Multiplication by
the finite principal-level index is harmless for this assertion.

Consider one irreducible cuspidal input and any smaller rational
member in the target tower. Removing a rational hyperbolic plane
removes one positive and one negative direction at every real
place. A nonzero value of a smaller global lift would give a
nonzero invariant pairing of the local oscillator with the
physical input at such a place. Indeed one may translate the
test to evaluate at the target identity, take pure data by
linearity, and hold every other local factor fixed. Right
invariance of the source integral gives this pairing. On
compact-finite vectors it yields a map to the contragredient
Harish-Chandra module of the source. A nonzero map is onto,
since this module is irreducible. The maximal source-isotypic
quotient of the oscillator would then be nonzero and would
have a nonzero irreducible quotient on the smaller target by
local Howe duality. Lemma~\ref{lem:rowreal809} rules this out.

Rallis's tower property now gives vanishing of all proper
unipotent constant terms of the lift at the specified target.
The hypothesis of that theorem is the vanishing of the lower
global lifts of the whole source representation, which was
just verified for arbitrary Schwartz tests. In particular the
conclusion applies to our prescribed point test and highest
polynomial. See \href{https://arxiv.org/pdf/1207.4709}
{Gan, Qiu and Takeda, \S11.2}. For a cuspidal target constituent
with nonzero projection, the projected theta map gives a quotient
of the local oscillator with that constituent as target and
the contragredient physical input as source. This follows by
pairing the lift with a target cusp form and fixing all other
places as above. Local Howe duality and
Lemma~\ref{lem:rowreal809} give its real parameters \(P_{j+1}\).
The fixed target test is compact-finite, and its finite part
is fixed by an open compact subgroup. At fixed infinitesimal
character, compact types and finite level the cuspidal space
is finite dimensional by automorphic finiteness. Thus the
argument applies to the sum of the resulting constituents at
the next step. Induction proves cuspidality all the way through
\(Y_m\), including the possibility of a zero output.

For the last assertion let \(R\) be the rational Witt index of
\(Y_{j+1}\). Scalar extension of a rational isotropic space to a
real place gives \(R\leq q_{j+1}\leq m/2<t\). Thus \(N-R>n\).
For the doubled unitary source of split rank \(n\) this is
Weil's convergence inequality. It also applies at later steps
where the Siegel Eisenstein series needs continuation at
\((N-n)/2\). See \href{https://arxiv.org/pdf/1207.4709}
{Gan, Qiu and Takeda, \S\S3.1, 11.4}.
\end{proof}

\begin{proposition}[Positive real factors for the specified tests]
\leavevmode\label{prop:rowrealgram809}
At each step and real place fix the nonzero joint highest polynomial
\(\Phi\) and the prescribed unit input line. Let
\(\eta=P_j^\vee\) denote the physical contragredient parameter,
obtained by negating entries and sorting in each compact block,
and write \(c_\eta\) for the coefficient of its unit highest line.
With \(A_\Phi(g)=\langle\Phi,\omega(g,1)\Phi\rangle\), the integral
\[
        \int_{U(n-q_j,q_j)}
               A_\Phi(g)\overline{c_\eta(g)}\,dg
\]
converges absolutely and is strictly positive. For the same real
tests and measures it is independent of the finite twist.
\end{proposition}
\begin{proof}
We specify the polynomial line to which the assertion applies.
Order the \(A\) positive centered entries of the lemma
decreasingly, followed by the alternating negative tail, and
call this list \((\xi_i)\). In the construction of
Theorem~\ref{thm:harmonic188} use the three lists
\[
 \left(\xi_i+i-A-\tfrac{t+1}{2}\right)_{i=1}^{A},\qquad
 \left(-\xi_{A+2v}+\tfrac{t+1}{2}\right)_{v=y,\ldots,1},\qquad
 \left(-\xi_{A+2v-1}-\tfrac{t+1}{2}\right)_{v=z,\ldots,1}.
\]
They are partitions. For the first list successive centered
entries differ by at least one, and its last part is at least
\(H-(t+1)/2>0\). The other two lists have decreasing parts
in the indicated reverse order and the same positive lower
bound when nonempty. In the three disjoint squares of respective
sizes \(A,y,z\), take the product of the leading minors whose
exponents are successive differences of parts, reversing the
two index orders in each dual square. The squares use, in
tensor coordinates, the positive source and target blocks,
the positive source and negative target blocks, and the negative
source and positive target blocks, respectively. Their occupied
rows and columns are disjoint. The determinant highest-weight
calculation of Theorem~\ref{thm:harmonic188} applies since it
only uses the ordered entries, their signs and the middle string.
It gives the two lowest compact highest weights after the
vacuum shifts computed above. At the frame where the three
squares are identity matrices the product is one. The highest
line in the matching space of joint harmonics has dimension
one by \href{https://arxiv.org/pdf/2002.09148}{Ichino,
\S7.4, Lemma~7.8}. Thus a different choice of nonzero vector
on the prescribed line multiplies this polynomial by a nonzero
scalar. In physical source coordinates its compact weight is
that of \(\eta\). In particular
\(\partial_{ai}\partial_{ai'}\Phi=0\) for every target row \(a\)
and every pair of source columns \(i,i'\) of opposite signs.

Let \(l=\min(n-q_j,q_j)\) and let \(u_1\geq\cdots\geq u_l\geq0\)
be the Cartan coordinates. The oscillator calculation in
Lemma~\ref{lem:oscillatorcartan710} works row by row for any
target dimension \(N\). For a paired boost its lowering operator
is a sum of the indicated mixed second derivatives, which
annihilate all compact source translates of \(\Phi\). Pairing the
creation-operator expansion with another such vector leaves the
diagonal factor \((\cosh u_i)^{-N-\deg_i}\), where \(\deg_i\)
counts variables in the two columns of that boost. Hence
\[
 |A_\Phi(k_1a(u)k_2)|\leq \|\Phi\|^2\prod_i(\cosh u_i)^{-N},
 \qquad
 A_\Phi(ka(u)k^{-1})>0 .
\]
The strict inequality follows by summing the positive diagonal
factors against the squared monomial coefficients of \(k^{-1}\Phi\).

The Cartan density in that same lemma is at most a constant,
depending on \(n\), times \(\exp(2(n-1)\sum_i u_i)\).
It follows by integrating the squared bound that
\(A_\Phi(g)\exp(\delta\sum_i u_i)\) is square integrable for
\(0\leq\delta<N-n+1=t+1\). The discrete series coefficient
\(c_\eta\) is square integrable since the real center is compact.
Cauchy--Schwarz proves absolute convergence, including at all
ranks \(N\leq 2n\) occurring here. When \(l=0\) both assertions
of integrability are immediate from compactness.

For the sign of the answer, take polar coordinates \(g=kz\)
with \(z\) positive Hermitian. Truncate to
\(\log\|z\|_{\rm op}\leq B\), which is compact modulo \(K\).
Let \(\rho\) and \(\rho_c\) be the full and compact positive
half sums in the order of \(\eta\), and put
\(\alpha=\eta+\rho,\ \mu=\eta+\rho-2\rho_c\).
Proposition~\ref{prop:polarcoefficient697} expresses \(c_\eta\)
as an average of its compact highest-line coefficient with
denominator \(m_{-\alpha}\), the weight product of the principal
minors. Applying compact Schur orthogonality to this average,
as in the proof of Theorem~\ref{thm:realgram710}, evaluates the
truncated integral as
\[
     \frac{1}{\dim V_\mu}
       \int_{\log\|z\|_{\rm op}\leq B}
                    \frac{A_\Phi(z)}{m_{-\alpha}(z)}\,d\nu(z).
\]
Here \(K\) has mass one and \(\nu\) is the polar quotient measure.
Both compact averages are over a finite measure space and the
radial cutoff is invariant under compact conjugation. Thus this
identity requires no decay estimate. All principal minors of
\(z\) are positive. The integrand is continuous and positive at
the identity, where its numerator is \(\|\Phi\|^2\).
Taking \(B\) to infinity gives a strictly positive value by
monotone convergence on this side and absolute convergence on
the original side. All quantities in this integral depend only
on the specified real data.
\end{proof}

\begin{corollary}[The terminal Gram form on the shorter intervals]
\leavevmode\label{cor:rowdefined809}
At every node above, the determinant \(Z_\vartheta\) of
Corollary~\ref{cor:remaininggram807} is defined for the actual
weighted row, and the identity there holds. Every real factor
in the successive Petersson pairings is strictly positive.
\end{corollary}
\begin{proof}
Apply Theorem~\ref{thm:rowcuspidal809} to the specified point
tests. The terminal coefficient is square integrable, which
is the analytic hypothesis in Corollary~\ref{cor:remaininggram807}.
Apply Proposition~\ref{prop:rowrealgram809} at each real place.
The common-field comparison for the integral projection still
requires the arithmetic size of the resulting global
determinant relative to the auxiliary determinants.
\end{proof}
\clearpage

\subsection{Change of the row determinant under a tame twist}
\label{sec:relativechain810}
Use the fields, weights, measures and test vectors of
Section~\ref{sec:actualrow809}. Write \(\mathcal M\) for the old
marks of \(F\), \(\mathcal Q\) for the new ones, and
\(S=\mathcal M\cup\mathcal Q\). If \(S\) occurs as a superscript
on an \(L\)-function over \(E\), omit all places above \(S\).
All \(L\)-functions below have their archimedean factors removed.
Set
\[
 \lambda_i=(\mu_i^0\vartheta^i)^{\rm BC},\quad
 \rho_{ai}=\lambda_a/\lambda_i\quad(a<i),\quad
 D_i=\upsilon^{-{\bf1}_{i\ {\rm odd}}},\quad z_0=t/2 .
\]
For a local cardinality \(Q=\#k(w)\) put
\([t]_w=\{Q^{(t-1)/2-h}:0\le h<t\}\), counted with multiplicity.
The common finite source in Theorem~\ref{thm:commonsource806}
has dimension \(D\). The notation \(D_i\) always denotes a
character.

\begin{lemma}[The spherical factor on a row constituent]
\leavevmode\label{lem:sphericalchain810}
Start with an irreducible bottom constituent and write \(B_{0,w}\)
for its ordinary standard base-change Satake multiset at \(w\)
above an unmarked place. Whenever the first \(i\) prescribed row
lifts have given a nonzero constituent spherical there, its multiset is
\[
 B_{i,w}=D_i(\varpi_w)
       \left(B_{0,w}\ \sqcup\
        \bigsqcup_{a=1}^i \lambda_a(\varpi_w)[t]_w\right).
\]
At the edge from \(n_{i-1}\) to \(n_i\), its unmarked doubling
product equals the corresponding product for the auxiliary
lift of the same bottom finite system from \(W_i\), multiplied by
\[
 \frac{d_s^S(z)}{d_{n_{i-1}}^S(z)}
 \prod_{a<i}\prod_{h=-(t-1)/2}^{(t-1)/2}
       L_E^S(z+\tfrac12-h,\rho_{ai}),\qquad
 d_n^S(z)=\prod_{c=1}^n L_F^S(2z+c,\omega_{E/F}^{\,t+c}).
\]
The identity first holds for \(\Re z\) sufficiently large.
\end{lemma}
\begin{proof}
In the tensor convention of Lemma~\ref{lem:rowreal809} the
characters attached to source and target at edge \(i\) are
\(\chi_S=\lambda_iD_i\) and \(\chi_T=\lambda_iD_{i-1}\).
The subscripts here refer to the two spaces, so the target
character twists the source representation in the doubling
integral. For a nonzero spherical theta pair the transfer is
\[
 B_{i,w}=(\chi_S/\chi_T)(\varpi_w)B_{i-1,w}
                    \ \sqcup\ \chi_S(\varpi_w)[t]_w .
\]
At an inert place this is Theorem~\ref{thm:inert175}, which
uses the two strata of the line Jacquet module and permits
nontempered spherical inputs. At a split place one can read
the conventions directly on the matrices \(hXg^{-1}\).
The physical source quotient is contragredient. Removing its
inverse determinant twist gives the dual bottom representation
at that edge tensored with \(\chi_T\). In the normalized matrix
correspondence the target support is its contragredient together
with the support of the trivial representation of the difference
in ranks. Restoring \(\chi_S\) gives the displayed transfer,
with \([t]_w\) the support of that trivial representation.
Equivalently one can use Proposition~\ref{prop:split175}.
The other split component inverts the multiset and character
values, and \([t]_w\) is invariant under inversion.
Induction from \(D_0=1\) gives the first formula.

For normalized unit spherical vectors and the standard lattice
vacuum, the unramified doubling integral is the standard factor
of \(B_{i-1,w}/\chi_T(\varpi_w)\) at \(z+1/2\), divided, over
the one place of \(F\), by \(d_{n_{i-1},v}(z)\). This is the
spherical doubling calculation in
\href{https://arxiv.org/pdf/1207.4709}{Gan, Qiu and Takeda,
\S11.6}, in the character convention of
\href{https://arxiv.org/pdf/2002.09148}{Ichino, \S6.3}.
Inverting the list and the twist simultaneously gives the same
product over \(E\), by conjugation. The finite framed lattices
are unimodular, the CM extension is unramified at finite places,
and off \(S\) the splitting characters are unramified and the
tests are the unit self-dual lattice vacua, as in
Lemma~\ref{lem:vacuumtail709}. These facts are unchanged by
increasing the matrix ranks. In particular the section is
spherical with value one at the identity.

On the auxiliary pair the target of rank \(s+t\) has torus
character \(\mu_i^0\vartheta^i\upsilon^{\lfloor(s+t)/2\rfloor}\).
Since \(s+t\) is even its space character is \(\lambda_i\).
The auxiliary source splitting acts on the target variable and
does not occur in its source spherical doubling numerator. Thus
the auxiliary list is \(B_{0,w}/\lambda_i(\varpi_w)\).
Theorem~\ref{thm:commonsource806} matches these bottom systems
with their dimensions. Dividing the row list by \(\chi_T\)
exhibits this same auxiliary list and the additional segments
\(\rho_{ai}(\varpi_w)[t]_w\). Multiplying their Euler factors
gives the asserted formula.
\end{proof}

\begin{proposition}[The Petersson pairing on every prescribed row step]
\leavevmode\label{prop:chainpairing810}
The weighted row lifts are injective on the successive images of
\(H_{Y_0}\). On the part arising from a single bottom spherical
system, the squared Petersson norm at edge \(i\), divided by the
squared norm of the input, is a positive scalar. Its unmarked
factor is the product in Lemma~\ref{lem:sphericalchain810}
evaluated at \(z_0\), including the auxiliary spherical factor.
At a finite marked place \(v\) its factor is \(B_{n_{i-1},t}(q_v)\)
of Lemma~\ref{lem:pointdoubling808}. Its other factors consist of
the specified real integrals and a positive constant depending on
the spaces and the Haar normalizations, common to the twists.
\end{proposition}
\begin{proof}
Apply the polarized Rallis inner product formula to a cuspidal
input at the given step. The smaller lifts vanish and \(n_i-R>
n_{i-1}\) by Theorem~\ref{thm:rowcuspidal809}, where \(R\) is
the target Witt index. For \(n_i>2n_{i-1}\) one has the ordinary
convergent Siegel--Weil formula. In the range
\(n_{i-1}<n_i\leq2n_{i-1}\) one uses
\href{https://arxiv.org/pdf/1207.4709}{Gan, Qiu and Takeda,
Propositions~11.1(ii) and~11.2}. This expresses the pairing
by the specified Laurent coefficient at \(z_0\) of the
global doubling integral. In the initial half-plane its
open-orbit unfolding is Eulerian on pure tensors. All proper
orbits vanish by cuspidality. In particular, with mixed inputs
the unfolded group coefficient is
\(\langle f',R(g)f\rangle_{\rm Pet}\), up to the common choice
of conjugations. If \(f,f'\) lie in orthogonal invariant
subspaces, it vanishes identically for arbitrary tests.
Continuation preserves that orthogonality. The conventional
constant in this formula is fixed by the global measures and
the pair of spaces. Changing the two unitary splittings
tensors the oscillator by determinant characters. Such a
twist changes the integrands by precisely the displayed
character twists. It changes neither the Haar measures nor
the doubled Schwartz identification, which pairs a test
with its conjugate. In Tamagawa measures the constant in
Proposition~11.2 is \([E:F]\).

We verify regularity at the evaluation point. This matters
because an Euler product initially evaluated on the boundary
of its absolute convergence need not represent a Petersson
norm by mere substitution. At any finite place take a flat
extension of the local Siegel--Weil section. On the open orbit
its restriction at \(g\) equals its value at \(z_0\) times
the positive Siegel height \(h(g)^{z-z_0}\). At an unmarked
place \(h(g)\leq1\). Indeed take a primitive top exterior
vector of the graph Lagrangian \(\{(x,x)\}\) in the direct
sum of the two source lattices. Under \((g,1)^{-1}\) its
projection to the second exterior summand is still a lattice
generator. Its lattice supremum norm is therefore at least one.
By decomposition into the Siegel parabolic times the lattice
compact, that norm is the inverse absolute determinant of the
Levi acting on the graph. Over a split algebra use the product
of the two component absolute values. This gives the asserted
bound for the height. It equals one on the source compact.
At a real place the same argument with a Euclidean exterior
norm bounds the height above by a fixed constant, the norm
of the graph vector divided by that of its second projection,
using the corresponding complex absolute value.

Let \(C_i^S(z)\) be the unmarked auxiliary spherical product
for the bottom system at this edge. On \(\Re z\geq z_0\) the
local integrand defining this product is dominated by the
absolute value of the rank-\((s,s+t)\) vacuum coefficient,
because a unit-vector coefficient has modulus at most one.
The integral over the maximal compact is one.
Lemma~\ref{lem:vacuumtail709} bounds the difference of the
local integral from one by \(128q_v^{-s-2}\). These errors are
summable over the field and are individually less than \(1/2\)
at all sufficiently distant nodes. The product is holomorphic
on \(\Re z>z_0\) and continuous and nonzero at \(z_0\), by
dominated convergence locally and uniform convergence of the
Euler product. At \(z_0\) each factor is real by inversion of
the two unitary matrix coefficients. It is thus strictly
positive. This argument requires no temperedness of a lift.

The denominators \(d_n^S\) have absolutely convergent nonzero
Euler products near \(z_0\), since \(2z_0+1=t+1>1\). For
\(\Re z>z_0\) the extra factors of the lemma have arguments
of real part strictly greater than one. At \(z=z_0\) their
arguments are precisely \(1,\ldots,t\). Each \(\rho=\rho_{ai}\)
and \(\rho^2\) has a nonzero angular exponent at infinity
because \(w_i-w_a=2(i-a)H\). By Hecke's continuation theorem
both \(L\)-functions are holomorphic at one. To check nonvanishing
there, for real \(u>1\) the Euler logarithm gives
\[
 \zeta_E^S(u)^3\,|L_E^S(u,\rho)|^4 |L_E^S(u,\rho^2)|\geq1 .
\]
Each prime-power summand in its logarithm has the sign of
\(3+4\cos v+\cos(2v)=2(1+\cos v)^2\), with \(v\) the
argument of the character to that power. A zero at one would
force the left side to tend to zero, since the zeta pole is
simple. Moreover \(S\) is conjugation-stable over \(E\), and
\(\rho^c=\rho^{-1}\). The Euler factors at a pair of distinct
conjugate primes have positive product for real \(u>1\).
At a prime fixed by conjugation an unramified value is real
of modulus one. Hence \(L_E^S(u,\rho)>0\) for \(u>1\), and
also at one by the preceding argument.

At the marks the row section has compact support on the source,
and Lemma~\ref{lem:pointdoubling808} evaluates it with the
principal weight. At the real places Proposition~\ref{prop:rowrealgram809}
gives integrable absolute majorants at \(z_0\) and strictly
positive values for the specified vectors. The upper height
bound gives dominated convergence from the right on bounded
subsets of \(\Re z\ge z_0\). Consequently the factored expression
is holomorphic in the right half-plane adjoining \(z_0\) and
has a finite nonzero limit from the right at \(z_0\) for
a unit pure input. It agrees with the global meromorphic
integral by continuation from large \(\Re z\). A pole of that
integral could not have this finite limit. Its prescribed
Laurent coefficient is therefore its ordinary value.

The marked polarized pairing is scalar on the entire
principal-fixed space. At infinity a prescribed output has
the target compact highest weight carried by its test, since
right translation acts on the test. In the discrete series
identified in Lemma~\ref{lem:rowreal809} the highest line of
that lowest compact type has multiplicity one. Off the marks
the fixed line in a spherical irreducible representation is
one dimensional. By multilinearity the same scalar norm
formula holds on the specified vectors of an irreducible
cuspidal summand, including linear combinations of pure
tensors. Between orthogonal source summands the mixed
pairing vanishes as shown above, also after allowing group
translates of their lifts. Within a common bottom spherical
system the formula is independent of the individual summand.
Starting at \(Y_0\), the positive formula proves nonvanishing
on each nonzero input, and Lemma~\ref{lem:sphericalchain810}
then supplies the spherical list for the next edge. This
induction proves the assertion on the successive images.
\end{proof}

\begin{theorem}[Exact relative determinant of the weighted chain]
\leavevmode\label{thm:detquotient810}
Let \(Z_\vartheta\) be the genuine terminal class-sum Gram
determinant of the row and \(A_\vartheta\) the product of the
\(m\) auxiliary class-sum Gram determinants, in the
normalizations of Corollary~\ref{cor:remaininggram807}.
Put \(Q_*=\prod_{v\in\mathcal Q}q_v\), and write
\(\rho_{ai}^0=(\mu_a^0/\mu_i^0)^{\rm BC}\). Then all four
determinants below are positive and
\begin{equation}
 \frac{Z_\vartheta A_1}{Z_1 A_\vartheta}
  =\left[
      Q_*^{-mst}
      \prod_{1\leq a<i\leq m}\ \prod_{h=1}^{t}
       \frac{L_E^{\mathcal M\cup\mathcal Q}(h,\rho_{ai})}
            {L_E^{\mathcal M}(h,\rho_{ai}^0)}
    \right]^D .                                      \label{eq:relchain810}
\end{equation}
The empty determinant is one. The Euler sets in the numerator
and denominator are those of their respective lifts.
\end{theorem}
\begin{proof}
Use Petersson-orthonormal bottom bases adapted to the unmarked
finite systems. The unitary module identifications of
Theorem~\ref{thm:commonsource806} give the same multiplicities
on every auxiliary source. Proposition~\ref{prop:chainpairing810}
makes the Gram form of the composite on each system the
product of the scalar edge factors, and makes the mixed
Gram forms of the orthogonal constituents zero. Converting
the terminal coefficient into the class form multiplies this
scalar by \(d_\infty I_\vartheta\), by
Proposition~\ref{prop:terminalmetric807}. The auxiliary rank
\(s\) integrals have the factor \(C_i^S(z_0)\) by the same
spherical calculation, or directly by Theorem~\ref{thm:globalgram711}.
Their fixed real factors are positive. Each auxiliary source
integral is unweighted as in Theorem~\ref{thm:auxdet805}.
Thus all \(C_i^S(z_0)\) cancel system by system in
\(Z_\vartheta/A_\vartheta\). The remaining real factors,
highest-line dimensions and measure constants are independent
of the finite twist. The extra \(L\)-values of
Lemma~\ref{lem:sphericalchain810} do not depend on the bottom
constituent, so their product occurs to power \(D\).

It remains to keep track of the normalizing \(L\)-functions.
Let \(d_n(z_0)\) include every finite prime of \(F\). At a
marked split prime \(v\), \(H_{n,t}(q_v)=1/d_{n,v}(z_0)\).
Hence replacing \(d_s^S(z_0)/d_n^S(z_0)\) by
\(d_s(z_0)/d_n(z_0)\) leaves the additional local multipliers
\(H_{s,t}(q_v)/H_{n,t}(q_v)\) at \(v\in S\). They are precisely
the multipliers in Corollary~\ref{cor:flagunit808}. That
corollary includes the terminal level index, the weights of
all source integrals, and the class measures of the auxiliary
targets. Its normalized local ratio at \(v\in\mathcal M\) is
\(|P_m(q_v)|\,|\mathrm{GL}_s(k(v))|^m\), independently of the
twist. At \(v\in\mathcal Q\) it is \(q_v^{-mst}\).
The complete finite denominators and all the old multipliers
are identical for \(Z_1/A_1\). Dividing the two determinant
ratios proves \eqref{eq:relchain810}. Positivity follows from
Proposition~\ref{prop:chainpairing810} and
Theorem~\ref{thm:auxdet805}.
\end{proof}

\begin{corollary}[An interval bound for the relative determinant]
\leavevmode\label{cor:detheight810}
At a node \(x\) let \(T_{\ell,\iota,x}\) be the left side of
\eqref{eq:relchain810} for the simultaneous twist attached to the
coefficient prime \(\ell\) and embedding \(\iota\). It is a positive
algebraic number. Given \(\delta>0\), the twists can be chosen
together with the bounds of Corollary~\ref{cor:actualdenom803}
on the same shorter interval so that, with \(v_\ell(\ell)=1\),
\[
 \sum_\ell(\log\ell)\sup_{x,\iota}
       \frac{\max(0,v_\ell(\iota T_{\ell,\iota,x}))}
            {\max(1,D_x)}\ \leq\ \delta ,
       \qquad D_x=\dim H_{Y_0}.
\]
\end{corollary}
\begin{proof}
Each quotient in \eqref{eq:relchain810} is
\(R_J(\rho_{ai}^0,h,(\vartheta^{\rm BC})^{a-i})\)
of Section~\ref{sec:samecritical731}. Indeed the quotient
character has tame part \(\omega_w^{a-i}\) at every old mark
and is unramified away from the old marks before twisting.
The prescribed order of \(\omega_w\) has both a 2-part and a
3-part greater than \(m\), by
Corollary~\ref{cor:quotientchainheight734}. Thus these are
exactly its old ramified pairs. The twist powers ramify at
every new pair. The two old infinite weights differ in
absolute value by \(2(i-a)H\geq t\), so every \(1\leq h\leq t\)
is one of the critical points of that corollary. Its algebraicity
statement applies in the fixed complex embedding.
The remaining scalar \(Q_*^{-mst}\) is rational and is a unit
at \(\ell\) since the new conductor primes avoid \(\ell\).

Choose positive \(e_\ell\) with
\(\sum_\ell e_\ell\log\ell\leq\delta\). At a node use error
\(e_\ell/(1+t\binom m2)\) per critical factor. These tolerances
can be inserted in the simultaneous specialization used in
Theorem~\ref{thm:actualselmer803}. The two distinct prime
divisors of the tame orders make the old term of
Theorem~\ref{thm:criticalheights734} zero for every quotient
character at every coefficient prime. Hence the positive
valuation of the product per source dimension is at most
\(e_\ell\) at every node and embedding. If the dimension
vanishes the left side of \eqref{eq:relchain810} is one.
Taking the indicated supremum and sum proves the bound.
\end{proof}
This computes the change with twist of the full complex Gram
determinant in the given level metrics. The passage to the
specified integral projected period still requires an estimate
for the untwisted comparison in its arithmetic lattice and for
the divisibility of the projected integral evaluations.
\clearpage

\subsection{Integral gluing at the old marked places of a single lift}
\label{sec:integralgluing811}
The factors \(\mathcal V_j(\beta)\) in Theorem~\ref{thm:auxdet805}
are periods on the definite target of the lift from \(W_j\).
To compare these periods as the character varies, we compute some
of their evaluations on products of coefficient vectors. Fix one
of the fields \(E/F\) of Section~\ref{sec:commonsource806}, an
index \(1\le j\le m\), and either \(\beta=1\) or
\(\beta=\vartheta^j\) at this same node. Split the standard
definite target into coordinate subspaces of ranks
\[
             p=s+9m+1,\qquad d'=t-9m-1,\qquad n=p+d'=s+t .
\]
Both subspaces have room for \(s\) coordinate vectors since
\(t=5s,\ s\ge x^5\) and \(m\le2\sqrt{x}\) at sufficiently
large nodes. Denote their integral
coordinate inclusions by \(e\) and \(z\), respectively, with
\((e,0)\) the inclusion prescribed by the original lift. The
choice of the \(s\) distinct rows of \(z\) is arbitrary.

\begin{theorem}[The lattice and the disjoint residue classes]
\leavevmode\label{thm:splitlattice811}
At an old marked place \(v\) write \(O=\mathcal O_{F_v}\),
\(\pi\) for a uniformizer, and \(q=|O/\pi O|\). Use its
designated split component over \(E\). Put \(L_+=O^p\),
\(L_-=O^{d'}\), and \(L=L_+\oplus L_-\). For an arbitrary
\(3m\)-by-\(p\) matrix \(\Lambda\) modulo \(\pi^2\), let
\[
 K=\{u\in L_+:\bar\Lambda\bar u=0\},\quad
 r=\operatorname{rank}\bar\Lambda,\quad
 r_e=\operatorname{rank}(\bar\Lambda\bar e).
\]
There exists an integral matrix \(\gamma\), invertible over
\(F_v\), with \(\gamma(e,0)=(e,z)\) and
\(v_\pi(\det\gamma)=r+d'-s\). For a matrix \(X=(X_+,X_-)\)
of \(s\) columns in \(F_v^n\), its translated point test satisfies
\begin{equation}
\begin{split}
 {\bf1}_{(e,0)+\pi L^s}(\gamma^{-1}X)
  =\sum_{Y\in\Gamma}
   {\bf1}_{\,e+(\pi K)^s}(X_+Y)\,
   {\bf1}_{\,z+(\pi^2L_-)^s}(X_-Y),\\[-2pt]
 \Gamma=\ker\bigl(\GL_s(O/\pi^2)\longrightarrow\GL_s(O/\pi)\bigr).
\end{split}
\label{eq:localpartition811}
\end{equation}
Use any integral invertible lift of \(Y\). At most one summand
is nonzero. Each restriction of \(\gamma^{-1}\) to a coordinate
block has only \(1\) and \(\pi^{-1}\) as elementary divisors,
up to units. The multiplicities of \(\pi^{-1}\) are \(r\) on
\(L_+\) and \(d'-s+r_e\) on \(L_-\).
\end{theorem}
\begin{proof}
Take for the column lattice of \(\gamma\) the submodule
\[
 \mathcal L=\{(ec+u,zc+\pi v'):c\in O^s,\ u\in K,\ v'\in L_-\}.
\]
It contains \(\pi L\). Modulo \(\pi L\) the graph
\((\bar e,\bar z)(O/\pi)^s\) meets \(\ker\bar\Lambda\oplus0\)
in zero, since \(\bar z\) is injective. These spaces span
\(\mathcal L/\pi L\), of dimension \(s+p-r\).
Thus \([L:\mathcal L]=q^{r+d'-s}\). Projection onto the
\(z\)-rows of the second summand restricts on \(\mathcal L\)
to a retraction onto the coordinates of \((e,z)O^s\). Its
kernel is free over the discrete valuation ring. A basis of
the kernel completes these \(s\) columns to a basis of
\(\mathcal L\). Place the specified columns at the positions
of \(e\) to obtain \(\gamma\).

The left side of \eqref{eq:localpartition811} is one exactly
on \((e,z)+\pi\mathcal L^s\). On this set write
\[
       X_+=e(I+\pi C)+\pi U,\qquad
       X_-=z(I+\pi C)+\pi^2 V,\qquad U\in K^s .
\]
Right multiplication by \(Y=(I+\pi C)^{-1}\bmod\pi^2\)
gives both required congruences, using \(\pi L_+\subset K\).
Conversely these congruences give the displayed expression by
right multiplication by a lift of \(Y^{-1}\). The square
submatrix formed by the \(z\)-rows of \(X_-\) must reduce to
\(Y^{-1}\bmod\pi^2\), which proves uniqueness. Both sides
vanish off \(L^s\), and changing the chosen lift of \(Y\) gives
the same indicators.

The map \(\pi\gamma^{-1}\) is integral. For \(u\in L_+\) the
vector \(\gamma^{-1}(u,0)\) is integral exactly when
\((u,0)\in\mathcal L\), or equivalently \(u\in K\).
For \(v'\in L_-\) the corresponding condition is
\(\bar v'=\bar z\,c_0\) for some \(c_0\in\ker(\bar\Lambda\bar e)\).
In fact this is the precise condition for \((0,\bar v')\)
to lie in \(\mathcal L/\pi L\). These integrality loci have
codimensions \(r\) and \(d'-s+r_e\) modulo \(\pi\).
The restrictions of \(\gamma^{-1}\) have integral left
inverses, obtained by composing \(\gamma\) with the appropriate
block projection. In Smith normal form their exponents
therefore lie between \(-1\) and \(0\). Each exponent \(-1\)
accounts for one dimension of the indicated codimension.
\end{proof}

\begin{corollary}[Evaluation against a principal-level source]
\leavevmode\label{cor:gluedperiod811}
At each old mark \(v\in\mathcal M\) choose its matrix \(\Lambda_v\)
and the matrix \(\gamma_v\) of Theorem~\ref{thm:splitlattice811}.
Complete these to a single finite unitary point \(\gamma\) by
using the inverse transpose on the conjugate components and
the identity elsewhere. Put
\[
       C_\gamma=\prod_{v\in\mathcal M}
         \chi_{W_j,v}(\det\gamma_v)
                    |\det\gamma_v|_v^{-s/2},\qquad
       N_{\mathcal M}=\prod_{v\in\mathcal M}q_v^{s^2}.
\]
Here \(\chi_{W_j,v}\) is the splitting character of the source
in the action of the target, fixed when \(\beta\) varies. In
the bilinear theta integral against any prescribed
principal-level source vector, translating the original
finite test by \(\gamma\) gives \(C_\gamma N_{\mathcal M}\)
times the integral using the following product test.
At each old mark its two factors are the cosets \(e+(\pi K)^s\)
and \(z+(\pi^2L_-)^s\). At each new mark the determinant
character and its invertibility test belong to the \(p\)-block,
and the \(d'\)-block is the lattice indicator. Elsewhere use
the two original lattice vacua. For an archimedean test
factored between the two blocks, the resulting kernel on the
block diagonal is the product of the two theta kernels with
compatible splittings. The identity holds at each node for
each of the stated twists.
\end{corollary}
\begin{proof}
At an old mark the split position model has additive measure
giving \(L^s\) volume one. The target acts by
\(\phi(X)\mapsto\chi_{W_j,v}(\det\gamma_v)
|\det\gamma_v|^{-s/2}\phi(\gamma_v^{-1}X)\).
Indeed the substitution in \(s\) columns has Jacobian
\(|\det\gamma_v|^s\) and the twist of the unitary action is
the source splitting character. The original tame determinant
on its point coset equals one. In \eqref{eq:localpartition811}
the first factor asks for \(X_+Y\equiv e\bmod\pi\) and
\(\Lambda_v(X_+Y-e)=0\bmod\pi^2\), exactly the two old row
conditions. There are \(q_v^{s^2}\) choices for \(Y\), obtained
by adding to the identity \(\pi\) times any residue matrix.
Each product coset has measure \(q_v^{-s(p+r)-2sd'}\).
Multiplying their total measure by \(|\det\gamma_v|^{-s}\)
recovers \(q_v^{-s(p+d')}\), the squared norm of the
untranslated test.

Each \(Y\) lifts to a source unitary element at this split
place, with inverse transpose at its mate, and identity away
from this place. It lies in the source principal subgroup.
In the position model its substitution \(X\mapsto XY\) has
unit modulus and its splitting character is trivial on this
principal subgroup because the character is tame. In the
global pairing one may move this right translation onto
the source vector by invariance of the quotient measure.
The vector is fixed. Thus every term has the same integral
as the identity term, giving \(N_{\mathcal M}\).
The only additional finite factor at a new mark is the tame
character of the determinant of the \(e\)-rows, extended by
zero at singular reductions, on integral matrices. All these
rows lie in the first block. This proves the asserted
factorization at finite places.

For the last assertion use a rational polarization adapted
to the orthogonal sum. With compatible splittings the two
Weil operators tensor on the diagonal source group and the
two target groups. Equivalently one uses the same source
character on both factors and multiplies the two target
splitting characters to get the full target character.
The rational position vectors are pairs, one from each
position space. For a tensor test the theta sum therefore
factors by absolute convergence of the Schwartz sums.
This argument includes products of the two archimedean
polynomials in disjoint rows times their Gaussians. The
integrals above are finite by rapid decay of the cuspidal
source against the theta kernels of moderate growth.
\end{proof}

\begin{lemma}[The integral Cartan inclusion on disjoint coordinate blocks]
\leavevmode\label{lem:dualcartan811}
For a list \(c_h\) of numbers of columns of heights \(h\)
in rank \(u\), denote by \(C_u(c)\) the integral span of
products of flag minors of those heights on a generic matrix.
A flag minor uses the first \(h\) columns and any \(h\)
distinct rows. Use the dual \(C_u(c)^*\) as the integral
Weyl lattice in this calculation. Suppose the positive
and negative lists \(c^+,c^-\) fit the \(p\)- and \(d'\)-blocks.
The tensor of two Weyl vectors included from these blocks,
the second in the contragredient action, belongs to the
mixed Cartan component. In the dual flag-minor lattices
the inclusion of this component into the tensor has an
integral left inverse. Consequently at a nonarchimedean
place any integer clearing the translated tensor clears
the corresponding translated mixed vector.
\end{lemma}
\begin{proof}
Let \(J_n\) reverse \(n\) coordinates. In functions on
\(\GL_n\) take the multiplication map
\[
 C_n(c^+)\otimes C_n(c^-)\longrightarrow
       (\det X)^{-u_-}C_n(c),\qquad
              F\otimes G\longmapsto F(X)G(X^{-{\rm t}}J_n),
\]
where \(u_-=\sum_h c_h^-\) and \(c\) has the positive columns
together with the complements \(n-h\) of the negative ones.
The inverse-minor formula, obtained by expanding the
determinant along complementary rows and columns, identifies
each negative minor with the complementary flag minor of
\(X\) divided by \(\det X\), up to sign. Conversely every
generating product in the target can have its columns
assigned to the prescribed two lists, even when some heights
coincide. Thus the map is surjective over \(\ZZ\).
Its target is finite free, being isomorphic to a finitely
generated subgroup of the integral polynomial ring. Lifting
an integral basis provides a section. Dualizing gives the
claimed split injection, equivariant for the standard and
contragredient actions by the substitution formula. The
left inverse need only be integral as a map of abelian groups.

Over any infinite field of characteristic zero the evaluation
functionals at invertible matrices span each dual, since
a polynomial zero at all invertible matrices vanishes
identically. Take evaluating matrices \(B\in\GL_p\) and
\(D\in\GL_{d'}\) for the two source blocks. The first \(p\)
columns of \(X=\operatorname{diag}(B,(D J_{d'})^{-{\rm t}})\)
are \(B\) included in the first block. The first \(d'\)
columns of \(X^{-{\rm t}}J_n\) are \(D\) included in the
second. The tensor of these evaluations is therefore
evaluation on the mixed target of the multiplication map.
By bilinearity and the spanning assertion this holds for
arbitrary vectors on the two blocks. The integral left
inverse takes an integral tensor in the image to its
integral preimage. It remains a left inverse after scalar
extension or localization, which proves the final assertion.
\end{proof}

\begin{theorem}[Denominators of the glued coefficient vectors]
\leavevmode\label{thm:gluedcost811}
Use the Weyl lattices of Lemma~\ref{lem:dualcartan811}. Let
\(f=[F:\QQ]\), \(H=s^8\), and \(M=p_0p_1\) for the two old
rational marks. For integral block coefficients of the
prescribed weights at any pair of block lattices, the mixed
coefficient at their block diagonal point translated on the
right by \(\gamma\) is integral after multiplication by
\(M^{E_j}\), where one can take
\[
       E_j=\sum_{\tau:F\hookrightarrow\RR}
                    \bigl(3m(w_{j,\tau}+k)+2Ht^2\bigr).
\]
On the original fixed-\((q_0,d)\) intervals one has, uniformly
in the matrices \(\Lambda_v\), row choices, \(j\), and the
prescribed twists,
\[
             \log(M^{E_j})\ll_{q_0,d} f m b_x\log b_x .
\]
\end{theorem}
\begin{proof}
In the gauge of the two block lattices the coefficient vectors
are integral by assumption. Passing from their block diagonal
point \(h\) to \(h\gamma\) first takes those coordinates at
\(h\) and then acts by \(\gamma^{-1}\). At the chosen component
the positive rectangular inclusion becomes
\(\gamma^{-1}\binom{1_p}{0}\).
Its minor of height \(h_0\) has denominator dividing
\(\pi^{\min(h_0,r)}\). Indeed in a Smith basis at most
\(\min(h_0,r)\) negative exponents occur in any exterior
product, and the changes of basis and their inverses are
integral. The negative inclusion is acted on by the
contragredient, so it uses \(\gamma^{\rm t}\binom{0}{1_{d'}}\),
which is integral here. At the conjugate component the
positive inclusion is integral, and a negative minor has
denominator dividing \(\pi^{\min(h_0,d'-s+r_e)}\), by
Theorem~\ref{thm:splitlattice811} on the second block.
At the remaining places both maps are integral.

For a rectangular inclusion \(A\), Cauchy--Binet expands
each flag minor of \(AZ\) as a sum of flag minors of \(Z\)
with coefficients given by minors of \(A\) of the same
height. Multiplying these expansions and evaluating with
an integral block vector shows that the exponents just
obtained add over its columns. Lemma~\ref{lem:dualcartan811}
transfers this bound to the mixed coefficient, including
the determinant twist. Every local uniformizer in question
divides \(M\) at that place. We may therefore use the integer
\(M\) throughout and add exponents over the \(f\) embeddings
that enter the weight.

We record the column count to keep the dependence on \(m\)
explicit. Fix \(\tau\), write \(w_i=\kappa_\tau+(2i+1)H\)
and \(b'=(t-1)/2\). Let \(\sigma_0,\sigma_1,\ldots\) start
with the \(k-1\) consecutive entries from \(k/2-1\) down
to \(1-k/2\), followed by \(H-w_1,H-w_2,\ldots\), in that
order. At height \(h_0\ge1\) the positive column count is
the positive part of
\(\sigma_{h_0-1}-\max(\sigma_{h_0},b'-w_j)-1\).
Put \(T=b'-w_j\). This count is at most
\((\sigma_{h_0-1}-T)_+-(\sigma_{h_0}-T)_+\), where
\(u_+=\max(u,0)\). The sum telescopes to at most \(w_j+k\).
The counts vanish for \(h_0>s\), as
\(\sigma_s=H-w_{m+1}<T\). Each positive column thus costs
at most \(r\le3m\).

For \(1\le h_0<s\) the negative count is the positive part
of \(\min(\sigma_{s-h_0-1},-w_j-b')-\sigma_{s-h_0}-1\).
At larger heights it vanishes because both indices in the
defining difference are zero under truncated subtraction.
The adjacent drops in the list are \(1\)
or \(2H\), apart from the drop to \(H-w_1\). At this exceptional
drop the truncated upper endpoint \(-w_j-b'\) is below
\(H-w_1\) since \(j\ge1\). Thus there are at most \(2H\)
negative columns at any height and the sum of their heights
is at most \(2Ht^2\). This proves the claimed integral
multiple.

Finally \(\kappa_\tau\in\{b_x,xb_x/x_{1,x}\}\) with
\(x_{1,x}\ge x\). Since \(s\le3x^5,\ t=5s,\ m\le2\sqrt x\)
and \(b_x\ge x^{400}\), one has \((2m+1)H\le b_x\),
\(k\le b_x\), and \(2Ht^2\le b_x\) at every sufficiently
large node. Hence \(E_j\ll f m b_x\).
Corollary~\ref{cor:actualmarks714} gives
\(\log(p_0p_1)\le22\log b_x\) on the same interval.
\end{proof}

The estimate concerns the coefficient vectors in these
period evaluations. To control the individual auxiliary
determinant one still needs product evaluations that span
the required projected image, together with an arithmetic
comparison of their values as \(\beta\) varies.
\clearpage

\section{Horizontal lifts and first occurrence}\label{nav:g-VI-3}

\subsection{Real first occurrence on each horizontal row}\label{nav:V-36}
Fix a real place of an original field \(E/F\) on an interval of
Theorem~\ref{thm:first250}. Use \(s=r+6\) for the horizontal rank
difference and \(t=t_{\rm r}\) for the vertical difference. Both
are odd. Set \(u_{ij}=r+is+jt\), \(q_j=\lfloor(m-j)/2\rfloor\) and
\(h_{ij}=-\tfrac12\mathbf1_{i+j\ {\rm odd}}\).
For \(A\in\RR\) write \(I_b(A)\) for the \(b\) numbers
\(A+(b-1)/2-l\), \(0\le l<b\). Define \(P_{ij}\) by adding
\(h_{ij}\) to each entry in the following signed list.
\begin{itemize}
\item The \(v_1,\ldots,v_{k-1}\) have sign plus. The numbers
\(-b_a\), \(1\le a\le m\), have sign plus for \(a\le j+1\),
and alternate starting with minus at \(a=j+2\).
\item Include \(I_t(-e_a)\) for \(1\le a\le j\) and
\(I_s(c_{m-b})\) for \(0\le b<i\), all with sign plus.
\end{itemize}
Here \(c_l=\lfloor(a_{l-1}+a_l)/2\rfloor\) as in
Theorem~\ref{thm:realcore253}. The \(a_l\) are the real core weights. The list defines a discrete
series on \(U(u_{ij}-q_j,q_j)\). Indeed, the first two kinds of
entries together with the strings of length \(t\) are distinct by
Theorem~\ref{thm:tower169}. All length-\(s\) strings lie above
\(v_1\), their centers are spaced by \(g_*>s\), and \(h_{ij}\)
puts the result in \((u_{ij}-1)/2+\ZZ\). There are exactly
\(q_j\) minus signs. Let \(\check P_{ij}\) denote the dual, formed
by negating the entries within the two compact blocks.

\begin{theorem}[Horizontal lifts and the cuspidality of their outputs]\leavevmode
\label{thm:realrow257}
For \(0\le i<m\) and \(0\le j\le m\), integrate the original
plain horizontal kernel from \(V_{i,j}\) to \(V_{i+1,j}\) against a
form of real type \(\check P_{ij}\). At the given real place its
quotient input has first occurrence at the prescribed target in
that Witt tower. The physical output has type \(\check P_{i+1,j}\).
Every nonzero horizontal lift of a cuspidal form of these types
at all real places is cuspidal. In particular the second lift
\(B_2\) with the prescribed nonzero restriction in
Theorem~\ref{thm:secondrow256} is cuspidal.
\end{theorem}
\begin{proof}
Write \(c=c_{m-i}\), \(u=u_{ij}\), \(p=u-q_j\) and \(q=q_j\).
The tensor for the kernel is the space of maps from
\(-V_{i+1,j}\) to \(V_{i,j}\). As for the first edge in
Theorem~\ref{thm:secondrow256}, use \(-iH\) on the smaller factor
and the Hermitian dual on the larger, with \(i=\sqrt{-1}\) in
this description. This exchanges the two compact blocks on the
smaller side. Integrating the plain kernel against a unitary
vector takes its conjugate for the quotient convention. The
entries on the smaller side are therefore those of \(P_{ij}\),
with the two blocks exchanged. The splitting rule attaches
\(\chi_{m-i}\) to an even-dimensional space and
\(\chi_{m-i}\eta^{-1}\) to an odd-dimensional one. Their half
angular exponents before the parity correction are \(c\), and
\(\eta\) has half exponent \(1/2\). Since \(u_{ij}\) is odd
exactly when \(i+j\) is even, the half exponents attached to the
larger and smaller spaces are respectively \(c+h_{ij}\) and
\(c+h_{i+1,j}\).
\srcloc{Development/AbelianVarietyTheory/Definitions/Divisors/LeastSplitAuxiliaryPrimeThetaChainSpec}{168--175,199--201},
\srcloc{Development/AbelianVarietyTheory/Definitions/Divisors/BasePlaceNotMemChiGridBadPlaces}{401--409,431--439}.

Subtract \(c+h_{ij}\) from the input. Every resulting entry is
negative, at distance greater than \((s+1)/2\) from zero. For
the cloud and the vertical strings their largest unshifted entry
is \(v_1\), since \(e_a=b_a+H_0\), \(H_0>(t-1)/2\) and \(b_a>0\).
Also \(c_m\ge v_1+g_*\) and \(c\ge c_m\). Every old horizontal
string has top at most \(c-g_*+(s-1)/2\). The bound on \(g_*\)
in Theorem~\ref{thm:first250} proves all the stated distances.
The ranks have opposite parity. In Atobe's criterion these
distances give \(k_\lambda=-1\). After the block exchange there
are \(p\) negative minus entries and \(q\) negative plus entries,
so \(r_\lambda=p,\ s_\lambda=q\). The target signature
\((p+s,q)\) has the form
\((r_\lambda+l+2T+1,s_\lambda+l)\) with \(l=0,\ T=(s-1)/2\).
A member of either counting set at \(T\) would have absolute
value at most \(T\), which is impossible. Subtracting any
positive number of hyperbolic planes forces \(l<0\), again
impossible. This proves the first occurrence.
\href{https://arxiv.org/pdf/2002.09148}{Ichino, Theorem 4.3(i),
after Atobe}.

For the parameter of the nonzero lift, the negative entries
exchange compact signs in Ichino's formula. This restores their
original blocks. The new block has signature \((s,0)\) and
precedes all old entries. More explicitly put \(N=u+s\).
For an old centered entry \(\beta\) in combined position \(k'\),
the inducing weight is
\(\beta-((N+1)/2-k')+c+h_{i+1,j}\). The weight on the new block
is \(-u/2+c+h_{i+1,j}\), since there are no positive centered
entries. Adding the ordered half sum gives the old unshifted
entry plus \(h_{i+1,j}\) in the first formula, and gives
\(I_s(c)+h_{i+1,j}\) on the new block. The separation just
proved puts these entries in strictly decreasing order across
Levi blocks. The blocks are compact, so cohomological induction
in stages gives the discrete series of this parameter. Pulling
back from the Hermitian dual of the target applies inverse
transpose and negates the entries in each block. The resulting
physical type is \(\check P_{i+1,j}\).
\href{https://arxiv.org/pdf/2002.09148}{Ichino, Theorem 4.1(i)}.

A global lift to a smaller rational member would yield a
nonzero local pairing at this real place by fixing the
other variables. Hence all smaller lifts of each relevant
cuspidal constituent vanish. The constant-term tower property
makes its lift at the prescribed member cuspidal and local
Howe duality gives the real type above. The assertion extends
to sums of constituents and to their smooth vectors by
continuity of the theta integral on the cuspidal realization.
Apply this first with \(j=i=0\) and then with \(j=0,i=1\)
to the two integrals defining \(B_2\).
\href{https://arxiv.org/pdf/2002.09148}{Ichino, Section 6.3}.
\end{proof}

\subsection{An explicit integral terminal coefficient for a restriction}\label{nav:V-37}
Use the original fields, intervals, finite tests and levels of
Theorem~\ref{thm:coregram258}. Equip the terminal quotient with
probability Haar measure, and take \(C_{\rm fix}\) in that
normalization. Fix a pure real unit vector \(v\) as in that
theorem. Write \(H(v)\) for its bottom finite multiplicity
space and \(T_v\) for its ascent with the selected real tests.
For the lists of Proposition~\ref{prop:schurframe259}, let
\(e_{c,a,b}\) be supported on level class \(c\), with value
\(l_a(\rho(k)v_b)\) on its compact real group. Here \(k\) is
the real variable, \(d_\rho=\dim\rho\), and \(d_0,a_0,L_0\)
denote the three quantities defined in that proposition.

\begin{theorem}[A lower bound for one integral terminal coefficient]\leavevmode
\label{thm:individual259}
For each complex linear functional \(\ell\) on \(H(v)\), write
\(z_{c,a,b}=\ell(P_vT^\dagger e_{c,a,b})\), where \(T^\dagger\)
is the reverse integral transform and \(P_v\) projects to \(H(v)\).
With \(B,H_*\) as in Lemma~\ref{lem:terminalclass259}, one has
\[
 \frac{B^{-4a_0} e^{-\delta} C_{\rm fix}}{h_m d_\rho}\|\ell\|_*^2
 \ \le\ \sum_{c,a,b}|z_{c,a,b}|^2\ \le\
 \frac{L_0^2 B^{4a_0}e^\delta C_{\rm fix}}{h_m d_\rho}\|\ell\|_*^2.
\]
In particular one of these integral terminal coefficients satisfies
\[
 |z_{c,a,b}|\ \ge\
 \frac{e^{-\delta/2}\sqrt{C_{\rm fix}}\,\|\ell\|_*}
      {H_*\sqrt{d_\rho}\,N^{d_0} B^{2a_0}}.
\]
This holds for the actual first restriction
\(\ell(A)=\mathcal B(A,T_0w)\). For each prescribed nonzero
smooth bottom row \(T_0w\), a pure \(v\) can be chosen with
\(\|\ell\|_*>0\). Uniformly on the same intervals there is a
polynomial \(P_{\psi,\eps}\) such that
\[
 \log\big(H_*\sqrt{d_\rho}\,N^{d_0}B^{2a_0}\big)
 \ \le\ P_{\psi,\eps}(x)+N^2\log\operatorname{Nm}\mathfrak m .
\]
\end{theorem}
\begin{proof}
On class \(c\) let \(\Gamma_c:W_c^*\otimes W_c\longrightarrow
\mathcal A_{m,c}\) be the complex coefficient map, using complex
fibers of the integral lattices. The target is the full right
\(\rho\)-isotypic space on that class. Schur orthogonality, with
the mass \(1/h_m\) of this class included, says that
\(\sqrt{h_m d_\rho}\,\Gamma_c\) is an isometry. It is onto by
the compact matrix coefficient decomposition in
Proposition~\ref{prop:terminalcore255}. Products of the two lists
have frame operator \(S_l\otimes S_v\), as follows by expanding
the sum of rank-one operators. Proposition~\ref{prop:schurframe259}
places its eigenvalues in \([B^{-4a_0},L_0^2B^{4a_0}]\).
Transport by \(\Gamma_c\) divides both bounds by \(h_m d_\rho\).
The terminal classes are orthogonal, so the same bounds hold
on their direct sum for the union of the transported lists.
This sum includes all indices, allowing repeated or zero vectors.

Let \(a_\ell\in H(v)\) represent \(\ell\) in the Petersson
metric. Moving the kernels across the pairings gives
\(P_vT^\dagger=T_v^*\) on the terminal space by
Theorem~\ref{thm:moment258}. Hence
\(|z_{c,a,b}|=|\langle e_{c,a,b},T_va_\ell\rangle|\).
Applying the frame bounds to \(T_va_\ell\) and using the
two inequalities for \(T_v^*T_v\) in Theorem~\ref{thm:coregram258}
proves the sum estimate. There are \(h_m N^{2d_0}\) summands,
so the largest square is at least the sum divided by this
number. Taking square roots and using \(h_m\le H_*\) proves the
individual bound. For the indicated first restriction the
existence of the required \(v\) follows from the last assertion
of Theorem~\ref{thm:moment258}.

We give the dependence on the arithmetic data. The proof of
Theorem~\ref{thm:first250} bounds the rank parameters and the
gap \(g_*\) by powers of \(x\), with \(v_{k-1}=0\) and
\(n^3\sqrt{v_1}\le g_*\). The terminal even core weights
are obtained from the core entries and their added strings in
Theorem~\ref{thm:realcore253}. The core entries have size
at most \(v_1+1+(m+1)g_*\), and the centers \(c_l\) add at
most a further polynomial in \(x\) by their defining formulas.
Each string displacement is at most its length, and passage from
the ordered parameter to the highest weight subtracts a half
sum of roots with entries of modulus less than \(N\). Thus
\(f,N,d_0,a_0\) are polynomially bounded. The polynomial Schur
spaces lie in tensor powers, so \(d_\rho\le N^{d_0}\). The
cyclotomic modulus \(M\) containing \(E\) in
Lemma~\ref{lem:dimension239} has \(\log M=o(\log b)\) with
\(\log b=x^\eps\). Its discriminant bound
\(\log |D_E|\le [E:\QQ]\log M\) follows either from the different
bound of Lemma~\ref{lem:lattices247} or from the abelian
conductor--discriminant formula. Substitution into the explicit
expressions for \(R,B,H_*\) gives the final inequality.
\end{proof}

The selected terminal input is integral in the coefficient
lattice. The bound for its descended restriction still contains
the real norm factor \(C_{\rm fix}\) and the dual norm of the
first restriction. Relating this complex period to the specified
integral projection after all remaining lifts is the comparison
required in Section~\ref{sec:missing}.

\subsection{Cuspidal ascent along the full horizontal row}\label{nav:V-38}
Write \(G_i=U(V_{i,0})\), \(u_i=r+is\), \(s=r+6\), and
\([G_i]=G_i(F)\backslash G_i(\mathbb A_F)\) on each original
field of Theorem~\ref{thm:first250}. The real signatures are
\((u_i-m/2,m/2)\). At the finite places use the tests of
Proposition~\ref{prop:rowfinite262} and denote their product
on edge \(i\) by \(\varphi_{i,f}\). For a polynomial Fock test
\(p\) at infinity define the plain theta integral by
\[
 U_{i,p}b(g)=\int_{[G_i]}
       K_{V_{i+1,0},\,V_{i,0}}(g,h;p\otimes\varphi_{i,f})\,b(h)\,dh .
\]
The quotient convention for the local correspondence therefore
takes the conjugate of the physical source form.

\begin{theorem}[Nonvanishing through the terminal horizontal group]\leavevmode
\label{thm:globalrow262}
For every nonzero compact-finite bottom cusp form \(b\) of the
prescribed real type and old level there are products of real
monomial tests \(p_0,\ldots,p_{m-1}\) such that
\(U_{j-1,p_{j-1}}\cdots U_{0,p_0}b\) is nonzero for every
\(1\le j\le m\). These are cusp forms at the given congruence
levels of real type \(\check P_{j0}\). In a bottom good Hecke
system \(B_0\), their good lists in quotient coordinates are
\(\mathcal B_{j0}\) of Theorem~\ref{thm:mixedvalues257}.
The ascent can also continue from the particular \(B_2\) obtained
in Theorem~\ref{thm:lowesttwo261}, keeping both of its initial
tests and its nonzero restriction.
\end{theorem}
\begin{proof}
Take first an irreducible cuspidal copy at a given stage, whose
good list is \(\mathcal B_{i0}\) for some original \(B_0\).
At stage zero this is the definition of the list. The smaller
rational targets in the Witt tower have zero lift by the real
first occurrence in Theorem~\ref{thm:realrow257}. Its tower
argument also gives cuspidality and the real type of any nonzero
output. Put \(u=u_i,\ N=u+s\) and let \(R\) be the rational
Witt index of the target. At any real place a rational isotropic
space has dimension at most \(m/2\). Hence \(N-R>u\), since
\(s>m/2\). This is the convergent range of the doubled
Siegel--Weil integral for unitary groups. For \(N>2u\) one has
the convergent doubling identity directly. For \(u<N\le2u\),
Propositions 11.1(ii) and 11.2 of Gan, Qiu and Takeda give
the constant Laurent coefficient of the doubling integral at
\(s/2\), up to a nonzero measure constant. Both formulas apply
to the mixed inner product for two Schwartz tests and two cusp
forms. The center has compact adelic quotient. The original
theta integral converges against a cusp form by rapid decrease
and the polynomial growth of a Schwartz theta kernel on Siegel
sets. \href{https://arxiv.org/pdf/1207.4709}{Gan, Qiu and
Takeda, Sections 11.4 and 11.5}.

Unfold on a right half-plane. The open orbit pairs the
restricted section with a global matrix coefficient. All other
orbits give constant terms of the source cusp form along
proper parabolics, hence vanish. For pure tensors in the copy
the normalized matrix coefficient factors. Outside \(S_0\) the
product is the partial standard value divided by the doubling
denominator, using Proposition~\ref{prop:rowfinite262}. By
Theorem~\ref{thm:mixedvalues257} with \(j=0\) this quotient
is holomorphic and nonzero at the theta point. At the marked
places the restricted section is compactly supported on the
small congruence group, where its flat extension changes it
by a bounded positive height to a complex power. These local
integrals are entire and give the nonzero stated scalars on
the whole fixed space. At infinity the discrete series is
tempered, so the local integrals converge locally uniformly
near \(s/2\). Their pairings on polynomial tests are nonzero
by the real first occurrence and Ichino's nonvanishing theorem
for the local Rallis form. Density gives the same conclusion
on Fock polynomials, using continuity of the local integrals.
The unfolded product is thus holomorphic at the required point
and its constant coefficient is its value.
\href{https://arxiv.org/pdf/2002.09148}{Ichino, Lemma 7.1 and
Proposition 7.2}.

To see that the tests detect the prescribed vector, let \(D\)
be the tensor product of the real source Harish-Chandra modules.
Those \(v\in D\) that annihilate the mixed real pairing for
every pair of polynomial tests and every other real source
vector form a submodule. Compact translations on \(v\) transfer
to the first test by equivariance. Differentiation gives the
same transfer for the complexified Lie algebra, whose oscillator
action preserves polynomials. The local nonvanishing just proved
makes this a proper submodule, hence it is zero by irreducibility.
In the fixed finite part of the global copy choose an orthonormal
basis \(e_j\) and write the input as \(\sum_jv_j\otimes e_j\).
The old bottom level contains the prescribed principal congruence
subgroup, and each subsequent output has the required invariance
by Proposition~\ref{prop:rowfinite262}. The marked pairings are
scalar on these finite spaces and the unmarked spherical lines
are one-dimensional. If \(v_j\ne0\), pairing the lift with that
of a suitable \(w\otimes e_j\) consequently gives a nonzero
mixed product. This proves nonvanishing for that input.

Distinct orthogonal invariant cuspidal copies give zero mixed
products even for different tests. Indeed their global matrix
coefficient in the open orbit is identically zero for every
group translate. The same observation for all tests makes
their invariant theta spans orthogonal. Fixed finite level
and the finitely many compact types in the input give finite
multiplicity by cuspidal admissibility. Thus a test detecting
one nonzero component detects the full input. Expanding a
polynomial into monomials selects one product monomial with
nonzero output. Such a test gives a compact-finite output.
At good places local Howe duality and spherical transfer give
the next list. In the split case pulling the target back from
its Hermitian dual inverts the list, and taking the conjugate
for the next input inverts it a second time. The resulting
recursion is precisely that of Theorem~\ref{thm:mixedvalues257}.
This proves the induction.

For the assertion with the first two tests fixed, the nonzero
restriction in Theorem~\ref{thm:lowesttwo261} gives \(B_2\ne0\).
Its Hermite tests give compact finiteness, and
Theorem~\ref{thm:realrow257} gives cuspidality and its real type.
At a marked place each individual coset of the integral
rectangle is fixed by the principal congruence groups. At an
unmarked place the tests in that construction are the lattice
vacua, also after the dual-lattice Fourier exchange of
Proposition~\ref{prop:oppline256}. Hence the same good transfer
identifies the lists of its constituents. Starting the argument
at \(i=2\) now proves the assertion, including the case \(m=2\)
where the ascent has already reached its endpoint.
\end{proof}

\subsection{The adjoint on consecutive line oscillators}\label{nav:V-39}
Consider \(i=2r+1\), \(r\ge1\), \(i+1\le m\). In the notation
of Proposition~\ref{prop:positions264} the row spaces decompose
as \(S=C\perp\mathbb H^r\),
\(T=(D\perp L_n)\perp\mathbb H^r\), with
\(C=C_i\), \(D=C_{i+1}\). The old Jacobi line \(L_o\) and
the middle line \(L_n\) of \(T\) have respective basis vectors
\(z_o,z_n\) of norm \(-1\). Identify the lines by these bases.
Set \(d=-\xi_0\), \(\zeta_o=d^r z_o\), \(\zeta_n=d^{r+1}z_n\),
and \(\mathfrak l=A/B\). The Jacobi period on \(S\) uses
\(\zeta_o\), and the Bessel period on \(T\) uses \(\zeta_n\).

\begin{proposition}[Transfer of the specified Jacobi test]\leavevmode
\label{prop:linkedline265}
There is a conjugate-linear transfer of Schwartz tests from
the oscillator on \(\operatorname{Hom}(C,L_o)\) to that on
\(\operatorname{Hom}(L_n,C)\). It identifies the new line theta
factor, evaluated on the stabilizer of the isotropic frame in
\(S\), with the conjugate of the old Jacobi theta factor.
The upper unitriangular character of that frame is treated
separately. The finite vacuum for the distinguished old Jacobi
line transfers to the full undeepened new line vacuum at each
place, up to a nonzero scalar. At a marked place the latter is
the sum of \((Nv)^{a_v\dim_E C}\) original coset indicators.
\end{proposition}
\begin{proof}
Let \(\dagger\) denote ordinary Hermitian adjoint. The map of
position-frequency spaces is \(w\mapsto w^\dagger\). Indeed
the map from the old to the new distinguished line carrying
\(\zeta_o\) to \(-\zeta_n\) is multiplication by \(\xi_0\).
In the opposite-row coordinates the inverse skew-adjoint
operation on \(v\) is \(-\xi_0^{-1}v^\dagger\), so their
composition on \(w\) is \(w^\dagger\). The trace of a product
in the Hermitian form on the map space is conjugated under
adjoint, by cyclic invariance of trace. Its alternating trace
after multiplication by \(\xi_0\) therefore changes sign.
The map preserves rational points and intertwines the two
Heisenberg groups upon inversion of the center.
Complex conjugation of a Schwartz function, followed by this
change of coordinates and Fourier transformation to the chosen
polarization, implements the transfer. With adelic dual measures
Poisson summation identifies the rational theta functionals.
Equivalently, in the realization as theta functions on the
whole Heisenberg group one composes conjugation with the
rational change of variables and evaluates at zero. This
description fixes their global scalar.

Here are the actions of the unipotent radical of the frame in
\(S\). Use its scaled hyperbolic columns \(e_j=d^j b_j\),
their dual columns \(e'_j\), and the residual space \(C\).
For \(n\) fixing the columns pointwise let \(B_n\) be its
block from \(C\) to \(\langle e_j\rangle\), and \(Z_n\) the
block from \(\langle e'_j\rangle\) to \(\langle e_j\rangle\).
In the following formulas \(x=\zeta_o e_r^{\,t}\),
\(y=-\zeta_n e_r^{\,t}=\xi_0 x\), where \(e_r^{\,t}\)
here denotes the last coordinate row on the index set.
The old Heisenberg displacement is \(xB_n\). In the new
opposite-row variable it is
\(-\xi_0^{-1}(yB_n)^\dagger=(xB_n)^\dagger\).
These formulas follow by composing the block from \(C\)
with the retained column in the first position and by taking
its inverse skew adjoint in the second position.

For the central terms write \(G_L\) for the line Gram matrix
and \(t(v)=\tfrac12\operatorname{Tr}_{E/F}(\xi_0 v)\).
The negative of the old argument of the additive character
and the new argument, respectively, are
\[
    t\!\left(\operatorname{tr}(x^*G_Lx\,Z_n)\right),\qquad
    t\!\left(\operatorname{tr}
      ((\xi_0\bar\xi_0)^{-1}y^*G_Ly\,Z_n)\right).
\]
For the first formula the shear of the opposite columns into
the retained columns has block \(Z_n\). Its Weyl phase is one
half of the alternating pairing of the retained displacement
with this shear, giving the displayed trace with minus sign
for the old argument. In the opposite row the retained
functional is \(-\xi_0^{-1}y^*G_L\). Inserting it on both
sides of the same shear in the opposite pairing gives the
second trace since the two scalar factors multiply to
\((\xi_0\bar\xi_0)^{-1}\). The formulas agree since
\(y=\xi_0x\). Upper unitriangular changes of basis have
determinant one and fix the last coordinate row on the
right. This accounts for the whole unipotent action.
\srcloc{Development/AbelianVarietyTheory/Proofs/DecoratedLatticeProofs/JacobiActionsMatchThroughLinkMap}{32--66}.

On the remaining unitary group \(U(C)\) the splittings give
conjugate actions. In split matrix coordinates the horizontal
character on \(U(C)\) is \(\gamma_i\eta^{-1}\), and the core
character there is \(\gamma_i\), leaving \(\eta^{-1}\). The old
line in the reversed position has factor \(\eta\) there. The
ratio of the two actions after adjoint and conjugation is a
character by Heisenberg irreducibility. It is trivial on
rational points by the theta functional, at split places by
this calculation, and on the full compact stabilizer at almost
all nonsplit places by the integral vacua. At those nonsplit
places an unramified character of the unitary group is trivial
by the hyperbolic-plane argument of
Proposition~\ref{prop:cancelsquare256}. Weak approximation at
the remaining finite set and the archimedean places gives
triviality there as well, as in that proposition. This proves
the asserted equality of theta factors.

Finally put \(M=\operatorname{Hom}(\Lambda_C,\mathcal O_Ez_o)\).
The old line displacement lattice is
\(\mathfrak A\mathfrak l^r(\mathfrak m/\bar{\mathfrak m})M\).
This is the line rule used in Theorem~\ref{thm:third263}, with
the indicated power for the \(r\)-th line. The new undeepened
line displacement is
\(\mathfrak A\mathfrak l^{-(r+1)} M^\dagger\), by the middle
ideal in \(T\). These are exchanged exactly by adjoint, since
\[
    \bar{\mathfrak A}\mathfrak l^{-r}
           (\bar{\mathfrak m}/\mathfrak m)
       =\mathfrak A\mathfrak l^{-(r+1)} .
\]
Both are self-dual for the trace form. Over a split place a
full translation and dual-frequency lattice has a one-dimensional
space of invariant tests, spanned in position coordinates by
its lattice indicator. This follows from the support and
translation equations. The same vacuum calculation in a
self-dual polarization applies at the unmarked nonsplit places.
Thus the transferred line is the claimed vacuum, whose marked
expansion partitions the full lattice at depth \(a_v\).
With full lattice indicators chosen at all finite places, one
may incorporate the resulting nonzero normalization scalar
in a single real line test.
\end{proof}

\subsection{Nonvanishing at the second horizontal lift}\label{nav:IV-35}
Retain the original fields, intervals, lattices and characters of
Theorem~\ref{thm:firstsquare255}. Use \([B]\) to denote the
automorphic quotient of \(U(B)\) and \(K_{D,B}\) for the kernels
of Proposition~\ref{prop:cancelsquare256}. For a fixed bottom row
realization \(T\) and a prescribed smooth vector \(Tw\) write
\[
 B_1(g)=\int_{[Y]}K_{S,Y}(g,h;p)\,(Tw)(h)\,dh,\qquad
 B_2(j)=\int_{[S]}K_{Z,S}(j,g;\Phi)\,B_1(g)\,dg .
\]

\begin{theorem}[A restriction after two horizontal lifts]\leavevmode
\label{thm:secondrow256}
For every nonzero smooth \(Tw\), the tests may be chosen so
that \(B_1\) is cuspidal and, with a descending core form
\(A_2\) on \(C_2\),
\[
            \int_{[C_2]} A_2(k)\,B_2(k\oplus1)\,dk\ne0 .
\]
One terminal integral class coefficient, all descending tests
above \(C_2\), and the entire second horizontal test \(\Phi\)
can be fixed for the given realization \(T\). At the marked
places \(\Phi\) is an individual integral coset indicator for
the actual lattice on \(Z\), including its negative line. Only
the real part of the first horizontal test \(p\) has to vary
with \(w\). Before selection of a single coset, the second
restriction for the Fourier-transformed test equals the
first projected seesaw period of Theorem~\ref{thm:firstsquare255}.
\end{theorem}
\begin{proof}
We first verify the tower property required on the intermediate
group. The dual of the bottom row has the signed parameter
\(\mu=(v_1,\ldots,v_{k-1},-b_1,\ldots,-b_m)\) of
Theorem~\ref{thm:first250}. Its cloud has plus sign and the
remaining signs alternate beginning with plus. Put
\(p_Y=k-1+m/2,\ q_Y=m/2\). The plain horizontal kernel is
integrated against \(Tw\), so its quotient input is the
conjugate, of parameter \(\mu\). In the tensor
\(\operatorname{Hom}(-S,Y)\) take the Hermitian form on the
dual of \(S\) and the skew form \(-iH_Y\) on \(Y\), up to
the same positive real rescaling as in
Theorem~\ref{thm:realcore253}. The latter form exchanges the
two compact blocks on the smaller group. The target angular
half exponent is \(c_m\), since this target has even rank.
All entries of \(\mu-c_m\) are negative and have distance
greater than \((d+3)/2\) from zero, since the largest is \(v_1-c_m\)
and \(c_m\) is within one of \((a_{m-1}+a_m)/2\).

The ranks now have opposite parity, so the absence of zero
gives \(k_\lambda=-1\). In the exchanged blocks there are
\(p_Y\) negative minus entries and \(q_Y\) negative plus
entries. Thus \(r_\lambda=p_Y,\ s_\lambda=q_Y\). The actual
target has signature \((p_Y+d+1,q_Y)\). In Atobe's criterion
this is \((r_\lambda+l+2t+1,s_\lambda+l)\) with \(l=0,t=d/2\).
Both counting sets are empty by the distance just proved.
For a smaller member of this Witt tower \(l<0\), violating a
necessary condition. Thus at every real place the given target
is the first occurrence. The global constant-term tower
property now gives cuspidality of \(B_1\) whenever it is
nonzero. It applies to the smooth row vector by continuity of
the theta lift on its smooth cuspidal realization.
\href{https://arxiv.org/pdf/2002.09148}{Ichino,
Theorem 4.3(i) and Section 6.3}.

Fix initially one nonzero lowest-type vector. By
Theorem~\ref{thm:firstsquare255} take its indicated \(p\),
its positive line test \(q\) (conjugating the negative line
test there), and its descending form \(A\) on \(S\).
Write \(A(g)=\int_{[C_2]}K_{C_2,S}(k,g;b)A_2(k)\,dk\).
Here \(b\) and \(A_2\) are the original test and the next
form of that descent. If \(C_2\) is terminal, \(A_2\) is
the class coefficient itself. The first period is
\(\int_{[S]}A(g)L_q(g)B_1(g)\,dg\) by
Proposition~\ref{prop:cancelsquare256}, so \(B_1\ne0\).
It is compact-finite on the target for the product Hermite
test and is a cusp form by the preceding paragraph. Substitute
the expression for \(A\) into this last integral. Both \(B_1\)
on \(S\) and \(A_2\) on \(C_2\) decrease rapidly. The line
theta function has polynomial growth, as does the second core
kernel on the product of the two Siegel sets by the Schwartz
estimate of Theorem~\ref{thm:seesaw245}. Hence this double
integral converges absolutely. Grouping the two kernels by
Proposition~\ref{prop:cancelsquare256} and integrating first
over \([S]\) gives the displayed restriction, with
\(\Phi=b\otimes Uq\). This order of integration uses a cusp
form on each noncompact integration group.

At a marked place the \(C_2\)-part of \(b\) is one core
rectangle indicator. Proposition~\ref{prop:oppline256}
expands the new line test as an average of \(D_L\) individual
residue indicators with coefficients of absolute value one
before division by \(D_L\). Tensoring with the rectangle gives
a single marked map coset per term, since
\(\Lambda_Z=\Lambda_{C_2}\oplus\mathfrak l z'\). All the
specified residues lift simultaneously to integral maps by
projectivity and the Chinese remainder argument in
Theorem~\ref{thm:markedcore253}. Keep the unmarked and real
reference factors common to all terms. The transformed period
is then a weighted sum of the individual periods with total
absolute weight one. At least one individual period has
absolute value at least that of the transformed period.

Fix this residue choice. At each marked place its inverse
Fourier transform in the line variable is supported on the
old integral position lattice \(M_v\) and is a character
there, constant modulo \(\varpi_v^{a_v}M_v\). This follows
by integrating over a translate of \(M_v^\vee\), exactly as
in Proposition~\ref{prop:oppline256}. It satisfies the old
depth equations. Thus for the selected term the reverse use
of the two kernel factorizations expresses its nonzero value
as the first restriction of a fixed compact-finite cusp form
on \(C_0\). Its first real test is still a product on \(Y,L\).
As in Theorem~\ref{thm:firstsquare255}, translate an arbitrary
nonzero smooth vector in the irreducible bottom row against
this fixed, bounded form. Transfer the translation to the
\(Y\)-test, and expand that one real variable in Hermite
functions. One product Hermite test gives a nonzero value.
The finite positions, the real line test, \(b,A_2\), and the
selected second horizontal test are all fixed. The same
factorizations and the cuspidality of the new \(B_1\) give
the displayed nonzero second restriction for this \(w\).
\end{proof}

\subsection{A nonzero period after the third horizontal lift}\label{nav:IV-38}
On the fields and intervals of Theorem~\ref{thm:lowesttwo261},
take \(m\ge4\). Keep the particular terminal integral coefficient,
descending tests and first two horizontal tests giving its
restriction \(J\ne0\). Write the latter as
\[
       J=\int_{[U(C)]}A_2(h)B_2(h\oplus1)\,dh,\qquad
       A_2(h)=\int_{[U(D)]}A_3(k)K_{D,C}(k,h;b)\,dk .
\]
Here \(C=C_2,D=C_3,S=V_{2,0},T=V_{3,0}\), and \(b\) is the
third descending core test. These cusp forms and kernels are
those in Theorems~\ref{thm:secondrow256} and~\ref{thm:downcore254}.
Use \(u=\dim_E S\), and the Haar conventions of
Lemma~\ref{lem:orbvol263} and Proposition~\ref{prop:unfold263}.
If a fixed scalar relates the original group measures to their
finite products, include it in the corresponding archimedean
quotient measure.

\begin{theorem}[A first positive-length Fourier--Jacobi period]\leavevmode
\label{thm:third263}
There is a third horizontal Schwartz test at the original depth,
with one individual coset indicator at each marked place, for
which the Fourier--Jacobi period \(P\) of its lift of \(B_2\),
paired with \(A_3\) at the distinguished character, satisfies
\[
  |P|\ \ge\ \frac{c_f\,\|q\|_2^2}{2(N\mathfrak m)^u}|J|
       \ \ge\ \frac{\|q\|_2^2}{2(N\mathfrak m)^{3u-1}}|J| .
\]
The same line test \(q\) may be used for every such pair of
forms on the given field. Its finite part is the integral
lattice test specified below. The real position test used
for this inequality has integral one on the norm shell and
may depend on the forms. For nonvanishing one may take a single
product of real Hermite functions as the entire third horizontal
real test.
\end{theorem}
\begin{proof}
Use the notation and position \(x=X(e)\) of
Proposition~\ref{prop:shell263}. In the residual variable on
\(\operatorname{Hom}(-D,S)\) take the test \(b\otimes q\)
for \(S=C\perp Ez\). At an unmarked place take for \(q\) the
lattice vacuum on the line summand. At a marked place use the
indicator of the zero coset modulo \(\pi^{a_v}\) in its
original integral position lattice. This is admissible jointly
with the original core indicator \(b\) on the other summand.
It also has the line depth prescribed for this Fourier--Jacobi
step. In fact the full, unmarked line displacement for the
summand would be
\(\mathfrak A\operatorname{Hom}(\Lambda_D,\mathfrak l z)\).
The separate line displacement multiplies this by
\(\mathfrak m/\bar{\mathfrak m}\). At a marked place this
multiplies positions by \(\pi^{a_v}\) and their dual
frequencies by \(\pi^{-a_v}\), whose common vacuum is exactly
the indicated zero-coset indicator. Take for its real test
a fixed nonzero Gaussian, with a fixed norm. This defines
\(q\) before choosing any row.
\srcloc{Development/AbelianVarietyTheory/Definitions/TwistedInnerProducts/IsAuxiliaryVectorWittDec}{177--184}.

Take \(a_f=1_{z+\mathfrak j\widehat L}\). The function
\(\mathcal R\) in Proposition~\ref{prop:unfold263} is constant
in its finite variable on the image of \(K_f\). The form \(B_2\)
is fixed by this group by Theorem~\ref{thm:globalrow262}.
At a marked place the residual tensor is a single coset in
the rectangle with target lattice \(\Lambda_S\). Principal
congruence acts trivially on its residue, and the determinant
character is trivial at this depth. At a good place the tensor
is the vacuum for the orthogonal sum. It is fixed by the full
stabilizer of \(\Lambda_S\) by the lattice calculation of
Proposition~\ref{prop:rowfinite262}. The larger residual space
\(D\) has even rank, so one may polarize \(D\) and act on
the odd space by substitution. At split places the same
statement follows from matrix substitution. The splitting on
the small group in \(\Omega\) uses the character attached to
\(T\), the same \(\gamma_2\) as for \(D\). It follows that
\(\Omega(k_f)(b\otimes q)=b\otimes q\) for \(k_f\in K_f\).
The formula for \(b_g\) gives the asserted invariance.

By Lemma~\ref{lem:orbvol263} the finite orbit integration is
now multiplication by \(c_f\). On the real orbit
\(H_\infty\backslash G_\infty\), write
\(R_\infty(\dot g)=\mathcal R(g_\infty,1)\). By
Proposition~\ref{prop:unfold263} it is continuous and has
value \(\|q\|_2^2J\) at the base point. Choose a product
neighborhood there on which its difference from this value
has modulus less than \(\|q\|_2^2|J|/2\).
The real orbit map is a smooth identification with the norm
shell at each real place. Its differential is onto the tangent
space by transitivity of the unitary group and its stabilizer
is \(H_\infty\). The shell is closed and smooth since its norm
value is nonzero. Choose nonnegative smooth bumps on the shell
with compact support in this neighborhood. In normal coordinates
extend them to ambient compactly supported smooth functions by
multiplication with a transverse bump equal to one at zero.
Normalize their product \(a_\infty\) to have shell integral one.
After multiplying the integral by \(\bar J/|J|\), its real part
is at least \(c_f\|q\|_2^2|J|/2\). This proves the lower bound
before returning to the original marked polarizations.

Proposition~\ref{prop:shell263} writes the marked hyperbolic test
as \((N\mathfrak m)^u\) original indicators with unit-modulus
coefficients. Tensoring each with the fixed residual indicator
gives a single marked coset on the entire edge. Linearity and
the triangle inequality give the first displayed inequality
for some one of these choices, and the measure estimate gives
the second. At the other places the unmarked vacua are unchanged
by dual-lattice Fourier transformation. The original \(A_3\)
and \(B_2\) are used throughout.

For this finite choice the period is a continuous linear
functional of the real Schwartz test by the absolute seminorm
bound in Proposition~\ref{prop:unfold263}. Finite sums of products
of Hermite functions are dense in the real Schwartz space,
also after the metaplectic change of polarization. If all
individual products had zero period, continuity would give
zero on the test just constructed. Thus some individual
product has nonzero period. Theorem~\ref{thm:realrow257}
gives cuspidality of the corresponding third horizontal lift.
The quantitative bound above uses the normalized smooth bump.
Quantitative Hermite selection would additionally require a
bound on the real seminorms of that bump.
\end{proof}

\subsection{Doubling values with both lengths of theta strings}\label{nav:VI-38}
Fix the original interval, \(E/F\), a word and its unitary splitting
characters. Write \(s=r+6,\ t=t_{\rm r}\), and put
\(\gamma_i=\chi_{m-i}\) for \(0\le i<m\). Write \(\mu_a\) for
the vertical character at step \(a\), including the word twist.
The angular half exponents of \(\gamma_i,\mu_a\) are respectively
\(c_{m-i,\tau}\) and \(-e_{a,\tau}\). The word twist has finite
order and thus contributes zero at the connected complex places.
Use the original marked set \(S_0\). One may also omit places
under the finite twist primes of the word. Write \(S\) for the
resulting set containing the finite ramification under consideration
and \(S_E\) for the places above it. The rational primes under
\(S_0\) split completely by Proposition~\ref{prop:tracelat253}.
At an additional twist prime both residue degree and ramification
index over \(\QQ\) are one by the definition of these places.
Since \(E/\QQ\) is Galois, that rational prime also splits completely.
\srcloc{Formalism/LocalToGlobal/IndefiniteReal/SubgroupCharExtensionOneApply}{91--93}.

At \(w\notin S_E\) let \(B_{0,w}\) be the rank-\(r\) multiset of
a unitary spherical bottom row in the quotient convention. Set
\(Q=Nw\), \(A_b(w)=\{Q^{(b-1)/2-l}:0\le l<b\}\), and
\(D_{ij}=\eta^{-\mathbf1_{i+j\ {\rm odd}}}\). Write
\[
 \mathcal B_{ij,w}=D_{ij}(w)\left(
   B_{0,w}\ \sqcup\!
   \bigsqcup_{a=1}^j\mu_a(w)A_t(w)\ \sqcup\!
   \bigsqcup_{h=0}^{i-1}\gamma_h(w)A_s(w)\right).
\]
These are the unramified parameter lists obtained by first
taking \(j\) forward vertical transfers and then \(i\) horizontal
transfers. They can equally be used to evaluate the Euler
products before constructing a global lift.

\begin{theorem}[The partial standard values at every horizontal step]\leavevmode
\label{thm:mixedvalues257}
Put \(z_*=(s+1)/2\), \(H=(t-s)/2\), an integer positive on the
given sufficiently distant intervals. For every \(0\le j\le m\)
and \(0\le i<m\) the partial standard product for
\(\mathcal B_{ij}\otimes(\gamma_i D_{ij})^{-1}\) is holomorphic and
nonzero at \(z_*\). Its value is
\[
\begin{aligned}
 &L^S(z_*,B_0\otimes\gamma_i^{-1})\\
 &\quad{}\times
 \prod_{h=0}^{i-1}\prod_{a=1}^{s}
        L_E^{S_E}(a,\gamma_h/\gamma_i)
 \ \cdot\
 \prod_{b=1}^{j}\prod_{a=1-H}^{s+H}
        L_E^{S_E}(a,\mu_b/\gamma_i).
\end{aligned}
\]
Write \(f=[F:\QQ]\). The absolute logarithm of the modulus of
the initial factor is at most \(2fr\log\zeta(4)\). The two
products of Hecke values are independent of the bottom
Hecke eigencharacter. This evaluates
the unramified numerator of the doubling formula on the original
fields, including all the long strings when \(j>0\).
\end{theorem}
\begin{proof}
For a step with underlying character \(\nu\), the characters
attached to the larger and smaller spaces are \(\nu D_{\rm old}\)
and \(\nu D_{\rm new}\), respectively. The quotient formula
multiplies the old list by \(D_{\rm new}/D_{\rm old}\) and
adds \(\nu D_{\rm new}A_b(w)\), where \(b\) is the rank
difference. This is Theorem~\ref{thm:inert175} and
Proposition~\ref{prop:split175} for the two kinds of place.
On a horizontal step the physical action on the target is
pulled back from its dual. At the next step the conjugate of
that physical unitary representation enters the quotient
convention. At a split place both operations invert the
standard multiset, restoring the quotient list. At an inert
place the standard list is reciprocal. This gives the stated
recursion from the vertical list of Theorem~\ref{thm:values171}.

Divide by \(\gamma_iD_{ij}\). On a right half-plane of absolute
convergence, disjoint union multiplies Euler products and a root
\(Q^\lambda\) replaces the Hecke argument \(z\) by \(z-\lambda\).
At \(z_*\) the strings of length \(s\) give the integers from
\(1\) to \(s\). Those of length \(t=s+2H\) give the integers
from \(1-H\) to \(s+H\).

We bound the initial factor at this shorter spacing. At a split
place the spherical bound gives \(|\alpha|\le Q^{(r-1)/2}\)
for each \(\alpha\in B_{0,w}\), by the split proof of
Lemma~\ref{lem:sphericalcore252}, which holds in either parity.
At an inert place write \(r=2a+1\). Besides \(e_l\pm e_k\)
of multiplicity two and \(2e_l\) of multiplicity one, the
relative roots include \(e_l\) of multiplicity two, from the
coordinates between the anisotropic line and the \(l\)-th
hyperbolic plane. Thus \(\langle\rho_F,e_1\rangle=
2(a-1)+1+1=r-1\) when \(a\ge1\). Macdonald's inequality in
that lemma applied to the Weyl orbit of \(e_1\) bounds the
reciprocal pairs by \((Nv)^{r-1}=Q^{(r-1)/2}\), with \(v\)
below \(w\). The remaining entry is one, also in rank one.
The difference \(z_*-(r-1)/2\) is four. Each local logarithmic
series is therefore absolutely bounded at \(z_*\), per root,
by \(\sum_{k\ge1}Q^{-4k}/k\). They converge locally uniformly
for \(\Re z>(r+1)/2\). Summing over the at most \(2f\) primes
of \(E\) above each rational prime gives the two asserted
bounds and holomorphic nonvanishing of the initial product.

For \(h<i\) the half exponents of \(\gamma_h/\gamma_i\) equal
\(-(i-h)g_*\) at the selected embeddings. These characters
have nonzero angular exponent and their positive-integer
values are nonzero by Lemma~\ref{lem:hecke171}. For
\(\mu_b/\gamma_i\) the half exponents are
\(-e_{b,\tau}-c_{m-i,\tau}\). Their absolute values exceed
\(H\) at every selected embedding. Indeed \(e_{b,\tau}>0\),
\(v_1\ge0\) and \(c_{m-i,\tau}\ge c_{m,\tau}\ge v_1+g_*\),
where Theorem~\ref{thm:first250} gives
\(g_*\ge2(k+m(t+s+1))^3>t\). These Hecke functions are entire
since their archimedean characters are nontrivial on the
circles. At positive integers apply the same lemma. At
\(a=-u\le0\) here \(0\le u\le H-1\). The complete gamma factors
are regular and Proposition~\ref{prop:splitvalue257} proves
nonvanishing of the partial values even at zero. This proves
both continuation and the formula.

The unramified doubling denominator at parameter \(s/2\) has
factors \(L_F^S(s+k,\omega_{E/F}^{\,s+k})\), \(1\le k\le
u_{ij}\). Their Euler logarithms converge absolutely near these
points because \(s+1>1\), so the same holomorphic nonvanishing
holds for the normalized good-place quotient. The marked
integrals and the prescribed global projection are additional
factors in a period calculation.
\end{proof}

\subsection{The relative norm on the full horizontal row}\label{nav:VI-40}
On one of the original fields let \(v\) be a pure unit vector in
the prescribed bottom real row representation, compact-finite at
each real place. Let \(H_Y(v)\) be the image of this fixed vector
tensored with the full multiplicity space in the old-level
cuspidal spectrum. Give this space and the terminal horizontal
cuspidal space their Petersson norms. Write \(f=[F:\QQ]\),
\(s=r+6\), and \(z_*=(s+1)/2\).

\begin{theorem}[A relative Petersson bound on the horizontal row]\leavevmode
\label{thm:rowgram262}
When \(H_Y(v)\ne0\), one can fix a single sequence of products
of real monomials in Theorem~\ref{thm:globalrow262} for the
entire space \(H_Y(v)\). Let \(U_v\) denote the resulting lift
through \(V_{m,0}\). There is a positive scalar \(C_Y\), common
to every bottom good Hecke system and its finite multiplicities,
such that
\[
  e^{-\delta_Y}I\ \le\ C_Y^{-1}U_v^*U_v\ \le\ e^{\delta_Y}I,
             \qquad \delta_Y=2mfr\log\zeta(s/2).
\]
In particular \(\|C_Y^{-1}U_v^*U_v-I\|_{\rm op}
\le e^{\delta_Y}-1\ll mfr\,2^{-s/2}\le\exp(-x^{\eps/8})\)
on all sufficiently distant original intervals, with the implicit
constant absorbed by increasing the starting point.
\end{theorem}
\begin{proof}
Start with a unit vector on the fixed real line of one cuspidal
copy. Theorem~\ref{thm:globalrow262} gives real tests with
nonzero output at every step. At each real place the joint
theta quotient for the specified source and target is unique up
to scalar by local Howe duality. Contracting the test with the
source vector therefore selects the same real target line in
every copy with that real input. This follows also by factoring
the global theta pairing at that real place after fixing all its
other arguments. On the reference copy the selected line is
nonzero. Choose a unit vector on it. Repeating at each step
determines the local real Rallis factors independently of finite
Hecke eigenvalues. They are nonzero by the positive norm of
each reference lift.

For another bottom system \(\beta\), let \(B_\beta\) denote its
initial standard Satake lists in quotient coordinates. On the
first step use the same real tests and input line. The mixed
formula in Theorem~\ref{thm:globalrow262} evaluates the
diagonal product with these real factors. The marked constants
are independent of the copy. At the other finite places all the
terms from earlier strings and from the doubling denominator
depend only on the step and the common field and characters.
Indeed the standard numerator in Theorem~\ref{thm:mixedvalues257}
at this step, after dividing out its horizontal string product,
is \(L^{S_0}(z_*,B_\beta\otimes\gamma_i^{-1})\).
That theorem proves its nonvanishing. The squared norm is thus
positive and the output has the real line already selected.
Induction gives the exact squared norm multiplier on this copy
in absolute-value notation as
\[
          C_Y\prod_{i=0}^{m-1}
             \left|L^{S_0}(z_*,B_\beta\otimes\gamma_i^{-1})\right|.
\]
The measure factor at a step comes from the fixed Haar measures
in the unfolded matrix coefficient, normalized by the squared
norm of the input. The real factors, good string factors,
denominators and marked volumes absorbed into \(C_Y\) are
therefore common to all copies. The marked scalar formula
extends to every vector in their finite fixed spaces by
polarization. Orthogonality of the invariant theta spans proved
in Theorem~\ref{thm:globalrow262} gives zero off-diagonal
products between distinct orthogonal copies through every
step. Consequently the displayed multipliers are eigenvalues
of \(U_v^*U_v\) on the entire multiplicity space.

At \(w\) outside the marked places every root of the actual
bottom list has absolute value at most \((Nw)^{1/2}\) by
Proposition~\ref{prop:basechange258}. The characters
\(\gamma_i\) are unitary. The absolute logarithm of the modulus
of each varying product factor is thus at most
\[
 r\sum_{w}\sum_{a\ge1}\frac{(Nw)^{-as/2}}{a}
          \ \le\ 2fr\log\zeta(s/2).
\]
Here the sum on the left may include all finite places of \(E\)
by positivity, and there are at most \(2f\) places over each
rational prime. Summing over the \(m\) steps gives the operator
inequalities. Finally \(\log\zeta(t)\le3\cdot2^{-t}\) for
\(t\ge2\) by integral comparison. The quantities \(m,f,r\) grow
at most polynomially in \(x\), and \(r\ge m\gg
x^{1/2-\eps/16}\) on the original intervals, as in
Theorem~\ref{thm:coregram258}. These estimates give the
asserted norm error.
\end{proof}

The designated Bessel coefficient involves applying the
particular Fourier projection of Corollary~\ref{cor:original200}
to a terminal form. The estimates here control its full
Petersson norm. A lower bound for that coefficient still
requires an estimate after the projection.

\subsection{Cancellation of the splittings through the second core}\label{nav:VII-40}
Use \(Y,C_0,S,C_2,Z\) of Proposition~\ref{prop:oppline256}.
Denote by \(K_{D,B}(g,h;\phi)\) the original theta kernel on
\(\operatorname{Hom}(-D,B)\), with the argument on the larger
space \(D\) written first. At the first core edge and horizontal
edge use the underlying character \(\chi_m\), and at the
second such edges use \(\chi_{m-1}\). To a space of even
rank attach that character, and to a space of odd rank attach
its product with \(\eta^{-1}\). This is the original prescription
in \srcloc{Formalism/LocalToGlobal/IndefiniteReal/HasInfExpF}{13--23}.
Let \(U\) denote the rational adjoint intertwiner on line tests
from Proposition~\ref{prop:oppline256}, normalized by adelic
Fourier inversion and rational Poisson summation.

\begin{proposition}[An identity with the original four splittings]\leavevmode
\label{prop:cancelsquare256}
For Schwartz tests \(p,q,b\) on the \(Y\)-part, the old positive
line part and the second core edge, respectively, one has
\[
\begin{split}
 K_{S,C_0}(g,h\oplus1;p\otimes q)\,K_{C_2,S}(k,g;b)
  =K_{S,Y}(g,h;p)\,K_{Z,S}(k\oplus1,g;b\otimes Uq).
\end{split}
\]
There is also a line theta function \(L_q(g)\), independent of
\(h,k\), such that the first kernel on the left factors as
\(K_{S,Y}(g,h;p)L_q(g)\), and the second kernel on the right
as \(K_{C_2,S}(k,g;b)L_q(g)\). These equalities hold for the
entire adelic groups.
\end{proposition}
\begin{proof}
The two total symplectic spaces identify over \(F\). The map is
the identity on the common \(Y\) and \(C_2\) parts and the
Hermitian adjoint on the line part. It intertwines the actions
of \(U(Y)\times U(S)\times U(C_2)\) on the underlying tensors.
The oscillator for a direct sum is the tensor product of the
oscillators. In a position model this follows by factoring the
translations and modulations and using the irreducibility of the
Heisenberg representation. A rational change of polarization
preserves the global theta functional under the normalized
intertwiner. For a Fourier exchange on a subspace this is
Poisson summation for its rational points. For a rational
change of basis it is reindexing. Consequently the ratio of
the two group splittings after transport is a continuous
unitary character, trivial on rational points. We check this
character at the places where the scalar is explicit.

At a split finite place write \(a,b_0,c\) for the determinants
on the chosen halves of \(h,g,k\), respectively. Evaluate all
characters on that half. For \(\operatorname{Hom}(-D,B)\)
the circle factor is the character attached to \(D\) evaluated
on \(\det h_B\), divided by the one attached to \(B\) evaluated
on \(\det g_D\). The other factor is the positive square root
of the Jacobian of the position change of variables. Fourier
exchange carries the unitary substitution on a position space
to the unitary dual substitution, preserving its circle scalar.
This can be checked by changing variables in the Fourier
integral. Thus the two circle products to compare are
\[
 \begin{aligned}
 &\chi_m(a)\chi_m(b_0)^{-1}
           \chi_{m-1}(b_0)\chi_{m-1}(c)^{-1},\\
 &\chi_m(a)(\chi_m\eta^{-1})(b_0)^{-1}
           (\chi_{m-1}\eta^{-1})(b_0)\chi_{m-1}(c)^{-1}.
 \end{aligned}
\]
They agree, including when these characters are ramified.
The equal positive Jacobians are forced by unitarity of the
same substitution. This proves triviality on each split finite
group. The determinant convention is specified in
\srcloc{Formalism/LocalToGlobal/IndefiniteReal/LocalNewformFunction}{51--75}.

Here is why those places suffice for the adelic equality. At
almost every nonsplit place the ratio is trivial on the
hyperspecial subgroup. Indeed the original splittings are
unramified there, and the rational isometry and its inverse
are integral there. The standard lattice vacua match under
the change of polarization, with any common scalar immaterial.
A character of such a unitary group trivial on that compact
subgroup is trivial on the whole group. To check this, choose
integral hyperbolic planes and, in odd rank, one remaining
unit line. Integral hyperbolic splitting follows by the
residue norm and Hensel argument of
Proposition~\ref{prop:vacuumcore254}. The Cartan decomposition
reduces the assertion to elements acting on one plane by
\(\varpi^b,\varpi^{-b}\) for a uniformizer of \(F_v\).
They lie in the special unitary group of the hyperbolic plane,
the split \(\mathrm{SL}_2(F_v)\). This group is generated by
its two elementary unipotent subgroups, each consisting of
commutators by conjugation with a diagonal element whose square
differs from one. Every character is trivial there.

The ratio is thus supported at finitely many places, including
infinity. A unitary group has weak approximation at every
finite set of places of \(F\). Indeed the Cayley transformation
\(X\mapsto(1-X)(1+X)^{-1}\) identifies open subsets of its
skew-Hermitian Lie algebra and the group. The group open subset
is locally dense because scalar norm-one units arbitrarily close to one
can move a given element off the divisor with eigenvalue \(-1\).
The Lie algebra is an \(F\)-vector space, so weak approximation
holds there, and then on the group by this parametrization.
This also applies to the product of the three groups. Triviality
on its dense rational image in the places of support forces
the remaining character to be trivial. Equivariance of the
theta functional now proves the identity of the two products.

For the separate factorizations, first restrict the splitting
of \(S,C_0\) to \(U(Y)\). Its action on the line tensor is the
identity up to a character. Comparison with \(K_{S,Y}\) gives
the same split scalar \(\chi_m(a)\) and hence exact agreement
by the preceding character argument. On \(U(S)\), divide out
the action on the \(Y\)-tensor. The resulting operators act
only on the line tensor, since they commute with the entire
Heisenberg action on the other factor. The group law defines
a line oscillator action there. On rational elements it
preserves the line theta functional by the product formula
for the rational theta functional. This gives \(L_q(g)\).
The already proved equality of the transported actions on
the total tensor identifies the action on the \(C_2\)-and-line
part with the product of this line action and the original
\(C_2,S\) action. Applying the rational theta functional gives
the second factorization as well.
\end{proof}

\section{Norms on the even core}\label{nav:g-VI-4}

\subsection{The spherical bound at the smaller core spacing}\label{nav:II-47}
Here \(n=k+m\) is the even rank of the first core. At a good place
\(v\) of \(F\) let \(q=Nv\), let \(w\mid v\) in \(E\) and write
\(Q=Nw\). In the normalized Satake convention the spherical character
of the trivial representation is the Weyl orbit of the half sum of
positive roots with their root-space multiplicities, exponentiated
with base \(q\). Use the standard base-change multiset of size \(n\)
at \(w\).

\begin{lemma}[The sharp general exponent for a unitary spherical character]
\label{lem:sphericalcore252}
For a spherical eigencharacter in the unitary bottom space of
Proposition~\ref{prop:heckecore252}, each entry \(\alpha\) of its
standard multiset satisfies
\[
                 Q^{-(n-1)/2}\le |\alpha|\le Q^{(n-1)/2} .
\]
The estimate also holds for the conjugate representation. In
particular, with \(d=n+4\), the difference between \((d+1)/2\)
and the upper exponent is exactly three.
\end{lemma}
\begin{proof}
We recall Macdonald's bounded spherical function theorem in the
normalization needed here. Let \(t\) be the normalized Satake character
of a hyperspecial spherical algebra, on a coroot translation lattice
of its semisimple root system. If the character is bounded on the
normalized double-coset averages, then for dominant \(\lambda\) in
that lattice and each relative Weyl element \(w_0\),
\[
           |t^{w_0\lambda}|\le q^{\langle\rho_F,\lambda\rangle} .
\]
Here \(2\rho_F\) is the sum of the positive relative roots counted
with their dimensions over the local field. This is the convex-hull
form of Macdonald's Theorem~4.7.1. See also
\href{https://www.maths.usyd.edu.au/u/pubs/publist/preprints/2005/parkinson-36.pdf}{J.~Parkinson,
Appendix~A, Remark~A.1 and Theorem~A.3}. In these conventions the
positive character on the right is the square root of the modulus
in the expanding direction. This follows by taking the determinant
of conjugation on the positive-root spaces. Our eigencharacter has
absolute value at most one on each normalized double coset since the
corresponding operator is a probability average of unitary translations.
For a unit simultaneous eigenvector, its normalized local matrix
coefficient is constant on each double coset and equals the corresponding
averaged eigenvalue. It is precisely the bounded spherical function
with this Satake character.
In applying the theorem on the derived group one may restrict a
spherical vector to its derived-group action. Its spherical
coefficient there has the restricted Satake parameter. Indeed the
normalized unramified induction model gives its spherical function
by the compact average of the unramified inducing section, and the
derived-group Cartan and modulus are the restrictions of those of the
original group. More explicitly, its hyperspecial group is the product
of the integral torus and the hyperspecial subgroup of the derived
group. At a split place this follows by taking a diagonal unit of the
required determinant. At an inert place every norm-one determinant is
\(u/\bar u\) for a unit in the unramified quadratic extension, by
Hilbert 90 and multiplication by a base-field uniformizer power. The
two entries \(u,\bar u^{-1}\) in one hyperbolic plane realize this
determinant. The inducing section is left invariant by these torus
units, so its compact average on derived-group elements equals the
average for the derived group. A positive integral
multiple of any coweight in the real span of the coroots suffices
when the coweight itself is outside the coroot lattice, since one can
take that root of the resulting absolute-value inequality.

At a split place write \(t=(t_1,\ldots,t_n)\) for the parameters
of \(\mathrm{GL}_n(F_v)\). The central character on the scalar
uniformizer has absolute value one, because that translation and its
inverse are unitary. Hence \(\prod |t_j|=1\). Take the dominant
coroot-lattice vector \(n e_1-\sum_j e_j\), where the \(e_j\) are the
standard coordinate vectors. The half-sum has entries
\((n+1-2j)/2\). The inequality for its Weyl conjugates gives
\(|t_i|^n\le q^{n(n-1)/2}\) for every \(i\). Applied to the unitary
contragredient, whose standard parameters are the inverses, it gives
the lower bound. The second prime over a split place has the dual
parameter, so both bounds hold there too.

At an inert place write \(n=2a\). In relative coordinates for the
quasi-split unitary group the positive roots are \(e_i-e_j,e_i+e_j\)
for \(i<j\), of dimension two, and \(2e_i\), of dimension one.
The paired off-diagonal entries in an isotropic basis supply one free
entry in the quadratic extension, and an anti-Hermitian diagonal entry
supplies one coordinate over \(F_v\), which gives these dimensions.
Thus \(\langle\rho_F,e_1\rangle=2(a-1)+1=n-1\). Up to multiples,
\(e_1\) is a dominant coroot-lattice translation, and its Weyl orbit
contains \(\pm e_i\) for every \(i\). The standard multiset over \(E_w\)
consists of \(t^{e_i},t^{-e_i}\). To see the normalization, write
dual Frobenius as \(s\vartheta\) with \(s\) diagonal and \(\vartheta\)
reversing and inverting diagonal entries. Squaring gives \(s\vartheta(s)\),
whose entries are \(s_i/s_{n+1-i}\) and their inverses. This ratio is
precisely evaluation on the translation with exponent \(1\) in the
\(i\)-th position and \(-1\) in the opposite position. The bound is
therefore \(q^{n-1}=Q^{(n-1)/2}\) in either direction. The trivial local representation attains the upper
exponent, as its Satake character is the indicated half-sum. The last
assertion follows by subtracting the two exponents.
\end{proof}

\subsection{Relative norms of the even core lift on its Hecke spaces}\label{nav:II-48}
On a fixed original field let \(f_E=[F:\QQ]\) and let \(D_0\)
be the bottom real representation. Fix a unit vector \(v\) in
it which is a pure tensor of compact-finite vectors over the
real places. In the cuspidal space with level \(K_{0,f}\)
and this real representation, write \(H(v)\) for the image
of \(v\) tensored with its full multiplicity space. Suppose
\(H(v)\ne0\). Give all spaces their Petersson norms. Keep the
finite indicators of Theorem~\ref{thm:upcore254}.

\begin{theorem}[Uniform relative norm comparison at fixed real input]\leavevmode
\label{thm:corenorms255}
There is a single choice of real monomial tests along the entire
even core giving a composition \(T=T_{m-1,p_{m-1}}\cdots T_{0,p_0}\)
such that for all nonzero \(a,b\in H(v)\),
\[
  e^{-4m f_E n\log\zeta(3)}
   \ \le\ \frac{\|Ta\|^2\|b\|^2}{\|Tb\|^2\|a\|^2}
   \ \le\ e^{4m f_E n\log\zeta(3)} .
\]
The same real tests work in every good Hecke space of \(H(v)\).
On each one the squared norm is multiplied by an exact scalar,
independent of the vector, including its finite multiplicity.
\end{theorem}
\begin{proof}
Choose a nonzero input on the prescribed real line in one
irreducible cuspidal copy. Theorem~\ref{thm:upcore254} gives
successive product monomials with nonzero outputs for this
input. At a real place the first test and the prescribed real
source vector determine one target vector up to scalar.
Indeed the local pairing into the irreducible source and target
is unique up to scalar by Howe duality. Contracting with the
same source vector in that pairing gives the same line in the
fixed target representation for every occurrence in a global
lift. It is nonzero by the choice of the reference input.
Use a unit vector on that line. Since the tests and real input
were pure tensors, this gives a pure real line for the next
step. Induction gives unit lines at every real place for all
the steps. The diagonal real Rallis factors on them are
nonzero, since their product for the reference lift computes
a nonzero squared norm. They depend only on the chosen tests,
unit vectors and local pairings.

Consider any bottom good Hecke system \(\eta\) on \(H(v)\).
Let \(B_\eta\) be its initial Satake lists in quotient
coordinates. The first norm formula of
Theorem~\ref{thm:upcore254} applies with the same real tests.
Its global measure constant is common because the open-orbit integral
uses the global Petersson matrix coefficient of the input,
which, on a pure tensor, is its squared norm times the product
of normalized local matrix coefficients. Passing to the
standard product Haar measures multiplies all such formulas
by the same measure factor. At marked places the product is
scalar on the entire space of finite fixed vectors, with
factors \(D_{i,v}\). At unmarked places there is a single
normalized spherical line. The varying part at step \(i\),
after extracting the factors \(C_i\) and the doubling
denominator of Theorem~\ref{thm:coreeuler252}, is precisely
\[
      L^{S_0}(z_*,B_\eta\otimes\kappa_i^{-1}),\qquad
                                  z_*=(d+1)/2 .
\]
The absolute value of this factor is nonzero. Thus the same
tests give a nonzero lift on this system and put its real
component on the previously determined line. This justifies
using the next norm formula for this input. Iterating, the
ratio \(\|Ta_\eta\|^2/\|a_\eta\|^2\) for one nonzero vector
on the line in a single global copy is
\[
       C_{\rm fix}\prod_{i=0}^{m-1}
            \big|L^{S_0}(z_*,B_\eta\otimes\kappa_i^{-1})\big| .
\]
Here \(C_{\rm fix}>0\) is common to all the copies. It consists
of the absolute real factors, measure constants, marked factors,
and the common good factors extracted above. Any phase in
writing the mixed-pairing formula disappears on taking absolute
values, as its diagonal left side is a squared norm.

For finite vectors that are sums of pure tensors in one copy,
the mixed formula gives the same scalar times their inner
product. This uses the marked scalar identity and the
one-dimensional good spherical lines. Mixed products from
orthogonal invariant copies vanish at the first step by
open-orbit unfolding. Their full cuspidal theta images for
arbitrary tests are orthogonal by that same calculation,
so their invariant spans are orthogonal. Applying this at
each further step gives the assertion for sums over copies,
including repeated ones. Thus on an arbitrary vector of \(H(v)\)
the squared norm ratio is a weighted mean of those positive
scalars, with weights proportional to its squared norms in the
orthogonal Hecke spaces.

Finally Lemma~\ref{lem:sphericalcore252} and the logarithmic
estimate in Theorem~\ref{thm:coreeuler252} put the absolute
logarithm of each displayed initial \(L\)-factor below
\(2 f_E n\log\zeta(3)\). For the ratio between any two
positive scalars the common factor cancels and the sum of the
\(m\) logarithmic differences has absolute value at most
\(4m f_E n\log\zeta(3)\). The same bounds hold between any
two of their weighted means. This proves the comparison on
the entire fixed real input space.
\end{proof}

This comparison controls variation among the bottom Hecke spaces
for one choice of real tests. An absolute estimate on the integral
projected period in Section~\ref{sec:missing} still depends on its
normalization in the arithmetic lattice.

\subsection{Exponential accuracy of the relative core norms}\label{nav:II-49}
Use the original field \(E/F\), ranks \(n,d=n+4\), and depth
\(\mathfrak m\) of Theorem~\ref{thm:upcore254}. Let
\(f=[F:\QQ]\). For a fixed unit pure real vector \(v\) that is
finite under the maximal compact, take \(H(v)\) as in
Theorem~\ref{thm:corenorms255}. Its elements have the prescribed
finite level. Write \(T_v\) for the composition of the \(m\)
even core lifts of that theorem restricted to \(H(v)\).

\begin{theorem}[The normalized Gram operator of the full core ascent]\leavevmode
\label{thm:coregram258}
For nonzero \(H(v)\), use the common real monomials and
original finite tests of Theorem~\ref{thm:corenorms255}. There
is one \(C_{\rm fix}>0\) independent of the bottom finite
Hecke system such that, in the order of Hermitian operators,
\[
  e^{-\delta}I\ \le\
  C_{\rm fix}^{-1}T_v^*T_v
   \ \le\ e^\delta I,\qquad
  \delta=2mfn\log\zeta(d/2).
\]
In particular
\[
 \big\|C_{\rm fix}^{-1}T_v^*T_v-I\big\|_{\rm op}
 \le e^\delta-1 \ll mfn\,2^{-d/2}
 \le \exp(-x^{\eps/8})
\]
on every sufficiently distant original interval. The
Hilbert--Schmidt norm squared of this difference is at most
\(\dim H(v)\,(e^\delta-1)^2\). Changing a group measure by a
constant changes \(C_{\rm fix}\) by the resulting common norm
factor. For the horizontal doubling values of
Theorem~\ref{thm:mixedvalues257} arising from an actual bottom
cuspidal row of rank \(r\), the absolute logarithm of the
modulus after division by the two string products is at most
\(2fr\log\zeta((r+6)/2)\) at every step.
\end{theorem}
\begin{proof}
On a good bottom core Hecke system \(\beta\), the proof of
Theorem~\ref{thm:corenorms255} computes the squared norm
multiplier exactly as
\[
  C_{\rm fix}\prod_{i=0}^{m-1}
      \left|L^{S_0}\big((d+1)/2,B_\beta\otimes\kappa_i^{-1}\big)\right|.
\]
Here \(B_\beta\) is the original bottom standard parameter,
and the \(\kappa_i\) are the prescribed unitary characters.
The common scalar includes the real local pairings, the
marked rectangle volumes, the string factors and the
doubling denominators. The mixed norm formula makes distinct
orthogonal cuspidal copies orthogonal under all these lifts.
Thus the displayed multipliers are eigenvalues of \(T_v^*T_v\)
on the whole finite multiplicity space at this real vector.

At every \(w\notin S_{0,E}\) Proposition~\ref{prop:basechange258}
bounds the absolute value of a bottom Satake root by
\((Nw)^{1/2}\). Multiplication by \(\kappa_i(w)^{-1}\) preserves
this modulus. At the indicated special value each local
Euler logarithm is bounded in absolute value by
\(n\sum_{a\ge1}(Nw)^{-ad/2}/a\). Above any rational prime
there are at most \(2f\) places and each norm is at least that
prime. Addition gives \(2fn\log\zeta(d/2)\). Summing over
\(i\) proves the two bounds in the order of operators.
For a self-adjoint operator, the largest absolute deviation
of an eigenvalue bounds the operator norm of its difference
from the identity. The sum of the squared deviations gives
the Hilbert--Schmidt bound.

For \(a\ge2\), comparison with the integral of \(x^{-a}\)
on \([2,\infty)\) gives \(\log\zeta(a)\le3\cdot2^{-a}\).
Here \(f,m,n\) are polynomially bounded in \(x\) and
\(n\ge m\gg x^{1/2-\eps/16}\) by
Theorem~\ref{thm:rangenew50}. Therefore \(\delta\) tends to
zero uniformly and the term \(2^{-d/2}\) absorbs the polynomial
factor with the displayed margin for \(0<\eps<1/52\).
For the horizontal assertion, the factor left after division is
the bottom standard function at \((r+7)/2\). Applying the same
logarithmic series bound with \(r\) roots and subtracting
\(1/2\) from this special argument proves the stated estimate,
independently of both row and horizontal indices.
\end{proof}

\subsection{The second horizontal restriction for a lowest-type row}\label{nav:II-50}
Use the test and measure conventions of
Theorems~\ref{thm:individual259} and~\ref{thm:secondbound260}.
Fix a pure unit vector \(v_0\) in the lowest compact type
of the bottom core and choose the even core tests for this
vector once on the given field. Let \(\phi\) be their first
real test in the downward convention. Write \(L_\tau\) for its
local energies, \(D=n(n+d)\), \(N=n_m\), and \(A(D,L_\tau)\)
as in Lemma~\ref{lem:weightedhermite260}. The constant
\(C_{\rm fix}\) is that of the complete even core ascent with
this choice and with terminal probability measure.

\begin{theorem}[An absolute restriction estimate for lowest-type rows]\leavevmode
\label{thm:lowesttwo261}
For every unit vector \(u\) in the lowest compact type of a
prescribed old cuspidal row, there is one terminal integral
coefficient \(e\), one product of normalized Hermite functions
\(h_\alpha\), and one integral marked residue \(\nu\) for which
the actual second horizontal restriction satisfies
\[
 \begin{split}
 |J_{e,\alpha,\nu}|\ \ge{}&
 \frac{\sqrt{C_{\rm fix}}}{\|\phi\|_2\prod_\tau A(D,L_\tau)}
 \frac{\operatorname{vol}(J)}{d_\sigma}
 \left(\frac{d_\pi}{\operatorname{vol}(K_f)}\right)^{1/2}\\[-1pt]
 &{}\cdot
 \frac{\|u\|_\infty}{(\operatorname{Nm}\mathfrak m)^{N^2}}\,
                        \exp[-P_{\psi,\eps}(x)] .
 \end{split}
\]
Here \(P_{\psi,\eps}\) is a single polynomial for the family
on all sufficiently distant original intervals. The factor
\(\|u\|_\infty\) is at least \(\operatorname{vol}([H])^{-1/2}\)
in the measure in which the row is a unit vector.
\end{theorem}
\begin{proof}
Take \(a\) and the functional \(\ell_a\) of
Proposition~\ref{prop:poincarenorm261}. With \(\delta\) from
Theorem~\ref{thm:coregram258}, Theorem~\ref{thm:individual259}
selects one terminal coefficient \(e\) for which
\[
 |\ell_a(P_{v_0}T^\dagger e)|\ge
 \frac{e^{-\delta/2}\sqrt{C_{\rm fix}}\,\|\ell_a\|_*}
           {H_*\sqrt{d_\rho}\,N^{d_0} B^{2a_0}} .
\]
We explain how to estimate this particular functional by the
actual horizontal periods of Theorem~\ref{thm:secondbound260}.
Use the Schwartz tests \(\Psi_j\) before Hermite truncation in
the proof of Proposition~\ref{prop:projecttests260}. Replace
them by \(\omega(a)\Psi_j\). Equivariance of the bottom group
gives \(F_e(\omega(a)\Psi_j)=R(a)F_e(\Psi_j)\). The latter
converge to \(R(a)P_{v_0}T^\dagger e\) in the translate by
\(a\) of one fixed finite-dimensional space of smooth cusp
forms. The first restriction against \(u\) is continuous
there. Hence its values converge to the functional in the
last display.

To bound the absolute Hermite coefficients of these translated
tests, use the integral formula for \(\Psi_j\) at a real place.
Before tensoring, the integrand after translation is
\(\omega_\tau(a_\tau k t)\phi_\tau\), with measures \(b_\tau\)
on the compact group and \(q_{\tau,j}(t)\,dt\) on the real
group. Now \(W(a_\tau k t)\le W(a_\tau)W(t)\). Lemma
~\ref{lem:weightedhermite260} and integration of the absolute
coefficients therefore bound their local sum by
\[
 \sqrt{\zeta(2)}\,A(D,L_\tau)\|\phi_\tau\|_2
        M_\tau(v_{0,\tau})\,\|b_\tau\|_{\rm TV} W(a_\tau)^{2s},
 \qquad s=\lfloor D/2\rfloor+1 .
\]
There is only one compact type \(\kappa_\tau\) in \(v_{0,\tau}\).
Its compact character projector has variation at most
\(\dim\kappa_\tau\), by orthogonality of the irreducible
character in Haar probability. The highest weight of the
lowest compact type is obtained from the Harish-Chandra
parameter by the difference of the compact and noncompact
half sums. Each coordinate changes by at most \(n\).
The compact Weyl dimension formula thus gives
\(\dim\kappa_\tau\le T_0^{C n^2}\). The lowest-type estimate
in Lemma~\ref{lem:weightedhermite260} bounds \(M_\tau(v_{0,\tau})\)
by \(T_0^{C n^3}\). Proposition~\ref{prop:poincarenorm261}
bounds \(W(a_\tau)\) by the same power, and \(2s=O(n^2)\).
Multiplying over the real places bounds the whole absolute
coefficient sum by
\[
             \|\phi\|_2\,T_0^{C f n^5}\prod_\tau A(D,L_\tau).
\]
This bound holds before and after any finite Hermite truncation.
Truncate with error tending to zero in the first restriction,
using Schwartz continuity as in
Proposition~\ref{prop:projecttests260}.

Each remaining Hermite product has first restriction equal to
the average of the second horizontal periods with the finite
Fourier phases, whose total absolute weight is one. This is
the exact two-kernel identity in Theorem~\ref{thm:secondbound260}.
The preceding coefficient bound times \(\sup_{\alpha,\nu}
|J_{e,\alpha,\nu}|\) therefore dominates
\(|\ell_a(P_{v_0}T^\dagger e)|\). Choose one pair attaining at
least half the resulting positive lower bound.

Insert Proposition~\ref{prop:poincarenorm261} for \(\|\ell_a\|_*\).
By Theorem~\ref{thm:individual259}, the logarithm of the
terminal denominator is bounded by a polynomial plus
\(N^2\log\operatorname{Nm}\mathfrak m\). The factors \(2\),
\(e^{\delta/2}\), and \(T_0^{Cfn^5}\) have logarithms bounded
by polynomials in \(x\) on the very same intervals, since
\(\log T_0\ll x^\eps\) and \(f,n,m\) are polynomially bounded.
Combining the polynomials gives the asserted lower bound.
Finally \(\|u\|_2^2\le\operatorname{vol}([H])\|u\|_\infty^2\)
proves the last assertion.
\end{proof}

This estimate accounts for the initial restriction functional
for all lowest-type rows. An absolute arithmetic normalization
of the even core theta factors, the comparison through later
horizontal lifts with the designated integral projection, and
the full residual local cohomology bound on the common intervals
are still needed for Section~\ref{sec:missing}.

\subsection{The finite position test on the first norm shell}\label{nav:II-51}
Suppose \(m\ge4\). At the third horizontal edge put \(C=C_2\),
\(D=C_3\), \(S=V_{2,0}=C\perp Ez\), and
\(T=V_{3,0}=D\perp(Ee+Ee')\). The unscaled grid columns satisfy
\(h_S(z,z)=-1\), \(h_T(e,e)=h_T(e',e')=0\), and
\(h_T(e,e')=1\). The distinguished isotropic column and its
prescribed image in the grid frames are \(-\xi_0e\) and
\(-\xi_0z\), respectively. The position variable in unscaled
coordinates is \(x=X(e)\) for \(X\in\operatorname{Hom}(-T,S)\).
Thus the prescribed position is \(z\).
\srcloc{Development/AbelianVarietyTheory/Definitions/TwistedInnerProducts/IsAuxiliaryVectorWittDec}{160--175}.

Keep \(\mathfrak A,\mathfrak m\) from
Proposition~\ref{prop:tracelat253}. Put
\(\mathfrak l=(\mathfrak A\mathfrak m)/
(\bar{\mathfrak A}\bar{\mathfrak m})\),
\(\mathfrak j=\mathfrak m\bar{\mathfrak m}\),
and \(L=\mathfrak l^{-1}\Lambda_S\). A hat denotes the
product of all finite completions of a lattice.

\begin{proposition}[A test at the distinguished position]\leavevmode
\label{prop:shell263}
In the model with position \(x\) and residual oscillator
\(\operatorname{Hom}(-D,S)\), the indicator
\(1_{z+\mathfrak j\widehat L}(x)\) is admissible for the
original finite depth on the hyperbolic plane. Here
\(L=\mathfrak l^{-1}\Lambda_C\oplus\mathcal O_E z\) is
unimodular. With the local position measures giving the original
unshifted lattice volume one at the marked places, the inverse
change to the original marked polarizations expresses this
test as a sum of exactly
\((N\mathfrak m)^{\dim_E S}\) distinct integral coset indicators
with coefficients of absolute value one. Those coefficients
are roots of unity.
\end{proposition}
\begin{proof}
The negative line lattice is \(\mathfrak l z\) by
Proposition~\ref{prop:oppline256}. Since
\(\mathfrak l\bar{\mathfrak l}=1\), multiplication of a
self-dual Hermitian lattice by \(\mathfrak l^{-1}\) preserves
self-duality. For the two new columns the staircase rule in
the original lattices uses
\(A=\mathfrak A\bar{\mathfrak m}^{-1}\) and
\(B=\bar{\mathfrak A}\mathfrak m^{-1}\). Its two ideals are
\[
 I_e=A^2/B=\mathfrak A\mathfrak l/\bar{\mathfrak m},
 \qquad I_{e'}=A/B^2=\mathfrak l\mathfrak m/\bar{\mathfrak A}.
\]
In particular \(I_{e'}=\bar I_e^{-1}\). The hyperbolic part
of the full map lattice has column lattices
\(L_e=\mathfrak A I_e^{-1}\Lambda_S=\bar{\mathfrak m}L\)
and \(L_{e'}=\mathfrak A I_{e'}^{-1}\Lambda_S\).
They are exact duals for the alternating trace pairing by
Proposition~\ref{prop:tracelat253}, since the two columns are
isotropic and pair to one. The depth displacements multiply
each of these lattices by \(\mathfrak m\). Translations of
the \(x\)-variable therefore form \(\mathfrak j\widehat L\).
The frequencies \(\mathfrak m\widehat L_{e'}\) have annihilator
\(\bar{\mathfrak m}^{-1}\widehat L_e=\widehat L\), since
the cross pairing is conjugate-linear on one coordinate.
The proposed coset is inside \(\widehat L\), is stable under
the translations, and all frequency phases are trivial on it.
The extra phase for a doubled mixed displacement is trivial
by this same duality. At a split place the two components
of a displacement supported on a single idempotent have
zero mutual pairing. This also checks the equations at two.
\srcloc{Development/AbelianVarietyTheory/Definitions/Divisors/CmHostImaginaryGeneratorNeZero}{120--145}.

At a marked \(v\) put \(a=a_v\), \(w=w(v)\), and choose a
uniformizer \(\pi\) of \(F_v\). The original positions are
\(X(e)_w,X(e')_w\). To obtain the mixed positions, Fourier
transform the second one to \(X(e)_{\bar w}\) by their trace
duality. Now \(L_{e,w}=L_w\) and
\(L_{e,\bar w}=\pi^a L_{\bar w}\). The desired indicator
on the two new positions is the product of the indicators of
\(z_w+\pi^a L_w\) and \(z_{\bar w}+L_{e,\bar w}\).
The inverse transform of the second is the indicator of
\(L_{e',w}\) multiplied by the additive character paired with
\(\pm z_{\bar w}\). Indeed the annihilator of \(L_{e,\bar w}\)
is \(L_{e',w}\) and both have volume one in dual measures.
This character is constant modulo \(\pi^a L_{e',w}\), since
\(\pi^a z_{\bar w}\in L_{e,\bar w}\). Partitioning the lattice
into these cosets produces exactly \((Nv)^{a\dim_E S}\) terms,
with the stated coefficients. At the other places the indicator
is the full position lattice. Partial Fourier transformation
exchanges the unmarked vacuum with the corresponding dual
lattice vacuum, using the unit normalization of that line.
Multiplying the marked counts gives the assertion.
\end{proof}

\subsection{Unfolding the Jacobi period on a full frame}\label{nav:II-53}
Let \(i=2r+2\), \(i+1\le m\), and put \(\ell=r+1\).
The two row spaces are \(S=(C\perp L)\perp\mathbb H^r\)
and \(T=D\perp\mathbb H^\ell\), with \(C=C_i,D=C_{i+1}\).
The middle vector \(z\) of \(S\) has norm \(-1\) and spans
\(L\). The scaled frame \(X=X_*\) in \(S\) ends with
\(\zeta=d^\ell z\), where \(d=-\xi_0\).
Write \(N_S,N_T\) for the flag radicals and \(G=U(S)\).
The Jacobi theta on \(N_T\rtimes U(D)\) uses this same
\(\zeta\) and a Schwartz test \(q\) on
\(\operatorname{Hom}(-D,L)\). The Bessel character on
\(N_S\) uses \(\zeta\) as its last vector.

For a Schwartz vector \(f\) on the residual oscillator of
\(\operatorname{Hom}(-D,C\perp L)\), write \(C_q f\) for
contraction against \(q\) in the line variable, linear in
\(f\). Use compatible dual Haar measures, as in
Propositions~\ref{prop:unfold263} and~\ref{prop:besorbit265}.
Use the finite position tests of
Theorem~\ref{thm:flagfinite264} and write \(\mathcal E(g;\Phi)\)
for the adelic integral in its two positions. The opposite
row position is zero on this edge.

\begin{proposition}[An orbital formula for the global Jacobi period]\leavevmode
\label{prop:jacorbit266}
Let \(B,A_D\) be compact-finite cusp forms on \(G,U(D)\).
The Fourier--Jacobi period of the horizontal lift of \(B\)
against \(A_D\) with line test \(q\) equals
\[
  \int_{R(\mathbb A_F)\backslash G(\mathbb A_F)}
       J_q(g;\mathcal E(g;\Phi))\,d\dot g,\qquad
       R=N_S\rtimes U(C),
\]
where
\[
 J_q(g;f)=\int_{[U(D)]}\int_{[U(C)]}\int_{[N_S]}
 A_D(h)B(nk g)\overline{\psi_{B,S}(n)}
             K_{D,C}(h,k;C_q f)\,dn\,dk\,dh .
\]
The functional is continuous in the real translated form
and test when \(g\) varies on a compact set, with fixed finite
data. In particular for a descending test \(b\) it satisfies
\[
     J_q(1;b\otimes q)=\|q\|_2^2\,
           P_{B,S}\big(B,\operatorname{desc}_b A_D\big).
\]
\end{proposition}
\begin{proof}
In the period defining the left side, expand the horizontal
theta in its first rational positions and multiply by the
conjugate of the line theta and the conjugate upper
unitriangular character on \(N_T\). On the center, the two
theta factors have the characters of \(x^*h_Sx\) and
\(\operatorname{diag}(0,\ldots,0,h_S(\zeta,\zeta))\),
respectively. A shear multiplies each fiber by the half
alternating trace pairing with its opposite-column shear,
so both characters use the same normalization. Orthogonality
and nondegeneracy of that pairing select precisely the frames
with the second Gram matrix. A marked position box has
independent columns after normalization in a split component.
Thus no rank-deficient rational tuple contributes. By Witt
extension the retained frames are a single rational orbit.
Right upper unitriangular changes preserve this Gram matrix.
The stabilizer modulo these changes consists of \(N_S\) and
the core isometries fixing \(z\). Indeed the action on the
first \(r\) columns is upper unitriangular, and the image of
the last column equals itself modulo their span. On the
middle quotient this means that it fixes the nondegenerate
line pointwise. Its remaining Levi factor is exactly \(U(C)\).
Conversely every element of \(N_S\rtimes U(C)\) has these
properties. Witt extension realizes every rational upper
unitriangular change of the whole frame.

Unfold the frames as \(\gamma^{-1}Xv_0\) for
\(\gamma\in R(F)\backslash G(F)\), \(v_0\in U_\ell(E)\), as
in Proposition~\ref{prop:besorbit265}. In changing the
noncentral and central coordinates by \(v_0\), the last row
of the noncentral block and the last diagonal component of
the transformed central matrix are unchanged. These are the
only such coordinates entering the line theta. The rational
sum in \(v_0\) again unfolds the upper-group quotient to the
entire upper group with weight \(\overline{\chi(v)}\).

On the fiber of \(X\), write the noncentral block as
\(\ell\) rows on \(D\). Composing its first \(r\) rows with
the first \(r\) columns of \(X\) gives all displacements into
the isotropic columns of \(S\). The line theta is independent
of these rows. In the second polarization their modulations
are paired with all the opposite rows by the perfect trace
pairing. Integration over their rational quotients therefore
selects precisely the zero opposite position. If \(r=0\)
there are no such rows.

The final noncentral row gives the maps \(w\mapsto
\zeta\,b_\ell(w)\) from \(D\) to \(L\), and hence ranges
through the whole line phase space. The two theta factors
have the same Heisenberg action there. In fact both use the
identical retained vector \(\zeta\), and their commutator is
the alternating trace pairing on these maps. The other rows
pair trivially with the line because the corresponding
columns are orthogonal to \(\zeta\). Inserting the final row
in the residual columns is translation, and its adjoint in
the opposite columns produces the dual modulation, with the
same half-commutator phase in the two kernels. The rational
linear map \(b_\ell\mapsto\zeta b_\ell\) identifies their
adelic additive quotients with their probability measures.

Integration of the product of the line theta and the
conjugate line theta gives the inner product of their two
Schwartz vectors. Explicitly choose rational position and
frequency coordinates. Frequency integration first makes
the two rational summation indices equal. The sum over that
common index unfolds the position integral to the full
adelic position space, giving the inner product. A rational
polarization for this computation exists on the underlying
symplectic vector space by ordinary symplectic linear algebra.
Product tests are dense and the identity extends to all
residual Schwartz tests by continuity. On \(U(D)\) the
two line actions coincide and are unitary. At a split
place the splitting remaining after removing the core pair
is \(\eta(\det h)\). The horizontal factor is
\(\gamma_i(\det k)/(\gamma_i\eta^{-1})(\det h)\)
and the core factor is \(\gamma_i(\det k)/\gamma_i(\det h)\).
The character comparison at the other places follows exactly
as in Proposition~\ref{prop:linkedline265}, from the integral
vacua, rational theta functional and weak approximation.
Consequently removing the line leaves the core kernel with
test \(C_q f\).

For \(n\in N_S\), write its action on the full frame \(X\)
as \(Xv_X(n)\). The upper left \(r\)-square block of \(v_X(n)\)
is the action on the first \(r\) scaled columns. Its last
column above the diagonal consists of the coefficients of
\((n-1)\zeta\) in those columns. Hence
\(\chi(v_X(n))=\psi_{B,S}(n)\), including the last simple
root. Substitution in the upper-group integral gives the
conjugate of precisely this character. On the opposite
position \(0\) the Heisenberg displacement and quadratic
phase from \(N_S\) are both zero. Indeed the displacement
is linear and the phase quadratic in that position. The
upper triangular block has determinant one. Thus the entire
residual action of \(n\) at that position is the identity.
Disintegrating over \(R\) now gives the formula. Its Haar
measure is the product on \(N_S\) and \(U(C)\), by the
unitary determinant argument of
Proposition~\ref{prop:besorbit265}.

Before unfolding the additional line theta has polynomial
growth on Siegel sets of \(U(D)\), uniformly on a compact
unipotent fundamental set, by summing absolute Schwartz
majorants. The two cusp forms therefore still give absolute
convergence. On the \(R\)-fiber the exterior-vector estimate
of Proposition~\ref{prop:besorbit265} applies with the
isotropic frame of \(S\) wedged with a rational flag in \(C\).
The extra orthogonal line is fixed in the Levi subgroup.
It follows that the translated \(B\) supplies arbitrary
power decay in the \(U(C)\) chamber, uniformly for \(g\)
in a compact set. The contraction against \(q\) is continuous
on Schwartz space by integration of Schwartz seminorms.
These bounds justify the disintegrations and continuity.
Finally \(C_q(b\otimes q)=\|q\|_2^2b\). Integration in \(h\)
gives the descent, leaving the old Bessel period.
\end{proof}

\subsection{Integral coefficient families with controlled metrics}\label{nav:IV-36}
Keep \(N,B\) and the original lattices of
Lemma~\ref{lem:terminalclass259}. Let \(\lambda_\tau\) be the
terminal highest weight at an embedding, let
\(j_\tau=\lambda_{\tau,N}\), and let \(d_\tau=\sum_i
(\lambda_{\tau,i}-j_\tau)\) be the tensor degree of its
polynomial part. Write
\[
          d_0=\sum_\tau d_\tau,\qquad
          a_0=\sum_\tau(d_\tau+N|j_\tau|),\qquad
          L_0=\prod_\tau d_\tau! .
\]
Give the Schur image its subspace metric from the Hermitian
tensor power. Use the induced exterior metric on the determinant
line, its dual metric on negative powers and the tensor metric
across embeddings. For a list \((v_b)\) in a Hermitian space
\(W\), its frame operator is the positive operator with
quadratic form \(w\mapsto\sum_b|\langle v_b,w\rangle|^2\).
We use the same definition on \(W^*\).

\begin{proposition}[Two integral lists in each terminal class]\leavevmode
\label{prop:schurframe259}
On class \(c\) there are lists \((v_b)\subset W_c\) and
\((l_a)\subset W_c^\vee\), each of length \(N^{d_0}\), with
frame operators \(S_v,S_l\) in the complex coefficient spaces
satisfying
\[
 \begin{aligned}
 B^{-2a_0}I&\le S_v\le L_0^2B^{2a_0}I,\\
 B^{-2a_0}I&\le S_l\le B^{2a_0}I .
 \end{aligned}
\]
All members come from the integral Schur image and determinant
model of Proposition~\ref{prop:terminalcore255}.
\end{proposition}
\begin{proof}
Work first in one factor, and suppress \(\tau\).
Use the standard orthonormal basis of \(\CC^N\) and its
ordered tensors \(e_I\) of degree \(d\). Let \(Y\) be the
unnormalized integral Young symmetrizer of the partition
\((\lambda_i-j)\) used to define the Schur image. The frame
operator of all columns \(Ye_I\) on that image is \(YY^*\).
It commutes with \(U(N)\). Its image is exactly the
irreducible Schur module, because in finite dimension
\(YY^*\) and \(Y\) have the same image. Schur's lemma thus
gives a single positive eigenvalue on that image and zero on
its orthogonal complement. The positive eigenvalue is rational,
being the integer trace of \(YY^*\) divided by the rank, and
is an algebraic integer since \(YY^*\) is an integer matrix.
It is therefore at least one. The norm of \(Y\) is at most
\(d!\). Indeed it is the product of a row sum and an alternating
column sum of permutation operators. The row and column
stabilizers intersect trivially, so the products of their
elements are distinct permutations, at most \(d!\) in number.
Thus \(I\le YY^*\le(d!)^2 I\) on the image. For the empty
partition this calculation uses the single tensor \(1\).

At \(\tau\) transport \(Ye_I\) by \(X_\tau^{\otimes d}\)
from Lemma~\ref{lem:terminalclass259}. This map commutes
with permutations, maps the standard Schur space onto the
new one, and takes the columns into its integral Young
image. On the tensor power its largest singular value is at
most \(B^d\) and the norm of its inverse is at most \(B^d\).
Both bounds persist on the Schur space. Transport of a frame
by an invertible \(A\) takes its frame operator \(S\) to
\(ASA^*\). The resulting bounds are \(B^{-2d}\) and
\((d!)^2B^{2d}\).

For the dual list, take all \(d\)-fold ordered tensors of
the integral \(\xi_i\) of the lemma and restrict them as
linear functionals to the Schur image. A restriction is
integral on this image because it was integral on the full
tensor lattice. On the full tensor space the sum of the
squares of the evaluations on any \(w\) lies between
\(B^{-2d}\|w\|^2\) and \(B^{2d}\|w\|^2\), by the singular
values of the tensor product of the dual column map.
Evaluating on \(w\) in the Schur subspace gives the same
bounds for the restricted covectors. By Hermitian duality
these are precisely the two frame inequalities on the
dual Schur space.

In the vector list use \(x_1\wedge\cdots\wedge x_N\) for
a positive determinant power, and
\(\xi_1\wedge\cdots\wedge\xi_N\) for a negative power,
taking \(|j|\) tensor copies. Reverse these two choices
on the dual side. Each wedge has norm in \([B^{-N},B^N]\)
by the singular values of its column map, and it belongs
to the indicated integral line or its integral dual.
A zero power contributes \(1\). This gives the factor
\(B^{2N|j|}\) in an upper frame bound and its reciprocal
in a lower bound. Taking all products of lists over the
embeddings tensors the frame operators, so their positive
bounds multiply. There are \(N^d\) indices in each local
list. Multiplication gives the stated lengths and estimates.
\end{proof}

\clearpage
\section{References}\label{nav:bibliography-XIV}
\begingroup\def\section*#1{}\begin{thebibliography}{99}\small
\bibitem{XIV-XI-IV-molev} A. Molev, \emph{Gelfand--Tsetlin bases for classical Lie algebras}, in \emph{Handbook of Algebra}, vol. 4 (2006), 109--170. \href{https://arxiv.org/abs/math/0211289}{Section 2, Theorem 2.1 and formulas (2.21)--(2.22)}.
\bibitem{XIV-XI-IV-gross} B. H. Gross, \emph{Algebraic modular forms}, Israel J. Math. \textbf{113} (1999), 61--93. \href{https://doi.org/10.1007/BF02780173}{Proposition 4.3}.
\bibitem{XIV-XI-IV-gan} W. T. Gan, \emph{Automorphic forms and the theta correspondence}, Arizona Winter School lectures. \href{https://arxiv.org/abs/2303.14918}{Lecture notes}.
\bibitem{XIV-XI-IV-remark} L. Alp\"oge, \emph{Remark}, manuscript supplied with this series, 20 November 2026. \href{https://alpo.ge/remark.pdf}{Author's note}.
\bibitem{XIV-XI-harrislabesse806} M.~Harris and J.-P.~Labesse, \emph{Conditional base change for unitary groups}, Asian J. Math. \textbf{8} (2004). \href{https://webusers.imj-prg.fr/~michael.harris/harlab.pdf}{Authors' text}, Theorem~2.2.2.
\bibitem{XIV-XI-innerunitary806} T.~Kaletha, A.~M\'{i}nguez, S.~W.~Shin and P.-J.~White, \emph{Endoscopic classification of representations of inner forms of unitary groups}. \href{https://public.websites.umich.edu/kaletha/unitary.pdf}{Authors' text}, Sections~1.3, 1.6 and 1.7 and Corollary~3.3.2.
\end{thebibliography}\endgroup
\end{document}
